\documentclass[a4paper, 12pt, openright]{ctexbook} 

\usepackage{geometry}
\usepackage{hyperref}
\usepackage{fancyhdr}
\usepackage{titlesec}
\usepackage{enumitem}
\usepackage{setspace}
\usepackage{xcolor}

% --- 1. 字体设置 (思源系列) ---
\setCJKmainfont{Source Han Serif CN}[ 
    UprightFont = Source Han Serif CN,
    BoldFont = Source Han Serif CN Bold,
    ItalicFont = Source Han Serif CN
]
\setCJKsansfont{Source Han Sans SC}[ 
    UprightFont = Source Han Sans SC,
    BoldFont = Source Han Sans SC Bold
]

% --- 2. 页面与行距 (保持原版宽松布局) ---
\geometry{left=2.5cm, right=2.5cm, top=3cm, bottom=3cm}\setlength{\headheight}{15pt} % 修复页眉高度溢出的警告
\linespread{1.5} 
% --- 3. 元数据与链接设置 ---
\hypersetup{
    colorlinks=true,
    linkcolor=black, 
    filecolor=magenta,      
    urlcolor=blue,
    pdftitle={在传统与革命之间},
    pdfauthor={曼弗雷德·里德尔},
    pdfsubject={黑格尔法哲学研究},
    pdfkeywords={欧陆哲学, 法哲学, 黑格尔哲学},
    pdflang={zh-CN}
}

% --- 4. 页眉页脚 ---
\pagestyle{fancy}
\fancyhf{}
\fancyhead[CE]{\leftmark}
\fancyhead[CO]{\rightmark}
\fancyfoot[C]{\thepage}

% --- 5. 标题格式 ---
\ctexset{
    % 设置“篇”的格式
    part = {
        format = \centering\Huge\bfseries,
        name = {第, 部分},
        number = \chinese{part},
    },
    chapter = {
        format = \centering\Huge\bfseries,
        number = \chinese{chapter},
        name = {,},
        beforeskip = 10pt,
        afterskip = 30pt
    },
    section = {
        format = \centering\Large\bfseries,
        beforeskip = 20pt,
        afterskip = 15pt
    },
    subsection = {
        format = \large\bfseries, 
        beforeskip = 10pt,
        afterskip = 10pt
    }
}

% --- 6. 引用环境 ---
\renewenvironment{quote}
  {\list{}{\rightmargin2\ccwd \leftmargin2\ccwd \itemindent=2\ccwd \listparindent=2\ccwd}%
   \item\relax\small\kaishu\color{black}}
  {\endlist}

% --- 书籍信息 ---
\title{\Huge \textbf{在传统与革命之间}\\\Large 黑格尔法哲学研究}
\author{\textbf{【德】曼弗雷德·里德尔}}
\date{}

\begin{document}

\frontmatter

\maketitle

\input{chapters/总序.tex}

\tableofcontents

\mainmatter

\chapter*{初版前言}
\addcontentsline{toc}{chapter}{初版前言} % 手动加入目录
\markboth{初版前言}{初版前言} % 设置页眉

本集中的论文作于1962至1968年间，旨在重构黑格尔法哲学与政治的自然法的传统之间的联系，同时也重构向现代世界以及与之相应的工业革命和政治革命的那种转变。在时至今日仍然影响深远的法、国家和社会等概念的崭新规定中，黑格尔法哲学也引领了这种转向。这一目标的前提，殊异于流俗理解之下的《法哲学原理》的效果史及其对马克思主义、自由主义和纳粹主义这些全球性意识形态的诸般影响所提供的历史观。如果我们与时下主流观点相反，想要像作者理解他本人一样去理解他，由此出发进行解释的话，那么在我们对法哲学进行历史——考证的解释时，就会发现自己需要面对黑格尔的问题，这些问题是前人在黑格尔自己的时代意识和问题意识下向他提出来的。这里，我试图通过一种允许我们以柏拉图和亚里士多德以及霍布斯、卢梭和康德的眼光去重新阅读黑格尔法哲学的视角，对其被歪曲的效果史进行校正。唯有经过此番校正，才能对自黑格尔、马克思以来决定了现代法哲学和国家哲学的那种概念变异与视角改换的影响做出评价。

1962年面世的《黑格尔〈法哲学原理〉中的传统与革命》一文最先提出了后来进一步研究的主线。此文是献给卡尔·洛维特(Karl Löwith)65岁寿辰的。本书的所有研究也都敬献给他。

\begin{flushright}
M. 里德尔\\
1969年7月\\
于海德堡
\end{flushright}
\chapter*{新版前言}
\addcontentsline{toc}{chapter}{新版前言} % 手动加入目录
\markboth{新版前言}{新版前言} % 设置页眉

眼前这部《黑格尔法哲学研究》于1969年面世，旋即出了第2版以及各种译本，至70年代中期即告售罄。现在，对人们经常提到的对于新版的期盼，我再也不能置之不理了，盖因我相信，它今天又有了新的读者。因为所有这些主题：在实定法与自然法、经济与政治、市民社会与国家、现状与历史之间的近代二分以及黑格尔对其所做的概念上的消解，至今分量丝毫未减。同时，《法哲学原理》还透露了黑格尔逻辑学的一些秘密。它揭露了辩证法的历史核心。

笔者在60年代初构思这个主题时，最初只是想对黑格尔的《法哲学原理》与康德的《道德形而上学》（1797年）和费希特的《自然法基础》（1798年）进行一个结构上的比较。[在黑格尔这里发生的]基本政治概念（如“市民”、“市民社会”、“国家”）的意义断裂不容忽视。这种断裂使我从康德出发，经克里斯蒂安·沃尔夫(Christian Wolff)、莱布尼茨和霍布斯回到经院哲学的亚里士多德传统，最终回到亚里士多德的《政治学》。正是概念与其历史之间解释上的差距，向我展示了黑格尔政治哲学的另一基础。唯有如此，方有可能证明，他并不是那种导致（从R. 海谋到K. R. 波普尔）对《法哲学原理》错误阅读的威权主义的非理性的无底深渊。《法哲学原理》的另一基础就是近代革命的自然法原则，即在公民生活的风俗、伦常和伦理的传统原则之后出现的、包含在基督教之中的所有人的自由及其在法中的制度化。

为了表明这两个原则在黑格尔的以历史为中心的辩证法中的共同作用，现将书名改为《在传统与革命之间》。无须对这些研究进行改动，对此我甚感满意。第一部分第二篇（“制度中的辩证法”）是新加的。它是我于1976年在帕多瓦大学所做的一次演讲，迄今为止，只有意大利文文本（Quaderni di verifiche 2）。第二部分第三篇（“自由法则与自然的统治”）和第四部分第八篇（“黑格尔历史哲学中的进步与辩证法”）取自（同样售罄的）论文集《体系与历史》（1972年）。它们也是60年代晚期为了演讲目的而写的。

\begin{flushright}
M. 里德尔\\
1982年1月\\
于埃朗根-纽伦堡
\end{flushright}

\part{法哲学的体系建构}
\chapter{客观精神与实践哲学}

客观精神学说为黑格尔解读提出了一系列难题，因此对属于客观精神的法哲学做出更进一步的阐释之前，任何解读都必须首先应对这些难题。我们可以将其归为三个问题域：(1) 这一概念的形成问题；(2) 哲学体系中的安排问题；(3) 特别难以处理的问题，即黑格尔思想发展的不同时期对这一概念的处理问题。第一个问题相应于这一概念自身所包含的多义性问题。我在这里只强调在黑格尔的语言使用中交替出现的几个最重要的意义成分。客观精神是超个体、跨主体的精神，即两个或更多的主体或者说主体统一体（家庭、等级、团体、民族等）之共同物。在这种意义上，客观精神就是普遍精神（《哲学全书》第1版，第399节）。依照普遍意志的模式——这种普遍意志并不是从个体意志中归纳出来的，而是贯穿于这种个体意志之中，并且统摄这些个体意志的——客观精神将个体的、“主观的”精神包含于自身之中。不同于普遍意志的公式——这种公式表达了一种抽离了时间和生成的“内在”意志关系——，普遍精神作为客观精神处于外部现象之中。就此而言，对于黑格尔来说，客观精神本质上就是历史精神。而外在显现，即客观精神的对象性或者现实性，又包含了许多方面，其中概念的意义繁多，模棱两可，时而指对象性的或者现实的精神，时而指对象化的或者现实化的精神，最后还指对象化了的或者现实化了的精神（或者用对这种模棱两可不满的尼古拉·哈特曼的话来说，就是“客观”精神和“客观化了的”精神）\footnote{《精神存在的问题》，第2版，柏林，1949年，第196页以下。}。

在涉及黑格尔思想中的体系安排及其发展演变的问题时，我希望自己限于如下主张。客观精神概念只是在1817年《哲学全书》中才与众所周知的主观精神、客观精神以及绝对精神的划分一道引入。早期作品中没有这种划分。尽管青年黑格尔的种种体系安排遵循了由笛卡尔和康德所规定的现代哲学的理性、自然与精神的三一体模式，但是精神哲学还没有在自身内部进行这样的区分，而是最初冠以“伦理体系”和“实践哲学”这样的名称。此时“精神”还完全是“绝对的”，以至于黑格尔所知的唯一区分涉及的是“绝对精神”自身：在一族人民中，精神是“存在着的绝对精神”\footnote{参见《耶拿实在哲学》第2卷（Jenenser Realphilosophie II），1805—1806年，J. 荷夫迈斯特出版，莱比锡，1931年，第272页。}。直到《纽伦堡哲学入门》之《哲学全书》（1811—1812年）中才出现了区分，尽管此时黑格尔似乎尚未拥有后来那样的成熟概念。这里，黑格尔区分了精神的实在化与其完成两个领域（第128节），但是前者并不是在“客观精神”而是在“实践精神”的标题下（第173节）加以探讨的\footnote{参见黑格尔：《全集》（Werke）第3卷，H. 格洛克纳出版，斯图加特，1927年，第200、215页。引文来自格洛克纳出版的《黑格尔全集》（百年纪念版），斯图加特，1927年之后，以下简称《全集》第1卷，等等。}。对于客观精神学说的概念形成与体系安排来说，这具有决定性的意义。《纽伦堡哲学入门》草稿仍然在概念上坚持了一种为黑格尔以及受其支配的19世纪哲学史研究所逐渐放弃的见解：“客观精神”的起源及其与传统称之为“实践哲学”的那个哲学部门之间的联系。要解释这种联系，黑格尔精神哲学中用来对概念进行处理的那些观点是不够的。当然我们可以说，这些观点对于评价实践哲学在近代思想中所起的作用来说，完全是致命的。[其后，]一直到《纽伦堡哲学入门》都没有分开的客观精神与实践精神开始分道扬镳。在1817年百科全书体系中，实践精神成为了主观精神的一个组成部分，即心理学的意志理论（第368-399节）。毫无疑问，实践哲学陷于遗忘根本上要归因于精神科学中的心理学理论，这与黑格尔体系中实践精神消融于主观精神以及自然哲学与精神哲学的对立有关。受这种对立的影响，狄尔泰创立了他的学派。他试图证明，早在17世纪，随着意识分析以及人类学情感理论的发现，这样一种对立就已经开始取代理论哲学与实践哲学的“错误划分”\footnote{参见《精神科学引论》（Einleitung in die Geisteswissenschaften）第1卷（1883年），收入《全集》第1卷，第4版，斯图加特/哥廷根，1959年，第225、378页。}。由此，狄尔泰基于精神哲学的体系构想对青年黑格尔的作品所做的解释也就时常误入歧途。在这种情况下，就完全不会提出“黑格尔用客观精神学说取代实践哲学到底意味着什么”这一问题。而且同样也不会问，这种概念的转变在历史上究竟意味着什么，以及为什么黑格尔只有在其哲学发展相对较晚的时候才会形成这个概念。为了对这些问题进行回答，我们就必须反其道而行之：从重构青年黑格尔的实践哲学观念出发，到其消融于“客观精神”学说以及作为其核心部分的法哲学之中。我首先尝试把黑格尔置入现代思想与传统之争的思想脉络之中，以便确定黑格尔得以介入这一争论过程的确切位置。

\section*{1}

由亚里士多德所创立的实践哲学学科及其中世纪和近代早期传统的突出之处，恐怕不是其理论内容，而是其在这些世纪中所拥有的独特形式和历史效力。黑格尔在《哲学史讲演录》中已经指出了这一点，并且指出了产生这一效力的原因。他说：“一般说来，一直到最近，我们所见到的还不是真正的思辨性的实践哲学。”\footnote{《哲学史讲演录》第1卷，C. L. 米希勒出版，《黑格尔全集》第17卷，第291页。（中译文引自〔德〕黑格尔：《哲学史讲演录》第一卷，贺麟、王太庆译，商务印书馆1960年版，第248页。——译者）}实践哲学之所以恒定不变，是由于其非思辨的特征。即是说，实践哲学的诸多公理独立于“第一哲学”的前提。对亚里士多德及其支配下的经院哲学来说，尽管实践哲学运用了形而上学的公理，但是两者之间并不存在直接的奠基关系。随着黑格尔所言的那种现代的到来，也就是17和18世纪自然法的发展，这一情况发生了改变。此时理论延伸到实践哲学的结构之中，实践哲学也提出了那种为现代科学方法论理想所推崇的普遍论证的要求。对于霍布斯而言，实践哲学的任务在于，如同认识图形的大小比例般，以同样的精确性认识人类行为的诸种关系。如同几何学一样，实践哲学现在也应当成为一门先天的、可以演证的科学。在霍布斯看来，实践哲学的对象非常适合于这一目的，因为对人来说，只有人自己产生的那些事物，才能有一种先天的、出于“第一根据”（ex causis）的科学\footnote{霍布斯：《论人》（De homine），第10章，第5节。}。

现代实践哲学的“思辨”特征在于这种对几何学认识理想的认同，而这同时也是批判实践哲学传统的第一步。第二步则是从实践哲学自身的经院形态出发，这种经院形态在18世纪初将它从外部通过先天演绎的方法论要求施加给它的批判原则纳入自身之中。沃尔夫对这一方法的运用导致实践哲学的体系分裂为一个“理性的”部分和一种“经验的”部分。理性部分对一般人类行动的区分以及权利和义务原则进行先天的证明，并且模仿笛卡尔的普遍数学（mathesis universalis）的观念，冠以普遍实践哲学（philosophia practica universalis）\footnote{《按照数学方法撰写的普遍实践哲学》（Philosophia practica universalis mathematico methodo conscripta），莱比锡，1703年。那篇莱比锡辩论论文——沃尔夫说是他写的，因为笛卡尔未曾研究过这一领域[《讲授纲要》（Ratio praelectionum），第6章，第3节]——构成了《普遍实践哲学》（第1-2卷，法兰克福/莱比锡，1738-1739年）的体系基础。}之名。这一新科学的引入为近代思想与实践哲学传统之争的第三步做好了准备。这一步由康德迈出。康德在“著名的沃尔夫之道德哲学导论”中清楚地意识到，这门新科学距离实践理性的批判有多远。这门新科学的引入源于对来自传统的概念和对象进行论证的需要，而普遍实践哲学则在方法和内容上仍然保持了与这些概念和对象的联系。依照康德的观点，这也是它没有进入一片“新领域”的原因之所在，因为它作为普遍的实践哲学仅仅把意志的普遍概念、“一般意愿”作为线索，因而并不考察一种“特殊”意志的可能，“比如一种无须任何经验性的动因完全为先天原则所规定，人们可以称之为纯粹意志的意志”\footnote{《道德形而上学奠基》（Grundlegung zur Metaphysik der Sitten，1785年，科学院版），第4卷，第390页。（中译文引自〔德〕康德：《道德形而上学奠基》，杨云飞译，邓晓芒校，人民出版社2013年版，第6页。——译者）}。

这种“特殊的”、从所有自然联系中摆脱出来的意志的发现，使得康德用实践哲学自身论证的观念取代了对实践哲学“特殊部分”（即传统的经院形态）的论证。同时，这种自身论证的观念也使得对于沃尔夫的普遍实践哲学来说具有构建意义的向传统流传下来的概念和对象的回溯成为了不可能的事情。康德《实践理性批判》（1788年）的主题在于，实践知识的最高条件不再像沃尔夫的“导论”所言的那样从自然规律中，而是从那个纯粹意志的理念中演绎出来，此纯粹意志自身在其自由中即是法则。自由的法则摧毁了实践哲学传统结构的一贯性，以及由沃尔夫再次建立起来的传统实践哲学的“普遍”部分与“特殊”部分之间的联系。政治学、家政学和伦理学作为明智和幸福的理论，完全被排除在实践哲学之外；对康德来说，它们都只是实用科学和技术性科学，属于经验性的经验世界、明智和技巧的活动领域。它们可能会从自然联系中借来规则，以便每次都能产生出按照因果律来说是可能的结果；然而它们并不能从自然联系中借来那些——如在旧的明智理论中——赋予伦理和政治行动以伦理义务的力量的目的。康德的自由法则终结了这种借用。依照康德的观点，只要它包含了所有行动的先天规定根据，那它就是一种“被称为实践哲学的特殊哲学”的对象，从而与那些违法地占据了“实践”之名义的（康德意义上的）实用学科和技术性学科相对立\footnote{《判断力批判》，导言。《全集》（科学院版）第5卷，第171页以下。}。

黑格尔在康德自由概念及其理论前提（即《纯粹理性批判》中发现的自我意识的自发性）的基础上，对实践哲学的经院传统进行了批判。这一批判在最大程度上超越了这个传统。辩证法完成了几何学方法、理性方法和先验方法未能完成的事情：传统的瓦解。这并不是因为黑格尔思想的尖锐批判和激进（在这个方面，霍布斯和康德都要明显胜过他），而是因为它们得以产生的条件。我们已经指出第一个条件：康德的意志自律的思想。黑格尔赞同这一思想，这使得他对传统的批判从一开始就超越了霍布斯和沃尔夫的批判。因为在霍布斯和沃尔夫那里，自由问题以及亚里士多德传统中与自由问题相关的主奴关系并没有得到解决。第二个条件是把现代国民经济学纳入实践哲学的构建中，因此也就是尝试将实践哲学的传统原则嵌入社会和工业世界的历史当下经验之中。由此产生的最重要后果就是黑格尔对制造（$\pi o\acute{\iota}\eta\sigma\iota\varsigma$）与行动（$\pi\rho\tilde{\alpha}\xi\iota\varsigma$）之间关系的思辨瓦解。尽管霍布斯已经将两者间的传统界线相对化了，但还没有真正克服，因为这种相对化还局限于政治领域，也就是人造的作品“利维坦”之内，还没有进入经济领域。在康德那里，这种情况再次出现。康德试图把当时刚刚确立的“国家经济学”从“特殊的”实践哲学中完全排除掉。最后，第三个条件，尽管它已经是黑格尔思想的成就，而不只是一个前提，但它涉及与历史维度的联系。这种联系使得实践哲学的范畴从完全无历史的经院传统中摆脱出来，使它们进入——将自己带入一种与这些范畴的内容完全改变了的关系之中的——“概念”的那种辩证运动之中。

同时，由黑格尔所引起的实践哲学史上的转折是从一个悖论开始的。如果我们想要理解这种转折的影响，那就必须认识并且说出这一悖论。这个悖论就是，一方面，黑格尔与传统保持了最遥远的距离，同时却又最先重新发现了一条直接通往这一传统的通道。这一点将他与传统的争论同霍布斯、沃尔夫和康德等人对传统的不同批判区别开来。就他们——如沃尔夫和康德——与亚里士多德主义经院哲学传统的一种特定的、相对较晚的阶段发生联系而言，他们与亚里士多德要么处于一种历史上相距遥远的关系中（霍布斯），要么处于一种仅仅是间接的关系中，青年黑格尔则研究了亚里士多德伦理学和政治学\footnote{黑格尔在斯图加特文理中学时期就开始研读《尼各马可伦理学》。参见罗森克兰茨（K. Rosenkranz）：《黑格尔传》（Hegels Leben），柏林，1844年，第11页。黑格尔对亚里士多德的真正研究（借助于16世纪的巴塞尔老版）则要追溯到在图宾根神学院的那些岁月。参见D. 亨利希：《洛伊特魏因论黑格尔：一份黑格尔传记文献》（D. Henrich, Leutwein über Hegel. Ein Dokument zu Hegels Biographie）一文中洛伊特魏因的证明，见《黑格尔研究》（Hegel-Studien）第3卷，1965年，第58页。《政治学》阅读的影响反复出现在耶拿时期的早期作品中。参见《全集》第1卷，第111、114、485页以下、494页以下、511页。《政治学与法哲学文集》（Schriften zur Politik und Rechtsphilosophie），G. 拉松出版，莱比锡，1913年，第417页以下、442页以下、460页以下、478页以下。当指亚里士多德的政治学和伦理学。——译者}，由此作为德国最早的人士之一，对传统获得了一副比那些批判者要恰当得多的形象，[反之，]对[霍布斯、沃尔夫和康德]这些批判者来说，这些渊源要么完全不存在，要么极端模糊不清。但是，这一过程中至关重要的倒不是黑格尔大体上重构了古典理论，而是他试图把古典理论同时运用到自己的思想以及当下历史经验向实践哲学所提出的任务上。黑格尔耶拿时期的著作即基于这种立场。他要求应用古典传统，从而不得不将一大堆材料放到一些概念之下，这些概念通常并不是来自这些材料，而是有着相反的来源。因而在黑格尔处理实践哲学的第一阶段（1801—1804年），概念与实在、形式与内容仍然是完全分裂的。尽管在第二阶段（1805—1807年）他逐渐对自己的哲学基础及其与实践领域的关系有了清晰的理解，但是只有在第三阶段——我们可以把两卷本的《逻辑学》以及《哲学全书》的出版（1816—1817年）作为这一阶段的开端——他才真正形成了“客观精神”这一概念。随着这一概念的形成，他克服了实践哲学传统的原则和体系形式。

\section*{2}

黑格尔在耶拿时期之初回归古典政治哲学并非什么意外之事。在18世纪90年代，黑格尔主要集中于神学和历史研究时就已经隐隐约约展现了这种倾向，但其意义有所不同。黑格尔早期作品的吊诡并不在于一般而言他又重新找到了回归古典城邦伦理的道路，而在于他同时按照康德《实践理性批判》的概念（自由、自律、自发性等）来对它进行衡量。在他看来，依据实践理性而行动和信仰的个体与古代的共和主义者都是自由的；前者如他的“精神性的、不朽的”理性本质所要求的那样来规定自己，后者则“在其人民的精神中”履行他的义务\footnote{《早期神学著作集》（Theologische Jugendschriften），H. 诺尔出版，图宾根，1907年，第70页。参见埃弗拉伊姆（F. Ephraim）：《黑格尔早期作品中的自由概念研究》（Untersuchungen über den Freiheitsbegriff Hegels in seinen Jugendarbeiten），见《哲学研究》（Philosophische Forschungen）第7卷，柏林，1928年，第59、63页。近著参见歌兰德（I. Görland）：《青年黑格尔的康德批判》（Die Kant-Kritik des jungen Hegel），美因河畔法兰克福，1966年，第4页以下。}。康德《道德形而上学》（1797年）的出版结束了这种混杂，因为康德的这部著作并没有为自律思想向政治世界这样一种拓展提供什么凭靠，并且遭到了法兰克福时期黑格尔（1798—1799年）的批评。黑格尔从古代实体性伦理、从一种人民宗教的完美自由中理性和自然的同一（Einssein）出发；但是古代的自由既没有构成他自己政治意愿的标准，他也没有用来解释康德的自律概念\footnote{《早期神学著作集》，第221页以下。}。与之相反，青年黑格尔的批判意图在于反对康德（同时还有费希特）通过将自由法则和自然规律割裂开来而未能为前者提供充分的范围，因为这种割裂导致实践理性自身饱受“差别”之害。由于黑格尔把康德和费希特哲学中的自由和自然之间的差别关系——黑格尔将这种关系解释为强制和压制的关系——置于批判的中心，并在一个政治整体的理念中寻求将两者结合起来的可能性，这就迫使他回到康德之前的政治哲学传统。

黑格尔在耶拿时期的早期著作中对康德和费希特的批判，就是从这一吊诡开始的。把自由概念扩展为一个政治整体，这一政治整体使个体能够“自由、理想地行动”，这就要求：现在，在伦理理念的建构中，那种遭到压制的“自然”应当重获它的“权利”。《论自然法的科学探讨方式》一文即以此为出发点，这篇论文的标题明白地提到经院哲学的体系形式：它“在实践科学中的地位”\footnote{参见拙文《黑格尔对自然法的批判》（Hegels Kritik des Naturrechts），载《黑格尔研究》第4卷（1967年），第177页以下，本书第42页以下。}。黑格尔在第一部分论述了这门科学的危机。自霍布斯以来，通过对其传统的批判，这一危机就存在了。自然法——同样也包括康德拒绝承认其为实践哲学的政治学和国家经济学——遭受这样的命运，即在现代思想历程中，人们仅将“哲学的哲学成分”归于形而上学，以至于这些科学“完全与理念无缘”\footnote{《全集》第1卷，第437页以下。}。黑格尔认为，对于将建立在经验和技术应用之上的那些科学从形而上学中解放出来，以及形而上学脱离自然、历史中现实物的现实性这样的做法，康德的批判哲学不仅无法逆转，相反，还把它们推向了无路可走的极端。因为批判哲学将自然作为非理性的杂多与作为纯粹统一性的理性对立起来，因而对康德的批判哲学来说，理性——与之一道，还有道德哲学的“纯粹意志”——成为了“一的无本质的抽象”；与此相反，理性则将自然蒸发为“杂多的无本质的抽象”。由此出发，黑格尔返回到了一种历史上先于这种对立的传统形态。这就出现了值得注意的现象：黑格尔试图重新恢复实践哲学中那些从形而上学中解放出来的科学，他不仅接受它们旧的概念形式，而且甚至为作为这些旧的概念形式之基础的“旧的、完全不一贯的经验”辩护。黑格尔在忆及那些将实践哲学的经验置于几何学的认识理想之下，因此必然发现实践哲学的概念形式“既不连贯又相矛盾”的批判家们时说道，一种伟大的、纯粹的直观能够“在其展现的纯粹建筑（在这个建筑中，必然性的联系和形式的统治并不是显而易见的）中将真正的伦理事物表达出来，各个部分以及那些自我修正的规定性的排列组合暗示了那种尽管看不见然而却是内在的、合理的精神”\footnote{《全集》第1卷，第453页以下。}。

在这种直观中，我们很容易发现亚里士多德实践哲学的自由建筑术的特征。通过返回到这种直观，黑格尔对现代自然法的原则和古典政治学的原则进行对照。个体的存在最清晰地表现了两者在出发点上的差异。在城邦伦理中，个体既不执着于自身，也不是通过其自然的特殊性（社交冲动、对死亡的恐惧、自我保存等）与他人结合为一个社会；他不是以抽象的方式——通过契约和个体意志的服从，而是自然地（$\phi\acute{\upsilon}\sigma\varepsilon\iota$），或者如黑格尔所言，以“活生生的”方式与伦理整体合而为一。与将“个体的存在设定为首位和最高原则”的现代自然法相对，古典政治学认为“个别性自身纯属虚无，它与绝对的伦理尊严纯为一物——唯有这种真正的、生动的、并非服从的同一才是个人的真正伦理”\footnote{同上书，第452页。}。我们可以在黑格尔早期实践哲学构想中观察到这种接近古典政治学概念的倾向。在他对伦理学的评价中，这种倾向也表露无遗。尽管伦理学作为德性论，必须是个人的伦理，但是“绝对”伦理也必须在个体上面表现出来，个体完全是否定性的东西，居于否定的形式之下；不仅现代主体性的反思性道德，而且外在地进行自我保证的、被嵌入伦理整体中的个体德性（黑格尔按照古典模型对其进行了自然描述\footnote{同上书，第513页。}）也是“否定的伦理”。这构成了“道德学”与自然法之间的形式区分，然而这种区分并非如康德和费希特所理解的那样，“仿佛它们是分离的，似乎道德学已经从自然法中排除出去了；毋宁说，道德学的内容完全就在自然法之中”\footnote{参见《伦理体系》，《政治学与法哲学文集》，第465页；《全集》第1卷，第510页。}。

这个根本而言为古典的自然法概念及其在实践哲学中“地位”的困境在于，它导致了个体存在以及作为消解和运动环节的否定物的消失。这种困境导致黑格尔无法对现代自然法做出合乎历史的恰当理解，而当黑格尔把亚里士多德关于劳动和行动的理论作为其自然法构建要素接受下来时，这种困境延续了下来。按照黑格尔的理论，伦理的保存要归功于那些个体的德行，个体的活动并不指向一件“产品”——亚里士多德的作品（$\acute{\epsilon}\rho\gamma o\nu$），而是直接在其规定性中摧毁之；它是“没有目的，没有需求，也不涉及与实践情感的关系，没有主观性，没有涉及占有和取得的主观性；相反，此种主观性的目的和产品与其自身一道终止”\footnote{《政治学与法哲学文集》，第467页，另见第463、466页。正是在这里，“步枪”作为一种普遍的、无人身的勇敢和“无人身的死亡”的特别现代的（“机械的”）发明，直接出现了（第467页以下）。}。德性的行动是“政治家和战士”的“绝对无差异的行动”，在《伦理体系》中，黑格尔想将其划给一个“绝对的等级”，我们很容易重新认出这就是古典意义上的市民等级、生而自由的、英勇善战的男性等级。与之相对的是从事“差异”事务（与自然关系相关的“经济”事务）的“正直等级”，这一等级从事“需要的劳动，占有、取得和财产”。这就是现代意义上的不参与战争、局限于私人生活的布尔乔亚（bourgeois）的市民等级。在黑格尔看来，随着奴隶关系的消失，产生了这一等级：“这里，对占有和特殊性的这种沉迷，不再受役于（‘绝对等级’的——笔者）绝对无差异。它尽可能地无差异化了，或者说形式的无差异，个人的存在（Personsein）在人民中反映出来，占有者通过其差异，没有没入整个共同体之中，因此没有陷入人身依附中；相反，其否定的无差异被思为某种实在之物，因此他是市民，是布尔乔亚，并且作为普遍物而得到承认。”\footnote{《政治学与法哲学文集》，第473页。}

另一方面，黑格尔又并不想像亚当·斯密在《国富论》（1776年）中所做的那样，去颠倒行动（Praxis）与制造（Poiesis）的传统关系\footnote{参见《国富论》（The Wealth of Nations）第1卷，第10章；第2卷，第3章。}。作为新兴的国民经济学理论家，斯密勾画了摆脱封建现代剥锁的现代市民社会的模式。在他这里，在古典政治看来保存了城邦的政治家、法律人与战士的行动，被降格为一种非生产性的活动。这种活动没有价值，因为它并没有在一件“作品”中，也就是在一个持续的对象或是一种商品中将自己实现出来。相反，对黑格尔来说，劳动和行动的关系本身不仅仍然有效，而且在其第一个实践哲学体系草案中，构成了他吸收现代国民经济学以及进行等级划分的出发点。这些草稿的大部分矛盾和晦涩可以归结为这种情况，即黑格尔将古典政治图式作为他分析国民经济学的指南。一方面，他将古典政治图式作为范本加以接受；另一方面，他又不得不对它进行修改和补充。而从古典政治图式本身来看，这些修改和补充并不就是一目了然的。

这些矛盾最为明显地出现在《论自然法的科学探讨方式》（1802—1803年）这篇论文中。这里，通过援引柏拉图和亚里士多德\footnote{参见《全集》第1卷，第495页以下。}，《伦理体系》中的“绝对等级”表现为“自由人等级”。不同于“非自由人等级”，这一等级天生就是独立的或者自为的，而非自由人等级则“按其天性不属于自己，而属于他人”\footnote{译自亚里士多德：《政治学》第1卷，第5章，1254 b 20-23；此外，参见同书第1卷，第4章，1253 b 27-30；《形而上学》第1卷，第2章，982 b 25-26。}。黑格尔继续论述道，第一等级的“劳动”并不限于个别的特定产品，而是指向“伦理有机体之整体的存在与维持”，并将自己展现为整体的“鲜活的运动和……神圣的自我享受”\footnote{《全集》第1卷，第494页以下。黑格尔直接提到了柏拉图和亚里士多德的政治理论：“亚里士多德为这个等级规定的事务，就是希腊人所言的‘$\pi o \lambda \iota \tau \epsilon \acute{\upsilon} \epsilon \sigma \theta \alpha \iota$’，在人民中生活，同人民生活，为人民生活，过一种普遍的、完全属于公共事务的生活，或是从事哲学思考。这两种事务，柏拉图以他更高的生命力认为它们是无法分开的，是完全连在一起的”（第495页）。}。因此，黑格尔赞同传统的实践优先的理论，实践的目的就是纯粹的运动（$\kappa\acute{\iota}\nu\eta\sigma\iota\varsigma$）和活动（$\acute{\epsilon}\nu\acute{\epsilon}\rho\gamma\epsilon\iota\alpha$），而不是一种产品的规定性\footnote{参见《形而上学》第9卷，第6章，1048 b 18-36；第8章，1050 a 21-31。}。精神概念最初在这个图式中得到了实现，并且也正好将表达了伦理行动的自我返回的（“内在的”）运动表达了出来。他在1803—1804年耶拿讲座中讲道：“人民的精神必须自身永恒地成为作品。或者说它只是作为一种朝向精神的永恒生成。”\footnote{《耶拿实在哲学》第1卷，J. 荷夫迈斯特出版，莱比锡，1932年，第232页。}然而，对黑格尔来说，在一族人民中，那种每次从自身中产生出一个产品的活动——现代意义上的与自然打交道的劳动（在这种交道中，运动超出自身，凝结于作品中）——也属于精神的这种“在作品中存在”（Am-Werke-Sein）。因此，黑格尔也在自我组织的伦理之内为“为了个人而劳动”的非自由人等级指派了一个（尽管是从属性的）地位。在黑格尔看来，自我组织的伦理必然被“划分为一个绝对接纳无差异之中的部分（公民的实践，其作品为维系城邦的政治活动自身——笔者），以及一个实在之物本身在其中是存在着的，因而是相对同一的部分（奴隶、手工艺人等生产出产品的制造活动[$\pi o\acute{\iota}\eta\sigma\iota\varsigma$]），这个部分只在自身中带有绝对伦理的反映”\footnote{《全集》第1卷，第494页。}。

不同于现代自然法理论的体系基础，在黑格尔的思想中是政治经济学体系的劳动概念提供了相对的辩护。这种辩护，我们在自然法论文中即已见到，在实践哲学的新建筑术中也同样清晰可见。一方面，黑格尔希望通过一种“伦理自然”法的建构，消除自然法与政治学之间的近代分裂以及私法范畴（比如说契约）的主导地位——特别是在国家法和国际法中，这些私法概念似乎已经变成了“自然法的家庭内部事务”\footnote{同上书，第410页以下。}。这样，在黑格尔试图建立起来的实践哲学体系中，自然法又重新与政治学和伦理学一道，落入根本而言需要从政治出发得到理解的“伦理”概念之下\footnote{同上书，第509-513页。}。同时，对黑格尔而言，政治经济学也是一门非常重要的“实践”科学。由此，它也就取代了旧的家政学，显现为体系的基础。其次，接踵而至的是与之相关的“形式”法。于此，我们主要看到的是旧的、处理私法材料的自然法，最后便是与前两者分离开来的、因成为“绝对”而摆脱了一切经济学和形式之“纠缠”的伦理：“这个领域的概念，作为实在的实践[领域]（Praktische），主观地看，就是感觉或者身体的需要和享受，客观地看，就是劳动和占有；这个实践[领域]（如按照其概念发生）被接纳到无差异当中，就是形式的统一性，或者法权（Recht），法权在这个实践[领域]中是可能的。第三个东西作为绝对物或者伦理，超乎二者之上。”\footnote{《全集》第1卷，第494页。}

毫无疑问，黑格尔由于采取了这样的一种划分方式，既居于传统的中心，又与传统相冲突。与18世纪对实践哲学的材料进行重组的企图相比，黑格尔将体系划分为政治经济学、法（即自然法）与伦理（即政治学和伦理学）的做法已经在自身中固定下来。而这又与自然法论文中政治经济学体系所发现的独特的承认的“等级”形式有关。就此而言，这种体系安排最清晰地表达了黑格尔建构其首个伦理体系的鲜明意图：将“绝对伦理的理念”与一个“贵族和自由人等级”的伦理现象混同起来，并将劳动等级的“相对伦理”与之对立起来。

黑格尔通过把劳动的领域从前述限制中解放出来，而超出了这种安排。这种超出发生在《伦理体系》（1802年）第一次阐述之后（此次阐述在许多方面都不充分），以及在1803—1804年和1805—1806年的“实在哲学讲座”中。在这两次讲座中，劳动与语言和行动（承认）一起，作为中心环节出现在精神的构成中。精神，并非如自然法论文所假定的那样\footnote{同上书，第509页以下；另见《信仰与知识》，同前，第424页。}，等同于“伦理自然”。同时，“伦理自然”也并不等同于一族“人民”的身位（在这一身位中，精神反思它自身，并且在理智直观的媒介中将自然带回到自身）；相反，取决于对这种带回[自然]本身进行反思，以便使精神的“教化”过程成为主题。因为黑格尔认识到，这种教化过程也同与自然打交道连在一起，伦理行动并不能穷尽其含义，由此，对他来说，“劳动”也就成为了实践哲学的问题。黑格尔勾画了一种劳动理论的基本特征，由此便将其自身的重要性给予了亚里士多德以“制造学”（Poietik）之名进行处理的那个哲学部门。自亚里士多德以来，对这个哲学部门的探讨便没有任何进步，在现代自然法理论（霍布斯、卢梭和康德）中，它也没有取得任何进展。新的制造学是把国民经济学与先验观念论联系在一起的一种结果，这种联系在青年黑格尔的思想道路中就已经有所准备，并在耶拿讲座中得到了实现。劳动本身构成精神的一种基本形式，作为这样的形式，它在实践哲学的体系中不再被贬低为一个次要的环节。

与哲学传统相对立，黑格尔既没有把实践行为局限于人与人之间行动的概念，也没有像康德和费希特那样将其局限于与自身感性相对立的道德主体性的内在活动中，而是将其扩展到人与自然或者（用先验哲学的普遍化术语来说）——主体与客体、自我与非我打交道的维度中。对劳动概念的这种新的解释导致实践哲学的原则和构造形式发生了根本性的转变。在其第一个实践哲学建构中，黑格尔通过个体的[伦理]自然法存在，以及通过将“个体的”劳动以等级的方式排除在伦理理念之外，消解了康德哲学的开端；其新的实践哲学则与他最初的实践哲学建构完全对立，以《纯粹理性批判》在统觉的原始综合统一的思想中阐发的先验自我的概念、以必然能够伴随我的一切表象的“我思”为出发点。然而，黑格尔并不是以康德的眼光，而是以——将这种伴随的自发性及其造成可能经验的统一性的成就解释为先验的“创造”的——费希特的眼光，来看待我思的功能。在从18世纪向19世纪过渡之际，费希特的这种解释使哲学卷入了一场撤除迄今为止各门科学之间所有教条主义樊篱的运动之中。青年黑格尔追随费希特和谢林，最先在耶拿讲座中制定出这门新的作为这场运动之表达的“意识的”先验“历史”的科学。黑格尔与费希特一样，执守先验哲学的“主体”的行为及其产物。这些产物规定了存在的意义，并且不再规定那些自笛卡尔以来即导致了那种没完没了的、在康德学派也还没有得到解决的争端的实体本体论的前提：质料或者形式、物或者自我、存在或者自我意识、主体或者客体是否拥有独立的持存。这些互相排斥的规定尽管一方面从古典的亚里士多德本体论的毁灭中产生出来，但是它们却无法逃脱作为一种独立或不独立、能动的或者被动的存在者的前提的实体本体论原则的魅力。对黑格尔来说，此类规定都是形而上学的抽象物，它们不仅错失了精神的教化过程，而且也错失了作为这些对立面之统一的精神的概念。诚如黑格尔所言，对于这样一种“非理性的争论”存在应该建立在意识的“被动性”还是“能动性”之中，“真正说来没有什么理性的东西可言”。实在论者主张形式、意识的被动性，观念论者把这种被动性归之于质料、存在。然而事实上，“真正说来”，这一争论的中心就是“中项的在自身中进行争论的潜能阶次（Potenz），其中诸规定性自身与关联项（Bezogensein）两者既设定为一又设定为有区别的东西”\footnote{《耶拿实在哲学》第1卷，第214页以下。}。黑格尔不是在其对立中思考对立的两个方面，而是将其思为诸对立的统一，但是他是在先验意识的基础上思考这种统一的。先验意识在其自身中融解了形而上学的僵化谓词，使之流动起来。因此，这种自我意识不仅仅是那种“伴随”表象的运动的“我思”的自然统一性，而是自身即处于运动之中的。黑格尔将这种运动表现为自我的历史，表现为自我与他者在语言、工具和作品这些中项之中的中介活动\footnote{同上书，第213页。}。在黑格尔这里，劳动——主体同对象打交道以及主体在作品中的对象化——也就成为意识的中介活动史的一个环节，而作品——作为工具和占有——则成为形式与质料、主体与客体、能动与被动等对立之间的中项。在这一中项中，意识（其自身就是打交道过程中的中项）自我组织起来，并且自我“教化”（bildet）。劳动和教化彼此共属一体。

黑格尔由此颠倒了旧实践哲学的基础。他解决了隐藏在劳动和行动关系规定中的困境。这一困境可以这样说明：如果活动者的活动（$\acute{\epsilon}\nu\acute{\epsilon}\rho\gamma\epsilon\iota\alpha$）为其自身的确证（也就是在伦理—政治实践的情形下）的话，那么这种活动就不能对象化；另一方面，如果活动对象化了的话，那么它就不在活动者当中，而是在其作品中。在作品中，活动同时就已经熄灭了。按照制造（poiesis）的古典模式，在劳动与劳动者之间并不存在任何反思性联系，也没有黑格尔的教化概念所表达的劳动者与作品、主体与客体之间的那种结合和摆动。原因在于活动本身的等级秩序，依据这一等级秩序，行动要高于劳动。实践的明智行动更加完善，因为它在实现出来时将行动者的完善表现了出来，而劳动者则似乎将形式沦丧于质料之中，因此也就不能导向劳动者的完善：“不是制造者的完善，而是作品的完善（non est perfectio facientis, sed facti）”\footnote{托马斯·阿奎那：《神学大全》第I卷第2册，第57题，第5节。}。在技术—实践的意义上与自然打交道时，制造（poiesis）生产出来的东西，并不是服务于作为生产者的生产者——比如说，使他在行动的过程中认识人类的事务——而是服务于这种实践对制造之物的使用。前工业社会的模式就是建立在这种使用之上，这就意味着，建立在那种按照定义剥夺了生产者的明智（phronesis），也就是实践智慧的能力之上。黑格尔通过他对国民经济学劳动过程的分析摧毁了这一模式。他的新模式取代了旧模式，它建立在经过先验哲学改变了的本体论基础之上，黑格尔的外化范畴则凝聚了先验哲学的这个本体论基础。劳动的过程无非就是进行制造，同自然打交道，对外物进行改造。然而黑格尔并不将这一过程理解为把一种外在于劳动者本人的形式转加到一种外在的质料之上，而是理解为劳动者的外化。而且，他还以双重化的方式进行理解：(1) 作为欲望的外化。欲望是“内在地扎根”于个体之中的东西，就此而言，也就是一种普遍的东西；(2) 作为同样方式的“内在”形式的外化，也就是意识的外化。黑格尔把劳动定义为“意识把自身造就为物”（Sich-zum-Dinge-Machen des Bewußtseins）\footnote{《耶拿实在哲学》第2卷，第214页。}。在劳动中，意识失掉了其主观的形式，以便在更高的层次上（在作品的客体化中）获得它，也就是说直观它。由于主体通过劳动而适应自然的联系，因此劳动过程就将主体外化为物；同时，因为主体也在加工的对象中使自己成为对象，也就从外化中摆脱出来。这种双重否定的运动不仅消除了质料和形式的简单对立，而且消除了在古典模式结构中发挥了重要作用的使用的优先性。概而言之，不同于亚里士多德以及追随他的前工业的制造学传统（技术理论或者工艺学），黑格尔并不是从劳动过程的终点而是从它的起点出发，对它进行解释：劳动源于欲望的否定物，在劳动的进展中，它一直都盯着欲望，同时也坚持要否定其所加工的对象：劳动是“受阻的欲望”，并且作为这样的东西而施与教化——否定之否定。

由此出发，黑格尔在《精神现象学》的那个著名章节阐述了主奴辩证法。在一定程度上，主奴辩证法也可视为劳动和行动的辩证法。从内容的阐述来看，这一章紧接在城邦伦理的结构之后。在城邦伦理中，相对于劳动等级及其陷入存在者的个别性之中的沉沦存在，“自由人等级”基于其“赴死的能力”——以及战士的实践——而获得了绝对的优先性。现在，黑格尔认为个别性环节已经在劳动中得到了扬弃，并且为该环节在对象中的持存进行了辩护，这就表明，黑格尔在此期间究竟在多大程度上超出了这种立场。在黑格尔看来，正因为主人对奴隶的统治与一种自身并不产生任何产品的行动相连，因此它就必须瓦解。尽管如此，古典意义上政治行动的环节也没有在此章中再度出现。因此，物品生活（Gut-Leben）不是解释为使用，而是解释为享受，由此，黑格尔就把他自己模式的出发点——欲望与劳动的联系——回植到亚里士多德的模式之中。这样一来，依照黑格尔的观点，在主人这一方面，那个按照古典模式为行动比劳动优越奠基的环节就具有一种缺陷：不能在一个产品中对象化。对于辩证的分析来说，这种观点证明它自己只是假象。因为物在劳动中保持了它的独立性，因此“对于物的非本质的联系方面”便落到了奴隶上面。对于主人来说，对象在享受中消失了；而对于奴隶来说，对于对象的否定联系便成为了“对象的形式并且成为了一种有持久性的东西，这正因为对象对于那劳动者来说是有独立性的”\footnote{《全集》第2卷，第156页。（中译文引自〔德〕黑格尔：《精神现象学》上卷，贺麟、王玖兴译，商务印书馆1979年版，第130页。——译者）}。劳动并没有消逝在加工的对象中，而是构成了欲望和对象之间的持久运动，“否定的中项”或者“赋形的行为”，对劳动者来说，这种行为就像古典观点中行动对于行动者一样内在。“这个否定的中项或赋形的行为同时就是个别性，或者意识的纯粹自为存在，意识现在在外在于它的劳动中进入持久的要素之中；由此，那劳动着的意识就直观到那个独立的存在就是它自身。”以及“由于形式被设定在外，所以对服务意识来说，形式并不是一个有别于它的东西；因为这形式正好就是它的纯粹的自为存在。在形式中，它的这个自为存在成为了真理。”因此，正是在似乎只具有异己意义的劳动中，服务的意识通过重新发现自己、通过自身获得了本己的意义\footnote{同上书，第157页。[这两段引文，贺译本译为“这个否定的中介过程或陶冶的行动同时就是意识的个别性或意识的纯粹自为存在，这种意识现在在劳动中外在化自己，进入到持久的状态。因此那劳动着的意识便达到了以独立存在为自己本身的直观。”以及“奴隶据以陶冶事物的形式由于是客观地被建立起来的，因而对他并不是一个外在的东西而即是它自身；因为这形式正是它的纯粹的自为存在，不过这个自为存在在陶冶事物的过程中才得到了实现。因此正是在劳动里（虽说在劳动里似乎仅仅体现异己者的意向），奴隶通过自己再重新发现自己的过程，才意识到他自己固有的意向。”见《精神现象学》上卷，第130、131页。——译者]}。用一句简短的拉丁语来说就是：factio perfectio facientis est（制造是制造者的完善）\footnote{参见罗森克兰茨：《黑格尔传》，第133页，另见《全集》第1卷，第495页以下；《政治学与法哲学文集》，第467页。}。

\section*{3}

如果我们想要理解黑格尔哲学中的客观精神的概念和安排，就必须要对劳动和行动之间关系的变化予以解释。正因为黑格尔在吸收实践哲学的过程中，不得不修正实践哲学的基础，因此他与传统的争论就使他远远超出了传统以外。这一点可以从1805—1806年的耶拿讲座中得到具体证明。在某些方面，这一讲座放弃了之前的“伦理自然”建构，同时也就舍弃了相应的所有实践哲学的划分尝试。在1801年到1804年之间，精神哲学根本而言与伦理理论一致，还没有进行内部划分。精神首先是一族人民的“绝对”精神，它的其余环节不是对民族精神开端的补充，就是对其终点的扩展。在这一时期，甚至宗教和哲学都属于伦理。依据柏拉图-亚里士多德模式，黑格尔将哲学比作一族人民的“绝对劳动”：正如牺牲的伦理行动的纯粹活动证明自己是为了政治整体的，哲学在对“个别性”的完全否定中将人民精神提升为“自然世界和伦理世界的精神”\footnote{《政治学与法哲学文集》，第465页。}。另一方面，在“伦理精神”和“绝对精神”同一的前提下，也就不可能对个体的（主观）精神进行独立的探讨。因此，在黑格尔看来，个体“以一种永恒的方式”直接存在于伦理整体之中，“其经验的存在和行为是完全普遍的存在和行为；因为个体并不是行动着的个体物，而是寓于他之中的普遍的绝对精神”\footnote{参见《耶拿实在哲学》第1卷，第230页；第2卷，第210页以下。}。

这里，对劳动概念的全新阐述导致了精神哲学体系安排上的深刻转变。此时，1803—1804年讲座所阐发的精神教化的思想，也就是“绝对伦理”从在自我行动的自我中被给予的意识阶段（“理智”和“意志”）“发展”出来的思想，开始发挥出全部效力。由此，也就抹平了自然法论文和《伦理体系》在“自然伦理”或者相对伦理（经济学）与“绝对伦理”（政治学，它对“政治”经济学进行不断的限制）之间划下的鸿沟。根本而言，有两个理由：(1) 借助国民经济学的劳动概念，“个体性”环节在与自由和平等的个体相互承认有关的行动过程中得到了辩护\footnote{《耶拿实在哲学》第2卷，第210页。}；(2) 在此之前，黑格尔将法权概念作为经济-实践领域的单纯“形式”统一性来处理，“绝对物或伦理物”毫不相干地漂浮于这种统一性之上。与此不同，现在黑格尔把法权概念作为“一般伦理”的环节从“承认的运动”中产生出来，同时，通过“承认的运动”，之前分离开来的体系各部分也得以联系起来。尽管如此，黑格尔并没有按照习惯的方式接受自然法的这个核心部分（作为使得从自然状态过渡到“公民状态”得以可能的契约论建构），而是对自然法范畴和政治经济学范畴进行对照和比较。他发现了对他自己和自然法理论家来说仍然隐藏着的东西：法权人格的概念，以及从“个体意志”到“普遍意志”的让渡概念中所包含的个人的自由和平等，是以通过劳动从自然中解放出来为前提的。黑格尔认为卢梭和康德的普遍意志不过是一种抽象物，因为它没有作品，其自身缺乏客观的联系和持久的存在。尽管承认的运动指向的是他人的自我认识而非一物\footnote{同上书，第113页。}，但是得到承认的意志的内容是经过了劳动和占有的中介作用的。只有这样，普遍的意志才是“现实的精神”。在其中，意志和理智如同劳动和行动一样彼此联系起来。最初包含在普遍意志概念之中的东西，并不是像意志的联合和契约这样的社会和历史的抽象构造，而是从人和自然的中介过程中产生出来的那些环节：劳动、产品和占有。这些环节先于所有的意志关系：“精神既不是作为理智也不是作为意志，而是作为本身就是理智的意志，才是现实的……在这种本身就是理智的意志中，抽象的意志必须自我扬弃，或者作为被扬弃了的，在承认（Anerkanntsein）的要素中、在这种精神的现实性中产生出来。就像之前劳动转变为普遍的劳动一样，占有也由此转变为法权；之前作为家庭财产的东西……成为了所有人的普遍作品和享受。”\footnote{《耶拿实在哲学》第2卷，第263页以下。}只有通过关联到一件作品，在这件作品中，普遍意志的概念、理智和意志的统一（“自我”）给予自身以定在，[只有这样，]普遍意志才是“现实的精神”。这种“现实的精神”还不是1817年的《哲学全书》引入的“客观精神”，但也是迈向客观精神的一步。黑格尔通过区分作为定在的教化的外化（这种外化同时就是对象化或客观化）和定在着的世界本身的外化完成了这一步。这一区分为1805—1806年讲座的结论奠定了基础。这个讲座的结尾不再是把伦理提升到宗教，而是把艺术、宗教和科学作为“绝对自由精神”的特殊阶段。绝对自由精神把迄今为止的所有规定都收回到自身，并且产生出“另一个世界”\footnote{参见上书，第212页，附释2。参见哈贝马斯：“劳动与相互作用——黑格尔耶拿精神哲学评注”（Arbeit und Interaktion. Bemerkungen zu Hegels Jenenser Philosophie des Geistes），载《自然与历史：卡尔·洛维特70大寿文集》（Natur und Geschichte: Karl Löwith zum 70. Geburtstag），斯图加特/柏林/科隆/美因茨，1967年，第132页以下、154页注释3。然而哈贝马斯忽略了耶拿精神哲学的现象学倾向。}。这里，一方面，差异似乎已经达到了“绝对精神”，另一方面，“现实的精神”和“主观精神”之间却没有清楚的界分。黑格尔之所以还没有认识到这一术语，与他在耶拿讲座草稿中对精神哲学的现象学倾向有着最为紧密的联系。尽管耶拿讲座从包含了认识（理论自我）和意志（实践自我）的意识概念中发展出了“现实的精神”，但由于现象学方式所要求的对象的相关性，后来所言的“主观精神”和“客观精神”的要素就在这里不断地混杂在一起：想象力、回忆与符号分别同自由意志、工具与技术的狡计一一对应。因此，在理论自我和实践自我之间也就并不存在真正的综合\footnote{《耶拿实在哲学》第2卷，第199页以下。}；实践自我的发展以极不清晰的方式汇入家庭之中\footnote{《哲学入门》（Philosophische Propädeutik）第3篇，第2章，第3部分，第128-129节、第173节以下（《黑格尔全集》历史考订版中对黑格尔纽伦堡时期这一精神学说的划分与里德尔此处引述的荷夫迈斯特的划分不尽相同。前者的划分如下：1. 精神的概念。A. 理论精神。B. 实践精神；第二章“实在精神”；第三章“在其纯粹阐述中的精神”。见G. W. F. Hegel, Gesammelte Werke. Band 10.1. Felix Meiner Verlag Hamburg 2006, S. 342, 353, 359, 362。中译本参见《黑格尔全集》第10卷，张东辉、户晓辉译，商务印书馆2012年版，第273、286、292、295页。——译者）。参见法尔肯贝格（R. Falckenberg）：《黑格尔客观精神的实在性·哲学及其历史论稿》（Die Realität des objektiven Geistes bei Hegel. Abhandlungen zur Philosophie und ihrer Geschichte）第25册，莱比锡，1916年，第62页。即便我们认为罗森克兰茨并没有认识到这些联系，并且人为地破坏了文本，但是这种模糊性依然存在。参见F. 尼科林：《黑格尔的主观精神理论著作》（Hegels Arbeiten zur Theorie des subjektiven Geistes），见《认识与责任》（Erkenntnis und Verantwortung），杜塞尔多夫，1960年，第367页（Th. 利特纪念文集）。}，以便再次开始新的为了承认的斗争。

同样的模糊性重又出现在《纽伦堡哲学入门》中。这里，“精神科学”分为三节：(1) 在其概念中的精神，(2) 实践精神，以及 (3) 在其纯粹阐述中的精神。第一节处理作为（理论的）理智的精神，并以“理性思想”结束（第170-172节）。它并没有实现与“实践精神”的综合，相反，实践精神单独包含了整个伦理学说，由此，伦理学说便以吊诡的方式从属于心理学体系之下。直到海德堡《哲学全书》（1817年）的概念构成才终结了这种追随18世纪90年代对个体灵魂学说的早期划分尝试\footnote{参见罗森茨威格：《黑格尔与国家》（Hegel und der Staat）第2卷，1920年，第83页以下。}而来的摇摆倾向。这部《哲学全书》产生了那种综合，并将“客观精神”定义为“理论精神和实践精神的统一”\footnote{《哲学全书》，第400节。《全集》第6卷，第281页。}。局限于意识之“自我”的现象学限制考察被取消了。“精神”是超越于思维、意志和对象之上的统一，唯有《逻辑学》的进程才能对这种统一进行阐述。由此造成段落安排上的根本差异：1807年《现象学》中的精神：精神的第一个阶段（尚未称作“客观精神”）是把“个体性”环节从自身中排除了出去的古代城邦的实体性伦理，“客观精神”的第一个阶段是“人格”的抽象法。在其中，通过作为普遍和个别之统一\footnote{《逻辑学》，第二部分，拉松出版，莱比锡，1951年，第219页以下。}的逻辑概念自身，它[指抽象法——译者]得到了绝对的辩护。“现象学的名称”与客观精神的意义的一致之处在于，精神在这里也意味着那种“普遍的作品”（而不仅仅是思维与意志的统一、“普遍意志”），“它通过所有人和每一个人的行动，作为他们的统一性和同一性而创造出来”\footnote{《全集》第2卷，第336页。（贺译本译为：“一切个人和每一个人通过他们的行动而创造出来作为他们的同一性和统一性的那种普遍业绩或作品。”见《精神现象学》下卷，贺麟、王玖兴译，商务印书馆1979年版，第2页。——译者）}。但是精神并不是一个自由意志的“客观化”的统一性概念，而是分化为意识的各个不同的“世界”。它在古代城邦中具有实体性，在现代教化王国中有外在的定在，而最终在普遍意志和道德中拥有知识和意志的自我确定性，但没有在一个积极的作品中得到实现。

《哲学全书》对客观精神的处理反其道而行之。它从道德的意志概念开始，这个概念乃是作为自由的起源本身的“普遍”意志，亦即在思维中把握自身的意志。就此而言，这种处理建立在由卢梭、康德和费希特开辟的现代法哲学和国家哲学的基础上。然而，对黑格尔来说，普遍意志既不是客观精神学说的起源，也不是这种学说的目的。之所以不是起源，是因为它是在与自然打交道的过程中获得的；之所以不是目的，是因为它自身必须在一个作品中得到实现。由于黑格尔把自由意志定义为自我规定的且给予其规定以外部实在性的精神\footnote{参见《纽伦堡哲学入门》（Nürnberger Propädeutik），第3篇，第2章，第3部分，第173节。《哲学全书》，第1版，第400节；第3版，第484节。}，所以他把自由意志重新纳入康德和费希特将其从中分离出来的客观定在的领域，由此，这种定在——也就是人的自然和历史（法、国家等）——即表现为由自由意志所“设定”和渗透的。自由只有在这种前提下，即它已经在历史中得到实现并且内在于实定法和国家之中，才能理解为“法的理念”。普遍说来，客观精神的核心范畴，即法，乃是意志在其自由的自我规定中自己给出的定在领域——实现了的自由的王国，精神的世界从它自身中产生出来\footnote{《法哲学原理》第4节。《全集》第7卷，第50页。参见同书第1节，补充，第38页以下。}。

这就是黑格尔的客观精神学说超越其18世纪前辈和传统实践哲学理解之处。自由的、自我客观化的意志的概念在自身中把那些在18世纪前辈和传统实践哲学处依然保持分裂的东西统一了起来：与伦理行动中众多个体意志彼此之间的关系的关联，以及与人类需要本性的关系及其通过劳动与外部自然物之间的中间活动之间的关联\footnote{参见《哲学全书》，第3版，第483节：“自由意志在它身上首先直接具有种种区别。自由是它的内在的规定和目的，并与一个外在的遇上的客观性相联系，这客观性分裂为：具有种种特殊需要的人类学东西，种种为意识而存在的外部自然事物，以及个体意志对个体意志的关系”；另参见《哲学全书》，第2版，第488节到491节。}。行为与活动的这种统一最终表明了精神自身的概念结构\footnote{参见《纽伦堡哲学入门》，第176节。}。尽管作为哲学的原则，精神是绝对的第一者并且构成自然的前提，但黑格尔还是用出自人与自然打交道的领域的那些范畴来解释精神的存在方式。精神不是静止的存在，而是过程、活动和劳动：“但是精神只有扬弃其直接的存在，方才存在。如果精神仅仅存在，那它就不是精神。因为它的存在正在于，作为自为存在着的精神，以自身为中介，自为地存在。”\footnote{《宗教哲学讲演录》（Vorlesung über die Philosophie der Religion）第1卷，G. 拉松出版，莱比锡，1925年，第70页。}在黑格尔这里，“精神的劳动”不是比喻，而是道出了它的本质：成为它自己的产品、成为自我对象化和作品。黑格尔在历史哲学讲座中说，精神“根本而言就是行动，它把自己造就为它自在地所是的东西，造就为它的行迹、它的作品；因此，它成为了自己的对象，使自己作为一个定在出现在自己面前”\footnote{《历史中的理性》（Die Vernunft in der Geschichte），J. 荷夫迈斯特出版，汉堡，1955年，第67页；另参见同书，第55页以下、123、135页。}。

在黑格尔体系内，这种按照劳动和行动相统一的模式对客观精神所做的解释导致了实践哲学的地位发生了根本的改变。在《精神现象学》之后，黑格尔在其体系构建中开始撰写《逻辑学》（1812—1816年），然后出版了作为整体纲要的《哲学全书》（1817年），《法哲学原理》于1821年出版，是他的最后一部著作，而且还是由他本人亲自出版的。这部著作对客观精神学说做出了真正的阐述。这里，黑格尔明确指出了继踵理论哲学与实践哲学的传统体系划分的意图，他在序言中说道：“关于思辨认识的本性，我在我的《逻辑学》中已予详尽阐述。”法哲学不仅建立在逻辑学的辩证方法上，而且还是逻辑学的对应学科，因为它阐述了“实践知识”的本性。黑格尔在柏林时期仍然能够并非毫无理由地按照理论哲学和实践哲学的“习惯划分”模式来解释它们的关系\footnote{参见《纽伦堡著作集》（Nürnberger Schriften），J. 荷夫迈斯特出版，莱比锡，1938年，第XIV页。}。不过，这种划分并不十分契合黑格尔体系的概念理解和结构。对黑格尔而言，逻辑学并不像旧形而上学那样，只是一门哲学的“理论”学科，而是作为“纯粹科学”构成了哲学的第一部分；这一部分要先于体系的其他两个部分，也就是自然哲学和精神哲学，由此，也就取消了逻辑学与自然法的并列\footnote{参见《逻辑学》，序言。《全集》第4卷，第13页。}。因为就事情本身而言，自然法在体系中的位置居于精神哲学当中。在黑格尔这里，“实践哲学”的概念和内容处于“客观精神”的名称之下。为此，体系就必须抛弃旧的“实践哲学”的名称，因为它既不契合于作为现代世界及其社会之基础的劳动，也不契合于精神客观化为其产品的本体论上的必然性。精神的主题是处于法和道德、家庭和社会、国家和历史诸领域之中的人类世界的行为。居于主观精神（人类学、心理学）与绝对精神之间的客观精神学说包含了那些在它产生出来的世界中既有其基础也有其目的之人的行动；此种行为就其意义而言也不是超越的，而是与人类世界及其“客观精神”不可分离的。在黑格尔这里，经院哲学所称的伦理学、家政学和政治学以及与自然法相关的东西，全都落入这个概念之下。由此，一方面设定了同这些科学之间的差别。对黑格尔而言，这些科学，就像旧形而上学、心理学、宇宙学以及自然神学一样，由于哲学的先验转向而“连根拔除了”\footnote{参见《哲学全书》，第1版，海德堡，1817年，第300节（原注似应为第300节，原文标为第299节。——译者）：“对于我们，精神以自然为它的前提，它是自然的真理性。在这种真理性中，自然在精神的概念面前消失了，而精神则作为理念产生了自己，理念的客体同主体一样，都是概念。”（《全集》第6卷，第227页）[中译文参见〔德〕黑格尔：《哲学科学全书纲要》（1817年版），薛华译，北京大学出版社2010年版，第155页。——译者]}。因为在其行动中实现了康德的那个“纯粹意志”的、自由的概念的客观精神，已经克服了自然（人的自然与人的外部世界的自然）的束缚，也就不同于经院哲学传统中的遵循自然规律的自由行为的实现\footnote{《哲学全书》，第3版，第24节，补充2，《全集》第8卷，第87节；《哲学史导论》（Einleitung in die Geschichte der Philosophie），第3版，J. 荷夫迈斯特出版，汉堡，1959年，第110页以下、223页以下；《历史中的理性》，第5版，J. 荷夫迈斯特出版，汉堡，1955年，第54页以下。}。对于黑格尔而言，“客观”精神之为精神，通过实现其“自在”存在着的自由，就不是一成不变的自然规律，而是使自然本身成为一种由精神所设定的、与精神一道得到实现的客体：自由就是在他在中自在地存在（Beisichsein im Anderssein）\footnote{《哲学全书》，第3版，第482节，附释，《全集》第10卷，第380页。（中译文参见〔德〕黑格尔：《哲学科学百科全书III：精神哲学》，杨祖陶译，人民出版社2015年版，第273-274页。——译者）}。对自然束缚的这种摆脱，就为客观精神不仅具有一种历史而是自身就是历史性的，提供了更深的理由。因而，就自由的理念来说，黑格尔能够说，它似乎还从未占领过（或者说尚未占领）非洲和东方的整片土地。这种说法同样适用于历史哲学：“[……] 希腊人和罗马人，柏拉图和亚里士多德，甚至斯多亚派都没有自由的理念；相反，他们只知道，人由于出生（作为雅典、斯巴达或者其他城邦的公民）或由于性格的刚强、教养、哲学……而是现实地自由的。自由的理念是通过基督教来到世上的。”

随着洞见自由理念的历史性（这一洞见对其客观精神学说产生了极深的影响），黑格尔就把传统实践哲学的整座大厦撂到了一边。对于现代思想而言，与青年黑格尔由以出发的那种对传统实践哲学的解构相比，这种洞见的影响要深远得多。19世纪都是从历史哲学出发理解黑格尔的客观精神学说的，由此也就必定忽视了客观精神学说发展起来的来源和背景。今天，哲学又要重新理解这些来源。这同时也就意味着，哲学可以从那种黑格尔彻底摧毁了的传统中获得教益。但是，所有的教益——无论人们怎样理解它，以及人们可以同这些教益走多远——都无法回避与黑格尔对传统的批判性解构之间的冲突。黑格尔的论证不仅尖锐，在许多方面令人信服，而且恰恰集中在经院实践哲学过于明显地与旧欧洲社会统治关系关联在一起的那些方面。问题在于，在当前社会历史世界条件下，实践经验的理论是如何可能的。没有人会说黑格尔给出了充分的回答；然而他的功绩在于，他为他的时代，同时也是我们的时代，提出了这个问题。
\chapter{制度中的辩证法}

黑格尔意义上的法哲学既不是一种实定法学理论，也不是阐述一系列个体自由权利的自然法的原则学说。一般说来，黑格尔法哲学的对象并不是任何历史的或者超历史的权利(Rechten)体系，而就是（单数的）法(Recht)——更准确地说：法的“概念”及其“定在”或者“现实化”（第1节）。黑格尔谈到了“法的理念”，因为正如这部著作开篇所言，哲学探究的是理念，而不是人们通常所言的“单纯概念”。

\begin{quote}
相反地，哲学应该指出单纯概念的片面性和非真理性，同时指出，只有概念才具有现实性，并从而使自己现实化。\footnote{《法哲学原理》，第1节，附释。《全集》第7卷，H. 格洛克纳出版（1927年之后），第38页。（中译文引自〔德〕黑格尔：《法哲学原理》，范扬、张企泰译，商务印书馆1961年版，第17页。——译者）}
\end{quote}

这种言说方式以对一些“逻辑学”的基本术语规定的理解为前提。按照黑格尔的理论，“定在”意味着“规定了的存在”，在其“形态”和“实存”的规定性中，定在并不等同于日常语言所言的“现实”。黑格尔把前述“形态”和“实存”称作过渡为他者的直接事物的诸形式。历史上的实定法，比如说一个地区或国家的法律，是现存的和有效的。也就是说，它们“在历史上实际存在”。但是，《逻辑学》引入的“实存”这一术语乃是“存在和反思的直接统一”，因而界定的并不是现实的结构，而是现象的结构。因此，尽管法也必须显现出来，尽管“伦理事物的现象世界”在法哲学的概念序列中占据了突出的地位，然而这里现象的语言在方法上的作用是次要的。每一种现象都是直接的，正如同黑格尔本人所言，它来自根据，又回到根据\footnote{“回到根据”德文原文为“geht zu Grunde”，意为“死亡、毁灭”。——译者}。只有那种不再“过渡到他者”的根据才对逻辑的东西和历史的东西“进行中介”。在黑格尔的术语中，被中介者的形式所特有的这个根据，就叫“现实”：

\begin{quote}
现实事物是上述那种直接统一的设定存在，是达到了自身同一的关系；因此，它得免于过渡……它的定在只是它自己本身的表现，而非他物的表现。\footnote{《哲学科学百科全书》，第一部分，逻辑学，第142节，附释，《全集》第8卷，第320页。（中译文引自〔德〕黑格尔：《小逻辑》，贺麟译，商务印书馆1980年版，第295页。按德文原注标明此句引文出自“说明”，疑误。——译者）}
\end{quote}

法哲学的对象是法的现实，依照黑格尔的观点，法的现实在“自由的理念”之中，后者建立在对立规定的统一之中。历史上设定的法所“表现”的，是那种自己赋予自己“现实性”的“概念”，它仿照自然法，在“其为自由的意志”（第4节）的自由的规范原则下构成了“哲学法学”的出发点。哲学法学在方法上的创新并不在于从规范性的法原则推演出超历史的法规范体系，而是把通达理念的道路理解为各种制度的历史形成过程（这是一个辩证的——矛盾的历程），并且把这条曲折的道路理解为认识逻辑概念自身辩证法的必要条件。

\section*{1}

在法哲学中，“制度”这个词并不是一个规范的术语，而是从语言的实际使用拿来并在不同的语境中加以使用的。我们可以对其作狭义(a)和广义(b)之分。狭义的制度(a)借助罗马法的规定，而黑格尔只把罗马法解释为私法（就历史而言，这是错误的）。按照这个词的原初法学语境，“制度”的例子有作为罗马私法概念之基础的“罗马家父权、罗马婚姻状况”\footnote{《法哲学原理》，第3节，附释，《全集》第7卷，第43页。（中译文参见《法哲学原理》，第5页。——译者）}。这些例子的语境为黑格尔所言的（以财产和契约为基本制度的）“抽象法”领域，在《法哲学原理》概念构建的体系进展中，它超出了这一语境。对于黑格尔而言，只有（在“道德”一章中）批判了这些制度之法的历史局限性并在伦理的概念和定在中扬弃了这些制度之后，自由的理念才是“现实的”。伦理作为“活的善，在自我意识中具有它的知识和意志，通过自我意识的行动而达到它的现实性；另一方面自我意识在伦理性的存在中具有它的绝对基础和起推动作用的目的”（第142节）。[由此，伦理]也需要历史的形态。伦理的语言不外就是“由于自由的理念而是必然的、因此是现实的那些关系在它们全部范围内即在国家中的发展”\footnote{同上书，第148节，《全集》第7卷，第229页以下。（中译文引自《法哲学原理》，第167页。——译者）}。

在伦理及其与“国家法”——公法——的关系的语境下，制度概念获得了其更加广泛的意义(b)。黑格尔在这里当然不再继续使用拉丁化的表达，而是代之以德语同义词“Einrichtung”。伦理事物的“固定的内容”，那些摆脱了个体的主观偏好和意见的“伦理力量”（第145节），乃是“自在自为地存在着的法律和制度”\footnote{《法哲学原理》，第144节，《全集》第7卷，第227页。（中译文参见《法哲学原理》，第164页。——译者）}。

如前述，逻辑的东西和历史的东西在法中的中介导向“逻辑学”的现实概念，而现实概念也为自然法、实定法以及国家学之间的关系且由此而为现在即将进一步研究的制度的辩证法之起源奠定了基础。“现实”一词每个人都很熟悉，它的每一用法都是开放的，遗憾的是，对于这个具有如此多意义的词语必然引发的误解，黑格尔并没有进行防备。尽管《法哲学原理》“序言”在其关于现实的东西的合理性和合理性的东西的现实性的著名的或者说声名狼藉的格言中已经涉及“现实”一词的逻辑规范意义，但由于它又同时影射了一种前哲学的——通俗的语言运用，引起了对他的双重误解：一方面，他为实定的东西特别是历史中现成的（“普鲁士”）国家进行明确辩护，另一方面，他明确地拒斥自然法原则。确实，黑格尔对“哲学法”的构想与其近代自然法前辈的设想没有多少共同之处，但他与那种他的后辈和批判者都想要读出来的复辟时期的“国家哲学”计划则更是格格不入。

根据《逻辑学》的概念定义，“现实”意味着“本质与实存或内与外所直接形成的统一”\footnote{《哲学全书》，第142节，《全集》第8卷，第320页。（中译文引自《小逻辑》，第295页。按德文原注标明此句引文出自“说明”，疑误。——译者）}。《法哲学原理》中有两个环节与这种把现实的东西的结构规定为中介的用法相符。这两个环节属于法概念的辩证法，同时也属于它的定在、伦理制度中的现实化。我们采用稍微不同于黑格尔的术语，称之为法的反思性和实定性环节。“反思性”这一术语在这里表示的是与法的“基地”的关联，这种关联属于法的现实性。对法来说，这种关联是一个“精神物”——也就是法主体——的构成性的自身相关，这个主体同时知道自己是那些承认和遵守规范和法律的义务的承担者。黑格尔的“自由意志”概念是那种反思性的基本表达。按照作为“外化”与自身反思之统一的一切现实的辩证法，自由意志在它的他者中意求它自身，因此也就是意求“自由意志”（第27节）。“自由意志”，即人之为人的普遍的法权能力（Rechtsfähigkeit），就是近代自然法意义上的“法”的基本概念。但是，正如辩证的重构所表明的那样，此种意志的内容是制度的各种历史形态。《法哲学原理》并不认为有从一种规范的普遍的法权原则中演绎出法的那些特定内容的可能性。黑格尔把这种思想斥为一种抽象的自然法假设：“可以有一个法体系和一种应当是纯粹理性的——仅仅是理性的——法状态。”\footnote{《法哲学讲演录》第2卷，伊尔廷（K.-H. Ilting）出版，斯图加特-巴德·坎施塔特，1974年，第89页。}在制度的特定发展阶段——诸如罗马法或中世纪封建法——上的法的内容是偶然的历史形态的产物。尽管对这些制度不仅要从历史的角度进行描述和解释，而且也要按照人之为“自由人”的法权能力原则进行评价（第3节），《法哲学原理》还是非常一贯地禁止从自然法原则出发进行演绎。

理解“哲学的法”之辩证法的关键在于“实定性”概念。众所周知，“实定性”这一术语在黑格尔青年时期作品中居于中心地位，而在这里，它表达的是法概念与其在制度中的定在、与作为道德规范和法规范之实定化的法律关系。黑格尔认为，这种关系对“一般的法”来说是建构性的。在一定程度上，可以通过分析从法的概念得出，法要实定化，进入定在和实存的外在性之中，并且作为“合法的东西”——法律——而具有效力。哲学的法和实定法彼此是有区别的，但是黑格尔认为自然法和实定法的截然对立是一个“巨大的误解”。他再次以罗马法为例\footnote{《法哲学原理》，第3节，附释，《全集》第7卷，第42页。黑格尔的比较是有问题的，因为不同于黑格尔法哲学教科书的“原理”，制度的原理部分并没有在严格的体系中建立起来。参见F. 舒尔茨（F. Schulz）：《罗马法学史》（Geschichte der Römischen Rechtswissenschaft），魏玛，1961年，第187、194页等处。舒尔茨中肯地对这些制度与亚里士多德的形而上学的松散构造进行了比较。}，拿自然法与实定法之间的相互联系跟《法学阶梯》与《学说汇纂》之间的关系进行比较：《法学阶梯》与《学说汇纂》共属于作为一个实定法体系的《民法大全》。法既由于其在某个国家中有效的形式，也由于其内容的要素，而必然是实定的。“内容的要素”包含了民族的特殊性、一个国家的历史发展阶段以及每一法律体系的组织和适用，这些要素绝不能从任何规范性的法原则中产生出来（第3节）。法必然要走向实定性，由此便产生了对于黑格尔的《法哲学原理》来说乃是奠基性的“自然法”与“国家学”之间的关系，而需要作为纲要来理解的《法哲学原理》一书的副标题也坚持了这种关系。从历史角度来看，这一标题描述的是哲学的两个不同学科。第一个学科（自然法）首先是在近代得到了发展，而第二个学科\footnote{即国家学。——译者}的起源则要追溯到古代世界、希腊的城邦理论。古典政治学体系的优越之处在于，它并不把自然法和国家学分开。柏拉图的《理想国》和《法篇》或者亚里士多德的《政治学》阐述的政治的理性理念，即最佳宪制，把自己理解为对自然法和一种正当地组织起来的（“公民的”）社会（$\pi o \lambda \iota \tau \iota \kappa \acute{\eta}\ \kappa \circ \iota \nu \omega \nu \acute{\iota} \alpha$, societas civilis）的陈述，这种公民社会与国家（$\pi \acute{o} \lambda \iota \varsigma$, civitas, res publica）是同一的。青年黑格尔追随这种从未间断的“国家学”的传统形态，首先在其中看出，对“自然法”进行科学探讨的任务在于，自然法应当在一个政治整体、在“国家”这一基本制度以内建构“伦理自然的法”\footnote{参见《论自然法的科学探讨方式，它在实践科学中的地位及其与实定法学的关系》（1801年），《全集》第1卷，第435页以下。以及本人与之相关的论文：《黑格尔对自然法的批判》，收入《黑格尔法哲学研究》（Studien zu Hegels Rechtsphilosophie），美茵河畔法兰克福，1969年，第42页以下。}。

然而不久之后，黑格尔就不得不认识到，在近代国家的特定历史形态的条件下，应用一种把个体行动埋葬于制度中，由此消除了矛盾的可能性的非辩证的政治概念，只会导致马基雅维利的立场上的政治绝对主义的实践。在与政治绝对主义的分歧中，从霍布斯到康德的自然法形成了与之相对的自然法的个体主义立场。这种自然法的权利在于，要么坚决限制国家对自身秩序井然的社会（“市民社会”）的“非法”干预，要么以一种新的自由概念和法概念革命性地克服已经固化为赤裸裸的权力展示的国家政治。这样，自然法和国家学的对立便引发了向政治启蒙的转向以及向革命的过渡，并且一直伴随着革命的进程。黑格尔的概念辩证法预设了自然法和国家学的对立，并且企图在哲学上和政治上对之予以克服。法哲学把自己理解为“法的哲学”，因为它试图在法理念的辩证法中对这种前国家的自然法与实定的国家法之间的对立进行调和。黑格尔说，法的理念就是自由，要真正把握法的理念，就必须在其概念和概念的定在中（《法哲学》第1节，补充）。也就是说，在各种制度的历史形成过程以及相应的法的现实性的过程中去认识它。

由于法概念与各种制度的定在之间的这种历史的辩证关系，自然法的语言当然就是成问题的了。黑格尔在1817年的海德堡《哲学科学百科全书》中就已经谈道：“迄今为止，自然法这一表达，对哲学法学来说已经成为了一种习惯表达。这一表达具有歧义：法是不是作为一种似乎由直接的自然植入的东西，还是指的是它是由事物的本性，也就是说，由概念来规定自身的。第一种意义是之前习惯所指的意义，以至于同时还虚构出一种自然法在其中应当有效的自然状态。”对黑格尔来说，就法并不是把它自身以及它的所有规定建立在自然之上，而是建立在“自由的人格性——一种为自然规定的对立面的……自我决定”（第415节）之上而言，事物的“概念”就是法。法的现实不是自然；法是“精神的世界，是从精神自身产生出来的”（《法哲学原理》第4节），一种“第二自然”，黑格尔称之为“实现了的自由王国”。在《法哲学原理》概要叙述的那些从最初抽象法的形态开始，一直到家庭、市民社会和国家的各个历史形态为止的制度中，法把自己建立起来。显然，这里处处都有“自然”关系发挥作用。问题在于，这些自然关系是否从自身中产生出了支配那些形态的法律。黑格尔坚决地拒绝这一问题以及对这一问题给予肯定性答复的自然法话语。必然性（比如说，在国家与个体的关系中，归之于“国家”这一制度的必然性）不再意味着，对个体来说，必须生活在国家中，是一条自然法则；相反，国家的伦理必然性以自由的法则为基础。这种自由的法则不是由一成不变的“自然”制定的，而是由——在其辩证法中依赖于各种制度的历史实现的一一概念自身制定的。对黑格尔而言，概念的自我实现乃是人的普遍法权能力的自然法原则，这里所言的人都起来反抗现存的支配和依附关系，并且冲破了迄今为止法的思想、自由都还是隐匿于其下的自然性的假象。

\section*{2}

因此，对黑格尔而言，“哲学法”之有别于国家的实定法，有如它之区别于自然法。在同时代的两种对立的法观念，即康德、费希特的理性法观念与萨维尼的“有机”法观念之间，黑格尔的法哲学处于中间的位置。黑格尔与萨维尼和历史法学派一道拒绝自然法的理性国家。但是，在黑格尔这里，历史解释不仅具有对实定法予以历史论证和辩护的意义，而且具有“概念的发展”的意义（[《法哲学》]第3节）。辩证的概念发展本身就是历史的，即自由在世界历史中的阶梯式进展，而《法哲学原理》的任务就在于，为不同时期的历史——实定法提供“自在自为地有效的辩护”（出处同上）。

这对于理解这部著作及其历史立场有着决定性的意义。黑格尔立足于现代革命及其制宪和立法的基地之上，现代革命的制宪与立法将自然法实定化了。对黑格尔而言，国家的合理性不再是一个公设，而是正在发生的历史现实，“哲学法”承认这一现实。黑格尔提出将“探究理性的东西”与“了解现在的东西和现实的东西”结合起来的哲学概念，反对在他的时代受到偏爱的“关于国家的哲学”，也反对这种国家哲学炮制出新的和特殊的理论的目标，

\begin{quote}
“似乎世界上从未有过国家和国家制度，现在也还没有，但是现在——这个现在是永远继续下去的——似乎应该从头开始”\footnote{《法哲学原理》，序言，第23页。（中译文引自《法哲学原理》，第4页。——译者）}。
\end{quote}

黑格尔承担了哲学的传统任务“理解存在的东西”，也就是对存在的思辨认识，以对哲学提出一种似乎将其颠倒了过来的解释：在时间发生的视域中思考存在。因为哲学作为对存在的东西\footnote{指哲学。——译者}的认识，同时又是“在思想中得到把握的它自己的时代”；哲学无法跳出它的世界，就像个人作为“时代之子”无法跳出罗陀斯岛一样。

黑格尔将“法的理念”规定为概念和定在的统一，就是基于时间和存在的这种交织。在解决将自由认识为这种理念这个问题的过程中，《法哲学原理》避免了自然法和法实证主义的双重片面性，它们要么仅仅执着于“自然”法的抽象概念，要么执着于实定法的定在。理念“为了得到真正的理解，必须在法的概念及其定在中来认识法”（第1节，“补充”）。其中包含了理念与历史的联系。自由只有在这样的前提下，即它已经在历史中进入定在并且内在于实定法之中了，才能作为法的理念得到把握。这发生在现代世界的国家中。现代世界的国家在自身中接受了法国大革命的自由原则，并且具有保存个体为自由的、具有法权能力的人格的任务。黑格尔把现代国家定义为“逻辑学”思辨意义上的“现实”，成为直接之物的内在本质与外在实存的统一：现代国家是“伦理理念的现实”（第257节）、“实体性意志的现实”（第258节）以及“具体自由的现实”（第260节）。这样，不同于从卢梭到费希特为止的观念论自然法，对黑格尔而言，法理念并不是思维的单纯可能性，而是一种经过历史的中介的现实，这种现实要求哲学放弃其抽象的自然法立场，法的理念要求哲学把当前的国家作为“其自身是一种理性的东西来理解和叙述”\footnote{《法哲学原理》，序言，第34页。（中译文引自《法哲学原理》，第12页。——译者）}。

黑格尔在《法哲学原理》序言中的格言：“凡是合乎理性的东西都是现实的；凡是现实的东西都是合乎理性的”，所指的便是这种现代世界的国家，而不是1821年的普鲁士国家。要认识这句格言的“进步”意图，我们不需要用第一句来反对第二句：因为这两句说的是一样的：在现代国家中，法的概念已经达到了定在，理性的东西——自由的思想——是现实的了，并且现实的东西——现代世界的国家——是理性的了。这并没有排除部分的不理性。黑格尔在最后一次法哲学讲演（去世前几天）对这句格言的一个解释中，简洁明了地指出：“现实的东西就是理性的。但是并不是所有实存的东西都是现实的，坏的东西是一个自身败坏的和虚无的东西。”\footnote{《法哲学讲演录》第2卷，伊尔廷出版，斯图加特，1974年，第923页。}

因而，黑格尔在分析法国大革命的自由原则时，提出了哲学的实现问题。那在法国大革命中起而反对已经成为无精神、无理性的封建制度的现实——“不法的旧制度”——并将其推翻的法的概念起源于哲学思想。而这也是这一事件无与伦比和史无前例之处。对于黑格尔而言，这一事件具有世界历史的意义：

\begin{quote}
“自从日立中天、众星绕行以来，还从未见过这样的事情，人立于头脑——也就是说思想——之上，并据此建造现实性。阿那克萨哥拉最先说出，努斯统治着世界；而只是到了现在，人才进而认识到，思想应当统治精神的现实性。”\footnote{《世界历史哲学讲演录》第4卷，G. 拉松出版，莱比锡，1920年，第926页。}
\end{quote}

但是哲学的实现并不意味着哲学的扬弃，并不是思想消解于一个历史上改变了的现实性之中，而是相反：要理解这种改变的根据，就需要哲学。因而哲学就成为了时代的理论，这是在直到黑格尔为止，哲学在其历史中从未有过之事；哲学取得了这样的任务——为了拯救法国大革命的自由原则：(1)反对其自身落入暴政之中；(2)反对复辟。复辟与当下的原则直接对立，因而哲学否定了复辟想要重新建立的历史连续性\footnote{参见里特尔（J. Ritter）：《黑格尔与法国大革命》（Hegel und die Französische Revolution），美茵河畔法兰克福，1965年，第24页以下。}。在这个意义上，黑格尔哲学是革命的哲学，而不是复辟的哲学或者普鲁士国家的哲学（海姆）。黑格尔哲学不是以改变的意志，而是以从法的现实性的事实中获得其原则的认识意志，面对时代及其“现实”。黑格尔在暗中提到培根的一句话时不无道理地说，半吊子哲学（尤其是哈勒、施莱格尔、德·迈斯特、博纳尔德的浪漫主义复辟理论属于其列）离开了国家，真正的哲学则导向国家\footnote{参见《法哲学原理》，序言，第36页。（中译文参见《法哲学原理》，第13页。——译者）}。这种“真正的哲学”的缺陷或许在于，它只是想要认识，并且通过认识，不是想要改变现实，而是想要“和解”。这种哲学的优越之处首先在于，它在概念——内在于完成了的历史之中的概念——的媒介中，对一种改变了的现实做出解释。

\section*{3}

因此，《法哲学原理》是那种青年黑格尔在《精神现象学》中提升为他的哲学规划的“被概念把握的历史”的一个教育剧本。对黑格尔而言，概念就是在其自由中的人的意志，自由构成了意志的“实体和规定”（第4节）。法在概念的自我展开中的辩证论证既涉及古典政治学和现代自然法的诸范畴，也涉及历史—精神的“现实”的形成过程，通过这一过程，这些范畴获得了新的内容。这一点从（“法”的诸制度的）体系纲要以及体系的三分法（一、抽象法；二、道德；三、伦理）可以看出来。抽象法是18世纪的自然法(jus naturale absolutum et hypotheticum)，探讨的主要是私法上的内容（财产、契约）；它是从一切支配和同业公会的束缚中解放出来了的抽象个体的法。黑格尔按照自然法的模式，把它放在《法哲学原理》的最前面，并且显然修正了他最初的立场，在不依赖于“实体性伦理”（家庭和国家）的前提下，对它进行阐述。对黑格尔而言，不是“具有一定身份而被考察的人”，不是作为一个不同于奴隶、农奴或“亲属”的市民，而是具有法权能力的“人格本身”（第40节）才是法和财产概念的基础。道德包含了经院哲学在伦理学的名义下进行探讨的那些概念和内容：古代意义上的德行论和义务论，但是黑格尔却将经由康德的主体性转向所形成的反思道德的原理纳入其中。

与经院传统相一致，黑格尔把抽象法（自然法）理解为“法的意志”的学说，把道德（伦理学）理解为“道德意志”的学说。在这两种学说中，随着现代革命——既是哲学革命也是政治革命——而立足于自身之上的人之为人，一方面在与外部的关系中，在与外部自然世界和周围人类世界的关系中，以财产（第41-71节）、契约（第72-80节）和不法（第84-86节）的形式，另一方面则是在与内在性的关系中，以故意和责任（第105-118节）、意图和福利（第119-128节）、善和良心（第129-141节）的方式，实现了自身，并将自己表达出来。尽管黑格尔承认这些范畴，并且把这些范畴吸收到了他的演绎中，但是这些范畴对黑格尔而言仍是抽象物，在《法哲学原理》体系中只具有有限的效力。被标明为“抽象法”的凝结为私法的自然法，乃是具有法权能力和财产能力的一般人格的法，借由康德-费希特哲学而成为主观性和反思性的伦理学也同样是与自身相关的主体性的道德。

如果说抽象法和道德与近代自然法和道德学(Moral)的范畴之间有一种显而易见的平行关系，那么如果把这种比较运用到《法哲学原理》的第三部分，就会产生不少困难。这些困难不仅仅在于家庭、市民社会和国家的划分，而且也在于伦理这个名称本身。因为在17和18世纪，自然法也划分为个体法和社会法，前者以财产和契约为内容，后者以家（家庭）和市民社会(societas civilis)或者国家(civitas)为内容。黑格尔的伦理理论仍然意味着把一个可以追溯到古典政治学的范畴重新引入《法哲学原理》之中。对于现代自然法而言，伦常(Sitte)并不具有一种独立形式的实践生活的效力，而是一般地从属于法概念的合理性之下\footnote{参见博比奥（N. Bobbio）：《黑格尔与自然法学说》（Hegel und die Naturrechtslehre），1966年，重刊于M. 里德尔（出版）：《黑格尔法哲学资料集》（Materialien zu Hegels Rechtsphilosophie）第2卷，美茵河畔法兰克福，1975年，第81页以下。}。这也适用于康德的法权学说。众所周知，法权学说是康德所称的《道德形而上学》的一部分。与之相反，黑格尔把伦理范畴与伦常的传统立场重新结合起来，在伦理范畴中，法国18世纪作家（孟德斯鸠）对“mœurs”（风俗）的发现，对柏拉图-亚里士多德城邦理论的重构以及早期浪漫派的“民族精神”的“有机”统一的理论相互交织在一起。黑格尔的“伦理”指的是个体与“伦理力量”（第145节）以及每一特定的、历史的民族和国家的“必然关系”（第148节）的统一。因此，伦理也就代表了道德学(Moral)与政治学(Politik)的关系，对于国家与“市民社会”统一(civitas sive societas civilis)的传统理论来说，这种关系是根本性的。黑格尔在这里提到“客观的”而非“道德的”（亦即“包括在道德主观性的空洞原则中的”）“伦理的义务论”与“第三篇中对伦理必然性领域”、国家内部的伦理关系和制度的“系统阐述”之间的同一性，并不是没有道理的（第148节）。

然而，制度论与伦常传统的结合仅仅触及黑格尔伦理理论的一个方面。另一个方面是，这个在形式和内容上超出近代自然法的法哲学范畴记录了从18世纪转向19世纪时在政治社会世界建构中已经完成的根本变化。毫无疑问，这里，对于各种制度的体系性重构遵循了《逻辑学》的概念辩证法。在《法哲学原理》中，黑格尔处处以《逻辑学》的概念辩证法为前提。尽管“概念”中包含的个体性、特殊性和整体性的三个环节已经为1817年的《哲学全书》中伦理的划分奠定了基础，但此时“民族”这一基本制度还是一种静态的统一，尚未按照特殊和普遍的环节进行区分，而是作为整体代表个别性范畴\footnote{《哲学全书》，[第三部精神哲学，]二、客观精神，3. 伦理：(1) 单个民族；(2) 外部国家法[在和平与战争状态中相互对立的（众多）“特殊”国家意义上的特殊性]；(3) 普遍的世界历史（即从特殊的“民族精神”的辩证法中产生出来的“普遍”精神意义上的普遍性）。参见《哲学全书》，第430-452节。}。黑格尔在《法哲学原理》中放弃了这种立场，而让“个别民族”的实体性伦理再度服从概念的辩证法：它分化为家庭（个别性）、市民社会（特殊性）和国家（普遍性）。

《法哲学原理》第三部分的这种划分涉及一个悠久的、令人崇敬的传统。黑格尔接受了这一传统，以便借助辩证法的手段消解这一传统，并在他的思想中克服之。黑格尔辩证地消解传统范畴的最重要理由之一就是要引入一个全新的社会概念，他由此开创了国家哲学历史的新纪元。直到今天，黑格尔借由这一概念让时代及其当下所意识到的，绝不亚于现代革命的历史成就：一个去政治化的、建立在个体的自由和平等之上的“市民社会”的产生，随着从英国开始的工业革命，市民社会的重心也从政治组织形式转移到了经济上面。自由主义自然法的社会思想并没有或者说仅仅是不充分地关注到国民经济学的语言，而在黑格尔这里，国民经济学则构成了一种从根本上发生了改变的概念演绎的参照系。按照自然法的观点，一切“社会”都是人格性地组织起来的，都是由许多个人构成的团体，他们通过理性的言论和由理性规定的行动形成一个对所有人都有法的约束力的共同意志，也就是一个“法权人格”的意志。这里，言和行都具有公共的意义，并在将每一个“社会”（而不仅仅是国家）建构为法律制度(Rechtsfigur)的那种公共领域的媒介中实现出来。相反，对黑格尔来说，“社会”按照定义是由通过需要和劳动彼此联系在一起的私人构成的。劳动是一种特殊的行动方式，需要是作为“私人”之人的自然基础。黑格尔在这一领域的哲学成就首先在于，将其“私人的”参照系理解为经过了公共中介的，并且将其自然基础理解为社会的常量。“市民社会”理论的定位系统并不是契约，即理性的、言行出众的个人之间的联合，而是“需要的体系”，也就是从需要、劳动以及满足需要的手段中产生出来的且在其活动中不断再生的“私人”（第187节）之间的关系网络。

黑格尔“辩证地”理解的私人目的与公共目的之间的交互关系，作为一种超越于“个人”行动之上的“社会联系”的基础，已经在英格兰-苏格兰道德哲学（休谟、斯密、弗格森）和法国百科全书派的利益理论中反映出来。在这里，通过将私人目的、对个人生命和财产的保护宣布为公共事务，契约模式就在其形式法功能之外获得了一种内容。功利和利益理论支配了欧洲晚期启蒙运动的思想，青年黑格尔研读了欧洲晚期启蒙运动的那些文本。它们宣扬的是市民私人彼此之间不受等级和地产关系的限制而进行交往的原则，这种交往通过“个人利益”（狄德罗、爱尔维修、霍尔巴赫）或者“自利”（富兰克林、边沁）规范了所有的社会关系。然而，它们和谐的社会图景现在被解释为冲突模式。黑格尔对那个不再与其他概念（国家、市民社会）同义的行为体系的逻辑-历史问题的态度是模糊的。一方面，“社会”是一个法的构造物(Gebilde)，是理性意志（“自由意志”）的定在的法；另一方面，意志只在私法领域才将自己概括为契约——社会，作为在家庭、市民社会和国家中实现其自身的自然-历史制度，并不是建立在契约的法律义务上的\footnote{《法哲学原理》，第181节，《全集》第7卷，第261页以下。黑格尔认为契约模式不足以从概念上把握家庭（第75节，第163节）和国家（第258节），而予以拒斥。}。

尽管“市民社会”也表现为法的构造物，但对黑格尔来说，在从理性意志（“自由意志”）的理念出发对其进行辩证的论证和辩护时，真正重要的是要表明这个制度落在了传统自然法的“社会”的起源以外，因为一方面它在体系上和历史上远远超出“家庭”这一制度之外，另一方面又没有达到“国家”这一制度。在[市民社会]这一阶段，被黑格尔普遍化为“法的理念”并整合到（客观）精神的辩证法之中的理性意志，将自己“特殊化”为许多异质的个人，“因为这些人本身在他们的意识里所有的和成为其目的的，不是绝对的统一，而是他们自己的特殊性和他们的自为存在，——这就是原子论的体系”\footnote{《哲学全书》，第三部分，第523节，《全集》第8卷，第401页。（中译文参见《哲学科学百科全书III：精神哲学》，第291页。——译者）}。

在这种制度的辩证法中构成的东西，是一种区别于国家、在需要和劳动的基础上成长起来的并且在其活动中不断再生的“私人”之间的关系网络——现代词义上的“社会”。它作为“市民社会”进入法哲学及其政治理论的中心，在这种地位中熔化了古典政治学和现代自然法的那些传统范畴。这首先表现在，社会并没有局限于划归给它的地方，而是扩展到了其他制度领域。因为，在黑格尔那里，市民社会一方面是抽象法的“定在”（第209节以下），另一方面又是主观道德的活动领域（对贫困领域来说）（第242节）、家庭的实体性基地和力量（第238节），最后还是国家的差异（第242节）。在伦理的建构中，市民社会以在历史上得到规定的方式成为了“中心”。在这里，市民社会表现为介于“家庭”与“国家”这两种制度之间的连接环节，这在古典政治学和现代自然法传统中是付之阙如的。市民社会这一“中心”不仅是在平衡与和解各种对立意义上的中介，而且更是形成了一个紧张的领域，这个领域把所有的制度都纳入其影响范围，并且改变了它们的传统结构。

这始于“家庭”这一制度，即“市民社会”的基本元素，按照自然法和国家学的传统理解，家庭乃是其自身构成为政治国家的中介。黑格尔之所以剥夺家庭的这一中介功能，不仅仅是因为市民社会和国家分离开来，而且也因为家庭自身按照其“概念”不能是国家团体的一个环节。家庭与国家的关系仅限于家庭是国家的“材料”（第262节）。也就是说，既在生命的自然繁衍过程中准备好作为“群众”的个人，也在教育中准备好直接的伦理情感和意识（第173-175节）。总而言之，在黑格尔之前，“家”之制度作为“整个家庭”，本身就是国家整体的一分子，现在，这个分子[不再是家庭，而是]在家庭的私人领域中成长起来并受到教育的个体。与此相应的是制度的构建过程中的一种与历史相平行的辩证法。一方面，与国家的关系要从家庭过渡到市民社会——它是黑格尔新颖地嵌入伦理体系之中的制度领域；另一方面，市民社会和家庭也必须将它们自己“设定”到一种制度上改变了的关系中。对黑格尔而言，家庭并不是一个由具有不同身份的成员所构成的“社会”，亦非由“更小的社会”（丈夫-妻子、父母-子女，主人-奴隶）组成，而是表现了一种“人格”，它在“一种所有权”中，而不是在家庭的法-经济上的统一中，具有其外在实在性（第169节）。在黑格尔这里，家庭概念的原初的“家政”特征已经消失殆尽，被一种现代的情感特征取而代之。这种情感特征完全建立在“私人”关系和情感的基础上，在18世纪末开始形成。因此，第158节这样定义家庭的概念：“作为精神的直接实体性的家庭，以爱为其规定，而爱是精神对自身统一的感觉。因此，在家庭中，人们的情绪就是意识到自己是在这种统一中，即在自在自为地存在的实质中的个体性，从而使自己在其中不是一个独立的人，而成为一个成员。”\footnote{中译文引自《法哲学原理》，第175页。——译者}黑格尔表述这一概念的时候完全认识到了《法哲学原理》体系中的市民社会的基本地位所引起的传统家政学形态的改变。因为，与迄今为止的政治学和自然法中的市民社会(societas civilis)理论——按照这些理论，市民社会以“家庭”社会为基础——相比，家庭在黑格尔所讨论的“市民社会”中是一个从属性的东西，仅仅是奠定基础而已：家庭不再具有如此“全面的影响力”。在超出了“整个家庭”界限的现代社会的经济再生产过程的条件下，个体已经从家庭中被“揪了出来”，并且如黑格尔所言，成为“市民社会的子女”；市民社会接过了经济的功能，“又用它自己的土地来代替外部无机自然界和个人赖以生活的家长土地，甚至使整个家庭的存在都依从它，而听偶然性的支配”（第238节）\footnote{中译文引自《法哲学原理》，第241页。——译者}。

正如黑格尔的市民社会概念一方面涉及改变了的家庭结构，另一方面它又涉及同样变化了的国家的地位。当我们把黑格尔伦理理论的第三章与传统政治学理论进行比较的时候，就可以确定，由于其种种前提，它不再是对“市民社会”及其政治制度（政府）的陈述，而是与其形成鲜明对照的“国家科学”。在黑格尔这里，传统通过市民社会(societas civilis)的关系概念融为一体的国家和市民社会第一次“设定”为一种关系，而且设定为分离、“差别”的关系（第182节，补充）。因而黑格尔拒绝“宪制的古老区分”，因为他认识到，在“差别”的前提下，奠定这种区分基础的社会和政治制度的“实体性统一”并不是“理性的”，或者按照黑格尔《法哲学原理》的前提所指的同样的东西，不再是“现实的”。它必然沦落为一种历史性的、只合乎“古代世界”立场的现象：“古代把国家制度区分为君主制、贵族制和民主制，这种区分是以尚未分割的实体性的统一为其基础的，这种统一还没有达到它的内部划分（一个在自身中发展了的机体），从而也还没有达到深度和具体的合理性。因此，从那种古代的观点看来，这种区分是真实的和正确的……”（第273节）\footnote{中译文引同上书，第287页。——译者}。

然而，[黑格尔]法哲学与传统的断裂不仅仅表现在，由于国家和市民社会的差别，以前的政治学概念和内容已经变成非理性的，成为了历史的逝去之物；通过为黑格尔的伦理理论奠定基础的国家与历史间的交互关系，这个决裂更是无比清晰地呈现出来。黑格尔的国家概念，往后看，以古代市民社会的“实体性统一”的瓦解和[新]出现的与市民社会的差异为前提；往前看，则与更广阔的世界历史领域相连。黑格尔使市民社会位居国家之下，与他使历史位居[国家]之上，两者形成一种互补的关系。如果我们看一下历史在迄今为止的政治学和自然法理论中的价值，就可以说，两者的共同做法是，都把法的问题或者“好生活”的问题化解为“最佳宪制”问题，也就是说建“国”(civitas)问题。古代人和现代人都认为，“最佳宪制”是与市民社会(societas civilis)、正当地组织起来的并且使得德行的施展得以可能的生活的建立同时发生的。为此，只需要回顾“先前的”科学，如形而上学、人类学、心理学，以及作为德行论和义务论的伦理即可。但是，所有这些科学都只不过构成了开端，作为市民社会的国家的学说基础。在18世纪，随着诸如国民经济学和历史哲学这些新学科的产生，前述科学在实践哲学上的紧密结合也就随之瓦解了。即便是最早在法学和国家学问题与历史哲学问题之间建立起联系的康德，仍然以传统自然法的方式把这个主题阐述为“达成一个普遍地管理法权的公民社会”\footnote{《在世界公民底观点下的普遍历史之理念》，[《全集》]第8卷，科学院版，第22页。（中译文参见〔德〕康德：《康德历史哲学论文集》，李明辉译，联经出版事业公司2002年版，第11页。——译者）}。而黑格尔则不同，在他看来，历史哲学把自身从自然法与“市民社会”理论的联系中解脱了出来。

国家的普遍理念在其完成的顶端下降为“众多国家”的特殊性，在历史的游戏中，这些国家彼此对立。国家理念是“作为类和对抗个别国家的绝对权力”的理念，是“在世界历史的过程中给自己以现实性的精神”（第259节）。在黑格尔这里，国家不再是传统的政治状态概念，不是古典政治学和现代自然法的无运动的、非历史的自然或契约模式，而是历史——作为实践哲学的新法庭——屹立于其上。在历史中，“伦理性的整体本身和国家的独立性”——鉴于家庭的伦理性整体所遭受到的同样命运，我们必须为黑格尔补充道：再次——“都被委之于偶然性”（第340节）。黑格尔在《法哲学原理》的末尾所引入的历史维度，是自然法学家在其体系的开端处设定的那个自然状态的理念的现实性。在自然法学家那里，运动乃是从自然状态走向“societas civilis”意义上的市民社会，并在这里进入静止状态；在黑格尔这里，则是当国家不再关系到市民社会，而是关系到其他国家的时候，这种运动便开始了。这种自然状态是一种现实的、绝非虚构的状态，是历史的运动，《法哲学原理》把这种历史运动接受到自身之中，借此再一次从抽象的法和社会的自然理论中解放出来。法概念的自我展开从“普遍精神”的定在要素开始，并结束于它。在世界历史中，普遍精神的定在要素就是“精神的现实性”，在其中人和物的“自然”并不是一种永恒的规律，而是黑格尔思为自由的概念，它“在其内在性和外在性的全部范围内”实现它自己。精神的运动与法哲学结束、历史哲学开始同时发生，在这一运动的普遍性之前，“只是作为观念性的东西……”的家庭、市民社会和国家，“并且在这个要素中，精神的运动就在于把这一事实展示出来”（第341-342节）\footnote{中译文参见：《法哲学原理》，第351页。除了“观念性的东西”外，引文后面的强调为里德尔所加。此外，引文来自第341节，而不是来自第341-342节。——译者}。

\section*{4}

如果对黑格尔《法哲学原理》的结构予以总结性地考察，那么我们可以说，那些具有历史意蕴的概念形式为逻辑上的各个建构原则奠定了基础。那些伦理事物的制度，特别是“社会”这个概念，还有“世界历史”的终结概念，都是逻辑的—历史的范例。这两个概念同时改变了古典政治学和现代自然法数世纪以来一以贯之地予以阐述的家庭和国家的体制。就像家庭处于作为成员而从属于它的个人与市民社会之间一样，国家也运行于市民社会和历史之间。在黑格尔这里，正是传统政治制度结构上的这种缝隙，推动了辩证地交织在一起的法概念的发展历程及其体系的精巧的结构形式。由此出发，《法哲学原理》尤其获得了其所特有的丰富的历史内容，正是这一点把它从原则上与近代自然法著作区分开来。

然而，《法哲学原理》作为制度辩证法的体系引入现代世界政治理论之中的这些新元素，也成为了黑格尔身后那些理论的出发点，这些理论动摇了他的体系，突破了他的那些中介。对《法哲学原理》的批判，正好就是挑战黑格尔本人在法的现实性的哲学重构中对逻辑的东西与历史的东西所提出的中介：(1)国家和历史的关系和(2)市民社会和国家的关系。对于黑格尔《法哲学原理》（1833年）的最早出版人爱德华·甘斯(Eduard Gans)来说，该书最重要的价值在于，黑格尔在书末为我们描绘的“壮丽景象”：他让国家“卷入历史的世界海洋”。甘斯不无道理地指出，黑格尔对于向历史的过渡所做的简要勾勒“不过预示了”这一领域中的“更为重要的关切”\footnote{《法哲学原理》，前言，《全集》第8卷（1833年），第VII—IX页。}。由此，如同他自己（以及后来的黑格尔主义者费迪南德·拉萨尔）已经用他的《世界历史发展中的继承权》（1824-1835年）向历史主义思维方式致敬一样，甘斯预见到了支配了19世纪的历史主义对哲学的未来批判。历史主义的本质在于，它不再把国家理解为自身的现实，而是理解为历史的产物。基于这种观点，阿诺德·卢格（1842年）对黑格尔的法哲学进行了批判。在卢格看来，《法哲学原理》并不关心时代和历史的进步运动，因为它在前述意义上定义国家。然而，国家实际上并不应当是普遍的现实，而是如同卢格从历史出发对其进行批判的1840年的普鲁士国家一样，是一个个别的—历史的实存。尽管《法哲学原理》中的国家令人想到“当下的国家”，是的，黑格尔甚至指出这个国家的名称，但他并没有让它从“历史过程”中产生出来，这也是为什么黑格尔对政治的生活意识和时代意识的发展毫无影响的原因所在\footnote{卢格：《黑格尔法哲学与我们时代的政治》（Die Hegelsche Rechtsphilosophie und die Politik unserer Zeit），《德意志年鉴》1842年卷，第762页以下。}。

批判的第二个出发点就是市民社会和国家的关系，这是青年马克思（1841或1843年）评论《法哲学原理》第261-313节的主题。马克思颠倒了黑格尔对市民社会和国家的关系的思辨解释：国家是市民社会的现象，而黑格尔贬为“伦理事物的现象世界”（第181节）的市民社会，那种远为坚实的、同时完全非思辨地理解的政治经济学的现实，对马克思来说，才是理解国家和历史的关键。马克思由此取消了黑格尔法哲学的中介。在这里，国家也就从一种现实贬低为一种产物，也就是经济的产物。政治经济学批判的问题不再是市民社会的“私人”原则与政治国家的中介，而是对这个市民社会及其私人—个人主义原则本身的批判。马克思抛弃《法哲学原理》，转向英国人和法国人的政治经济学。这并没有妨碍他对那种在黑格尔市民社会理论中开始浮出水面（第243-245节）的贫富辩证法予以批判。通过将[贫富]这两个极端及其基础（黑格尔发现的市民社会领域）独立出来，马克思将辩证法驱至一场[新的]运动，它超越了《法哲学原理》从个别的个体经过家庭和市民社会再到国家的概念运动。按照将马克思与圣西门、孔德、约翰·斯图亚特·密尔和斯宾塞联系起来的社会学的思维方式，国家是社会运动的产物，而不是黑格尔赋予它的那种伦理“理念”的现实性。

就此而言，《法哲学原理》实际上处于一个终结点，但它并没有那种产生于工业社会早期阶段的“国家式微”的思想。它只是通过把国家和市民社会分离开来，为这一思想的出现做好了准备，但却未能设想到这一思想。黑格尔哲学概念还包含了限制的智慧，因此他不可能通过其哲学概念做到这一点。因为哲学既不想教导世界应当怎样，也不能预言世界将会怎样。在黑格尔看来，哲学作为世界的思想，只有在历史形成过程完成和“结束”的时候，才会显现在时间之中。因此，法哲学执着于一种“旧”的生活形态，这种形态可以认识，但是不能使之变得年轻或者进行改变。认识的任务就在于，不要让哲学超出它自己的时代界限。黑格尔在《法哲学原理》序言的最后说道，密涅瓦的猫头鹰只有在黄昏来临之后才起飞。但这并不意味着，哲学对我们而言就属于过去，其本身已经成为“历史的”了。时下，东方和西方之所以同样都对《法哲学原理》进行深入的研究，是因为工业社会的现实在转入其对立面之前，绝对不会兑现圣西门和马克思的那些超越了它的发展的思想。在现代社会的形成过程中，国家变成了一种神话般的存在，它在自身中将权力和智慧统一了起来，并且剥夺了具有矛盾可能性的个体的自由，而在黑格尔看来，自由乃是思想的源泉。意识形态和伟大预言的时代并没有阻挡住其所反对的政治统治的神话，反而是促进了它。利维坦的回归是向哲学提出的一项挑战，哲学必须接受这一挑战，以便能够重新找回它的经典志业：通过理性的力量摧毁神话。黑格尔的《法哲学原理》就是这种力量的见证。绝非偶然的是，人们已经再一次从黑格尔开始，在一种“自身”从政治统治中解放出来的基地上对自由、法和国家的关系进行讨论。今时今日，澄清这种关系要比以往更加紧迫。

\part{自然法与实定法}
\chapter{自由法则与自然的统治}

在《法哲学原理》的序言中，黑格尔把这部著作中阐述的“伦理世界”理论与自然哲学进行了对比。在涉及自然的时候，大家都承认，哲学要如其所是地认识自然，自然“自在地”是理性的，并且科学要概念把握自然中的现在的、现实的理性，然而由人所产生的且依赖于人的意志的伦理—历史世界则处于任意和偶然的摆布之下，就此而言，就应该是无法则的和无理性的。\footnote{《法哲学原理》，序言，《全集》第7卷，H. 格洛克纳编，1927-1930年，第23页以下（《百年纪念版》）。}

关于自然规律和法权法则之间的区别这一哲学史古老主题，黑格尔自己并没有在《法哲学原理》的文本本身中，而是在柏林时期的讲座导言中进行了处理。而在甘斯出版的第二版《法哲学原理》（1833年）的序言中，紧接在这个对比之后，则有一个较长的补充。\footnote{甘斯编辑的注释以H. G. 霍托1822-1823年冬季学期讲座笔记中的一段话为基础，参见《法哲学：依据黑格尔教授先生演讲（1822-1823年）》（Philosophie des Rechts, Nach dem Vortrage des H. Prof. Hegel, im Winter 1822/1823），柏林，手写笔记第2册，第1-4页（柏林普鲁士文化遗产基金会国家图书馆黑格尔遗稿）。进一步参见《1824-1825年冬季学期法哲学讲演录》（Vorlesungen über Philosophie des Rechts WS 1824/1825），格里斯海姆笔记，四开本德语手稿第545号，第13页以下。[这两本笔记现收入《黑格尔全集》（历史考订版）第26卷第2分册和第3分册（G. W. F. Hegel, Gesammelte Werke, Band 26, 2-3. hrsg. von der Rheinisch-Westfälischen Akademie der Wissenschaften, Felix Meiner Verlag Hamburg 2015）。——译者]}按照这里所做的区分（这种区分同样依据一种古老的概念分类），有“两类规律或法则(Gesetzen)，自然规律和法权法则”。自然规律确实存在，它们的存在就使它们有效，由此我们即把“一般自然”思维为由规律所规定的。存在(Sein)和效力（应当）在自然规律领域是一致的：“因为这些规律是准确的，只有我们对这些规律的观念才会错误。这些规律的尺度是在我们身外的，我们的认识对它们无所增益，也无助长作用，我们对它们的认识可以扩大我们的知识领域，如此而已。”\footnote{参见《法哲学：依据黑格尔教授先生演讲（1822-1823年）》，第1页，以及《法哲学原理》序言，第24页，附释。（中译文引自《法哲学原理》，第14页。——译者）}就此而言，自然规律与法权法则没有区别，我们也是从“外部”、作为某种给定的东西认识法权法则的，而法律人也必然停留在这种给定存在的立场上。按照这种立场，法则(Gesetze)之所以有效，乃是因为它们是被设定起来的(gesetzt)。然而，从法权法则的差异性和可变性就可以非常明显地看出，它们并不像自然规律那般“绝对”，因此存在和效力在这里就是彼此分离的。法权法则在字面意思上是“被设定的东西”，一种“源出于”人的东西，并依赖于人的意志和意识。自然规律直接规定物的存在，其效力根本而言是不可变的，法权法则由于这种依赖性而处于不断的变化之中。法权法则具有历史，因此人不会让它们直接存在和有效；相反，它们的存在和效力都要以人自身即人的意志和意识的主体性为中介。我们无法像服从自然规律的必然性那样服从由一位立法者所颁布的法律的权力和权威的必然性，因为自然规律调整的乃是人之外的物的运行，而法权法则的尺度不再在我们之外，而是就在我们之中。

在自然界中，一条规律的最高确证是，它一般地存在。而在法权领域，一条法则的确证并不是在其存在中，而在于它被认知和意求，因为只有在这里存在和应当之间的冲突才是可能的。在自然规律和法权法则的那种对比中，最重要的立场并不是自霍布斯、休谟和康德以来出现在自然法基本概念中的二元论。黑格尔感兴趣的是这种概念体系中更深层次上的二元论——自然与精神之间的对立，它通过激化存在和应当的对立来扬弃这一对立。黑格尔的对比进一步化为这种尖锐的对立：

\begin{quote}
如果我们对这两类法则（或规律）的这种差别进行考察并追问，法权法则所归属的基地是什么，那么我们就会看到：法只产生于精神，因为自然没有任何法。因此就有一个存在着的自然世界和一个与自然相对立的精神世界。\footnote{《法哲学：依据黑格尔教授先生演讲（1822-1823年）》，第2页。（本段引文见《黑格尔全集》（历史考订版）第26卷第2分册，第772页。——译者）}
\end{quote}

在这种意义上，黑格尔在《法哲学原理》第4节中说，法的基地一般说来是“精神的东西”，并且其出发点是自身意求着的“自由意志”，自由意志在法的体系中有其历史形态、“实现了的自由的王国”，即“从精神自身产生出来的、作为第二自然的那个精神的世界”\footnote{《法哲学原理》，导论，第4节，第50页。（中译文参见《法哲学原理》，第10页。——译者）}。

作为第二自然的精神世界，并不是亚里士多德的那种以礼法($\nu\acute{o}\mu o\varsigma$)和习惯($\check{\epsilon}\theta o\varsigma$)为基础的素朴伦常和城邦伦理的第二自然($\delta\epsilon\acute{\upsilon}\tau\epsilon\rho\alpha\ \phi\acute{\upsilon}\sigma\iota\varsigma$)，而是一种由人产生并设定于作品之中的自然。就此而言，较之于亚里士多德的《尼各马可伦理学》或《政治学》，这种第二自然就更接近于霍布斯的《利维坦》。霍布斯在现代自然法理论的开端处即借助“自然”社会和“人为”社会的经典对比，削弱了亚里士多德的城邦的自然理论以及传统的智者派的礼法和自然的二分理论。实际上，几乎所有黑格尔借以区分自然规律和法权法则的那些谓词，都可以在这种经典对比中重新找到。“自然社会”中的和谐（霍布斯和黑格尔都认为，这种和谐是动物式的和谐）是“以自然为中介的神的作品”；与之相反，人的和谐则是“人为的”，它以意志的协调一致（契约）为中介。\footnote{霍布斯：《法律原理》（Elements of Law）第1卷，第19章，第5节；《论公民》（De cive）第2卷，第5章，第5节；《利维坦》第2卷，第17章。}与黑格尔一致，霍布斯首先也把自然物视为是持存的和必然的东西，而人所产生的法律和契约则是人为的作品，由于依赖于人的意志，它们便被视作任意的和可变的。个中缘由完全在于自然(Physis)和礼法(Nomos)这一对概念主题：自然性的东西不依赖于契约和行为（在黑格尔看来，则是不依赖于意志和意识），而人为的东西则始终受制于它们。但是霍布斯从这一对比中得出了一个不同于智者的二分以及柏拉图和亚里士多德对这种二分消解的结论。霍布斯不再想要表明，唯有一个自然社会（也就是以自然尺度为准的社会）才能对个体施加义务，一个人为的社会则让个人保持自由；相反，他想要证明，唯有一个严格意义上的“人为”社会才可以向公民施加服从的义务——在涉及自然的时候，已经不再能够有意义地使用“义务”概念了。就此而言，与礼法和自然的二分相对，义务的思想恰好从“自然”颠倒为“人为”。\footnote{对此，参见沃尔夫（F. O. Wolf）：《霍布斯的新科学：论近代政治哲学的基础》（Die neue Wissenschaft des Thomas Hobbes. Zu den Grundlagen der politischen Philosophie der Neuzeit），斯图加特-巴德·坎施塔特，1969年，第82页以下。}

这种颠倒对近代自然法理论的意义，可以从圣经中造物主的功能上看出来：造物主把自然性的东西贬为一种被造物，自然性的东西不再在自身中，而是在上帝的意志有其根据。对于义务的概念来说，这就意味着：如果现在还要宣布自然性的东西对人具有约束力，那么这种约束力的根据就不再是自然，而只在与自然相分离的上帝的意志之中。这种模式引起了更进一步的转变：一种完全不同的约束力与一般“自然”的约束力相对立，同时也抛弃了从圣经的造物主的意志中推导出这种约束力的做法。\footnote{参见上书，第84页。沃尔夫明确地阐述了这一观点。}在黑格尔看来，迈出这一步的不是霍布斯而是卢梭。他在《法哲学原理》第258节中把这算作卢梭的功劳：“一个原则，它不仅在形式上（诸如社会性冲动、神的权威），而且在内容上也是思想，并且是思维本身。这就是说，他提出意志作为国家的原则。”\footnote{《法哲学原理》，第258节，附释，第330页。（中译文参见《法哲学原理》，第254页。——译者）}卢梭把“人的最内在的东西”、作为“与自身相统一的”自由提升为法的基础，并且借此给了人以对抗自然或者说转移到自然上的上帝意志的“一种无限的力量”。\footnote{《哲学史讲演录》第3卷，《全集》第19卷，第527页。（中译文参见〔德〕黑格尔：《哲学史讲演录》第四卷，贺麟、王太庆译，商务印书馆1978年版，第233页。——译者）}同样，自然、自由与法之间的对立已经出现在霍布斯的作品中。在黑格尔看来，现代自然法的真正历史开始于霍布斯，因为霍布斯最先“试图把维系国家统一的力量、国家权力的本性回溯到内在于我们自身的原则，亦即我们承认是我们自己所有的原则”。\footnote{参见《哲学史讲演录》第3卷，《全集》第19卷，第442页。（中译文引自《哲学史讲演录》第四卷，第157页。——译者）}霍布斯与这种假设决裂：存在着一种超出人的意志之上的、支配着万物的自然目的法则(lex naturalis)，理性要成为“正确的理性”(recta ratio)并使“正义的”行动成为可能，就必须摹写、效仿这种法则。在霍布斯那里，就自然的“规律”包含了使得抛弃自然法成为必然条件而言，法和社会的理论仍居于自然的条件之下。通过霍布斯，在自然法的传统名称中产生了一种“歧义”，对黑格尔来说，这种歧义至关重要，他在霍布斯讲座中也把注意力放在了这种歧义上面：“‘自然’这一名词具有歧义，一方面，人的自然指他的精神性、合理性而言；另一方面，他的自然状态则指人按照他的自然性而行动的状态而言。”\footnote{同上书，第443页。（中译文引自《哲学史讲演录》第四卷，第159页。译文稍有改动。——译者）}在黑格尔看来，卢梭加剧了霍布斯所进行的与传统自然法理论的决裂，这种歧义在卢梭那里进一步升级。卢梭把人的“精神性、合理性”理解为人的自由，自由把人从自然中提升出来，并且成为区分人与动物的根本标志。霍布斯在基于契约的制度性的或者“人为的”统治团体之外，还让主奴关系的“自然性的”统治关系继续存在下去，而卢梭则一般地追问统治的绝对合法性。卢梭既不是将实定法，也不是将霍布斯作为基础的自然规律，而是将自由作为合法性的原则。卢梭并不考虑他的17世纪先行者，在他看来，他们与为奴役人、压迫人进行辩护的亚里士多德和实施这种行径的卡利古拉同属一类。同时，卢梭也不考虑欧洲国家现行有效的法权法则：“他对上述问题做出这样的回答：人是有自由意志的，因为‘自由是人的品德。放弃自己的自由，就是放弃做人。放弃自由，就是放弃一切义务和权利’。”\footnote{《哲学史讲演录》第3卷，《全集》第19卷，第527页。（中译文引自《哲学史讲演录》第四卷，第233页。——译者）黑格尔引用了《社会契约论》第一卷第四章的著名段落，参见第一章、第七章、第十一章；第三卷，第九章，注释；《论人类不平等的起源》（Discours sur l'origine de l'inégalité parmi les hommes），K. 维甘德编，汉堡，1955年，第106页以下。另参见《世界历史哲学讲演录》第4卷，第920页以下。}诚如黑格尔所言，卢梭的原则：人是自由的，并且建立在“公意”基础上的国家是这种自由的实现，“乃是正确的”。在卢梭这里，只有当他让公意由众多个体的意志、它们的自然性的自由倾向来“构成”的时候，才出现了“歧义”。由此，“自然的自由”——个体意志的自发性——又取消了“作为完全绝对物的自由”，卢梭在其公意概念中思考的正是这种自由。\footnote{参见上书，第528页：“这些原则表达得很抽象，我们必须把它们正确地找出来；可是歧义立即就发生了。人是自由的，这当然是人的实质本性，这种本性在国家里不但没有被扬弃，事实上倒是开始被建立起来了。本性的自由，自由的禀赋并不是现实的；因为国家才是自由的实现。”（中译文引自《哲学史讲演录》第四卷，第234页。——译者）此外，参见《哲学全书》，第163节，《全集》第8卷，第360页；《法哲学原理》，第258节，第530页以下。}同时，卢梭也使我们意识到，自由乃是完全的绝对物、“人的概念”：

\begin{quote}
恰恰就是思维本身；要是抛开思维来谈自由，就不知道自己说的是什么东西。思维的自身统一性就是自由，就是自由意志……意志只有作为思维的意志才是自由的。现在自由的原则出现了，它把这种无限的力量给予了把自己理解为无限者的人。\footnote{参见《哲学史讲演录》第4卷，《全集》第19卷，第528页以下。（中译文引自《哲学史讲演录》第四卷，第234页。——译者）}
\end{quote}

意志于其自身之所是，既不能通过与自然及其“规律”的类比，也不能通过意志与自然或自然性地得到规定的自由任意之间的单纯对立来把握。意志——卢梭诉诸“公意”，作为国家的基本原则，使我们想到这一点——必须超出这些对立。意志只有在这种情况下才是自由的，即它“意求的不是他者、外在的东西和陌生的东西，因为如果这样它就是依赖性的，而只是意求自身——意求意志。绝对意志就是意求成为自由”。卢梭用这种自由意志的理念，完成了由霍布斯开启的意求和思维自身的个体的主体性对[传统]自然目的论的颠覆。这种自然目的论构成了传统自然法的自然法则(lex naturalis)的基础。取代自然的约束力以及霍布斯对造物主意志的如此模糊的诉求的，乃是绝对化了的意志的约束力。黑格尔认识到了卢梭对霍布斯模式的这种转变，并且承认这一转变乃是他自己的法概念的最重要的前提：

\begin{quote}
绝对意志就是意求成为自由。意求自身的意志乃是一切权利和一切义务的根据，因此也就是一切法权法则、义务要求和所施加的约束力的根据。意志自由本身就是一切法的原则和实体性基础，它自身就是绝对的、自在自为的永恒的法和最高的法。就此而言，其他的、特殊的法被置于次要的地位，它甚至是使人之成为人的东西，因此也就是精神的基本原则。\footnote{《世界历史哲学讲演录》第4卷，第921页以下。}
\end{quote}

由此出发，黑格尔推出了向康德哲学的过渡。康德通过将自然的立法与自由的立法、经验的意志与“自由的和纯粹的意志”彻底分离开来，为终结在卢梭那里尚且存在的自然法概念中的“歧义”做好了准备。\footnote{《哲学史讲演录》第4卷，《全集》第19卷，第529页以下，第552、590页。}在先验哲学的立场上，意志的“概念”达到了对其自身的理解。“自我意识的简单统一”、“自我”是从一切给定的自然秩序中摆脱出来的、“密不透风的、完全独立的自由和一切普遍规定（也就是思维规定）的渊源，——理论理性，同样也是一切实践规定的最高的东西，——作为自由的和纯粹的意志的实践理性；而意志的理性在于，自身保持在纯粹的自由之中，在一切特殊物中只意求这种纯粹自由，仅仅为了法而意求法，仅仅为了义务而意求义务”。\footnote{《世界历史哲学讲演录》第4卷，第922页。}

由此，“对精神性的东西的意识”将成为法哲学的基础；它发现了一种“国家的思想原则，现在，这种原则不再是诸如社会性冲动、财产安全的需要之类的任何一种意见的原则，也不是像君权神授这样的虔诚原则，而是与我的自我意识同一的确定性原则[……]”\footnote{同上书，第924页。参见拙文，《黑格尔对自然法的批判》，《黑格尔研究》第4卷，1967年，第178页以下。现在也收入《黑格尔法哲学研究》，美茵河畔法兰克福，1969年，第42页以下。（见本书第85页以下。——译者）}

\section*{2}

如果要把握《法哲学原理》体系中自然和自由的关系以及《法哲学原理》在现代自然法终结处的独特的矛盾地位，那么我们就必须注意到黑格尔对霍布斯、卢梭和康德的法原则的这种评价。黑格尔在下述两种意义上站在霍布斯预备好的自然法理论的基地上：一方面，是在对于那种自然法理论具有构成性意义的自然概念的基地上；另一方面，则是在一种从既有的自然性和历史性权力中解放出来的意志的前提下，这种意志必须通过它自身的运动才与一切给定性发生关系，并且使这些给定的东西与自身和解。《法哲学原理》与17、18世纪自然法共享了一种自然的理念，这种自然本质上是否定性的，是通过消除“自然性”目的以及这些自然性目的对“最终”目的的隶属而得到界定的。作为精神哲学的一部分，它并不承认经院哲学中的那种自然的阶梯式进展。按照经院哲学的观点，这种阶梯式进展一直会进到“恩典的王国”(regnum gratiae)，其每一阶梯都有各自的自然地位。《法哲学原理》也不承认[莱布尼茨的]那种“自然王国”与“目的王国”之间的前定和谐学说：[按照这种学说，]“自然王国”与“目的王国”服从不同的原因和法则，并由于两者之间的这种和谐，自然的道路会从自身通向恩典。\footnote{莱布尼茨：《单子论》（1714年），第88节。}《法哲学原理》阐述的是“精神王国”的实现。“精神的王国”以一种目的论为前提，这种目的论内在于自由意志自身之中，并且就是渗透到一切实在性并在其中获得其自由的那种“概念”的运动。精神的阶梯式进展并不处于自然概念之下；它基于它自身的概念。而对黑格尔而言，这也就意味着：精神乃是基于自由之上。只有不断地摆脱自然，并且产生一种“第二自然”即精神的历史世界，自由才拥有其定在。精神的王国是自由的王国，它作为这样的王国不是自然性的等级制度，或者无世界的精神王国，而是“从它自身产生出来的精神的世界”。\footnote{《法哲学原理》，第4节；《哲学全书》，第381节。}

在这一世界——“第二自然”——中，自然与自由就像精神对自然那样互不相容。“精神的自然可以在其最完全的对立中认识到。我们把精神与物质对立起来。就像重量构成物质的实体一样，因此我们必须说，自由乃是精神的实体。”\footnote{《历史中的理性》，第5版，J. 荷夫迈斯特出版，汉堡，1955年，第55页；《哲学全书》，第381节，补充，《全集》第10卷，第19页以下。}在自然法理论中，同一种对立支配了对于自然法理论为构成性的自然的概念。这个概念在自然状态学说中扮演了一种既根本又不幸的角色。对黑格尔而言，这个概念具有一种“重要的、会导致绝对错误的歧义。一方面，自然意味着自然性存在，正如我们发现自己在不同的方面是直接造成的那样，是我们存在的直接方面。与之相对并与之相区分，自然亦是概念，事物的自然[本性]也就是事物的概念，就是事物的合乎理性方式的存在，这种事物完全不同于单纯自然性的东西”。\footnote{《历史中的理性》，第117页。}当“自然”指的是事物的概念时，那么自然状态中的法就必须被把握为那种“按照人的概念、按照精神的概念应当属于人的法”。然而，这不能与精神在其自然状态、在不自由与依赖于自然的状态中所是的东西混淆起来。因此，黑格尔与霍布斯、斯宾诺莎一起说道：必须走出自然状态(Exeundum est e statu naturae)。

在18世纪现代自然科学背景下，这句话的去目的论意义获得了一种独特的尖锐性。只要自然还是被规定为一种由目的所推动的秩序，那么自由就会与自然关联起来，并且人的行动也会被思维为对自然的延续或者模仿。随着自然被构建为一种可以以数学的方式加以描述和以力学的方式组织起来的因果关系，返回到这种自然就成为不可能的了；康德的《纯粹理性批判》为整个近代时期阐明了自由的思想，这种自由思想在自然中再也找不到类似于目的的东西，而是立足于自身之上。这对于黑格尔的《法哲学原理》具有决定性的意义。尽管黑格尔的《法哲学原理》也在重量概念中为其自由的基本概念找到一种自然式的类比，然而这里所涉及的是一个奠定了自从牛顿物理学以来系统理解自然的基础的范畴。《法哲学原理》第4节的补充指出，意志自由最好“参照”物理自然得到解释。自由和重量是意志和物体的“根本规定”。也就是说，这些规定对于意志和物体来说不是偶然的，而是必然的和本质的。它们的基本规定是：“当我们说物质是有重量的，我们可能认为这个谓语只是偶然的。然而并非如此，因为没有一种物质没有重量，其实物质就是重量本身。重量构成物体，而且就是物体。说到自由和意志也是一样，因为自由的东西就是意志。意志而没有自由，只是一句空话；同时，自由只有作为意志，作为主体，才是现实的。”\footnote{《法哲学原理》，第4节，补充，第50页。（中译文引自《法哲学原理》，第11-12页。——译者）}在重力中，自然物体追求一个中心点，也就是由自然物体在空间中的位置和关系所决定的物理体系的重心，但是这个重心对于这些自然物体而言始终是外在的中心。与之相反，在作为意志的基本规定的自由中，意志始终是自身相关的，这种自身相关正好就是意志的基础或者实体。意志尽管也追求一个中心点，但是追求的是一种它已经在自身中拥有的中心点：“它不在自身之外拥有统一；它始终在它自身中发现这种统一，它在自身中且依于自身；物质在它之外有其实体。与之相反，精神则是自在地存在，而这正是自由。”\footnote{《历史中的理性》，第55页。}

黑格尔与现代自然法所共享的第二个前提，涉及的是“自由意志”及其“概念”的起源。《法哲学原理》的基础正是“自由意志”及其“概念”的这种起源，而不是亚里士多德和经院哲学的自然法理论所刻画的那种具有复杂层级结构的统治和社会团体的目的论的阶梯式进展。尽管黑格尔明确地与霍布斯、卢梭和康德的理论保持距离，然而《法哲学原理》的整体安排在一个根本点上仍然是自然法式的：《法哲学原理》的概念发展始于与自然物和其他个体意志发生关系（财产和契约）的“一个主体的个体意志”\footnote{参见《哲学全书》，第400节以下；《法哲学》，第34节以下。——黑格尔一再强调这个属于自由意志的直接性的个别性环节，参见第13节附释、第39、43、46、52节。}，也就是始于“抽象”法的法，这种法重演了自然法学说的前政治状态。我们常常忽视了从个体意志出发的构造贯穿了整个体系，甚至延伸到对凝聚在“国家”中的意志的演绎。在黑格尔看来，这个意志\footnote{指凝聚在“国家”中的意志。——译者}也同样必须是一个“个体”的意志（君主的意志）。\footnote{参见《法哲学原理》，第279节，及附释：“任何科学的内在发展，即从它的简单概念到全部内容的推演[……]都显示出这样一个特点：同一个概念（在这里是意志）开始时（因为这是开始）是抽象的，它保存着自身，但是仅仅通过自身使自己的规定丰富起来，从而获得具体的内容。”（第381节）（中译文引自《法哲学原理》，第296页。——译者）另参见第190节，附释，以及《哲学全书》，第2版，第514节，《全集》第10卷，第397页。}

尽管如此，这个个体意志并不是自然法意义上的“自由”意志，并不是每个人的偶然的、个别的任意；相反，法、普遍意志的概念正是产生于对每个人的偶然的、个别的任意的交互限制。因为只有意志概念在主观精神学说中从一切偶然的自然规定性中被纯化之后，黑格尔才引入作为《法哲学原理》的对象的意志概念，所以《法哲学原理》的概念起源就是从包含在个体意志中的“普遍意志”出发的。因此，对黑格尔而言，自由意志是“理性意志和个体意志的统一，个体意志是理性意志的活动的直接的和独特的要素”。\footnote{参见《哲学全书》，第2版，第485节，第382页以下；另参见《法哲学原理》，第2节，第39页；第4节，第50页以下。}然而，这并不意味着取消自然法思想，实际上毋宁是对它的一个巨大的加强。\footnote{参见罗森茨威格：《黑格尔与国家》第2卷，慕尼黑/柏林，1920年，第106页以下。}借此，个体不再作为其自身，在其自然性中，而是作为理性存在者，而成为法学的出发点和对象。当康德和费希特把他们所发现的自由概念局限在伦理学领域并把自然法建立在每个人的（由自然所规定的）自由任意之上时，黑格尔也把这种自由概念运用于法学中。法不限制自由意志，而是自由意志的“定在”，即“作为理念的自由”。对黑格尔来说，这一措辞说的是：自由不再仅仅是规范和公设（康德意义上的理念），而是事实，它存在于社会历史世界之中，而不仅仅是某种没什么指望的东西。

为了思维这种理性的事实，就需要黑格尔在《法哲学原理》导言中将自然法的意志“概念”应用于其下的那种辩证法。依照黑格尔的观点，康德尽管在先验哲学的自我原则中发现了意志的自我决定的奠基性环节，即“纯粹的无规定性或自我在自身中的纯粹反思”的环节，在其中，所有的限制和每一种由自然“给予的”内容都消解无余了，但是自为地看，意志的自由仿佛只是消极的；在“政治”中，意志自由（如同在以卢梭理论为范本的法国大革命进程中所看到的那样）导致了“一切现存社会秩序的破坏”和“对任何企图重整旗鼓的组织的消灭”。\footnote{《法哲学原理》，第5节，第54页以下。（中译文参见《法哲学原理》，第14页。——译者）}因此，自我作为意志的纯粹概念，必然在自身中包含了第二个环节，即“规定和设定一种规定性作为一种内容和对象”，借此自我将自身设定为有规定的东西，也就是说进入定在。这种对意志予以规定和限制的特殊物，绝不是任何给定的外在限制，而是（在黑格尔看来，康德和费希特忽略了这一点）内在于自我决定的行为之中的。在这两个环节（即意志的无规定的普遍性和有规定的特殊性）的统一中，作为意志的“概念”的自由将自身实现为第三个环节，——意志在他在中自在地存在，或者说“自我的自身规定，在一中，它设定自己为[……]被限制的东西；它留在自己那里，即留在与自己的同一性和普遍性中；又它在这一规定中只与自己本身联结在一起”。\footnote{同上书，第7节，第58页以下。（中译文引自《法哲学原理》，第17页。译文稍有改动。——译者）}

\section*{3}

自由意志的辩证运动将普遍性的环节与特殊物\footnote{指普遍性的环节与特殊物。——译者}中介起来，也将这两者与它自身中介起来，而成为“具体的”个体性、“自由意志的定在”（《法哲学原理》，第29节）。自由意志的这一辩证运动，就是我们在黑格尔《法哲学原理》中作为“概念”而见到的那种“事物的自然[本性]”。概念把自然的意义完全把握在自身中；概念指的无非就是把它自身以及它的所有规定都建立在“自由的人格性”上面的法，也就是“一种毋宁说是自然规定的对立面的自我规定”。\footnote{《法哲学原理》，第7节，第58页以下。（里德尔原注疑误。该书第7节没有这段引文。——译者）也参见《纽伦堡哲学入门》，《全集》第3卷，第55、71页。特别是《纽伦堡哲学入门》，第三篇，第2章；《哲学全书》，第三部，二之1 法，第181节：“精神作为自由的、具有自我意识的存在者，是自相等同的自我，这种自我在其绝对否定关系中首先是排斥性的自我，即个别的、自由的存在者，或人格(Person)。”[本段引文见《黑格尔全集》（历史考订版）第10卷第1分册，第355页。中译文引自《黑格尔全集》第10卷，第289页。译文稍有改动。——译者]}概念的辩证运动开始于摆脱了所有“自然性”规定的“人格”的自我规定，也就是开始于作为法的唯一原则的自由意志的无限的、不受任何外物限制的“概念”。黑格尔在柏林时期的法哲学讲座中反复讲述了概念（即自然）、自由和法之间的这种联系。第一次柏林法哲学讲座在1818-1819年、《法哲学原理》出版之前。这次讲座第3节说道：“法的原则不在于自然，既不在外在自然，也不在人的主观自然，因此时人的意志是自然地被规定的，也就是说，是欲望、冲动和偏好的领域。法的领域是自由的领域，在其中，自由外在化自身并赋予自身以实存，此时尽管自然会出场，但却是作为一种没有独立性的东西[而出场的]。”\footnote{《黑格尔教授1818-1819年冬季学期讲授的自然法和国家法》（Natur- und Staatsrecht nach d. Vortrag des Professors Hegel im Winterhalbjahr 1818/1819），霍迈尔(G. Homeyer)笔记，四开本德语手稿第1155号，第9页。[本段引文见《黑格尔全集》（历史考订版）第26卷第1分册，第239页。——译者]}正如法哲学所述，从抽象法的现象形式开始，直到家庭、市民社会和国家为止，自由的各种外化就是社会制度和行为方式的所有环节。毫无疑问，在所有这些环节中，处处都有“自然性”联系的身影。但问题在于，这些“自然性”联系是否从自身产生出一种支配那些形式的法律。在黑格尔看来，必然性——比如说，在与个体生命的关系中，国家的定在所具有的那种必然性——不再意味着，对于个体而言，必须在国家中生活乃是一种自然的规律。毋宁说，国家的必然性建立在自由自身所立的法律基础上。黑格尔对第3节的解说毫不含糊地指出：“这一节是由自然法之名引发的。自由与必然性的统一并不是通过自然，而是通过自由产生出来的。自然事物一成不变，自己并不能从规律中摆脱出来，以便自己制定法律。然而，精神则使自身从自然中挣脱出来，并且自己产生它的自然、它的法则本身。因此，自然不是法的基地。”与康德和费希特的看法一样\footnote{参见康德：《反思录》，第7084条，康德：《全集》（科学院版）第19卷，第245页。费希特：《自然法权基础》，《全集》第3卷，第148页。}，黑格尔也认为自然法的名称“不过是习传下来的”，按照事物[本身]、法的“概念”，这一名称甚至是不正确的，因为“自然可以理解为：(1)本质、概念。(2)无意识的自然（真正的意义）”；“真正的名称”必须是“哲学法学”。\footnote{《自然法和国家法》，霍迈尔笔记，第9页。[本段引文见《黑格尔全集》（历史考订版）第26卷第1分册，第240页。——译者]}

因此，在起草法哲学体系的时期，黑格尔的思想发展完全回到了《论自然法的科学探讨方式》一文想要通过“伦理自然”法的建构予以克服的那种自由哲学的基础上。在结束了对“哲学法学”基本概念的漫长理解过程（这一过程同时也澄清了哲学法学与传统自然法的历史关系）之后，现在对此也就无须赘言了。在黑格尔的体系中，自然与自由、自然规律与法权法则分离开来了。因而黑格尔决然地拒绝传统观念：就只有人可以认识和遵守自然规律而言，“自然规律”对人来说可以成为法的“范本”；人的历史性定在所唯一具有的法则，就是自由法则，不是“自然”，而是“概念”自身制定了这一法则。\footnote{参见上书，第3页：“对我们而言，事物的概念并不来自自然。”[本段引文见《黑格尔全集》（历史考订版）第26卷第2分册，第774页。——译者]}概念的自我立法——对于黑格尔而言，这意味着18世纪发现的人的普遍法权能力的观念\footnote{参见上书，第5页：“这须被尊重为某种伟大之事：现在人因其为人，必须认为拥有权利，以至于他的人的存在要高于他的身份。”}，哲学通过这种观念去反对现成的东西，并且打破了自然性的假象，在这种假象背后，“法的思想”、自由都还一直隐而不彰。这种思想将现成的法和对它的知识变为同一个东西。黑格尔在1824-1825年的讲座导言中说，尽管在以自然、自然需要和目的等为基础的“习惯的自然法”中，自由的思想并没有被有意排除掉，但“事实上，由于没有领会到这两种原则的独特性就接受了它们，从而怠慢了自由”。自由不能作为自然，而是必须在“概念”的规范性中去思维，在自由与其自身以及自由与自然之间，概念起着中介的作用。当自由以“概念”的要求出现时，这种要求就包含了“自由仅仅指与法和伦理有关的自由，这样自然法之名就陷入了动摇之中。”\footnote{《法哲学讲演录》，格里斯海姆笔记，第9页以下。[本段引文见《黑格尔全集》（历史考订版）第26卷第3分册，第1054页。——译者]}只有当我们“承认”，自由就是自然法所公设的“事物的自然[本性]”——也就是“概念”——的时候，我们才能保住自由的概念。然而，这种表达是“不恰当的”，因为不同于自由，自然“是无拘束的，并不与某个它物相对立。而自由则相反，同时显得是好战的，并且具有众多的对立物，而它的头一个对立物就是自然本身”。\footnote{同上书，第10页。[本段引文见《黑格尔全集》（历史考订版）第26卷第3分册，第1054页。——译者]}

\section*{4}

对于黑格尔法哲学的体系安排和概念处理来说，自然与自由的对立具有决定性的意义。这一对立不仅决定了这部著作两个标题中“哲学法学”和“自然法”的区分\footnote{黑格尔在1824-1825年讲座中明确提到了这一点。参见格里斯海姆笔记，第3页；以及拙文“黑格尔《法哲学原理》中的传统与革命”见《黑格尔法哲学研究》，前揭，第100页以下。（见本书第170页以下。——译者）}，而且也决定了体系最后也是最重要的部分随着在家庭与国家之侧引入“市民社会”概念而经验到的双重化，因为区分市民社会和国家的体系根据在于自然概念和自由概念的分离。黑格尔在历史中思考终结这种分离。在近代自然法的发展进程中，国家(civitas, res publica)与市民社会(societas civilis)的同一瓦解了。这种同一曾在传统的“封建”国家和社会理论上面打下烙印，并且在哲学上通过自然规律(lex naturalis)的目的论而得到合法化：这种目的论让个体的“自然性目的”（冲动、需要等）与整体的伦理目的彼此协调一致。然而，这种同一性的瓦解并没有导致一种持久的差别，而是再次被克服了。吊诡的是，其原因恰恰在于自然法的契约理论。这种理论既是分离的催化剂，同时又是克服这种分离的手段。一方面，自然的目的论的、超出个体和市民社会之定在的运动秩序遭到了抛弃，而另一方面，自然又在自然法学说的出发点上，作为多的“任意”的环节，进入契约的建构之中。如果每个人的任意都应当能够按照一条普遍的法则（这一普遍法则取代了被废弃的自然法则）而与所有人的任意相一致的话，那么法概念所必须限制的，就正是包含在“自由任意”这一说法之中的那些需要、冲动和爱好。或许在这里可以找到黑格尔之所以反对康德和卢梭之引起争议的法原则与自由原则的界分的真正动因，对此，黑格尔在《法哲学原理》中也非常重视。与此同时，黑格尔自己也对卢梭—康德与旧自然法之间在自然观上所具有的区别非常清楚。传统的探讨方式（在这种探讨方式中，自然概念的“两种含义”还没有发生分离）以人的“自然性存在”即“爱好、需要、冲动”为基础：首先是肉体的需要及其满足的必然性，然后是“社交的冲动、社会的冲动，这种冲动一方面涉及两性关系，另一方面也扩展到普遍的市民社会”，最后自身“发展为国家”。\footnote{参见《1824-1825年讲演录》，格里斯海姆笔记，第5-8页。《法哲学原理》导论，第19节，附释。}而对黑格尔而言，这里正好表现了“这种事情的不足”。并不存在自然的目的论，可以对个人的冲动和需要的本性与“国家”的伦理历史性定在进行中介；相对于国家和市民社会按照“概念”之所是，所有从中推演出来的规定都是“抽象的”：

\begin{quote}
如果我们必须说，市民社会、国家是从一种冲动中产生出来的，这马上就表明了从一种冲动到国家这种事情的不足。之前在哲学中被假定为国家之基础的社交冲动，乃是某种缺乏规定的、抽象的东西，对于具有丰富层次的国家来说，它只能提供少量的规定，而且对它所要奠基的东西来说，它也显得太过贫乏了。\footnote{参见格里斯海姆笔记，第8页。[本段引文见《黑格尔全集》（历史考订版）第26卷第3分册，第1053页。——译者]}
\end{quote}

霍布斯最先抛弃了这种自然法原则，其后卢梭与康德费希特哲学又把它与自由原则对立起来。但是，正如黑格尔在1824-1825讲座中所言，[卢梭、康德和费希特]“所采取的方式并非我们的方式，也没有让法学成为一种全面的和一以贯之的科学”。\footnote{同上书，第11页。〔本段引文见《黑格尔全集》（历史考订版）第26卷第3分册，第1054页。其中“我们的”在历史考订版中为“真正的”。——译者〕也可参见《1823-1824年讲演录》（Vorlesungen von 1823/1824），霍托笔记，第7页以下。}黑格尔同意康德，认为在自然规律之下自由是不可能的，因为自然和自由彼此对立。因此他和康德一样，拒绝传统目的论自然观中所包含的对“国家的历史起源”\footnote{参见《法哲学原理》，第258节，第330页。（中译文参见《法哲学原理》，第254页。——译者）}的追问，或者用政治学和自然法的传统术语来说，对“市民社会的开端”的追问，因为法哲学只是涉及“国家的理念”、“被思维的概念”。而在康德和费希特那里（以及他们之前的卢梭那里），他们作为法权学说基础的体系原则是与他们对体系原则所做的历史阐述相矛盾的。因为“思维本身”是作为意志，国家的原则“不仅在形式上（诸如社会性冲动、神的权威），而且内容上也是思想”，因此一方面这种原则把自然从自身中排除了出去，而另一方面又通过他们把仅仅在历史规定的形式上的“个体意志”中的作为“特殊个体”的意志置于顶端，而将自然重新引了进来。\footnote{参见《法哲学原理》，第258节，以及第29节，第79页以下，第182节，补充；《历史中的理性》，第111页以下。}结果，他们就不是把“普遍意志”、国家的“概念”理解为“意志的自在自为的合乎理性的东西，而是理解为从这种个体意志中产生出来的共同的东西”。由此，在国家和市民社会的“概念”方面，近代自然法和启蒙的古典自由主义又一起重新落到了之前的自然法学说的立场上。霍布斯、卢梭和康德作为前提来瓦解国家(civitas)和市民社会(societas civilis)的古典同一性的契约建构重复了那种对两个概念领域的语言体系上的“混淆”，在《法哲学原理》第182节中，黑格尔拿这种混淆去反驳“国家”的自然法“看法”。\footnote{同上书，第182节，补充：“如果把国家想象为各个不同的人的统一，亦即仅仅是共同性的统一，其所想象的只是指市民社会的规定而言。许多现代的国家法学者都不能对国家提出除此之外任何其他看法。”（第262页以下）（中译文引自《法哲学原理》，第254页。——译者）}

这一反驳之所以可能，前提是黑格尔自己从近代自然法的出发点所得出的结论。通过把自然规律和自由法则彻底对立起来，黑格尔不仅抛弃了传统自然法名称，而且也放弃了市民社会(societas civilis)这一说法。与抛弃传统自然法概念相比，后者所造成的影响要深远得多。随着黑格尔引入国家和市民社会的分离，一个问题视域被打开了。在这一问题视域之内，各个个体的生命和意志的关系、他们彼此之间的联系以及他们与自然的联系就能在理论上获得独立且相对隔绝开来，作为“社会”联系得到研究和把握。恰恰因为个体意志的自然规定性、“特殊的个体”及其需要、冲动和爱好与普遍的、在国家中自身显现着的意志的演绎并没有任何直接的论证关系，所以黑格尔才有可能发现那种作为现代市民社会的经济基础的“需要的体系”。而对自然法理论的自由主义信徒们来说，这个需要的体系却被社会契约模式及其政治意涵遮蔽了。当然，在黑格尔这里，自然和自由的关系问题在《法哲学原理》中也并没有得到满意的解决，而是在一个新的层次上重新出现了。“作为理念的自由”在概念上要求市民社会从国家中脱离出来，由此，自由也就在[市民社会]这一层面从现实性坠入表面现象、“伦理的现象世界”。在这个世界中，伦理的同一性不能作为自由，只能作为必然性而得到思考。\footnote{《法哲学原理》，第186节，第267页；参见第184-185节。}这里，集中于自由意志领域的概念运动已经蓄势待发，准备转变为一种社会的运动法则，这种运动法则的基础不是作为概念之前提的自由，而是自然——“不平等的要素”。“不平等的要素”表达了精神的特殊性的客观法。其中，不仅自然设定的不平等没有被扬弃，反而“从精神中产生出来”。\footnote{同上书，第200节，附释，第278页以下。（中译文参见《法哲学原理》，第211页。——译者）此外，黑格尔在对世袭君主进行演绎（第280节以下）和对战争进行辩护（第324节）时，都援引了概念所无法化解的自然作为依据。}这里构成了突破黑格尔法哲学的地方之一。在这些突破黑格尔法哲学的地方，概念的辩证起源达到了它的界限——也就是黑格尔的最重要的、最具有历史影响力的学生马克思想要通过把这种辩证法“唯物主义地颠倒”为一种生产力和生产关系的运动而加以超越的那种界限。
\chapter{黑格尔对自然法的批判}

在1802-1803年发表于《批判的哲学杂志》的论文中，黑格尔阐述了他对迄今为止的“自然法的科学探讨方式”的第一次大批判。在黑格尔学界，这篇论文在个别段落上争议极大，但在整体上却获得了近乎一致的评价。这种评价大体上概括如下：如果我们忽略这篇论文里面的一些晦涩之处，则可以认为，这篇自然法论文已经包含了1821年《法哲学原理》的纲要。如同在其他地方一样，这里，这种评价的文献来源是罗森克兰茨的《黑格尔传》一书。在罗森克兰茨看来，黑格尔在其成熟时期的法哲学中只不过是在一个精妙的体系中更为专门、详细地展示了所有的概念，相比之下，其构思的原创性已经在将其表达得“更美、更新鲜、部分地也更为真实”的杂志论文的“青年形态”中呈现出来。\footnote{罗森克兰茨，《黑格尔传》，柏林，1844年，第173页以下。}罗森克兰茨的判断在代表了1910年到1945年之间的新黑格尔主义的海林(Theodor Haering)和格洛克纳(Hermann Glockner)的黑格尔学述中延续了下来。格洛克纳在这一点上走得最远。他精炼地指出，黑格尔在其辩证方法最终形成之前，即已凭借自然法论文，作为伦理学家“完全”出现在我们面前了。\footnote{格洛克纳(H. Glockner)：《黑格尔》(Hegel)第2卷，第2版，斯图加特，1958年，第302页以下。也见海林(Th. Haering)：《黑格尔：他的意愿和他的著作》(Hegel. Sein Wollen und sein Werk)第2卷，莱比锡与柏林，1938年，第389页以下、404页以下。}

我们想问的是：情况真是这样吗？毫无疑问，就材料的丰富及其相对清晰的体系安排而言，自然法论文超出了青年黑格尔所有以前的著作。然而争议在于，黑格尔在基本概念的处理上是否有所超越，有多大程度的超越。由于对“绝对伦理”的那些概念（人民、政府等）的内容上的兴趣，新黑格尔主义跳过了这个问题。为了能够对这一问题做出决定，就需要反思为黑格尔自然法的伦理—政治范畴奠定基础之理论上的基本概念，首先就是“自然”概念。由于黑格尔自己认为需要对——正如他在《自然法》这一标题中表达的——自然和法的联系进行反思，因此之前的黑格尔学界居然没有提出这个本身就很明显的问题，就更是令人惊讶了。对这一联系进行考察，我们可以得知黑格尔法哲学的发展历程，以及黑格尔法哲学在近代自然法历史上所占据的地位。下面将会表明，这两个方面都依赖于黑格尔对其哲学本体论基础所获得的洞见。只有当辩证方法的论证已经进展到能够在自然法领域中为解决与传统自然法概念的争论提供指南时，黑格尔才作为伦理学家“完全”出现在我们面前。

\section*{1}

自然法论文在其批判的部分限于分析近代自然法，这一点对于论文的体系价值和历史价值来说具有核心的意义。在黑格尔看来，自然法在其历史发展过程中形成了两种彼此对立的科学探讨方式，也就是17世纪的“经验主义的探讨方式”和18世纪晚期的“形式主义的探讨方式”。前者从它认为是构成了实践经验之本质的众多自然关系和伦理关系中分离出一些个别因素（诸如自我保存的冲动、社交性、不合群性等），赋予它们概念统一性的形式，然后最终把它们提升为它所追求的科学体系的基本原理；后者则从经验中分离出统一性形式（对它来说，这种统一性形式乃是实践理性的本质），将它在自为的纯粹意志概念中固定下来，并把这种意志概念与杂多的经验关系对立起来。尽管两者间也存在着一种特别的区分，也就是说，其一的原则“是经验直观与普遍物的关系和混合，而另一种方式的原则则是绝对对立和绝对的普遍性”，但在黑格尔看来，这两种体系的“成分”（直观和概念）及其探讨方式（对规定性进行分离和固定，这些规定性的对立或不完善的连接），都是一样的。\footnote{参见《论自然法的科学探讨方式》，第328页以下、332页。黑格尔的这一著作以及《信仰与知识》和《费希特与谢林哲学体系的差别》等著作的页码，请参见《全集》第1卷，柏林，1832年。[这也包含了《全集》（百年纪念版）第1卷，H. 格洛克纳出版，斯图加特，1927年。因为这个版本也同时包含了第1版的页码。]}两者的共同基础是道德法则(lex moralis sive naturalis)与经验自然的分离，“自然法的概念及其在整个伦理科学中的处境的最新规定，均依赖于”这种分离。\footnote{同上书，第359页。}

黑格尔以17世纪的“经验主义的”探讨方式为例来阐明这一点。由于在纯粹经验中，自然的每一个因素似乎与其他因素一样，都是原初的，因此在追求使实践经验服从于概念统一性的过程中，纯粹经验就在方法论上被消解掉了，否定性的自然概念、“混沌”遂上升为体系的出发点。取代肯定性的、由“纯粹”经验把握到的自然（在这种自然中，“杂多之物的地位”并未遭到动摇）的，是一种否定性的自然概念。它要么在“存在的样态下被想象”为自然状态，要么在“可能性和抽象的形式”下表现为对于人类学上发现的能力、人的自然的列举。\footnote{《论自然法的科学探讨方式》，第333页以下。}这些规定、自然状态和人的自然的否定统一性是从伦理状态的肯定统一性中抽象出来的，然而后者并不能从前者推导出来。因此，经验主义的探讨方式就必须预设其他的自然规定性（比如社交的冲动和对暴死的恐惧），作为[从自然状态向社会状态]过渡的基础，或是寻求实定的、历史的规定性（比如强者对弱者的征服）的帮助。\footnote{同上书，第337页。}

对黑格尔来说，标示了17世纪自然法特征的体系基础，即对自然的否定，在康德—费希特哲学中臻于完成。一方面，黑格尔在那里承认，伦理总体在霍布斯、斯宾诺莎、普芬道夫的体系中达到了一种即便只是“没有形式的、外在的和谐”，另一方面，黑格尔在这里意在表明，观念论自然法的形式主义如何“徒劳地想要达到一个肯定的有机体”。\footnote{参见上书，第329、337页。}使得经验主义的探讨方式回溯到自然物的个别规定性成为可能的自然和概念的不完全分离，让位于一种完全的分离、“绝对的”对立；在[康德、费希特的]观念论体系中，自然显现为[作为]杂多的无本质的抽象物，而与一的同样的无本质的抽象物（也就是实践理性）相对立。\footnote{《论自然法的科学探讨方式》，第345页。}黑格尔在《费希特与谢林哲学体系的差别》（1801年）中就已经指出了从这种开端出发对于自然法原则和“人类共同体”建构必然得出的结论。自然法论文没有重复他在《差别》中得到的这种结论，而是以之为前提。就其实质而言，这种结论构成了批判自然法的经验主义探讨方式的反面。在黑格尔看来，在康德那里，特别是在费希特那里，自然就其依赖于概念而言只具有一种意义。费希特把这种依赖性提到概念支配自然的极端，但是通过将自我的这个自身无意识的产物本身\footnote{即自然。——译者}理解为自我，这种依赖性也并无改变。在费希特的自然法中，只是为了解释自我的“被规定的自然”，即冲动的自然[本性]，自然才被作为外在的自然（“规定性”）演绎出来。\footnote{参见《费希特与谢林哲学体系的差别》，第226页以下。黑格尔批判中的自然法概念的特殊价值来自于其哲学上的出发点，也就是从谢林那里接受过来的先验直观的客观设定(Objektivsetzen)方法。“由于坚持先验直观的主观性，自我仍然是一个主观的主体客体，这最鲜明地表现在自我与自然的关系上：一方面表现在自然的演绎上，另一方面表现在由此建立起来的许多科学上。”（第226页）（中译文参见黑格尔：《费希特与谢林哲学体系的差别》，宋祖良、程志民译，商务印书馆1994年版，第55页。——译者）}黑格尔如此总结他的批判：“在自然法中，理智的行动只是把自然作为一种可塑的质料产生出来，因此这就不是自由的、观念的行动，不是理性的行动，而是知性的行动。”\footnote{出处同上，第247页以下；也见《信仰与知识》，第140页以下。}

黑格尔的自然法概念试图重建法和自然的原始关系，从而与由费希特体系中“特别激烈地”表达出来的概念所造成的对自然的肢解相对立。\footnote{参见《费希特和谢林哲学体系的差别》，第234页：“在《自然法体系》所做的对自然的演绎和陈述中，自然和理性的绝对对立[……]表现得特别激烈。”如同自然是“由自我所设定的被设定存在”（第226页），自由因而也就是一个“纯然的否定物”（第237页）、绝对无规定性，其规定只有在对其自身不断地进行限制的情况下才是可能的，——“而限制的概念构成自由的王国，在这个自由的王国里，由于生命物被撕裂为概念和物质，由此，生命的每一个真正自由的、自为无限和不受限制的，也就是美好的相互关系就被毁灭了，自然就成了被统治的东西”（第236页以下）。黑格尔一再强调，费希特并不把自然理解为活的、自立的、发挥作用的东西，理解为“物体的内在充盈和力量的表达”（第248页）。参见第230页：“因此，不仅在理论方面，而且在实践方面，自然都是一个本质上被规定的东西，是一个死物”；第232页：理性在费希特那里不过是“自我的无依赖性以及自然的绝对被规定存在的理念，自然被设定为一个要被否定的东西、绝对依赖的”；第233页：“先验的观点[……]从显现为自然的东西中抽出了主体—客观性，只为自然保留了客观性的死壳”；第235页：“自然的基本特征”是“成为绝对对立的东西”。参见《信仰与知识》，第136页以下、140、142页。}黑格尔认为自然法的任务在于，“按照其名称”建构起“真正的肯定物”，——“自然法应该建构出名副其实的伦理自然”。\footnote{《论自然法的科学探讨方式》，第397页。格洛克纳没有注意到这个地方，他错误地主张，黑格尔之所以在自然法论文中把法哲学称为自然法科学，是因为在18世纪时并没有其他的名称，不是“因为他自己支持今天我们称之为自然法的思维方式的那种立场”。参见格洛克纳：《黑格尔》第2卷，第304页。}自然法名称中所包含的肯定物，导致黑格尔走向一种绝对设定自然和伦理的方法，这一方法是黑格尔法哲学的第一个发展阶段所特有的。\footnote{参见《论自然法的科学探讨方式》，第372页：为了能够规定自然法的真实概念及其与实践科学的关系，就必须突出“无限性方面”。也就是说，要预设“肯定的东西”即“绝对伦理总体就是一族人民”。进一步参见《伦理体系》，第415、406页以下、462页（页码根据黑格尔：《政治学与法哲学文集》，G. 拉松出版，莱比锡，1913年）；《差别》，第242页。}也就是说，黑格尔在与奠定近代自然法体系之基础的否定性的自然概念发生争论时，不得不重新回到近代自然法体系以前的法和社会的自然理论上。这就是黑格尔在术语和思想上与谢林发生关联的意义之所在，我们可以在他耶拿早期的作品中看到这种关联。黑格尔从谢林的自然哲学借来的自然概念，是传统的目的论自然概念，黑格尔卓有成效地将其重新用来为自然法奠基。\footnote{参见《信仰与知识》，第141、148页。黑格尔在这里把“服从自然和神圣必然性的永恒规律”称作真正伦理的原则。在这里，自然就是“活生生的东西自身，它在规律中同时设定自身是普遍的，并在人民中成为真正客观的”（第149页）。参见《伦理体系》，第460页以下。}这种自然理论的独特之处在于，自然不是以个别规定性的形式，而是作为整体自然进入伦理体系的；它既不与一个无论怎样理解的绝对概念（自我、自由、理性）相对立，也不作为否定性的自然状态理论与法权状态的实定权力相对立。它不与前者对立，是因为“伦理自然”在自身中统一了（纯粹个别性和普遍性的）概念；\footnote{《论自然法的科学探讨方式》，第395页。}不与后者对立，是因为伦理自然包含了“作为彻底同一的自然状态和威严”。而法权状态的威严对黑格尔而言，“自身不过就是绝对的伦理自然”。\footnote{同上书，第338页。}与古典自然法传统相一致，自然规律和道德法则彼此直接地关联在一起，公民状态被重新视为自然存在的实现，而非其克服。黑格尔在影射霍布斯及其后继者的自然状态理论时说道：“然而，那种在伦理关系中必须认为是要予以抛弃的自然物，本身就不是伦理，因而在其本来的形式中也就完全不能表现伦理。”\footnote{同上。除此之外，这个在原则上为目的论的自然概念还促使把前述“伦理自然”与（理性或自然的）“有机体”的概念等同起来，并将这个概念与“机械”的概念、自然法的“知性国家”进行对比。参见《费希特与谢林哲学体系的差别》，第242页。}

黑格尔对近代自然法的批判恰恰就是从这一点出发的：思维一种自然概念，这种自然概念不是一种“要予以抛弃的东西”，而是自身就是“伦理的”，是“伦理自然”。\footnote{这一措辞在《伦理体系》和自然法论文中一再出现，而且在某种程度上，可以认为它对黑格尔早期的自然法观念具有建构性意义。参见《伦理体系》，第443、471、473、478页；《论自然法的科学探讨方式》，第338、347页以下、369、393、397页。}按照黑格尔的观点，把伦理事物思维为“自然”，也就意味着：思维其一切潜能阶次的持存，思维这些潜能阶次的实在性以及它们与必然性的统一。因为“第一位的东西不是个体的个别性，而是伦理自然的生命(Lebendigkeit)、神性，对于伦理自然的本质来说，个别的个体太贫乏了，无法在其全部实在性中把握它的自然。”\footnote{《伦理体系》，第470页以下。对于黑格尔而言，个体只能“短暂地”（也就是说在时间中）、个别地表达伦理事物，作为否定物的个体通过时间的否定而被设定起来，以至于伦理事物的统一性（“无差异”）在个体中只不过是“形式的”而不是“实在的”；由此可以得出它（指伦理事物的统一性）与自然的等同：“[……]伦理事物必须把自身把握为自然，把握为一切潜能阶次的持存，并在其生动形态中把握每一潜能阶次，它必然与必然性合一，并作为相对同一性而持存。”（第471页）}这个具有多重含义的措辞\footnote{当指“自然”。——译者}应当把伦理整体对其部分、众多个体的优先性以及构成伦理整体一切环节的独立持存都表达出来。显然，这里自然概念用于两个方面：(1)作为总体，在斯宾诺莎“神或自然”(deus sive natura)的意义上，和(2)作为“本质”，在亚里士多德的城邦“按照自然先于”个人的学说的意义上。自然法论文与此相应，说伦理“如果不是个人的灵魂，就不能在个人身上表达出来。只有当它是一个普遍物和一族人民的纯粹精神时，它才是个人的灵魂。按照自然，肯定的东西先于否定的东西。或如亚里士多德所言，人民按照自然先于个人”。\footnote{《论自然法的科学探讨方式》，第396页。}这一基本原理的双重表达方式表明，黑格尔在进行其称作“伦理自然”的自然法构建时，同时以亚里士多德和斯宾诺莎为基础。\footnote{伊尔廷着重强调青年黑格尔的法和国家学说的这一重要的、迄今尚未得到足够重视的出发点，这是他的功劳。参见伊尔廷：《黑格尔对亚里士多德政治学的分析》（Hegels Auseinandersetzung mit der aristotelischen Politik），载《哲学年鉴》（Philos. Jahrbuch）第71卷(1963-1964)，第38-58页。}黑格尔把斯宾诺莎关于有限物的本质的一般教导\footnote{参见斯宾诺莎：《伦理学》，第一部分，命题八，附释一。黑格尔在《信仰与知识》第64页引用了这段话。}——有限物是否定物，它不能独自存在——与亚里士多德相对具体的也就是政治的主张——孤立的个人不是自足的，而是“在自然上”依赖于城邦——等同起来。我们在这里感兴趣的并不是那些在本体论上引发或者激发这种等同的理由。重要的是，黑格尔诉诸于斯宾诺莎，并不是考虑到他与亚里士多德迥然不同的自然法和国家法，而是只考虑到他的那种在某种程度上可以与亚里士多德的《政治学》或柏拉图的《理想国》相提并论的形而上学学说。是的，我们可以说，黑格尔之所以在批判自然法时利用斯宾诺莎的形而上学，主要是为了能够更为纯粹地阐明“绝对伦理”的结构，而其构造要素则由他借自古典政治学。\footnote{黑格尔自己并没有非常清晰地呈现奠定这种借用的基础的理念，谢林则在《学术研究方法讲座》（1802年）中对其进行了陈述：“就其内在生命而言，国家的科学建构并不会在后来的时代找到任何相应的历史要素。”《全集》第3卷，施罗特(M. Schröter)出版，慕尼黑，1927年，第335页；也见第315页。[谢林]这些地方[的说法]很可能已经考虑到了黑格尔的重构努力，或是参考了与黑格尔的谈话。同时也不排除黑格尔的论述也受到了谢林的影响。两人进一步的一致，比较谢林，同上，第314页（关于古代和现代国家之差异的原则），以及黑格尔，《论自然法的科学探讨方式》，第383页以下。}黑格尔同时也赋予[斯宾诺莎的]无限实体“按照自然先于”其分殊\footnote{参见《伦理学》，第一部分，命题，定义三和定义五。（中译文参见〔荷兰〕斯宾诺莎：《伦理学》，贺麟译，商务印书馆1983年版，第5、3页。——译者）}的学说以一种政治-实践意义。对斯宾诺莎来说，就像被理解为对“绝对伦理”的“理智直观”的“对上帝的理智的爱”的学说一样，无限实体学说并不具有政治实践意义。黑格尔明显影射斯宾诺莎道：“按照哲学关于世界和必然性的观点，万物都在上帝之中，没有任何个别性存在。对经验意识而言，由于每一个别的行动或思维或存在唯有在整体中才具有其本质和意义，因而这种观点便完全实现了。”\footnote{《伦理体系》，第465页，以及前一句：“因此在伦理中，个体以一种永恒的方式存在，其经验的存在和行动是完全普遍的存在和行动。因为不是个体，而是居于个体中的普遍的绝对精神在行动。”}黑格尔通过把斯宾诺莎吸收到对近代自然法的批判中，就在形而上学上升华了作为这种批判之尺度的古典政治学的那些基本原理。为黑格尔的自然法概念奠定了基础的“肯定的”自然概念，这个表达伦理整体绝对优先于其部分以及这种关系之必然性的概念，只会赋予个体一种完全“否定的”意义。\footnote{当在近代道德哲学和法哲学中表达出来的伦理意识在“普遍物和特殊物的同一”（其中普遍物为根据）之间“插入任何一个其他的个体性为根据”时，黑格尔就径直把它称为“非伦理的意识”。参见《伦理体系》，第466页；以及《论自然法的科学探讨方式》，第349页。}由此，黑格尔就在其法哲学的第一个发展阶段从古典政治学的立场出发对近代自然法进行了毁灭性的批判。黑格尔对斯宾诺莎形而上学的运用，使得这种批判彻底化了，更有甚者，它首先阻碍了黑格尔对近代自然法获得一种恰当的历史性理解。

\section*{2}

黑格尔批判各个自然法原理的起点和终点是，在这些自然法原理中，“个人的存在被设至首位和最高位置”。现代“主体的绝对性”在一种“比较低级的抽象”层次上出现在启蒙运动时期的幸福论当中；在“那些被称作反社会主义的”（霍布斯、卢梭）体系中，自然法把这种绝对性提升为理论。在康德和费希特的哲学中，这种主体的绝对性达到了对自身的理解，也就是对它的“概念”的理解。\footnote{《论自然法的科学探讨方式》，第343页以下。}在17世纪自然法中，这一原则把绝对伦理的环节歪曲为两种特殊的“本质性”，即自然状态和法权状态，在这两种状态中，个体性以不同的方式固定起来：在自然状态中，作为绝对自由；在法权状态，作为臣民的绝对服从。预设的自由以服从于一种对个人而言是外在的，并且就此而言其本身也是个别的、特殊的权力为条件。然而，个体性的取消不过是一种假象；实际上，它通过“服从的同一”关系不断地再生出来。对于结合起来的众多个体来说，伦理的结合(unio civilis)仍是一个异己之物，自然法“以社会和国家的名义”表象出来的一个无形式的、外在的统一。\footnote{同上书，第337页以下。}

黑格尔把这种“空名”与“绝对的伦理理念”进行比较。在绝对的伦理理念中，个别性并没有被固定起来，处于一种相对的同一性，[也不是]作为一种服从的关系（在这种关系中，个体性是某种完全被设定起来的东西），相反，“个体性自身纯属虚无，它与绝对的伦理威严纯为一物，——唯有这种真正的、生动的、并非服从的同一才是个人的真正伦理”。\footnote{《论自然法的科学探讨方式》，第338页。}就像在所有地方一样，在这里，个体也是否定物，其持存必须在绝对理念中“流动”起来。黑格尔很少像这一时期那样如此清楚地表达出古典城邦伦理的立场。他对[个人]与伦理整体的生动同一与近代自然法的国家理论的“服从的同一”进行了鲜明的对比，并把彻底否定自然法理论的出发点、“个人的虚无”作为前一种同一的前提。他在批判康德和费希特的自然法时又重复了这一点。对黑格尔而言，诸如自由、纯粹意志、人权这些概念都不过是“更为纯粹的否定物”\footnote{同上书，第340页以下。这种批判同时也是对他自己伯尔尼和法兰克福时期草稿所持立场的批判。青年黑格尔的这种立场，与卢梭、康德和费希特一道，想要让“从内心出发自我立法的不可转让的人权”抽象地生效。参见《早期神学著作集》，图宾根，1907年，第212页以下，另见第42、53、71、139页以下、173、188页以下、365页。参见《黑格尔来往书信集》(Briefe von und an Hegel)第1卷，J. 荷夫迈斯特出版，汉堡，1952年，第23页以下。对此立场的最早修正已经出现在《论符腾堡最近的内部情况》（1798年）和《德国宪制》（1800-1802年）中，然而这两个文本在我们这一研究的考察范围之外。}，它们的长处在于把最新伦理体系原则的“自为存在和个体性”更加清楚地表达了出来。在17世纪自然法中，“服从的同一”关系有时候前后不一，有时候又成为了简单的统治关系。\footnote{参见《论自然法的科学探讨方式》，第337页。}黑格尔以费希特为例表明，这种“服从的同一”关系会成为一种强制的体系，其全面性(Universität)使得“个别性与普遍性的同一”完全不可能。\footnote{同上书，第362页以下。}尽管观念论哲学的“伟大方面”在于，把权利和义务的本质与思维主体和意求主体的本质设定为同一，但是观念论哲学并不忠实于这种同一性原则。不仅这个原则的有效性由于限于行为的主观道德而被打破，分裂也在合法性中同样被绝对地设定，而且观念论哲学认为[道德和合法性]两者具有同等价值，并且没有丝毫关系。主体与权利和义务的本质的同一的可能性（即道德性）与不同一的可能性（即合法性）相对立，自由的体系与强制的体系相对立。\footnote{参见《论自然法的科学探讨方式》，第360页以下。}

因此，尽管黑格尔把道德的思想尊为康德和费希特哲学的“伟大方面”，但是对他而言，道德只不过以更为纯粹的形式表现了个体性的否定物。因为作为单纯的可能性来经验的并且与合法性分离开来的意识和义务的同一把道德主体抛回到其个体性中。诚如黑格尔所言，构成了道德和合法性两者之本质的，以及两者在其中是一个东西——“伦理自然”——的，并不是“肯定的绝对物”，而是以前述的形式与之完全对立的否定的绝对物，或者说“绝对概念”。现在，黑格尔同时也把观念论自然法体系内部的这种对立以及道德原则中对否定物的固定与“绝对伦理”的立场进行对照。在方法论上，这是通过诉诸“理智直观”而发生的。在黑格尔看来，在理智直观中并不存在可能性和现实性、概念和存在（“自然”）的分离。因为就像古代城邦理论一样，这种历史性地在其诸环节中重建了古代城邦理论的伦理自然的直观，包含了现实东西的当下性，一个作为“活的关系和绝对当下”的“这一个”。\footnote{参见上书，第357、359页。此外，这里我们还可以注意，黑格尔对从谢林那里接受过来的合乎历史地占有古典伦理学和政治学的方法进行了多大程度的创造性发挥。}在这里，黑格尔为了对他把古典政治学和现代自然法进行对比的方法作辩护，转向了语言：“表示‘伦理’的希腊词，还有德语词，都卓越地”暗示了语言的一般预设，即在“绝对伦理的自然”中存在着一种“普遍物或伦常”，“而最近的伦理体系由于以一种自为存在和个别性为原则，就无法做到使用这些词语而不言明它们的联系；这种内在的暗示如此有力地证明，为了标明其本质，这些最近的伦理体系就不能滥用这些语词，而是采用了道德一词[……]”。\footnote{《论自然法的科学探讨方式》，第396页。}将伦理的理智直观与伦常的“这一个”的联系起来，其结果就是必然向古典自然法靠拢。

这种靠拢首先表现在黑格尔对道德学的评价上。尽管绝对伦理作为德行论直接就是个人的伦理，然而对黑格尔来说，个体乃是完全的否定物，[由此]绝对伦理就以否定的形式显现在个体中；不仅现代主体性的反思性道德，而且渗入伦理整体的联系之中，并且表现于外的个体的德行（黑格尔想要按照道德学的古典概念模式将这种个体的德行理解为对道德学所做的自然描述\footnote{参见上书，第399页。}），也是“否定性的伦理”。道德学与自然法（在一种构建“伦理自然”法意义上的政治学）的形式差别就在于此；我们不能像在康德和费希特那里那样这样理解这种差别：“仿佛它们是分离的，道德学被自然法排除掉了，而是它们的内容完全就在自然法之中。”\footnote{参见《伦理体系》，第465页；另见《论自然法的科学探讨方式》，第396页。}正如自身否定的东西的领域、作为普遍物的单纯可能性的个体的伦理归于道德学之下一样，归于自然法下面的是“真正的肯定物”的领域、它的现实性。在这一时期，黑格尔把这种现实性叫作“实在的绝对伦理”或者“伦理自然”。在它的由一族人民设定起来的法中构建这种“自然”，就是黑格尔为自然法规定的任务。如果像在近代[自然法的]科学探讨方式那里那样，当“自然法的规定由”否定的东西、道德的抽象物、“纯粹意志以及个体意志的抽象物以及诸如强制、通过普遍的自由概念对个人进行限制等抽象物的综合表达出来时”，那么自然法的概念就只能是一个否定的概念，“一种自然的非法[Naturunrecht]，因为当这些作为实在性的否定物为[自然法]奠基时，伦理的自然即已处于彻底的败坏和不幸之中了”。\footnote{《论自然法的科学探讨方式》，第397页。[引文中的“自然的非法”(Naturunrecht)，德文原文误作“自然法”。兹据黑格尔原文（见G. W. F. Hegel. Gesammelte Werke. Band 4. Felix Meiner Verlag Hamburg 1968, p. 468）改正。——译者]}

\section*{3}

对这种自然法概念以及与之相联的将伦理设定在一族“人民”（对人民的理智直观只不过是非常明显地模仿了雅典城邦的轮廓）中的方法论立场，黑格尔并未坚持多久；他在自然法论文发表之后不久就修正了自然法概念，最后在耶拿末期完全放弃了它。黑格尔学界只是偶尔注意到这种概念上的转变。\footnote{参见F. 达姆施泰特(Darmstaeter)：《作为社会权力的自然法与黑格尔的法哲学》(Das Naturrecht als soziale Macht und die Rechtsphilosophie Hegels)，载《智慧》(Sophia)第5卷(1937年)，第215页，注释13。达姆施泰特只是发现黑格尔主张过各种不同的自然状态观，但没有对这种不同进行追问。海林也注意到黑格尔“在非常不同的意义上”使用“自然”的概念，一种无助的状态（海林：《黑格尔》第2卷，第392页以下）。认为黑格尔后来在概念的使用上仍然摇摆不定（第393页），则是错误的。}这种引人注目的概念转变与黑格尔在1803-1804年与1805-1806年的耶拿讲座之间完成的思想基础上的一般转变有关。\footnote{对此，参见罗森克兰茨：《黑格尔传》，第178页以下；F. 罗森茨威格：《黑格尔与国家》第1卷，慕尼黑/柏林，1920年，第183页。}在这几年间，黑格尔似乎重新研究了费希特哲学，同时从谢林的术语和方法中摆脱了出来。这也意味着同时放弃迄今为止以亚里士多德和斯宾诺莎为准的自然法构想。就在细节上对黑格尔体系和概念形成上的这种转变的前提和结论进行考察而言，那就要一方面在与自然、“概念”和法的联系中，另一方面在它们与“个体性”以及“个体性”在体系中的价值关系中去寻求。黑格尔之所以在伦理体系草稿和自然法论文中把个体的存在规定为完全的否定物，并且沉没在整体的“伦理自然”之中，是因为他在这一时期仍然主要把否定物等同于虚无，等同于趋于纯粹的毁灭。这样，为近代自然法奠定其“社会和国家”演绎之基础的个体的存在，就在第一个体系草稿中出现在“否定物或自由或犯罪”这一标题下面。\footnote{参见《伦理体系》，第446-460页。}它的主题是：物理的毁灭、抢劫、偷盗、征服、谋杀、报仇、斗争、战争。体系的这一部分位于(1)“自然伦理”（按照关系的绝对伦理）和(3)（绝对）伦理之间，而不触及两者中的任何一个。倘若伦理作为“肯定物”直接设定起来，那就必须事先拒绝一种过渡的可能性。在后来的那些体系草稿中，这一部分获得了一种完全不同的功能。在1803-1804年讲座中，这一部分尚处于自然伦理（家庭）和人民之间，但是已经失去了其自己的体系地位。\footnote{参见《耶拿实在哲学》第1卷，J. 荷夫迈斯特出版，莱比锡，1932年，第226-232页。这一章没有自己的标题。对此，参见伊尔廷（本篇注释21）非常中肯的观察，第56页。}在第一个体系草稿和自然法论文中，由于“人民”的自然设定而被排除在外的那种向绝对伦理的过渡，已经在这里出现。最后，在1805-1806年讲座中，体系的这一部分压缩为短短几页，而且还是出现在完全不同的位置，或者说页边。\footnote{《耶拿实在哲学》第2卷，莱比锡，1931年，第210-212、219-221页（即荷夫迈斯特对第237页注释3的页边补充）。}

然而，黑格尔在自然法论文中就已经暗示了从个体的否定性存在过渡到伦理的肯定性存在的最初可能性。个体本身所是的否定物，可以由个体所设定，对其各种规定性的否定也可以被绝对地否定。这种从一切限制中抽身出来的能力，尽管也是一种否定物，但却是一种这样的否定物，通过它，个体存在成为“被绝对吸收到概念之中的个体性、否定的绝对无限性、纯粹自由”。这样理解的自由、个体的“概念”，作为否定的绝对物，是绝对物自身的环节。也就是说，包含在绝对伦理中。绝对伦理出现在“自由人等级”。自由人等级在生死斗争的活动中证明了纯粹的自由。但一个肯定物作为其条件要先于这种否定的活动，[这就是]“绝对的伦理事物，亦即属于一族人民之物[……]个人与它的同一唯有在否定物之中、通过死亡的危险才能得到确凿无疑的证明”。\footnote{参见《论自然法的科学探讨方式》，第370页以下；《伦理体系》，第465页以下。}这里表现出来的局限是黑格尔当时的立场；过渡的可能性只保留给了自由人，而非自由人等级，由于其劳动并不合乎于“概念”\footnote{参见《伦理体系》，第473页：“他们并不在无限的概念之中，通过无限的概念，这个对他们的意识而言不过是作为外在物而被设定之物，就会完全是他们绝对的、自身的、推动他们的精神，这个精神克服了他们的一切规定性。他们的伦理自然达到了这种直观，这种好处是第一等级提供给他们的。”}，则囿于个体性之中。此外，第一等级的否定性活动不能从自身中产生出任何恒常之物（“肯定物”）；相反，伦理的肯定的持存总是已经预设好了的。

黑格尔在1803-1804年讲座中还在做出这种预设。在实现从个体性到普遍性过渡的地方，讲座还几乎与《伦理体系》和自然法论文逐字对应：“个人作为人民的成员是一个伦理的存在，他的本质是普遍伦理的生动实体，他作为个别的东西，仅仅作为一种被扬弃了的东西，[是]一种观念性形式的存在者。伦理在其生动的多样性之中的存在就是人民的伦常。”\footnote{《耶拿实在哲学》第1卷，第232页。}这仍然完全是古典城邦伦理的立场，以及前述所有方法论的和体系与内容的结果。很明显，只要坚持这种立场，与近代自然法理论的关系就会一直都是否定性的。黑格尔在边注中写道：“没有和解，没有契约，没有任何默示或明示的原始契约。个人[不仅必须]放弃他的自由的一部分，而且必须完全放弃[自身]。”\footnote{同上。}尽管如此，过渡不再局限于“自由人等级”。否定性的活动不再是为了保存伦理整体而斗争的活动，而是个体为了相互承认的斗争。\footnote{在《伦理体系》，这个主题尽管没有特定的价值，但已经在发挥作用了：一方面，它出现在论述交换、财产、货币和贸易的承认领域（即个体性成为普遍）之后，另一方面它位于家庭关系之前。在家庭关系的“无差异”中，为了承认的斗争的结果（主奴关系）得到了扬弃。参见《伦理体系》，第440-443页；另见第495页以下，在这里，法权体系事后被解释为承认的形式体系（第497页）。}只有到了现在，那个在体系草稿中只占据了一个漠不相干的位置，并且在自然法论文的过渡理论中束缚在一个等级的偶然性上面的“否定物”，方才获得了那种构成黑格尔法思想最为重要的基础之一的功能：个别性与普遍性之间的中介。在这里，为了承认的斗争首先构成了家庭的“自然”伦理与人民的“绝对”伦理之间的中介，同时，人民的伦理一如既往地来自于对亚里士多德和斯宾诺莎的范畴的改写。\footnote{参见《耶拿实在哲学》第1卷，第232页：“一族人民的绝对精神是绝对的普遍要素，是将所有单个的自我意识交织于自身的以太，是绝对简单的、活生生的、唯一的实体。”}这样，个体得到承认，也就包含了它被接纳到绝对的伦理实体之中，[从而]意味着个体自身的消失\footnote{同上：“因此，这种绝对意识（在这个阶段，是为了承认的斗争的结果——笔者）是作为个体的意识的被扬弃了的存在[……]。它是普遍的、持存着的意识；它不是无实体的个体性的单纯形式，而是这些个体不再存在；它是绝对实体。”}。

尽管如此，却由此迈出了克服古代城邦理论的抽象方法论立场以及重新评价现代自然法的第一步。这一点可以从1805-1806年讲座的相关段落中很好地看出来。在这些段落中，黑格尔从体系、方法和内容上得出了这种出发点的结论。谢林的方法及其术语的含义完全消失不见了；黑格尔现在使用的范畴不再是“肯定物”和“否定物”、“伦理自然”及其法、“直观”和“概念”的对立等，而是起源于“自我”的“理智”和“意志”。1803-1804年讲座中隐现在“承认”的中介功能中的东西，现在变得完全清楚了：黑格尔在耶拿时期后半段重新对费希特进行了研究，并大约从1804年开始在他的精神哲学和自然法的讲座中回到了费希特。\footnote{对费希特“主体性形而上学”的积极继受（这种继受规定了黑格尔耶拿讲座一些章节的结构和安排）肯定在耶拿逻辑中发挥了重要作用。在出版者（H. 埃伦贝格、G. 拉松）看来，耶拿逻辑属于黑格尔最早的、时间定为1801-1802年的讲座草稿。然而，其理论基础“自我”作为个别性和普遍性的综合的观点（参见《耶拿逻辑学、形而上学和自然哲学》，G. 拉松出版，莱比锡，1923年，第161、163页以下、178页以下），在1801到1803年间撰写的论文中以及部分地在1803-1804年的讲座中都还绝对没有达到，这种情况在前面探讨过的黑格尔自然法理论第一个文本（即《论自然法的科学探讨方式》）中已经表述出来了。只有到了1805-1806年讲座才出现了转向。这一讲座的安排也与《耶拿逻辑学》的相关章节（第三章 主体性形而上学，第161页以下）明显一致。}现在，其探讨方式暴露了费希特的直接影响的承认主题，\footnote{比较费希特：《自然法权基础》（1796年），《费希特全集》第3卷，第44页以下、85页以下，与黑格尔：《耶拿实在哲学》第1卷，第226页以下；《耶拿实在哲学》第2卷，第194页以下、205页以下。对于黑格尔而言，费希特对“承认”作为（“现实的”）行动与“概念”的综合解释提供了一种结合的可能性，在这种综合中，各个个体彼此交互地“设定”起来。}被明确地引入近代自然法的提问方式中：“这就是习惯称之为自然状态的关系：各个个体之间的自由、平等的存在，自然法应该回答，按照这种关系，个体彼此间具有怎样的权利和义务，对他们[作为]按照他们的概念而且是独立的自我意识而言，他们的行为的必然性是什么。”\footnote{《耶拿实在哲学》第2卷，第205页。}黑格尔自己用众所周知来自霍布斯的论点：必须走出自然状态(exeundum e statu naturae)，进行回答。他已经在耶拿教授资格论文的第九个命题中运用过这一论点，尽管论证意图有所不同。\footnote{参见黑格尔：《最早出版作品集》(Erste Druckschriften)，G. 拉松出版，莱比锡，1928年，第404页：“自然状态是不义的，因此必须走出这种状态。”(Status naturae non est injustus, et ob eam causam ex illo exeundum.) 众所周知，霍布斯和斯宾诺莎恰恰论证的是相反的观点：应当离开自然状态，因为它是不义的。黑格尔在注释53中所提到的地方赞同这一说法。}康德和费希特也持有相同的论点，即个体“在自然上”不具有任何权利和义务，只有法权状态、“国家”才要求权利和义务具有现实性（效力）。

[在他们这里]“设定起来的”，不再是一种无论怎样理解的“自然”，而是“概念”，正是这种概念应当达到它的“法”：“彼此自由的自我意识概念被设定了起来，但也恰好只是概念而已；概念，正因为它是概念，毋宁必须实现它自身，也就是说，扬弃与其实在性相对立的、处于概念的形式中的自我。”\footnote{《耶拿实在哲学》第2卷，第205页。之前的黑格尔学界想要用在许多方面都不充分的主要与《精神现象学》发展阶段相连的公式“伦理的去浪漫化”去把握我们所描述的这种黑格尔自然法构想的转向。参见海姆(R. Haym)：《黑格尔及其时代》(Hegel und seine Zeit)，柏林，1857年，第207页以下；梅茨格(W. Metzger)：《德国观念论伦理学中的社会、法和国家》(Gesellschaft, Recht und Staat in der Ethik des deutschen Idealismus)，海德堡，1917年，第310页以下。}自身实现的概念无非就是承认的运动；这个运动没有直接过渡到在1803-1804年讲座中依然还是唯一的肯定物的伦理的“绝对实体”，而是过渡到“一般伦理”，并且是过渡到它的直接性，也就是法。而伦理观上的这种值得注意的转变意味着对黑格尔迄今为止的自然法构想基础的颠倒。承认的运动得以产生的“概念”并不是那种将字面意义上的“无”从自身中释放出来的[作为]个人的“纯粹自由”的否定物；因为概念的运动自身产生出了肯定物的要素，在其中个体性达到了一种恒常物，也就是法：“法是人格(Person)彼此对待的关系，是人格的自由存在的普遍要素或者说规定，[也就是]对它的空洞的自由的限制；我不必为我自己策划和产生这种关系或限制；相反，对象自身就是一般的法，亦即得到承认的关系的这种产生[……]被承认者通过他的存在而被人承认为直接有效，然而，这种存在也正好是从概念产生出来的。”\footnote{《耶拿实在哲学》第2卷，第206、212页。}概念的运动打破了自然、伦理和法迄今为止的（“直接的”）联系，并且与自然相对，在法中将自身固定下来。黑格尔继续展开的与自然法的自然概念的争议也不能掩盖这一事实，即这场争执现在已经取得了一种新的标准。这个标准本身归功于观念论晚期的自然法，并且恰恰与迄今为止的自然法相对立。这一标准就是普遍的法权能力，作为纯粹概念（即人的自我）的“人格”。“人格”的发现以与一切既定的自然秩序及其“规律”决裂为前提。“法包含了纯粹人格，纯粹的被承认存在。因而人格并不在自然状态，而是沉入定在，由此他是一个在其概念中的人；但在自然状态中，他并不在他的概念中，而是作为自然存在，在其定在中。问题直接地自相矛盾——我考察的是在其概念中的人，而不是在自然状态中的人。”\footnote{《耶拿实在哲学》第2卷，第205页（黑格尔的边注）。荷夫迈斯特并没有完全正确地读懂第二句；“[……]沉入定在”后面似乎应当打逗号，因此读作“借此，他是人，[他是]在他的概念中”。}

向费希特靠拢首先表现在我们这里感兴趣的讲座段落的体系安排上。在1803-1804年讲座中，承认的运动还是意识的一个“潜能阶次”，它在伦理的“绝对意识”中扬弃其自身。而在1805-1806年讲座中，意识的基础则是自身设定为“理智”和“意志”的“自我”的自发性。这种自发性就是上面所言的概念的运动，承认表达了这种运动的“肯定的”形式。黑格尔在这一章把“自我”等同于“概念”，实非偶然：“[……]自我，概念自身[就是]运动。自我通过他者而运动，将他者转化为自身运动，并且反过来把这种自身运动转变为他在，自我的对象化，以及自为存在。”\footnote{同上书，第204页（边注）。}黑格尔在耶拿《逻辑学》中已经把自我理解为个别性和普遍性的统一，并且与承认范畴联系起来。\footnote{参见《耶拿逻辑学》，第164页以下，第171页。}这里，自我不再显现为居于自然的“生动”统一之下的空洞的统一，也不显现为对立环节的单纯统一，而是显现为这些对立环节的一种通过进行设定的活动创立起来的中介，概念的自身中介。\footnote{对此，参见施瓦茨(J. Schwarz)：《黑格尔的哲学发展》(Hegels philosophische Entwicklung)，美茵河畔法兰克福，1938年，第339页以下。施瓦茨指出，要对《伦理体系》第431页、《信仰与知识》第56页以下与《耶拿实在哲学》第2卷第190页和195页注释5进行对照。}因此，黑格尔最初赋予传统“自然法”之名的积极意义，必须被抛弃。与这一出发点相应，黑格尔在耶拿《逻辑学》中已经把精神抬高到自然之上，并在自然法论文中把“物理自然”和“伦理自然”区分开来。\footnote{参见《耶拿逻辑学》，第192页以下；《自然法》，第395页。}但在涉及自然法的奠基和概念时，他却并没有利用这种区分，其原因要在片面的古典政治学和斯宾诺莎的实体形而上学的定位上面去找；由此，黑格尔就为自然法赢得了“伦理自然”的“法”的意义，并将其等同于一族“人民”的实体性伦理。\footnote{《论自然法的科学探讨方式》，第346页以下。}直到在1805-1806年讲座，在对费希特和——有很大可能——卢梭重新进行研究之后，\footnote{参见《耶拿实在哲学》第2卷，第213页以下、244页以下，这里尽管没有出现卢梭的名字，但是明显涉及《社会契约论》的表述。在之前的草稿中，“公意”的概念完全没有起作用。}黑格尔才超出了这种立场，并与一种以现代理论的“否定性的”自然概念为基础的自然法概念达成一致。他构建的自然法并不是一种无规定的“伦理自然”的法，而是承认运动的法。在这种承认运动中，个体性的持存、“法”得以产生出来。这种运动的必然性在于个体的“概念”，正如黑格尔所言，这种运动“是概念自己的运动，而不是与内容相对立的我们思维的运动。作为承认，概念自身就是运动，而这种运动正好扬弃了概念的自然状态：概念是承认；只有自然的东西存在，它并不是精神的东西”。\footnote{《耶拿实在哲学》第2卷，第206页。这个地方必须与第242页对照：“精神是众多个体的自然，是他们直接的实体，以及这一实体的运动和必然性。”}进行承认的运动——它的“概念”构成法——在个体性的方面把一种“意志”作为结果。这种意志是“认识着的”意志，或是一个“普遍物”，——卢梭的“公意”，它的直接现实性就是具有法权能力的“人格”。\footnote{《耶拿实在哲学》第2卷，第212页：“现在，这个认识着的意志是普遍物。它是得到承认的存在。它在普遍性的形式下与自身对立，它就是存在、一般现实性，而个体、主体是人格。个体的意志就是普遍的意志，普遍的意志就是个体的意志，一般的伦理，而它直接就是法。”相反的观点，参见拉伦茨(K. Larenz)：《德国观念论的法哲学和国家哲学》(Die Rechts- und Staatsphilosophie des deutschen Idealismus)，见《哲学手册》(Handbuch der Philosophie)，第四部分，慕尼黑和柏林，1934年，第153页。他认为意志的出现始于《精神现象学》。}把这两个范畴纳入伦理体系，表明黑格尔与古典政治学构想及其自然概念的决裂，而在1802-1804年间，他是把它们同现代自然法抽象地对立起来的。随着对作为伦理事物的直接性的“法”的承认以及将法纳入“概念”之中，个体性的环节，也就是现代自然法的出发点，便得到了辩护。由此，在经历了耶拿前期的动摇之后，黑格尔最终重新找回了他曾在1790年代追随过的卢梭、康德和费希特的自然法立场。\footnote{参见《书信集》第1卷，第18、23页以下；《早期神学著作集》，第42、53、62、72、139页以下、173、188页以下、365页。黑格尔于1798年研习了康德的法权论，1796年研读了费希特的自然法。参见罗森克兰茨：《黑格尔传》，第48、87页；梅茨格：《德国观念论伦理学中的社会、法和国家》，第331页以下。}黑格尔在耶拿末期获得了这样的洞见：“一切”都取决于“不仅把真实的东西理解和表述为实体，而且同样理解和表述为主体”。\footnote{参见黑格尔：《精神现象学》，J. 荷夫迈斯特出版，汉堡，1952年，第19页。（中译文引自《精神现象学》上卷，第10页。——译者）}通过“概念”（即自我）——在概念中，个别性和普遍性得到了中介——对个体的存在所做的这一辩护，就是这种洞见的结果。这尤其适用于自然法的观点；自然法的真理是（再次引用《精神现象学》序言，黑格尔在其中保证了他的哲学的新方法论视域）：只有当主体“是设定自身的运动时，或者说，只有当它是自身转化与其自己之间的中介时”，它才是“现实的”。\footnote{《精神现象学》，第20页。（中译文引自《精神现象学》上卷，第11页。译文稍有改动。——译者）}自然法不再以“伦理自然”和黑格尔最初将之等同于希腊城邦的实体性伦理的“伦理自然”的法，而是以在其运动中渗透了伦理实体的法的“概念”为对象。

\section*{4}

设定自我的概念打破了自身“自然的”直接性和实体性伦理的“自然”。黑格尔只有在洞见了这种自我设定的概念的运动之后，才作为完成了的“伦理学家”出现在我们面前。这种双重运动相应于这样一种基本立场，在这种立场中，黑格尔第一个自然法构想的困境及其克服历史性地发展出来。现在，在耶拿末期，由于黑格尔也从自然法的立场出发批判古典政治学，从古典政治学的立场出发来对现代自然法进行批判便颠倒了过来。这第一次明确地出现在1805-1806年讲座中。一方面，黑格尔把对从卢梭那里接受过来的从个体意志中构建起来的公意理论的解释追溯到亚里士多德的“整体按照自然先于部分”的思想，由此得出，对于个体意志而言，公意\footnote{即普遍意志。——译者}乃是“第一位的东西和本质”。\footnote{参见《耶拿实在哲学》第2卷，第244页以下。}另一方面，尤其是由于吸收了那种自然法理论，黑格尔又把分裂、作为“近代的更高原则”的“个体的自身绝对认识”与古典城邦伦理中的“普遍物与个体的直接统一”对立起来。\footnote{《耶拿实在哲学》第2卷，第251页。}其后，黑格尔即用这个产生了“概念”双重运动的标准去衡量[古典]政治学和自然法，并由此理解到两者的局限。在哲学史讲座中，我们又遇到同一个标准。这种一致绝非偶然。因为众所周知，黑格尔在1805-1806年第一次讲授了这门课程，并且后来也利用了耶拿手稿。在关于柏拉图的讲座中，奠定国家(Politeia)的基础的“实体性的东西、普遍的东西”与“就自然而言的法，[……]在个体上面与为了个体的法”形成了对照。\footnote{黑格尔：《哲学史讲演录》，《全集》第14卷，柏林，1832之后，第269页以下。这里“伦理自然”这一术语还在使用，但显然意思已经转变为“在其合理性中的自由意志”（第269页）。（中译文参见〔德〕黑格尔：《哲学史讲演录》第二卷，贺麟、王太庆译，商务印书馆1960年版，第243页以下。——译者）}黑格尔把这种近代自然法称为一种“对实在的实践本质、对法的琐碎抽象”\footnote{参见《全集》第14卷，第290页以下。}；而另一方面，他又强调柏拉图城邦理念的局限，其对个体性、个体意识、直到近代自然法才使之有效的“人格”\footnote{同上书，第295页。（中译文引自《哲学史讲演录》第二卷，第266页。译文略有改动。——译者）在对亚里士多德《政治学》的讲座（第400页）中，再次重复了黑格尔批判的双重方向。比较对柏拉图《理想国》的评断（第275页）与《耶拿实在哲学》第2卷，第251页：“柏拉图并没有提出一种理想，而是内在地把握到了其时代的国家。但是这种国家没落了[……]，因为它剥夺了绝对个体性的原则。”（边注）}的排斥：“与柏拉图的原则正相反对的是个人的自觉的自由意志原则，这一原则近来特别被卢梭提到很高的地位：认为个人本身的意志、个人的表现是必然的。”\footnote{参见《全集》第14卷，第295页。}

对古典政治学和现代自然法的局限的认识，同时也是以一种对近代自然法本身的特定历史理解为前提的，因为黑格尔的批判标准并不是从随便哪个作者那里取来的，而是取自作为批判的背景的卢梭。《哲学史讲演录》和《历史哲学讲演录》都表明，在黑格尔这里实际上存在着这样的理解。尽管黑格尔也认同18和19世纪广为传播的意见，即自然法的“启蒙运动”是从格劳秀斯的《战争与和平法》（1625年）开始的，但同时黑格尔也指出，支配格劳秀斯著作的仍然是斯多亚—西塞罗的自然目的论。\footnote{参见《哲学史导论》(Einleitung in die Geschichte der Philosophie)，第3版，J. 荷夫迈斯特出版，汉堡，1959年，第159页以下；《哲学史讲演录》（《全集》第15卷），第439页以下、445页以下、502页；《世界历史哲学讲演录》，拉松出版，莱比锡，1920年，第917页以下。黑格尔指出的古老自然法的“naturaliter insita”，显然指的是西塞罗：《论法律》(Legg.)第1卷，6.18。这里，自然法(lex naturalis)作为“根植于自然的、指挥应然行为并禁止相反行为的最高理性”(“ratio summa insita in natura, quae jubet ea, quae facienda sunt, prohibetque contraria”——中译文引自〔古罗马〕西塞罗：《国家篇 法律篇》，沈叔平、苏力译，商务印书馆1999年版，第158页。——译者)出现。也参见《论法律》第1卷，12.33。（中译文参见同书，第165页。——译者）}按照这种自然目的论，权利意识对人而言是“自然的”或者说与生俱来的(“natura insitum”)，因此，无论是格劳秀斯、普芬道夫和沃尔夫，还是苏格兰学派的道德哲学，都未曾与法传统决裂。而在黑格尔看来，摧毁了格劳秀斯仍然接受的“自然”与法之间的奠基关系的真正“革命性的”自然法理论来自霍布斯的著作，因为他最先“试图把维系国家统一的力量、国家权力的自然回溯到内在于我们本身之中的原则，亦即我们承认是我们自己所有的原则”。\footnote{《全集》第15卷，第442页。（中译文引自《哲学史讲演录》第四卷，第157页。译文略有改动。——译者）}在霍布斯那里，法和社会的理论设定一种自然秩序(ordo naturalis)的条件是，这种自然秩序需要规定使得离开它成为必然的那些条件。由于霍布斯，自然法的传统名称中产生了“歧义”。对黑格尔来说，这种歧义至关重要，并在解释《论公民》时注意到了这种歧义：“‘自然’这一名词具有歧义，一方面，人的自然指他的精神性、合理性而言；另一方面，他的自然状态则指人按照他的自然性而行动的状态而言。”\footnote{《全集》第15卷，第443页。（中译文引自《哲学史讲演录》第四卷，第159页。译文略有改动。——译者）}在黑格尔看来，卢梭、康德和费希特推进了这种歧义，加剧了霍布斯与传统自然法理论的决裂。卢梭把人的“精神性、合理性”理解为人的自由，自由使人从自然中走出来，被视为人与动物区分开来的绝对标志。\footnote{同上书，第527页。（中译文引自《哲学史讲演录》第四卷，第233页。——译者）黑格尔引用了《社会契约论》第1卷第四章的著名段落；参见第一章、第八章、第十一章；第3卷，第九章，注释；《论人类不平等的起源》，K. 维甘德出版，汉堡，1955年，第106页以下。此外，参见《世界历史哲学讲演录》，第920页以下。}如黑格尔所言，我们必须认为卢梭的原则：人是自由的，并且建立在“公意”基础上的国家乃是这种自由的实现，“是正确的”。当卢梭从个体意志、个体自然性的自由倾向“共同设定”公意时，他的“歧义”方才产生。由此，“自然的自由”又取消了卢梭在公意概念中所思考的“作为完全绝对物的自由”。\footnote{参见《全集》第15卷，第528页：“这些原则表达得很抽象，我们必须把它们正确地找出来；可是歧义立即就发生了。人是自由的，这当然是人的实质本性；这种本性在国家里不但没有被扬弃，事实上倒是开始被建立起来了。本性的自由，自由的秉赋并不是现实的；因为国家才是自由的实现。”（中译文摘自《哲学史讲演录》第四卷，第234页。——译者）另参见《哲学全书》，第163节，补充一（《全集》第6卷，第360页）；《法哲学》，第258节。}尽管如此，同时卢梭也使人意识到了，自由应是“人的概念”。\footnote{参见《全集》第15卷，第528页。}即对黑格尔来说，从中可以得出“向康德哲学的过渡”。康德哲学通过把自然的立法与自由的立法、经验的意志与“自由的和纯粹的意志”彻底分开，终结了自然法概念中的“歧义”。\footnote{参见《全集》第15卷，第529、552页以下、590页。}

在先验哲学的立场上，人的“概念”才达到了对其自身的理解。“自我意识的单纯统一”、“自我”乃是从一切既定的自然秩序中摆脱出来的、“不可穿透的、完全独立的自由和一切普遍的规定也就是思维规定的源泉”。\footnote{《世界历史哲学讲演录》，第922页。}“精神的东西的意识”成为了自然法理论的基础；自然法理论发现了一种“国家的思想原则[……]这种原则不再是诸如社会性冲动、财产安全的需要之类的任何一种意见的原则，也不是像君权神授这样的虔诚原则，而是与我的自我意识同一的确定性原则”。\footnote{同上书，第924页。}

要理解黑格尔《法哲学原理》的形而上学基础及其在终结自然法中独特的模糊立场，我们就必须要知道他对近代自然法的这种历史评价。这种评价在根本观点上偏离了自然法论文。就像现在自身表明的那样，自1805-1806年起就构成了黑格尔法思想的视域的“概念”的辩证运动，绝不是抽象的、空无内容的公式；相反，这种运动的内容就是意志的绝对“自由”。自从卢梭、康德和费希特以来，这种自由就被提升为一切法的原则。卢梭早就认识到，这种法原则隐含了自然法理论的标准乃是对自然的否决。他区分了两种自然法：(1)“本来意义上的自然法”(droit naturel proprement dit)，这种自然法以一种自然情感为基础（同情的自然法），并使自然状态中的人能够组成社会；(2)经推理而来的自然法(droit naturel raisonné)，这种自然法在国家建立之后，在市民状态中起支配作用。\footnote{参见《社会契约论第一版》(Première Version du Contract Social)第1卷，伏汉(C. E. Vaughan)编，曼彻斯特，1918年，第493页。}卢梭拒绝一种“理性的”即以自然秩序为范本的自然法，也就是格劳秀斯及其18世纪的后继者们都从中获得国家法上效力的那种自然法。卢梭和康德、费希特在这一点上是一致的，尽管康德、费希特还比卢梭更进了一步。康德偶尔也会指出黑格尔在霍布斯讲演中提到的这种自然概念的“歧义”。康德在其“道德哲学反思录”中写道：“没有市民秩序，全部自然法就不过纯然是一种德行论，并且只是作为一种可能的外在强制法律的计划，而具有一种法的名称[……]因为‘自然法’这个词曾如此歧义地使用，因此我们必须小心避免这种歧义。我们要区分自然法(Naturrecht)和自然性的法(natürliche Recht)。”\footnote{反思录，第7084条（约1776-1777年），康德，《全集》（科学院版）第19卷，第245页。}自然法是内在于市民秩序的理性法，而自然性的法则是建立在自然的个别规定上面的法，依照其内容要么是私法要么是公法。不同于卢梭和康德，费希特不仅把这个传统的名称劈分为两个方面，而且完全抛弃了这个名称。因此，费希特毫不怀疑“除了在一个共同体中，除了在实定法之下，根本就不可能存在人们经常运用的意义上的自然法”。\footnote{费希特，《自然法权基础》（1796年），见《全集》第3卷，第148页；参见第55、92页以下。《遗著集》(Nachgelassene Werke)第2卷，第498页以下。}

在其法哲学思想的第三个阶段，黑格尔从这种特定的现代自然法问题中引出了概念历史上必然形成的结论。他在多大程度上向霍布斯、卢梭、康德和费希特的立场靠拢，或许最明显不过地表现在，他在这一时期赋予前述自然和法的关系上的“歧义”以何种意义。他在1817年的《哲学全书》中写道：“迄今为止，自然法这一表达，对哲学法学来说已经成为了一种习惯表达。这一表达具有歧义：法是不是规定为似乎由直接的自然植入的，还是指的是，它是由事物的本性，也就是说由概念来规定自身的。第一种意义是之前习惯所指的意义，以至于同时还虚构出一种自然法在其中应当有效的自然状态。”\footnote{黑格尔：《哲学全书》，海德堡，1817年，第415节。（中译文参见〔德〕黑格尔：《哲学科学百科全书》，薛华译，上海世纪出版集团2002年版，第301页。——译者）对此，参见福莱希曼(E. Fleischmann)：《黑格尔的政治哲学》(La philosophie politique de Hegel)，巴黎，1964年，第17、113页以下。}通过法学术语“事物的[自然]本性”的中介作用，“概念”完全把“自然”的意义集中于自身。这里，概念的含义与耶拿末期是一样的：法，就其把自身及其一切规定都建立在“自由的人格性”之上而言，乃是“一种自我规定，这种自我规定毋宁是自然规定的对立面”。\footnote{《哲学全书》（1817年），第415节。（中译文参见《哲学科学百科全书》，第301页。——译者）也见黑格尔：《纽伦堡著作集》，J. 荷夫迈斯特出版，莱比锡，1938年，第156、170页以下。特别是第284页[《哲学百科全书》(Philosophische Enzyklopädie)，第181节]：“精神作为自由的、具有自我意识的本质，乃是自身等同的自我，在其绝对否定性关联中首先进行排斥的自我，个体的自由的本质或人格。”}现在，在黑格尔最终理解了近代自然法的历史意义之后，他就比耶拿时期更为明显地把概念的抽象自我运动等同于摆脱了所有“自然性的”规定的“人格”的自我决定。自由，观念论自然法所阐发的人的无限的、不为任何外物所限制的概念，乃是法的唯一原则。

黑格尔在柏林时期的法哲学讲座中反复讲述了概念（即自然）、自由和法之间的这种联系。第一次柏林法哲学讲座在1818-1819年、《法哲学原理》出版之前。其中没有收入正式出版文本中的第3节说道：“法的原则不在于自然，既不在外在自然，也不在人的主观自然，因此时人的意志是自然地被规定的，也就是说，是欲望、冲动和偏好的领域。法的领域是自由的领域，在其中，自由外在化自身并赋予自身以实存，此时尽管自然会出场，但却是作为一种没有独立性的东西[而出场的]。”\footnote{《自然法和国家法》(Natur- und Staatsrecht)，G. 霍迈尔根据黑格尔教授先生1818-1819年冬季学期演讲所作笔记，四开本德语手稿第155号，第9页。这一笔记以及下面引用的笔记藏于普鲁士文化遗产基金会柏林国家图书馆。[本段引文见《黑格尔全集》（历史考订版）第26卷第1分册，第239页。——译者]}正如法哲学所述，从抽象法的现象形式开始，直到家庭、市民社会和国家为止，自由的各种外化是人的历史—社会的（“客观的”）现实性的所有环节。毫无疑问，在所有这些环节中，处处都有“自然性”关系的身影。但问题在于，这些“自然性”关系是否从自身产生出一种支配那些形式的法律。在黑格尔看来，必然性——比如说，在与个体存在的关系中，国家的定在所具有的那种必然性——不再意味着，对于个体而言，必须在国家中生活乃是一种自然的规律。毋宁说，国家的必然性建立在自由自身所立的法律基础上。黑格尔对第3节的解说毫不含糊地指出：“这一节是由自然法之名引发的。自由与必然性的统一并不是通过自然，而是通过自由产生出来的。自然事物一成不变，自己并不能从规律中摆脱出来，以便自己制定法律。然而，精神则使自身从自然中挣脱出来，并且自己产生它的自然、它的法则本身。因此，自然不是法的生命。”与费希特的看法一样，黑格尔也认为自然法的名称“不过是习传下来的”，按照事物[本身]、法的“概念”，这一名称甚至是不正确的，因为“自然可以理解为：(1)本质、概念；(2)无意识的自然（真正的意义）”；“真正的名称”必须是“哲学法学”。\footnote{《自然法和国家法》，第9页。[本段引文见《黑格尔全集》（历史考订版）第26卷第1分册，第240页。——译者]黑格尔也在《世界历史哲学讲演录》的导论中指出了“自然”的双重含义。参见《历史中的理性》，第5版，J. 荷夫迈斯特编，汉堡，1955年，第117页：“如果‘自然’这个词描述的是一种事物的本质、概念，那么存在自然状态，自然法就是此种状态。这种法，按照其概念、按照精神的概念，应当归属于人。但这绝不能与精神在其自然性状态中之所是混淆在一起。”进一步参见《法哲学讲演录》，1824-1825年冬季学期，格里斯海姆笔记，四开本德语手稿545号，第5页：“[……]我们必须注意，‘自然’这一表述马上就具有一种歧义，一种重要的、会导致绝对错误的歧义。一方面，自然意味着自然性存在，正如我们发现自己在不同的方面是直接造成的那样，是我们存在的直接方面。与之相对并与之相区分，自然也是概念，事物的自然[本性]也就是事物的概念，就是事物的合乎理性方式的存在，这种事物完全不同于单纯自然性的东西。”[本段引文见《黑格尔全集》（历史考订版）第26卷第3分册，第1052页。——译者]}

我们已将黑格尔法思想发展的第三个阶段确定在1816年到1820年之间。这一阶段完全回到了《论自然法的科学探讨方式》一文想要通过“伦理自然”法的建构予以克服的那种自由哲学的基础上。在结束了对“哲学法学”基本概念的漫长理解过程（这一过程同时也澄清了哲学法学与传统自然法的历史关系）之后，现在对此也就无需赘言了。在黑格尔这里，自然与自由、自然规律与法权法则分离开来了。黑格尔在1822-1823年冬季学期的柏林法哲学讲座导言中说：“在自然中，有着一般说来一条规律存在的最高证明；而在法的法则中，事情之所以有效，并不是因为它存在，而是因为每个人都要求，它们应当符合一种自身的标准[……]如果我们对这两种法则[规律]的这种差别进行考察并追问，法的法则所归属的基地是什么，那么我们就会看到：法只产生于精神，因为自然没有法。”\footnote{《法哲学：按照黑格尔教授先生的讲演》(Philosophie des Rechts. Nach dem Vortrage des H. Prof. Hegel)，1822-1823冬季，霍托笔记，手写笔记第二册，第2页。在甘斯版的《法哲学》中，霍托所记黑格尔法哲学讲演导言部分的部分笔记作为序言的附释刊印（《全集》第8卷，第8页以下）。[本段引文见《黑格尔全集》（历史考订版）第26卷第2分册，第772页。——译者]}因而黑格尔决然地拒绝传统观念：就只有人可以认识和遵守自然规律而言，“自然规律”对人来说可以成为法的“范本”；人的历史性定在所唯一具有的法则，就是自由法则，不是“自然”，而是“概念”自身制定了这一法则。\footnote{参见上书，第3页：“对我们而言，事物的概念并不来自自然。”}但概念的自我立法——对于黑格尔而言，正是这种人的普遍法权能力的思想\footnote{参见上书，第5页：“这须被尊重为某种伟大之事：现在人因其为人，必须认为拥有权利，以至于他的人的存在要高于他的身份。”}，反抗现成的东西并打破了自然性的假象，在这种假象背后，“法的思想”、自由都还一直隐而不彰。卢梭、康德和费希特的观念论自然法首先把握到的这种思想将现成的法和对它的知识变为同一个东西。黑格尔在1824-1825年的讲座导言中说，尽管在以自然、自然需要和目的等为基础的“习惯的自然法”中，自由的思想并没有被有意排除掉，但“事实上，由于没有领会到这两种原则的独特性就接受了它们，从而怠慢了自由”。自由不能以自然的形式，而是必须以“概念”的形式来思维，在自由与其自身以及自由与自然之间，概念起着中介的作用。当自由以“概念”的要求出现时，这种要求就会主张：“自由仅仅指与法和伦理有关的自由，这样自然法之名就陷入了动摇之中。”\footnote{《法哲学讲演录》，格里斯海姆笔记，第9页以下。[本段引文见《黑格尔全集》（历史考订版）第26卷第3分册，第1054页。——译者]}

现在，黑格尔在这次讲座中明确地区分自然法和法哲学；讲座开篇这句话就包含了对《法哲学原理》这本书的“双重标题”的提示\footnote{参见上书，第3页。[见《黑格尔全集》（历史考订版）第26卷第3分册，第1051页。——译者]}：《自然法和国家学或法哲学原理》。我们的阐述表明，这个标题绝非偶然。它在一种非常具体的意义上保存了构成这部著作前史的那些环节。副标题“自然法和国家学纲要”涉及黑格尔法思想的出发点，以及他与古典政治学和现代自然法的争执。在这一过程中，黑格尔逐渐认识到两者的局限和克服两者的条件。这种理解的结果就是标题：法哲学，这一标题摆在整本书最前面，并且把自然法和国家学在自身中统一了起来。应当指出的是，由黑格尔引入的“哲学法”\footnote{黑格尔，《法哲学原理》，第3节。}的名称以及它所描述的事物，即自由的“概念”，在历史上产生于观念论自然法，在法哲学中，黑格尔知道他自己就是观念论自然法的继承人。

\part{政治经济学与政治哲学}
\chapter{对国民经济学的接受}

如果我们在近代思想的发展脉络中考察黑格尔与古典政治学和近代自然法的分歧，那么他似乎首先是在实践哲学领域中进一步反对柏拉图—亚里士多德传统，借由理性的批判力量，作别和消解传统的对象与原则。然而，当政治学与自然法的对比限于卢梭《社会契约论》与柏拉图《理想国》的原则上的对照、限于“个别性”和“普遍性”的环节时，真正说来，这算不上是真正具有原创性的对比。如卢梭本人即已进行了这种对比，其尖锐程度远在黑格尔之上。而且卢梭还确实从他所占有的“正确”原则推出了错误的结论：《社会契约论》以柏拉图的《理想国》和亚里士多德的《政治学》为典范，将国家思作“市民共同体”，从而将其等同于“市民社会”(société civile)，由此，其撰构也就重归于传统的外表之下。

与此相反，黑格尔同古典政治和近代自然法的分歧所产生的一个最重要的成果是，他洞见它们的原则的双重不足，由此得出了斩断两者在这个关键点上的传统联系的结论。霍布斯和康德之间的自然法理论设定“个体”意志与“普遍”意志（或者“市民社会”与“国家”）是同一的。与此不同，黑格尔在法国大革命历史经验的影响下，从这样的洞见出发：国家和市民社会的分离应当作为差别的原则得到承认，并应使之在理论上具有效力。古典政治学与近代自然法之间的对照使他认识到，对国家与市民社会的同一性的坚持只是一种假象，同时也与历史上决定两者的前提相冲突。

尽管黑格尔只在《法哲学原理》（1821年）中才论述了现代的、有别于国家的“市民社会”概念，但他在1805-1806年耶拿讲座中即已言及那种“差别”的必然性。由于黑格尔与古典政治学和近代自然法都有分歧，结果这里就为法哲学同时系统地批判柏拉图的《理想国》和卢梭的《社会契约论》做好了准备。近代国民经济学的社会模式使自然法的契约理论遭到了彻底的质疑，黑格尔对近代国民经济学的分析恰好与前述批判形成对照，并弥补了它的不足。

历史上，个体意志与普遍意志的直接同一出现在古代城邦；\footnote{参见《耶拿实在哲学》第2卷，第249页以下：“这是如此令人羡慕并将一直令人羡慕不已的希腊人的美的、幸运的自由。人民同时融入公民之中，人民同时就是个人、政府。它只与自己交往。同一个意志既是个人又是普遍物。然而，必然会出现一种更高的抽象、一种更加伟大的对立和教化，一种更深的精神。”也见罗森克兰茨：《黑格尔传》，第136页。}卢梭（以及青年黑格尔本人）在理论上、罗伯斯庇尔在现实中\footnote{参见《耶拿实在哲学》第2卷，第245页以下。}所做的恢复这种直接同一的努力不仅证明这与从革命释放出来的社会条件不合，而且也与一般现代国家生活的结构形式和要素不合。黑格尔将那种“新科学”，即18世纪产生的、尽管方式不同但与自然法理论一样以个体性原则为基础的国民经济学，也归在一般现代国家生活的结构形式和要素之下。

在当时德国观念论哲学中，黑格尔对从詹姆斯·斯图亚特到亚当·斯密和（在1821年《法哲学原理》中）大卫·李嘉图这些英国经典作家的最先进的国民经济学理论的接受是绝无仅有的。康德基本上将（作为家庭经济、农业和国家经济的）经济理论排除在实践哲学之外，因为它只包含了与社会概念无关的、按照因果性的自然概念才有可能产生结果的技术—实践规则。\footnote{《康德全集》（科学院版）第5卷，第172页。}尽管康德的《法权论》（1797年）第31节接受了亚当·斯密的货币定义（“货币是一种其转让同时是劳动的手段和尺度的物体”），但却鲜明地以个体主义的方式关联到契约上的“‘我的’和‘你的’的转移(commutatio latio sic dicto)中的法概念”，而不是关联到社会的契约建构本身。费希特对国民经济学的看法也是如此。他是重农学派的信徒，比康德还要落后。谢林对这个主题从未表现出更大的兴趣。尽管当时伟大的哲学似乎很少关注这个领域，然而黑格尔对国民经济学的接受也并非孤立事件。从历史上看，它与人们贬为“通俗哲学”的德国启蒙哲学对进步的西欧社会理论的广泛接受浪潮有关。作为新兴的市民社会的科学，国民经济学是那些[启蒙]哲学家们（伽尔韦[Garve]、施洛瑟[Schlosser]、阿普特[Abbt]、伊泽林[Iselin]）想要优先普及的对象。通过他们的介绍，黑格尔熟悉了斯图亚特、斯密、弗格森和休谟的作品。然而，黑格尔由法国大革命及其对德国的影响所产生的实际政治兴趣超过了对公众私人目的进行启蒙的想法。与其18世纪的哲学前辈不同，青年黑格尔研究国民经济学不再出于启蒙公众的需要（在普及科学的过程中，公众确信其私人目的与公共目的实现了同一）；相反，他的研究恰恰是要削弱这种需要，消除那种通过革命进程保证实现这种同一的可能性。在伯尔尼和法兰克福时期，黑格尔试图通过文章和小册子积极介入革命过程，他对现代国民经济学感兴趣的根源似乎就在这里。他认识到，法国大革命对自然法“最佳宪制”理想充满矛盾的实现，其责任并不在于革命的意识形态标准，而在于无法与自然法理想保持一致的社会条件。黑格尔通过分析这些条件，扩展了其革命经验的视野：在法国旁边，矗立着英格兰，这个工业革命的发源地。对黑格尔法哲学的基本历史地位来说，这一点具有决定性的意义。正如卢卡奇中肯地指出的那样，黑格尔不仅“在德国最深刻、最正确地洞见了法国大革命和拿破仑时期的本质，同时还是唯一认真地探讨了英国工业革命问题的德国思想家”。\footnote{G. 卢卡奇：《青年黑格尔和资本主义社会问题》(Der junge Hegel und die Probleme der kapitalistischen Gesellschaft)，柏林，1954年，第25页。}

然而，对于黑格尔法哲学的发展来说，决定性的并不仅仅是那种背景，而且还有这种事实：他将这种由现实的政治兴趣激发起来的对国民经济学的接受从历史的角度与古典政治学和现代自然法关联起来，正因为如此，他就不得不越出它们的界限。诚然，这种关联与青年黑格尔从“应当的东西”发展到认识“存在的东西”的一般转向大致平行，然而他一开始即为这种接受划定了界限。只有在经历多次思想转变后，他才逐渐超越这些界限。本研究将对这些转变进行描述，直到他成功实现了从简单接受国民经济学到阐明现代社会和国家理论的突破为止。与广泛流传的理解——青年黑格尔似乎即已清楚洞见“市民社会的本质”\footnote{参见卡尔·罗森克兰茨：《黑格尔传》，柏林，1844年，第51页。}或者发现他自己面临“资本主义社会的问题”（G. 卢卡奇）——相反，我们坚持认为，直到耶拿时期末，黑格尔都还没有一种允许他把经济研究与政治—历史研究的材料融贯起来的理论。我们手头掌握的，首先不过是这种理论的一些片段，黑格尔处理国民经济学范畴时之所以出现很多矛盾和晦涩，不仅是因为德国社会的落后\footnote{参见G. 卢卡奇：《青年黑格尔》，第127页以下等处。}，而且还因为下面即将谈到的接受过程中一些结构上的特点。

这一过程可以追溯到黑格尔的伯尔尼时期。在其历史—政治的研究和构思的过程中，黑格尔以其特有的细致方式（通过摘录报刊、议会报告等）熟悉了英国的经济和社会状况。根据罗森克兰茨的记载，他对济贫税的辩论尤其感兴趣。\footnote{参见《黑格尔传》，第85页以下。}在对詹姆斯·斯图亚特爵士的《政治经济学原理研究》（1767年）一书德文版（1769年）的评注中，黑格尔将这里获得的对需要和劳动、劳动分工与财富、慈善和赋税（国家管理等）的实践—社会意义的洞见集中地表达了出来。可惜这一评论现已佚失。\footnote{参见上书，第86页，据罗森克兰茨的陈述，[黑格尔的这一评论]写于1799年2月19日至5月16日之间。此外参见夏穆富(P. Chamley)：《斯图亚特和黑格尔的政治经济学和哲学经济学》(Economie politique et philosophique chez Steuart et Hegel)，巴黎，1963年。}除斯图亚特外，在耶拿时期初，亚当·斯密的经济理论又对他产生了重大影响。他通过伽尔韦的翻译（1794-1796年）了解了斯密的《国富论》（1776年）。\footnote{参见《耶拿实在哲学》第1卷，其中黑格尔引用伽尔韦版（布雷斯劳，1794年）第1卷，即斯密。关于讨论黑格尔与国民经济学的分歧的基本文献为G. 卢卡奇：《青年黑格尔》，第208页以下、369页以下。另见K. 赫内(Höhne)：《黑格尔与英国》(Hegel und England)，载：《康德研究》，第36卷，1931年，第301页以下。}黑格尔此时著述的卓异之处在于，他一边研究经济学，一边重构古代城邦理论。就像他在与自然法发生争执时诉诸亚里士多德的《政治学》一样，现在他也以古典经济学原理为背景，处理英国人提出的国民经济学问题。然而，亚里士多德路向的经济学和政治学结构与黑格尔纳入其最早的体系构思当中的那些现实之间的对比之鲜明，远远超过了自然法批判的情形。

在黑格尔最早的政治体系草案《伦理体系》（1802年）中，他与亚里士多德的一致最突出地表现了出来。第一部分（“依照关系的伦理”或“自然伦理”）的关键段落相应于亚里士多德在《政治学》第一卷讨论的那种经济学或家政学(Okonomik)。与在亚里士多德那里一样，黑格尔对人的自然的、出自需要和冲动的活动与团体（劳动、工具的使用、占有、交换；夫妻、父母与子女、主人与奴隶之间的关系）的分析都汇入与自然有关的相对伦理通过支配和家庭经济所能达到的最高“总体”即“家庭”之中：“支配的差别是表面上的。丈夫是主人和管理者，不是与其他家庭成员对立的财产所有者。作为管理者，他只有自由处置的外表。劳动同样按照每个成员的天性进行分配，而产品则是共有的；正是通过这种分工，人人都有剩余，但并不是作为财产。产品无须交换，它直接地、自在自为地就是家庭共有的。”\footnote{参见黑格尔：《政治学与法哲学文集》，第444页。}对自然不同方面（土地、植物、动物）劳动的讨论，都在追踪亚里士多德的主题\footnote{参见《政治学》第1卷，第8章，1256a19—1256b7。}，对具体经济现象的分析（交换、商业、价格、货币）亦然。\footnote{参见《政治学》第1卷，第9章，1257a5—1257b30。另一方面，同样值得注意的是，对亚里士多德暗示的“自动”工具（《政治学》第1卷，第4章，1253b33—1254a1），即机器，黑格尔指出了它的现代形式；他直接从“机械的死板劳动”推演出机器：这种劳动使得“水、风、蒸汽等的运动”作为劳动力进入抽象化的劳动过程成为可能（黑格尔：《政治学与法哲学文集》，第433-434页）。}对黑格尔来说，这个“相对伦理”的所有活动和团体都是通往“人民”这个自足的绝对伦理路途上的前站，而人民则是他按照城邦的模范构建起来的共同体，在这个共同体当中，人类活动的各种可能性都得到了完全的实现。\footnote{黑格尔：《政治学与法哲学文集》，第460页以下。}

然而，黑格尔对古典家政学范畴的部分地、非常紧密的依赖并不能掩盖这一事实，即这些范畴已经再也不能为工业和社会经济解放的事实提供坚实的基础，因为他已经通过研究国民经济学把握了这种解放的意义。在亚里士多德那里，“经济”活动是与“整体”脱离开来的。也就是说，它们要么束缚于一个为同样“个别的”主人而劳动的奴隶的个别性，要么通过其他的方式被排除在作为政治共同体(Koinonia politike)的城邦以外。对黑格尔来说，它们一开始就具有一种特别的“社会的”特性，因为它们是由社会形成的、不依赖于政治的统治活动：“因为劳动以这种方式成为一种普遍的劳动，由此，由于劳动按照其内容不是指向全部需要，相反，按照它的概念，一种普遍的依赖关系便因为自然需要的满足而设定起来。”这样，与亚里士多德的前提完全矛盾，经济学就在《伦理体系》中作为市民的实践，出现在政府的名称之下，表达了公共—政治地位。与在亚里士多德那里不同，在黑格尔这里，经济学取得了这种地位。\footnote{这里，黑格尔即已言及“需要的体系”，这是他后来在《法哲学原理》中分析现代市民社会的基础。然而，我们可以看到，黑格尔尚未将此“体系”设定到“需要”的多样性之中（《法哲学原理》，第三部分，第189-208节），而是从个别“需要”的体系和总体出发的。也见《政治学与法哲学文集》，第490页。}

如果我们追问：黑格尔为什么必须要把经济—政治活动纳入伦理体系之中，那么很快就可以在他那里找到答案。其理由就是：对现代的、自身成为“体系”的政治经济学的接受。在政治经济学中，从需要和劳动产生出来的众多产品的“诸规定性”和“各种实在性”已经获得独立，创造出了一个个人之间的普遍的社会依赖的领域：“在实践领域中，这些处于它们纯粹内在的无形式和单纯性之中的实在性，或者说这些感觉，出于差异，进行自我重建，并且通过消灭诸直观、扬弃了无差异的自我感而得到了重新恢复；——身体的需要和享受重新在总体中自为地设定起来，在它们无限复杂的关系中服从于一种必然性之下，形成了身体需要、劳动以及劳动积累方面的普遍的相互依赖体系，以及——这个体系作为科学——所谓的政治经济学体系。”黑格尔的思想在这个发展阶段必须解决的基本矛盾是，他一方面试图遵照柏拉图—亚里士多德的典范，将人伦政治的（“伦理的”）行动摆到劳动的制造过程之上，同时又要依据现代政治经济学，承认经济制造活动所具有的铸造社会的功能。由此，黑格尔一方面在伦理总体下“纯粹否定地”对待“实在性体系”，让它“接受伦理总体的主宰”，\footnote{参见《政治学与法哲学文集》，第487页以下。}而另一方面，又由于这个体系在其自身得到如此深入的发展，从而“必然摧毁它与那些关系混为一团、完全没有分别开来的自由伦理”。换言之，“有意识地接受这个体系，承认它的权利，在高贵等级之外承认它是一个独特的等级，作为它的领域”。\footnote{参见《政治学与法哲学文集》，第499页。}

从黑格尔最早体系草案对政治经济学体系的这种承认，产生了进一步的矛盾。这个矛盾与上面提到的那个矛盾密切相关。随着将经济—社会领域限制在一个特定等级，对“绝对伦理”的建构也就要从历史哲学的角度进行修正。在黑格尔看来，这是近代世界与古代世界之间的对立所要求的。\footnote{谢林的耶拿时期的《学术研究方法讲座》（1802年夏季学期讲授，1803年印行）也为此提供了动因：“想要作为自由人去理解法和国家的实定科学的每一个人，他的第一个追求必然是：通过哲学和历史获得对晚近世界以及其中的公共生活必然形式的生动直观。”[《谢林全集》第5卷，第315页]也见同书第314页对近代“市民”国家与古代城邦之间区别所作的评论。黑格尔在自然法论文[《全集》第1卷，第497页以下]中的历史哲学论述即由此出发。}然而，这种《伦理体系》尚且陌生的历史哲学修正本身产生了一系列困难。这些困难既与黑格尔建构时从古典政治哲学那里借来的那些材料有关，也与其所接受的近代政治经济学成分有关。对于历史上与近代政治经济学成分相应的那个等级，黑格尔说，它“自身处于一般的占有以及对于占有的可能正义之中，同时，这个等级还构建起一个连贯的体系，并且通过将占有的关系接纳到形式的统一性之中，每一个个人，因其潜在地能够占有，便直接作为普遍物或者作为市民（在作为布尔乔亚[bourgeois]的意义上）与所有人发生关系”。\footnote{《全集》第1卷，第499页。}然而，黑格尔在历史哲学上将这个似乎被他归在近代市民阶层下面\footnote{其更加具体的表现是，他把近代才产生的贵族市民关系也纳入这种论述之中。参见《全集》第1卷，第495页。}的等级解释为古代城邦伦理衰落的产物。随着“绝对伦理的沦丧”，随着城邦的特殊形态消亡于罗马帝国的普遍性之中，古代的自由“等级”和不自由的“等级”成为平等的了，同时“随着自由消失，奴隶制也必然消失”。\footnote{同上书，第496页。}在吉本《罗马帝国衰亡史》的影响下，黑格尔认为，鲜活的城邦伦理僵化为罗马法无生命的形式，作为“私人”之人的普遍性和平等的原则获得了抽象的效力，就像在基督教中，灵魂在上帝面前获得了抽象的平等一样。从历史上看，在一个关键点上，这种说法是错误的（随着人格概念的形成，罗马法绝没有消除黑格尔所言的等级的“分别”，而是在自由人身份(status libertatis)的名义下巩固了这种分别\footnote{参见佐姆(R. Sohm)：《罗马法制度》(Institutionen des Römischen Rechts)，第4版，莱比锡，1891年，第100页以下、105页以下。}）。从其本身来看，如果这种说法不涉及黑格尔如此戏剧性地提出的“伦理悲剧的上演”（这种悲剧在何时何地上演，对这个问题争论不休，绝非偶然）的话，那也并没有大不了的。黑格尔早期著作的解释者们以在其他地方罕见的方式异口同声地对这种悲剧的“晦涩”或“神秘”提出了指控。\footnote{参见卢卡奇：《青年黑格尔》，第465页；H. 格洛克纳：《黑格尔》第2卷，第330页。}实际上，这种“晦涩”或“神秘”的真正原因在于，黑格尔此时对近代“市民社会”的概念和起源还不甚了了。对罗马法的错误解释（这种错误在《法哲学原理》中才得到纠正\footnote{参见《全集》第7卷，第40页以下。}）是他将市民的形象放到古代世界终结处的外部动因，这种投放允许黑格尔在其称赞为所有时代都是有效的埃斯库罗斯的悲剧中去说明历史上随着近代（“市民”）社会的政治和经济解放而产生的作为公民(citoyen)的市民(Bürger)与作为布尔乔亚(bourgeois)的市民之间的冲突。\footnote{参见《全集》第1卷，第500页以下。}伦理悲剧尚未表达出在这种冲突中设立起来的“市民社会”与“国家”之间的差别\footnote{这与G. 洛尔摩泽尔相反。见氏著：《主体性和物化》(Subjektivität und Verdinglichung. Theologie und Gesellschaft im Denken des jungen Hegel, Gütersloh 1961)，第86页以下。他在这一点上追随卢卡奇的主张，大致说来忽略了黑格尔的发展历史。}；相反，它表达的是“第一等级”与“第二等级”之间的差别。前者依照其指向伦理整体的活动，与其自身处于一种无对立的（“无差异的”）关系之中，后者则沉入生产的规定性之中，它的“无差异”只是相对的，受到“现有对立”的束缚。\footnote{参见《全集》第1卷，第500、505页。黑格尔的前提，即绝对伦理通过“贵族和自由人等级”得到表现，更加清楚地出现在1801—1802年第一个等级体系草案中。也见黑格尔：《政治学与法哲学文集》，第472页以下。拉松证实自然法论文已经完成了绝对伦理与第一等级之间的等同（参见《政治学与法哲学文集》，导言，第XXVII页），却又因此而将《伦理体系》的时间放到后面，并且认为其撰构更加成熟，这是不可理解的。}那种“伦理悲剧的上演”，即绝对物所做的永恒的自身游戏，应当保证这种见解：这种对立（第三等级对第一等级的相互依赖以及它们对劳动产品的依赖）就是伦理理念自身的“命运”和“必然性”；因为绝对伦理为了能够在自身深入、巩固的国民经济学体系的历史条件下得到思考，并且作为“形态”得到保持，就必须要这样建构：它通过“牺牲它自己的一部分”，“有意识地承认”一种独立的存在，它也就“从中得到净化，保存了它自己的生命”。\footnote{《全集》第1卷，第537页。}

伦理悲剧的结局，即古代城邦伦理与政治经济学体系的和解，作为该体系通过“分离”和“抵抗”而保存其自身的“命运”，再一次将我们前面讨论过的所有矛盾都保留了下来。这些矛盾构成了黑格尔法哲学和国家哲学第一阶段的特征。它们产生于黑格尔对近代政治经济学和古典政治理论行动与劳动之间关系的同时接受，并且在他引入“伦理体系”结构当中的第一等级与第三等级的独特“分离”中固定下来。显然，真正说来，黑格尔只是在突破了古典樊篱的近代国民经济学基础上重复了政治学和家政学的古典分离，而我们前面讨论过的他的第一个体系和概念构建的困难也与此相连。

其后，黑格尔的思想克服了这些矛盾。当我们从自然法论文的第一次构划来到1803-1804年和1805-1806年耶拿讲座时，就会发现，此时体系和事实似乎完全割裂。基于柏拉图和亚里士多德的划分已经被抛弃。这一时期草稿的特征不仅仅是“材料的任意堆砌”。\footnote{见罗森克兰茨：《黑格尔传》，第194页。}这种堆砌至少可以由此得到部分说明，即这里见到的不是为了出版而写的草稿，而是用于讲座的笔记。而且更加值得注意的是在黑格尔的概念和体系构造的这一阶段所达到的对国民经济学的接受范围，它完全对立于以前的时期。他这一时期阐述的那些范畴既没有受到特殊对象领域的限制，也不固定在一个“等级”上，而是渗入体系的不同部分。其中两点特别有意思：(1)黑格尔将经济学范畴和最近的自然法范畴联系起来，使得政治经济学和（此前一直等同于罗马私法的）“形式”法之间的关系以一种不同的方式出现；(2)他不再将“与需要和劳动有关的相互依赖体系”局限于一个“等级”，而是认识到其中表达了一种关于社会的概念，这个概念使得继续坚持传统的等级学说成为完全不可能的了。遗憾的是，1803-1804年讲座手稿没有对黑格尔政治哲学的这个至关重要的变化为我们提供更多的说明。手稿恰好在黑格尔比较广泛地探讨了国民经济学问题之后，准备向“财产权”和“人格”概念过渡的地方突然中断了。当然，这可能纯属偶然。然而我们只要考虑一下前面段落的处理，就很容易看出，这里，黑格尔消解了其耶拿著作矛盾的一个方面，然而却以极端的形式重复了另一方面。也就是说，一方面，他同先前一样将“绝对伦理”理解为“人民”——个人的存在同它处于纯然否定的关系中；作为“总体”，个体是“一个单纯可能的、并非自为存在的个体，在其现时存在中只是这样的，即个体时刻做好了赴死的准备，舍弃自己。它不是作为个别的总体，作为家庭并且在占有和享受中；相反，这种关系对它自己来说是一种观念性的关系，并且证明为一种自我牺牲的关系”；\footnote{《耶拿实在哲学》第1卷，第231页。}另一方面，黑格尔将政治经济学体系作为绝对伦理的自然基础融入其中。\footnote{参见上书，第234页。在这种联系中，有意思的是，黑格尔在这里将经济活动和政治活动与精神的同样的“必然性”——必须在一族“人民”中工作——联系起来。当然，有别于德性行为的“伦理工作”，第一个等级的工作乃是“否定性的工作，[精神]对准不同于它自己的他者的现象，或是它的无[机]自然”（第234页）。}并且如在自然法论文当中一样，它是伦理的“无机自然”，但同时，它不再只是归于一个“等级”，而是属于“人民”。此前，等级限制了它对于整体的重要性。因此，尽管手稿关于等级学说的段落付之阙如，但是我们还是能够根据黑格尔的论述得出结论。黑格尔在诸如劳动、占有、价值、货币等经济学范畴上面发现的“普遍性”方面本身就是“人民”的定在环节。在这种存在中，这些范畴“直接成为有别于它们自身在其概念中所是的东西”。于是，满足个人需要的劳动——其直接运动形式为个别形式——就转化为一个普遍物：劳动的规则，或是为了所有人的需要的劳动。\footnote{关于劳动的普遍化，参见《耶拿实在哲学》第1卷，第236页：“对劳动本身来说，现在要求同样出现了：[它]要得到承认，具有普遍性形式，正是一种普遍的方式、所有劳动的规则，构成了某种自为存在的东西，表现为一个外在之物、无机自然，并且人们必须要学会它。对劳动来说，这个普遍性是真正的本质。”关于需要的普遍性，参见同书，第237页以下。}在黑格尔看来，这种过程的经典例子就是亚当·斯密所言的劳动分工，制造业中劳动过程的分化。这种分化，作为“劳动的个别化”，无限地增加了生产商品的数量，由此也就无限增加了“个别劳动与完全无限的需要量”之间的相互依赖。现在黑格尔认识到，政治经济学体系在一族“人民”内部构成了一个“团结和相互依赖的体系”，它的“治理”需要不同于等级“区分”和“划界”的其他手段。\footnote{同上书，第239页。另外，也见《论自然法的科学探讨方式》，《全集》第1卷，第498页，这里经济体系的扩张与前述一个等级的构成联在一起。}

1805-1806年讲座对这个体系进行了阐述，它不仅否定了等级区分，而且否定了从自然状态向社会状态过渡的自然法模型。这里，社会劳动的过程取代了自然法模型。在社会劳动过程中，个人将其自然和意识外化为物，以在外化中与其自身进行中介。承认的运动和被承认的存在(Anerkanntsein)，作为普遍的社会中介形式，进入经济过程的运动之中。这份讲稿的伟大在于：一直到费希特为止，现代自然法的社会概念都是抽象的，这份讲稿则通过将现代自然法的社会概念回放到政治经济学的基础上（或者用一个公式来表达：将亚当·斯密和卢梭结合起来），将它从它的抽象中解放了出来。

对于“社会”的创构来说，不仅个人的意志和权利的关系，而且他们在需要和劳动中建立起来的同物的关系，都是构成性的：“被承认的存在是直接的现实性，并且在其要素中，首先就是作为一般自为存在的个人；他享受着、劳动着。唯有这里欲望才有了出场的权利，因为它是现实的了。也就是说，它自身具有了普遍的、精神的存在。所有人为了所有人的劳动，以及享受——所有人的享受。”\footnote{《耶拿实在哲学》，第213页。}对于社会劳动过程来说，物态化(Dingverfassung)或者“意识的自身物化”(Sich-zum-Dinge-machen)\footnote{同上书，第214页：“自我的无机自然，作为存在着的东西，同样与作为抽象的自为存在的自我相对立；自我否定地对待它，并[在]二者的统一[中]扬弃它，然而方式却是这样的，即自我像塑造对象一样塑造其自身，直观它自己的形式，以同样的方式消耗它自己……而这种加工本身就是多方面的；它是意识的自身物化。”}就像个人的劳动一样根本。但在社会中，劳动和需要总是普遍的——“抽象劳动”和“抽象需要”。这种抽象在于“自为存在”、个人在社会中实现的个别化。自为存在和社会两者都是“在其运动中的抽象物”，两者都建立在个人的自为存在着的运动以及个人的劳动之上：“他的劳动的内容超过了他的需要；他为许多人的需要而劳动，每一个人都是如此。因此，每个人都要满足许多人的需要，而他许多特殊需要的满足则来自许多人的劳动。”\footnote{《耶拿实在哲学》第1卷，第214页以下。}劳动和享受过程中的物化同时意味着，物本身进入社会的中介化过程，同时在这些物的上面，被承认的存在的环节、“普遍意志”开始出现。黑格尔在物的经济学存在方式（价值、货币、交换）和法律范畴（如契约）上证明了这一点。在“各种抽象加工”之间必然产生一种“运动”，由此那些为了抽象需要而加工的物便重新成为“具体的[物]”。同时在劳动的抽象中存在着这种运动的可能性，即物与物之间是可以相互比较的。社会劳动是“普遍性要素”，物在其中成为“普遍”和“等值”的了。被加工之物的普遍性就是“它们的等值或价值。在价值上，它们是一样的。这种价值本身作为物，就是货币。回到具体，回到占有，就是交换。在交换中，这个抽象的物表现为其之所是，即从物性(Dingheit)回到自我的变化，并且是以这样的方式：它的物性在于，成为他人的占有……这种物的等值，作为其内在物，就是它的价值，价值完全出自我的同意和他人的意见——我的积极主张，同样还有他的主张，我的意志和他的意志的统一，我的[意志]作为现实的意志、具体存在着的意志而有效；被承认的存在就是定在”。\footnote{同上书，第215页以下。黑格尔边注：“在他的抽象劳动中，他直观到他自身的普遍性、他的形式的普遍性，或者他是为了他人的存在。因此，他想要同他人共享这种设定，或者这些他人自身应当在其中得到直观：第二种运动，它发展并包含了第一种运动的诸环节。自我[是]在与另一个自我的关系中的行动（并且是得到他的承认了的），而另一个自我也关系到我的占有，然而他只有征得我的同意才能拥有我的占有，正如我只有征得他的同意才能[指向]他的占有一样。——两人作为得到承认者的平等。——价值、物的意义。物具有联系到他人的意义。”（第215页以下）}劳动和交换就是将“社会”建构为意志关系和权利关系的中介形式；正如劳动构成了个人与其自身以及作为“物”的自然之间的中介，交换便构成了个人劳动与所有人劳动之间的中介。对黑格尔来说，这些中介形式就是被承认的存在的定在（“事实”），个人在其中直观到其在物的要素中形成的普遍性。在传统自然法中，先占(Okkupation)作为传统自然法财产权的权利资格，在从自然状态向公民状态的过渡中使得社会契约成为必然。在黑格尔这里，与近代国民经济学相一致，先占不再起这种作用。不是先占，而是劳动和交换，构成财产权取得的真正资格，而“社会”的外化形式和中介形式，即被承认的“定在”，则总是构成了财产权取得的前提：“同一个普遍性构成财产权上的中介，作为认知的运动，因此就是直接拥有，它经过了承认的中介，或者说它的定在就是精神的本质。这里，据有的偶然性被扬弃了：通过劳动和交换，我在被承认的存在中拥有了一切。……这里，财产权的渊源和起源是劳动、我的活动本身的这种[普遍性]——直接的自我和被承认的存在和根据。”\footnote{《耶拿实在哲学》第1卷，第217页。参见同页对这一部分的总结。这里首先需要注意的是，黑格尔不仅将交换和契约而且将劳动理解为外化(alienatio)的形式：“我已经在交换中表达意志，将我的物设定为价值。也就是说，内在的运动、内在的活动，正如劳动是沉入存在之中的活动一样；[在两种情况下都是]同一种外化。（在劳动中，我使自己直接成为物，[成为]形式，即存在。）以同样的方式，我外化我的这个定在，使之成为我的异己的存在，并在其中保存自我……因此，这里，我就将我的得到承认的存在直观为定在，并且我的意志就是这种效力。”（第217页）}

由此，黑格尔将劳动理解为现代社会的解放形式，在现代社会，个人被“教化”为法律人格的自由。无论是在古典政治学中，还是在近代自然法中，劳动的这种“社会”性都隐而未显：前者将伦理行动置于制造之上；后者则着眼于意志关系和契约关系。[而在黑格尔这里，]契约的“社会”事实则是这种本源性的承认(Anerkanntsein)的结果，个人通过劳动和交换而获得这种承认。这种确定性，即个人之间的关系和交付通过契约而获得保证，使得建立在劳动和交换之上的社会的偶然关系成为可预测的和理性的了：“我的话必定有效，这不是出于我内心诚实，我的意向、信念等不可更改的道德理由”，而是因为意志只有作为“被承认的[意志]”或者“社会的[意志]”才有其定在：“我不仅自相矛盾，而且与我的意志要得到承认这一点相矛盾。人们不能相信我的话。也就是说，我的意志只是我的，只是单纯的意见。”\footnote{参见《耶拿实在哲学》第1卷，第219页以下。此外，参见H. 马尔库塞：《理性和革命》(Vernunft und Revolution, Neuwied, 1962)，第81页以下。}契约把个人当作自由、平等的个人对待，个体意志不是个别的意见，而是可以预见的普遍意志。诚然，在契约（黑格尔将契约称为“观念性的交换”、主张的交换，而非物的交换）中，“概念”和“定在”彼此分离；由于个体坚持其个体意志，而非普遍意志，这样就可能会违反契约。只有履约才是“定在”或者“存在着的普遍意志”。\footnote{同上书，第218页以下。}

因此在契约中，普遍意志也就“隐藏在特定的物中”\footnote{参见上书，第220页。}，这样，个人便在外化的行为中彼此结合起来。在黑格尔看来，普遍意志必须要走出这种隐藏状态；因为契约很好地规范了整个社会内部个人之间的关系，但是并没有规范社会内部的运动。这里，在对这种运动之规律的认识中，黑格尔完全超出了自然法的视域。他在“被承认的存在”这一部分（这一部分至少形式上坚持了从自然状态到公民状态过渡的理论）之后，在讨论“伦理”（政治学，这里包含了国家宪法和行政理论[宪政]）之前，进一步论述了“掌握权力的法律”，在法治的视角下对前面讨论过的材料进行了重新处理。这里很容易看到，现在，政治经济学在何种程度上更加广泛、系统地进入政治哲学当中。\footnote{1803-1804年讲座绝没有讲得这般清楚。尽管黑格尔在这里讨论“绝对伦理”时重复了语言、工具、家庭、家产等“自然阶次”，但是他并未改变“人民”伦理的体系地位。相应地，作为第一部分之结束的承认过程，并不以一个独特的“得到承认的存在”的领域为结果，而是直接过渡到对伦理的阐述。1805-1806年的“得到承认的存在”构成“现实精神”的第一阶段(甲)，它与(乙)契约、(丙)犯罪与惩罚和(丁)掌握权力的法律一道，构成了体系的第二部分，并且作为独立的领域，插入体系的第一部分（理论自我和实践自我的发展史）和第三部分（宪法，即国家理论和行政理论）之间。对此，参见F. 罗森茨威格：《黑格尔与国家》第1卷，第178页以下。事实上，他并未对（紧随耶拿早期论文之后的）1803-1804年讲座体系做出中肯的判断。}黑格尔在前面的部分将“普遍性要素”描述为经济法律范畴的“社会”环节，此时这些范畴还带有个人行为所具有的偶然性，这里，他明确地直呼其名：现在，正是社会本身在“掌握权力的法律”的标题下找到了它的位置：“普遍物……对个别劳动者来说是一种必然性。他在普遍物上面有其无意识的存在；社会是他的本性，他依赖于社会的根本的、盲目的运动，这种运动在精神上和肉体上维护他，或是毁灭他。”\footnote{《耶拿实在哲学》第1卷，第231页。}政治上得到解放、按经济规律运行的“社会”以双重方式出现：(1)以肯定自为存在的个人的偶然性和权利的形式，由此社会将自身作为整体重新产生出来，并且表明自己优越于传统政治社会的实体性伦理。这种肯定的表达就是黑格尔所言的将人格、“个人的自我包含于自身”的法律。\footnote{《耶拿实在哲学》第1卷，第225页。参见同书第226页以下对法律各环节的描述。}然而，(2)这个社会处于“自然必然性、个人的偶然性”之下，以至于当它独立自为时，就会产生许多对立，使它和个人的定在面临碎裂的危险。个人的劳动，维持其存在的可能性，“处于完全交织在一起的整体的偶然性之下。因此大量的人口注定要在工厂、制造业和矿山中从事完全机械的、不健康的、不安全的低技术水平的劳动。由于外国发明等造成的[新的生产]方式或者改良，养活了一个大的阶级的工业部门一夜之间便垮掉了，所有这些人随之陷入他们自己无能为力的贫困之中。巨贫与巨富之间的对立出现……万般无奈的贫困”。\footnote{同上书，第232页。}

尽管这种对立摧毁了“人民的绝对纽带、伦理”，但《伦理体系》仍然“直接通过等级自身的构成”来保证整体的维系。\footnote{参见黑格尔：《政治学与法哲学文集》，第492页。尽管对立产生于“需要的体系”内部，但因为黑格尔此时在其等级学说的视角下思考“需要的体系”，因此对立便出现在“营利等级”内部。盖营利等级“又重新分为许多特殊的营利等级，并将这些等级区分为不同财富和享受的等级”（第491页）。关于政府强加的外部限制，参见同书，第493页以下。与等级的建构相比，这一部分的论述相对简略，基本上只涉及税务制度。}1803-1804年讲座不再将这个对立简化为一个特殊等级，而是将其理解为那个设定在“人民”自身中的“巨大的团结和相互依赖的体系”的一个环节。但由于讲座在讲到这个体系时中断了，\footnote{参见《耶拿实在哲学》第1卷，第239页以下。}因此对此问题我们也就不得而知了。1805-1806年讲座则有所不同，它在社会的现代名义之下，将该体系作为“国家”与人民的“宪法”外在地分开。它不仅花了大量的篇幅分析社会的对立，\footnote{参见《耶拿实在哲学》第2卷，第231-233页。}而且还比之前更纯粹、更清楚地论述了这些社会对立的“必然性”的辩证法。社会整体的根本必然性渗入社会个体的意识之中，同样也关联到偶然性的基本形式。通过偶然的运动，社会整体维持其自身和个人。这一法则在肯定个人偶然性的同时，又必然再度否定之。也就是说，社会的盲目运动自己摆脱偶然性，使它自己的定在归于毁灭。对立的社会本身无法摆脱这个魔咒；为此就需要一个救急神(deus ex machina)，在法则中将必然和偶然结合起来。在黑格尔看来，这个平衡社会对立并使整个社会运动与其自身和个人的运动实现和解的法则，就是国家：“然而，这种构成个人存在的完全偶然性的必然性，同样是维持个人存在的实体。国家权力进来，并且必然关心每个领域的维持，或者进行干预，在海外寻找出路、销售渠道等……职业自由[是必需的]，干预必须尽量不要表现出来，因其属于任意的领域。”\footnote{同上书，第233页。}因此，在黑格尔从对“社会”的经济运动的论述过渡到国家（三、宪法）之前，国家便出现在“掌握权力的法律”的概念之下。尽管黑格尔明确提到了社会的名称，个人的“法律”也已经区分为家庭（“个人直接存在的持存”\footnote{参见上书，第227页：“A. 法律[作为]其直接定在的持存。个人在其中直接作为自然的整体，个人作为家庭。他作为这个自然整体，不是作为个人(Person)；他必须首先成为个人。”另见同书第228页以下。}）的法律和家庭解体后的人格与财产领域的法律，然而在其思想的这一阶段，他仅将处于家庭和国家之间的个人存在叫作“国家”。\footnote{《耶拿实在哲学》第2卷，第231页：“B. 这种个人直接定在的法律，作为法律，构成他的意志，或者它在偶然存在的消失中维持此意志本身。[通过]父母离世，它成为肯定的；它作为父母之前所是的定在出现——国家。”另见罗森茨威格：《黑格尔与国家》第1卷，第180页以下。}

由此，这里概念和得到表现的实在性仍然是分开的。但同时概念（“国家”）越来越具有了讲座大量涉及的材料的历史外貌。在最后的段落，国家是“作为财富的国家”，而在“宪法”部分，国家成为了作为“普遍力量”和“绝对权力”的国家。\footnote{参见上书，第242、244页以下。}黑格尔似乎在修改这一部分时利用了1799至1802年间所写的关于德国宪法的手稿。这份手稿的第一部分提出的“国家的概念”在许多方面同讲座的国家概念是一样的\footnote{见黑格尔：《政治学与法哲学文集》，第17页以下；《耶拿实在哲学》第2卷，第246页以下。这种一致首先来自主权的特征以及自保的规则方面，为此黑格尔在这两部作品中都引用了马基雅维利的《君主论》。——海姆在《黑格尔及其时代》第167页正确地主张，《德国宪法》中的国家概念差不多完全跟早期哲学草案中“构建的”国家概念矛盾。然而，海姆没有提到这一事实，即这种看法并不适用于1805-1806年讲座。}；无疑，黑格尔的“国家”同现代政治概念（从马基雅维利到霍布斯阐发的自保和权力行使的规则）之间的联系，与前面观察到的进一步将国民经济学纳入政治哲学之中，是联系在一起的。由此而在政治学和自然法的传统概念和体系架构中构建起来的现代社会概念，一方面使人再也不能死抱住传统等级学说不放，另一方面也使得国家概念再也不能片面定位于古代的“实体性伦理”。这样，黑格尔便在本次讲座中最终与其青年时代的城邦理想诀别。他看出了这个理想与社会契约理论之间的隐秘联系：“有人如此设想普遍意志的构成：所有公民都参加集会、商议、投票，于是多数便形成了普遍意志。”\footnote{《耶拿实在哲学》第2卷，第245页。}黑格尔指责自然法理论，因其假定，普遍意志唯有出自个体意志方能形成，而另一方面又预设，甚至在人们集会成为国家之前，普遍意志即已存在于他们心中。在黑格尔看来，这种契约模式与“国家”（这里理解的国家为“近代”国家，既包括绝对主义国家，也包括革命以后的国家）的历史和社会基础不合。讲座非常清楚地指出了这一点。普遍意志不是国家的原因，而是国家的结果。国家的起源来自外部力量：“国家是他们的自在的本质(Ansich)，也就是说，是对他们进行强制的外部力量。因此，所有国家都是由伟大人物的非凡力量缔造的……”\footnote{同上书，第246页。}由此，黑格尔转向其时代的国家，对他来说，就是被法国大革命动摇了的、其基础得到了重新改造的近代欧洲国家。革命实现了社会契约，社会契约认为，普遍意志由所有人的意志构成，它是“真正的、自由的国家的原则”。\footnote{参见上书，第245页。}然而革命并没有实现城邦理想，而是实现了其在近代国家的经济社会条件下已经有所不同的摹本。契约理论的出发点是普遍意志与个体意志之间的直接同一。从历史的现实性来看，这种直接统一已经分裂为“自身即为个体性、[作为]政府的普遍物的端项，[政府]不是一个国家的抽象物，[而是]以普遍物自身为目的的个体性，以及以个人为目的的另一端。而这两种个体性是一样的。同一个人既关心自己和家庭；他劳动、订立契约等，同时也为普遍物劳动，以普遍物为目的。按前一方面，他叫市民(bourgeois)；按后一方面，他叫公民(citoyen)。所有人都要服从作为多数意志的普遍意志，[它]通过众多个人的特定表达和投票形成”。\footnote{《耶拿实在哲学》第2卷，第249页。黑格尔边注：“小市民(Spießbürger)和帝国国民(Reichbürger)，两者形式上都是市民。”}

个体性和普遍物的统一在历史上得以“定”在的“双重方式”已经暗示了1821年《法哲学原理》的结构：市民社会理论（人作为市民[bourgeois]）与国家理论（人作为公民[citoyen]）的分离。然而，此时（1805-1806年）黑格尔自己还对这个由法国革命提出的人[homme]作为公民与人作为市民的区分问题并不是很清楚。上述引文之后，黑格尔仍然直接过渡到对古代民主的描述，从而我们也就无从知道，法国的、为民主的革命实现奠定基础的这个对立与“希腊人美的、幸福的自由”究竟有何干系。然而，讲座对一点，即现代国家优越于古代“民主”的原则，是给出了明确答复的。这里显然，正是近代国家的原则自身产生了市民和公民的对立，并在思想和现实中将这种对立推向顶点，由此而能够承受它：“因此，更高的分裂就是，每个人都完全回到自身，知道他自己本身就是本质，以至于坚持自己与定在的普遍物分离，而绝对存在。”在黑格尔看来，近代国家的原则在于，个人让普遍物获得自由，普遍物也同样让个人获得自由。通过彼此都让对方获得自由，由此两个领域都获得了独立。普遍意志不仅由个体意志构成，而且还摆脱了个体意志——“摆脱了所有人的知识”。\footnote{同上书，第249页以下。}黑格尔由此推出了世袭君主制是与近代世界相符的国家形式的结论。从法哲学和国家哲学的角度来看，这个结论并不像如下事实那样具有决定性意义：他用这种原则摧毁了从契约推出国家概念的基础，由此也就摧毁了政治（“市民”）社会概念的基础。现在，随着同“社会”的分离，国家才首次获得了一个独特的、与其现代概念相符的地位，并成为后来法哲学体系的出发点。国家在两方面获得了自由：相对于个体意志的自由，个人落入“市民”独立性的要求之中；相对于君主属性的自由（君主只是一个“组织起来的”整体的“空洞的节点”，整体则生活于制度之中，并且知道承受它在其内部形成的那些对立[“极端”]）。\footnote{参见《耶拿实在哲学》第2卷，第251页：“现在，在那些宪法依旧的国家，统治和生活已经发生改变——这是随着时间逐渐改变的……每一个领域：城市、行会在管理其特殊事务时[都是实际的]。但人民构成政府时，就是糟糕的——像非理性一样糟糕。然而整体就是中心，就是自由的精神，它摆脱了那些完全固定的极端，也摆脱了君主的个性；[他]是空洞的节点。”另见同书，第252页。类似思想已经出现在《论德国宪法》（1801-1802年）中，载黑格尔：《政治学与法哲学文集》，第27页。}
\chapter{“市民社会”概念及其历史起源问题}

如果比较一下1820年至1850年间黑格尔同时代人对（当时即已引发不少赞叹的）《法哲学原理》第三部分所阐述的市民社会的评论与当今解释者的评论，那么对这两拨人来说，这个概念的历史地位似乎不是一个问题。对于像甘斯这样的黑格尔的亲炙弟子来说（甘斯1833年在“逝者友人版”《黑格尔全集》总体框架下编辑出版了《法哲学原理》第2版），不言而喻的是，这里“与国家有关的一切，包举无遗”，因为“甚至国民经济的科学”都能“在市民社会中找到恰如其分的地位和讨论”。\footnote{黑格尔：《法哲学原理》，序言，引自《全集》第8卷(1833年)，第VII页。下文引用依据此版。}魏塞(Chr. Hermann Weisse)、卢格(Arnold Ruge)、康士但丁·勒斯勒尔(Constantin Rößler)和布鲁诺·鲍威尔(Bruno Bauer)都有同样的论断；至于问题的实质，即将问题归结到国民经济学，则从黑格尔法哲学出发的那些关于“自然法”和“国家学”的最早努力：贝塞尔(Konrad Moritz Besser)的《自然法体系》（1830年）和艾泽伦(J. Fr. Gottfried Eiselen)的《国家学体系手册》（1828年），都在甘斯的框架下对市民社会进行了描述。埃德曼(Johann Eduard Erdmann)在其《国家讲座》（1851年）中对市民社会地位的理解也差不多。在他看来，市民社会的“本质”在于“明智的计算、理性的考量、工业等”以及“信用”，并且它的这些突出特征首先体现在“地地道道的市民”即“实业家”身上。\footnote{《国家讲座》(Vorlesungen über den Staat, 1851)，第28-29页，及《学术生活与学术研究讲座》(Vorlesungen über akademisches Leben und Studium, 1858)，第97-98页。}但由此埃德曼使市民社会更加靠近中产阶级；这样，黑格尔对“市民社会”一词的使用便与17、18世纪以来欧洲市民等级的解放运动平行。

马克思以及当时的其他人，如里尔(Wilhelm Heinrich Riehl)、洛伦茨·施泰因(Lorenz Stein)、J. C. 布伦奇利(Bluntschli)和瓦格纳(Hermann Wagener)，基本上都是这样理解市民社会的概念的。如马克思在《〈政治经济学批判〉序言》中所言，黑格尔将“物质的生活关系的总和”，按照“18世纪英国人和法国人的先例，概括为‘市民社会’”。\footnote{马克思：《政治经济学批判》，序言，引自1951年柏林版，第12页。（中译文引自《马克思恩格斯选集》第二卷，人民出版社2012年版，第2页。——译者）}提到18世纪英国人和法国人，显然是要强调这个概念是一个特别现代的概念。甚至在马克思之前，就有人（如1850年之前德国最重要的国家学家之一毛伦布雷歇尔[Romeo Maurenbrecher]的做法\footnote{《当代德国国家法原理》(Grundsätze des heutigen deutschen Staatsrechts)，第2版，未修改版（1843年），第22页以下。}）为了在历史上对黑格尔的市民社会进行归类，即由此出发将其与自然法的个体主义观念联系起来。同样，在拉萨尔(Ferdinand Lassalle)看来，霍布斯的“一切人对一切人的战争”(bellum omnium contra omnes)，作为这种自然法的版本之一，不过就是“人们相较于国家，称之为进行自由竞争的市民社会领域”。滕尼斯(Ferdinand Tönnies)也并不是最后一个如是理解的人。他在《共同体与社会》（1887年）中从与现代自然法相区分的概念对立出发，在黑格尔的意义上将“市民社会”与现代的契约、交换和劳动社会联系起来。\footnote{拉萨尔：《巴斯提亚特-舒尔策·德里奇先生》(Herr Bastiat-Schulze Delitzsch, 1864)，引自《演说与著作全集》(Gesammelte Reden und Schriften, ed. Eduard Bernstein)第5卷(1919年)，第69页。滕尼斯：《共同体与社会》，第2版（1912年），第25页。（滕尼斯的著作为《共同体与社会》，正文原文作《社会与共同体》。现改正。——译者）}

在近来的解释家中，罗森茨威格在其著作《黑格尔与国家》（1920年）中追问的不是该词的历史起源，而是其传记上的发生学起源。其结果是指出了黑格尔早期对舒尔策(Sulzer)《一切科学简论》（1759年）以及黑格尔对弗格森《文明社会史论》（1767年）的知识，但这并不能令人完全信服。只有在洛维特这里，黑格尔法哲学的市民社会似乎才第一次得到全面考察，因为他将市民社会与国家的区分（按照他对他援引的《法哲学原理》第260节的具体解释方式）与古代和基督教的二元论、卢梭的《社会契约论》和柏拉图的《理想国》联系起来。与之同时，马尔库塞也能够在其《理性和革命：黑格尔与社会理论的起源》（1941年）中指出内在于黑格尔法哲学之中的社会理论及其对晚近社会学的影响。当然，这必然导致对黑格尔本人及其时代生活于其中的历史维度的论述着墨不多。也正是这种简化处理首先导致了卢卡奇的《青年黑格尔与资本主义社会问题》（1938年写成，1948年出版）一书中解释者与其对象总是不成比例，这个问题唯有成功的个别分析才能解决。当卢卡奇将“市民社会”的表达引入黑格尔的分析（尽管这一表达在他解释的著作中一次也没有出现）时，这一点就最清楚地表达了出来。最后，也能在里特尔的关于《黑格尔与法国大革命》（1957年）的研究中指出同样的问题。此著在现实问题的强有力推动下，强调黑格尔的市民社会第一次从概念上阐述了现代劳动社会。\footnote{参见罗森茨威格，《黑格尔与国家》第2卷(1920年)，第118页；卡尔·洛维特，《从黑格尔到尼采》（1941年），第二部分，第1章，第2节；卢卡奇，《青年黑格尔》，柏林1954年版，特别参见第152页以下、227页以下、280页以下、368页以下以及435页以下；约阿希姆·里特尔，《黑格尔与法国大革命》（1957年），特别是第35页以下、38页以下。此外，参见以下对市民社会的同样重要的解释：魏尔(E. Weil)，《黑格尔与国家》(Hegel et l'Etat, 1950)，第71页以下；罗西(M. Rossi)，《马克思与黑格尔的辩证法》(Marx e la dialettica hegeliana)第1卷《黑格尔与国家》(Hegel e lo Stato, 1960)，第3篇，第543页以下。}

然而，这里出现的问题是，一般说来，这种事后的解释是否抓住了黑格尔的概念的真正历史内容。当我们从后来的社会学角度出发解读黑格尔的“市民社会”时，难道不应该反问一下：它与在康德和沃尔夫那里反复遇到的这个词的以前的用法有何关联？还有：超出康德和沃尔夫以外，这个概念与亚里士多德的政治共同体(koinonia politike)和在西塞罗那里、在中世纪经院哲学和近代自然法当中的市民社会(societas civilis)有何关系吗？因为历史理解的可能性正好在于：将发问背后变动不已的现实——在我们这里是：现代社会现象——一起纳入对过去的理解之中，同时悬置过去自身的视域，并对这些视域进行修正。因此，首先在欧洲政治哲学的概念传统之下对黑格尔的“市民社会”进行打量，就是正当的了。

\section*{1}

在“市民社会”概念史中，一直都有一个值得注意的现象，这就是，这个词差不多同时（1820年左右）不仅以全新的含义出现在黑格尔的《法哲学原理》中，而且也作为政治的核心概念出现在哈勒的《国家学的复兴》中。对于公认的复辟时代理论家哈勒来说，“市民社会”是启蒙哲学的一个革命性概念。启蒙哲学从罗马共和主义的语言中接过了这个词语及其古代含义。这个古代政治概念的人为设立构成了18世纪大革命及其颠覆一切自然—社会关系的真正原因：铲除人与人之间的统治和依附关系、人们服从于国家、消灭世俗和教会的“统治”、消灭自治的团体和同业公会。因此，哈勒在其六大本复辟著作的纲领性前言中宣称，革命时代“所有错误的由来和根源”首先就是“人们将其照搬到完全不同的社会关系当中的罗马市民社会(societas civilis)的不幸观念”。\footnote{参见《国家学的复兴》(Restauration der Staatswissenschaft, 1816-1820)第1卷，前言，第XXVIII页。}哈勒不厌其烦地重复，由于将人们视为一个“市民社会”的自由、平等的成员，因此，“社会联系”也就瓦解了。在哈勒看来，只有自格劳秀斯和霍布斯以来的近代自然法才知道一个与“自然社会”不同的“所谓的市民社会”，他主要是在近代欧洲学术界的拉丁语使用中看见了这种区分的萌芽。\footnote{参见《国家学的复兴》第1卷，第89-90页：“因为拉丁语只有共和主义的言说方式和名称，或者至少在言及国家时，这是最常见的方式和名称，于是这些表达便用到了完全不同的事物和关系上面。同样，正如罗马的市民彼此构成了一个共同体、一个市民团体、一个真正的市民社会：因此，所有其他的人类结合和关系也就必须同样称为市民社会或市民联合体。”}

黑格尔也在《法哲学原理》第182节界定“国家”和“市民社会”时对近代自然法持批判立场。他指责近代自然法没有将这两个概念明确分开：“如果国家仅仅被设想为不同的个人的统一体，亦即仅仅是一个同质的统一体，则其所指就只是市民社会的规定。许多近代国家法学家都不能对国家提出任何别的见解。”针对国家和市民社会的这种混淆和相互混同，黑格尔将市民社会规定为处于国家和家庭之间的“差别”：“市民社会是处于家庭和国家之间的差别[阶段]。尽管其形成要晚于国家。因为它作为差别，以国家为前提，它要存在，前面就必须要有一个独立的国家。”\footnote{《全集》第8卷，第182节，补充，第246页以下。}这里，在市民社会与国家以及家庭与市民社会之间的差别建立起来之后，哈勒所言的“社会联系”也同样瓦解了。但与狂热的复辟理论家哈勒不同，黑格尔肯定地将这一过程规定为国家从社会、社会从国家解放出来的过程，它使国家和社会两者第一次处于它们的真正关系之中。因为他对近代自然法的批判反对显然意味着：当国家与市民社会总是混淆在一起时，国家就不是国家；同样，当市民社会是“政治”社会（即国家）时，它也不是“社会”。

然而，黑格尔关于市民社会及其所谓的与国家混淆的观点仍然停留在哈勒的视域中；两人均将这个概念的起源和历史意义局限于近代自然法，特别是18世纪自然法，而不知道其古老的欧洲传统及其真正的发源地，首先是亚里士多德的《政治学》以及中世纪和近代对亚里士多德的接受。然而，黑格尔自己都不知道，根本而言，他反对的正是这种在近代自然法中依然有效地将国家(polis或civitas)界定为与市民社会(koinonia politike或societas civilis)同一的古典政治哲学传统。因为按照古典政治哲学传统对政治规制的人类世界的理论理解，国家既不自身包含社会，也不预设社会为其存在的前提；相反，国家就是“社会”，然而却是“市民社会”“政治社会”。

“社会”(Gesellschaft)一词古老的原始含义为：结合、联合——希腊语的“koinonia”，拉丁语的“societas”或“communitas”。这样，19世纪以来的、在我们今天看来不言而喻的国家和社会的分裂以及二者的并称，在古代欧洲政治哲学中并不存在。直到18世纪中叶，国家和社会在政治上由同一个概念连在一起，被概括为一个社会整体。这个概念就是古典政治的基本概念“koinonia politike”或“societas civilis”（以前的市民社会）。在亚里士多德《政治学》开头，它将古代城邦的政体表达了出来。这个市民社会从一个自身就是“市民的”社会（也就是政治的社会）的社会实体中接过了其公共政治结构。为此，我们摘引由J. G. 施洛瑟翻译的亚里士多德最早德译版《政治学》第1卷（《政治学》第1卷，第1章，第1节，1251a1—7）译文。尽管译文显然想要在革命时代的经验之后重铸这部政治学的影响力，\footnote{施洛瑟为其翻译写了一个前言，而非导言，它以这个典型句子开始：“值此人人相信自己皆有责任议论和反驳国家形式与革命、市民权利与君主义务之时代，于我而言，将留存至今的几千年前亚里士多德的政治著述译成德语，昭之于世，诚非无益之事。”（亚里士多德：《政治学及经济学残篇》，施洛瑟译自希腊文，附注释和文本分析，第一部分[1798年]，第四页。）}但“社会”一词仍然原封不动在它以前的、与“koinonia”相一致的意义上使用：

\begin{quote}
Es ist offenbar, das ein jeder Staat aus einer Gesellschaft besteht. Eine jede Gesellschaft hat aber, wenn sie sich verbindet, die Absicht, einen gewissen Vorteil zu erreichen. Denn alle Menschen handeln bloß, um das zu erreichen, was ihnen nützlich scheint. Es ist also kein Zweifel, daß die Gesellschaften alle in dieser Absicht zusammentretten, und daß die wichtigste und vortrefflichste, nämlich der Staat, oder die bürgerliche Gesellschaft, auch auf den höchsten und vortrefflichsten Vorteil hinzielt.\footnote{《政治学及经济学残篇》第1页。同样参见施洛瑟之后不久克里斯蒂安·伽尔韦的《政治学》译本（《亚里士多德的政治学》，克里斯蒂安·伽尔韦译，G. 许勒伯恩[Fülleborn]出版、注释并撰文，布雷斯劳，1799年）。此译本同时使用了“市民社会”和“市民联合”。唯有在后来从经济上进行界定的社会概念（尤其是滕尼斯区分社会和共同体）的影响下，晚出的政治学翻译（亚里士多德：《政治学》，E. 罗尔费斯[Rolfes]翻译并撰写导言和注释，1912年；以及O. 纪贡[Gigon]翻译的亚里士多德：《政治学和雅典人的国家》，1955，第1、55页）用共同体(Gemeinschaft)翻译“koinonia”，其后便鲜有例外，而“koinonia politike”也被译为“国家共同体”。在此中立化翻译中，施洛瑟和伽尔韦仍受其约束的古代传统概念的影响力及其典范性已然失效。}（显然，每一个国家都由一个社会组成。然而，当每一社会联合起来时，都是想要达到一定的利益。因为所有人行动，都只是为了达到他们所认为的利益。因此，毫无疑问，所有社会都为此目的而集合在一起，并且最重要、最卓越的社会，即国家或者市民社会，也就指向最高的和最卓越的利益。）
\end{quote}

国家或者市民社会（$\pi\acute{o}\lambda\iota\varsigma\ \kappa\alpha\grave{\iota}\ \kappa o\iota\nu\omega\nu\acute{\iota}\alpha\ \pi o\lambda\iota\tau\iota\kappa\acute{\eta}$；civitas sive societas civilis sive res publica）——这正是古老欧洲政治哲学的经典公式。从亚里士多德到大阿尔伯特、托马斯·阿奎那和梅兰希顿，以及从博丹到霍布斯、斯宾诺莎、洛克和康德，在国家和社会的现代分离之外，这个公式一直有效。大阿尔伯特在他的《亚里士多德评注》中说：“国家、市民沟通或者政治沟通，属于人的自然，亦即那些对人来说乃是自然的事物。”(De numero eorum, quae sunt natura homini, id est, quae sunt naturalia homini, civitas est, et communicatio civilis, sive politica.)\footnote{大阿尔伯特(Albertus Magnus)：《亚里士多德政治学八卷评注》(Commentarii in Octo Libros Politicorum Aristotelis)第1卷，第1章，引自《全集》第8卷(1891年)，第6页。}博丹同样将国家(res publica)规定为市民社会：“由于国家就是市民社会，它无须行会和同业公会即能自立，但它却不能没有家庭。”(Est enim res publica civilis societas, quae sine collegiis et corporibus stare per se ipsa potest, sine familia non potest.)\footnote{《论国家》(De Republica)第3卷，第7章，引自第4版(1601年)，第511-512页。}梅兰希顿和斯宾诺莎只是在形式上偏离了“civitas”的规定，从而将“civitas”和“societas civium”等同起来：“像这样坚实地建筑在法律上和自我保存的力量上的社会就叫作国家，而在这国家法律下保护着的个人就叫作市民。”(Haec autem societas legibus et potestate sese conservandi firmata civitas appellatur, et qui ipsius iure defenduntur cives.)\footnote{《伦理学》，第四部分，命题三十七，附释二。（中译文引自《伦理学》，第200页。译文略有改动。——译者）}梅兰希顿在其《亚里士多德政治学评注》中的说法高度相似：“国家就是一个建立在以彼此利益和最大安全为宗旨的法律之上的市民社会，市民就是那些在这个社会中能够担任公职或者司法职务的人。”(Civitas est societas civium iure constituta, propter mutuam utilitatem, ac maxime propter defensionem. Cives sunt qui in eadem societate magistratus aut judicandi potestatem consequi possunt.)\footnote{梅兰希顿(Melanchthon)：《亚里士多德政治学评注》(Commentarii in Politica Aristotelis, 1530)，引自《梅兰希顿全集》(Ph. Melanchthonis Opera quae supersunt Omnia)第16卷，1850年，布雷特施奈特、宾德塞尔(Bretschneider/Bindseil)编，第435页。}国家和市民社会的传统等同至少在观念上仍然出现在霍布斯的“市民结合”(unio civilis)中，通过让市民服从于国家，市民结合终结了自然状态。对此，他说道：“这样形成的结合叫作国家或者市民社会。”(Unio autem sic facta appellatur civitas sive societas civilis.)\footnote{霍布斯：《论公民》，1642年，第5章，第9节。}最后，在这一脉络中，我们还可以提到洛克的《政府论》（1689年）和康德的《道德形而上学》（1797年）的第一部分。洛克此书第7章的标题是“论政治社会或市民社会”，而康德也在第45节将国家界定为“civitas”，并在第46节中将其视同于被解释为“societas civilis”的国民社会(staatsbürgerliche Gesellschaft)：“这样一个社会(societas civilis)（亦即一个国家）的成员，为了立法的目的集合起来，就叫国民(cives)。”康德意义上的“国民”(Staatsbürger)——这个词被他的同时代人当作一个新词，一开始还遭到拒绝——事实上就是旧的市民社会的具有法律能力和政治权利的市民，即城邦的市民(politai der polis; cives der civitas)。他们具有特别突出的国民品质，“了解市民事务和审慎”(“rerum civilium cognitio et prudentia”)（西塞罗：《演说家》第1卷，14，60），参与公共生活、立法和国家管理。因为直到18世纪，只有生活在市民状态或政治状态(status civilis sive politicus)中的那些人才是“市民”，市民社会建立在这种状态之上。它不是“社会”观察方式的主题，而仅仅是广义上的政治学的主题。沃尔夫说道：“对生活在国家或者市民状态下的人进行思考的哲学部门，叫政治学。因此，政治学就是指导市民社会或者国家中的自由行动的科学。”(Ea philosophiae pars, in qua homo consideratur tamquam vivens in Rep. seu statu civili, Politica vocatur. Est itaque Politica scientia dirigendi actiones liberas in societate civili seu Rep.)\footnote{克里斯蒂安·沃尔夫(Christian Wolff)：《理性哲学或逻辑学》(Philosophia rationalis sive Logica, 1740)，预备性讨论，第3章，第65节。}

因此，我们可以说，“市民社会”在古老的欧洲意义上是一个传统的政治概念，是政治世界的基本核心范畴，其中“国家”和“社会”尚未发生分离，反倒构成了市民—政治社会自身同质的统治结构。这种统治结构建立在家庭奴隶劳动、奴隶制度或者农奴制度和工资制度的“经济”领域之上，并与之形成鲜明对照。因为在政治学的这个古典传统中，并非共同体的所有居民都拥有市民身份；所有类型的非自由人：那些居于公共政治的市民领域之下、在家庭的私人领域中为了基本生存需要而劳动的人，以及另一方面，同样从事“经济”活动的、束缚于家庭作坊的手工艺人以及妇女，都不属于市民社会或国家(societas civilis sive res publica)，因为他们是“Oikos”（即“家庭社会”）的一部分，没有赋予他们市民身份的政治地位。亚里士多德政治学的原理，同时构成了其经济学的前提：“奴隶没有城邦”（《政治学》第3卷，第9章，1280a32）。正如康德在其时代（18世纪末）所言，这一原理也依然适用于佣人、日工以及作为佣工的手工艺人。他们生活在真正意义上的市民社会之外，市民社会参照他们得到界定，并与他们形成对照。除了从属于政治以外，这一点构成了旧市民社会的第二个基本结构要素：往下与经济、与家庭分离——市民社会(societas civilis)与家庭社会(societas domestica)的古典政治对立。在托马修斯、沃尔夫和康德那里，都还能找到这种对立。托马修斯说：“然而人的社会本身要么是市民社会，要么是家庭社会。家庭社会构成市民社会的基础，因为这里市民社会的意思不外乎许多家庭社会及其人口的联合，而处于一个普遍的统治之下。”\footnote{托马修斯(Chr. Thomasius)：《政治明智短论》(Kurtzer Entwurf der politischen Klugheit, 1725)，第204-205页。市民社会与家庭社会之间的这种对比也同样适用于康德，因此，尽管康德提出了国家革命的、着眼于市民自由的新原则，但仍然保持了市民社会与国家的古老同一，因为他仍然知道旧欧洲的“家庭法”（《道德形而上学》，第一部分，1797年，第24节以下。引自科学院版，第6卷，第276页以下），以及与之相联的“物的方式的人格权”（同书，第22节以下，第276页以下），由此导致了他的全部困难，即要规定一个“构成他自己的主人(sui juris)”的“人的等级”（见“论俗语：这在理论上可能是正确的，但在实践上是行不通的”[1793年]，第二部分，引自《全集》[科学院版]，第8卷，第295页）。}

\section*{2}

历史政治生活的细节复杂多变，从亚里士多德到康德的伟大政治著作对它们的书写也各不相同。如果我们不考虑这些，那么就得承认，就国家与社会的现代关系而言，直到黑格尔以前，古老的欧洲学术传统中，没有它的“概念”，而根据传统的“市民社会”概念，这一点也绝非偶然。这个概念的连续性，它传承的漫长历史及其对于政治学理论所具有的典范意义（在黑格尔这里，这种典范意义似已不再存在），也并不能阻止黑格尔进一步使用这一概念，与此同时又与这个传统以及古老的政治世界决裂。黑格尔完成了这个传统和这个古老的政治世界，表示与之决裂，并且提出了一个新的市民社会概念。只有在1800年左右这个伟大的欧洲传统发生断裂的大背景下才能够说明，为什么黑格尔这个欧洲形而上学中的历史思想的代表，在历史性地反对市民社会及其与国家的“混淆”时，仍然停留在哈勒的视域中。就事物本身而言，我们也可以反过来说，要对这个概念做出新的界定，就必须跟古老的欧洲政治传统一刀两断。18世纪末的革命开始了这一断裂，如在哈勒处所见，其中也存在着激烈反对这种趋势的复辟思想。唯有在这种语境下阅读《法哲学原理》中黑格尔市民社会概念的引入，并像这里一样反对从马克思和洛伦茨·施泰因出发对它所做的流行的历史性解释。\footnote{这一方面最典型的是福格尔(Paul Vogel)：《黑格尔的社会概念以及L. 施泰因、马克思、恩格斯和拉萨尔对它的历史续造》(Hegels Gesellschaftsbegriff und seine geschichtliche Fortbildung durch L. Stein, Marx, Engels und Lasselle, 1925)。参见历史学家布伦纳(Otto Brunner)对现代社会概念问题的（哲学上有关的）与之相对的阐述，载《社会史的新道路》(Neue Wege der Sozialgeschichte, 1956，尤其是第7页以下、30页以下和207页)；康策(Werner Conze)：《德国早期革命时代的国家和社会》(Staat und Gesellschaft in der frührevolutionären Epoche Deutschland)，载《历史杂志》(Historische Zeitschrift)第186卷(1958年)，第1页以下，此文扩展后收入文集《德国三月革命前的国家与社会》(Staat und Gesellschaft im deutschen Vormärz, 1962)，第207页以下，我从此文中获益良多。}换言之，首先从政治传统和前革命世界出发进行解读，才能衡量它对于19世纪所具有的真正意义和历史现实性。

从亚里士多德直到康德的政治形而上学大传统将国家称为市民社会，因为在市民的、康德依然用拉丁语对其进行解释的“cives”（市民）的法律能力中，社会生活本身就已经是政治的，而且这个如此理解人类世界的政治状态(status politicus)，同时也将处于统治—家庭的或者等级的分层中真正“经济的”和“社会的”要素涵括于自身当中；反之，黑格尔则将政治领域、国家同一——现在已经成为“市民的”(bürgerlich)——“社会”领域分开。由此“市民的”这一表达便对立于其原始含义，获得了一种首要的“社会”内容，从而不再像18世纪那样，在与“政治的”相同的意义上使用。当国家在政治上成为绝对，并由此出发保证社会拥有它自己的重心，将其作为“市民社会”释放出来的时候，这种表达便仅仅指称在这种国家中私人化了的市民的“社会”地位。当人们就像自从马克思、拉萨尔、滕尼斯以来所做的那样，在近代自然法中追溯黑格尔的[社会]概念的起源，就会将这一历史过程变成自然状态(status naturalis)和市民状态(status civilis)的自然法模式，而没有正确对待近代自然法本身以及在其中仍然发挥作用的传统。因为譬如说霍布斯的市民社会(societas civilis)（《论公民》，第1章第2段；第5章第9段）也并不是那种“bellum omnium contra omnium”（一切人对一切人的战争）、彼此之间互相契约的交换社会（如亚当·斯密和重农学派最先设想的）；相反，它依然终究是一个政治秩序概念、自然状态(status naturalis)的对立物。在德国，直到施洛策(August Ludwig Schlözer)为止\footnote{参见施勒策(August Ludwig Schlözer)：《国家法总论》(Allgemeines Stats-Recht, 1794)，其中市民社会成员被称为“宗子，因而是自由的家长”（第一部分：《元政治学》[Metapolitik]，第17节，第64、65页）。施勒策用传统公式表达了它的结构：“Societas civilis或者civitas”，以区别于作为“Societas civilis cum imperio”的“国家”（同书，第44节注释）。他在这种仅仅暗示了黑格尔的“差异”的对立中，把握到了“国家与市民社会”之间的区别。}，亚里士多德的独立家长的观念都一直保存在自然状态(status naturalis)中，这种家长在家中自由地发号施令，因此属于市民社会的成员。在施勒策这里，市民社会第一次与“国家”形成鲜明对照，这一事实乃是独立于自然法模式、基于前述国家与社会之间的历史“差异”而发生的。随着政治集中到国家（行政、宪法、军事）、“市民事务”转移到社会，便同时产生了这种历史差异。正是黑格尔命名的这种差异规定了1821年《法哲学原理》“市民社会”部分的内部结构和历史内容。这里，从其政治—法的意义中解放出来的市民概念便与同时得到解放的社会结合起来了。当旧的市民社会(societas civilis)经典表达的政治实体消解为社会功能时，“市民”(Bürger)和“社会”两者，通过18世纪末对欧洲传统的突破，随着工业革命和政治革命而获得了这种社会功能。至此，市民才作为布尔乔亚，成为政治哲学的根本问题，这是与现代社会的形成同时发生的。由于现代社会自身基本上接过了“经济”的功能，便一天天消解了旧的家族实体。以前，家庭是旧的市民社会的社会细胞，现在，则是市民社会的变化形态构成了现代国家的社会基础。

在政治概念的内在转变中，这里简单勾勒的这个历史过程更加具体地显露出来。1800年以后，“市民社会”的结构似乎在一个汲取了旧的意义成分的新的中心上面建立起来。因此，只要我们注意到自然法理论在人作为人与人作为公民之间所作的区分（北美革命和法国革命的人权与公民权将这种区分法典化了），那么黑格尔的概念与18世纪的概念之间的区别及其内在地属于19世纪也就一目了然了。按照自然法思想，作为人(Mensch)，他是“societas generis humani”（人类社会）的成员，同时是类存在(Gattungswesen)和个体，服从伦理学法则（这种法则在其普遍性中是没有规定的）。作为公民，他属于市民社会、国家及其法律，服从政治学所要求的规则。\footnote{对此，参见克里斯蒂安·沃尔夫：《理性哲学或逻辑学》（1740年），预备性讨论，第3章，第64-66节；康德：《道德形而上学》，第一部分（1797年），法权论导论；卢梭：《爱弥儿》（1762年）第3卷，见《全集》第3卷（1823年），第371页，等等。}相反，在《法哲学原理》的“市民社会”中，自然法上作为人的人，也就是类的代表，消融于他的自然需要、消融于黑格尔在第190节解释中嘲讽般地谈到的“观念的具体物、人们所谓的人”及其本身归于“需要的立场”之下的东西。\footnote{《全集》第8卷（1833年），第256页。参见接下来的限制（我们现在必须参照马克思对需要体系的普遍化以及将人固定在生产和消费之间的运动上，阅读这种限制）：“因此，只是在这里，并且真正说来也只在这里，才在这种意义上谈到人。”（第256页）}这个具体物，第190节基于人和动物的自然区分将其规定为抽象—普遍的需要者，因而现在就是市民社会的对象，即“市民”——黑格尔用法语在括号中标明为“bourgeois”。作为单纯的（即自然的）人，人就是需要者，作为需要者，他就是私人，即作为“bourgeois”的市民。人和市民不再像在18世纪中那样彼此对立；相反，在现代市民社会，市民将人包含于自身。

在历史概念的这个改变了的结构中，还隐藏了一个进一步的联系：欧洲革命与古老的古今之争之间的联系。这里，这个联系只是通过暗示得到了表达。18世纪末的一个真正历史性标志就是，在政治学中，正是通过接受古人的政治概念，开始意识到与古代、与希腊人和罗马人的古老的市民自由之间的距离。因为在欧洲政治传统的历史中，这种接受第一次无法“应用”\footnote{这个概念在这里的含义，在伽达默尔1960年出版的《真理与方法》一书中所指出的解释学的传统和本质环节的意义上使用。}于时代的历史实际：现代领土国家、“臣民”、城市居民、现代社会等。这种变化明显出现在18世纪的市民概念的问题上，也同样出现在我们即将看到的新“社会”给政治思考所带来的问题上。骤然间出现的旧概念无法应用的情况，其最典型的表征可能是卢梭的建议：将诸如“公民”(citoyen)或是“祖国”(patrie)之类的词语完全从语言中抹去。\footnote{参见卢梭，《爱弥儿》（1762年）第1卷（引自《全集》第3卷，1823年，第15页）。（中译文参见〔法〕卢梭，《爱弥儿》上卷，李平讴译，商务印书馆1978年版，第11页。——译者）}那时，德国的思想没有这样激进，人们也不知道“bourgeois”（市民）和“citoyen”（公民）之间的区别，而是满足于这样的定义：市民概念在一种“普遍的意义”和一种“特殊的意义”上得到界定。市民作为“国家的成员(civis)”，属于前者，由此，与传统模式一致，“君主自身”出现为“民族的第一市民”。\footnote{对此，参见当时影响很大的哥廷根法学家谢德曼特尔(Scheidemantel)，《德国国家法和采邑法》(Teutsches Staats-und Lehnrecht, 1782)，第一部分，第2节。同样，阿亨瓦尔(Achenwall)也说，邦君并不完全在市民社会以外，而是“要把自己当作国家最杰出的市民(civis eminens)”，见《欧洲帝国的国家宪法》(Staatsverfassung der europaischen Reiche, 1752)，第2页。两个人的真正问题是市民和臣民的关系，因此就是市民的政治概念，因为“我们发现有许多臣民不是市民，比如说外邦人或者奴隶；其他一些人既是臣民也是市民；相反，也有不是臣民的市民，如君主。”（《德国国家法和采邑法》，第一部分，第439页）。}人们所理解的“特殊的意义”下的市民，则是城市成员、聚居区居民(Ortsbürger)，因此就是有别于贵族和农民的市民等级。现在，决定性的是，由于现代政治国家几乎铲除了市民(civis)的社会基础、市民社会(societas civilis)当中被统治者的自治，因此也就将市民概念的“普遍意义”减损至最低。黑格尔在1800年左右就已经罕见地理解了这一过程；因为只有从他开始，才直接出现了公民(citoyen)与市民(bourgeois)、私人与扩展到所有“臣民”的国民(Staatsbürger)的并列。现在，关心“普遍物”的政府、国家与“以个人为目的”的另一“端”即“市民”之间的现代差异同“古代伦理”及其政治实体形成了鲜明对照：“同一个人既关心自己和家庭、劳动、订立契约等，同时他也为普遍物劳动，以普遍物为目的。按前一方面，他叫市民，按后一方面，他叫公民。”\footnote{黑格尔：《耶拿讲座》（1805-1806年），J. 荷夫迈斯特编，《耶拿实在哲学》第2卷（1931年），第249页。黑格尔还加了边注，将这些词语转译到特殊的德国语境：“小市民和帝国国民，两者形式上都是小市民。”}青年黑格尔已经在希腊城邦共和国的衰落中瞥见这一过程，并且想要在罗马法中确立市民的法律地位。毫无疑问，这里，在晚期古代的那些现象中，他尖锐地道出了他的时代问题。

在晚期柏林讲座中，黑格尔在分析亚里士多德《政治学》的最后，以史论政，比较了现代国家和唯一与之相符的现代社会现象：工厂。他在这里谈到，“我们近代国家的抽象权利把个人孤立起来，允许他如此”，继而建立起一种必然的联系，以至于“真正说来，没有人具有整体的意识，或者为了整体而活动；他为整体而工作，但是不知道他是怎样为它工作的，他只关心他个人的保护”。这里，黑格尔指出了近在眼前的现代工厂形象。在现代工厂，就像在国家和社会中一样，整体活动的生命分散开来，分成不同的部分：“这是一种分工的活动，在这种活动中，每个人只占一份；正如在一个工厂里面，没有人自己单独制造一件完整的产品，每个人都只制造产品的一部分，不懂得制造其他各部分的技能，只有少数几个人才把各部分装配成一件产品。”只有古代产生的那些自由人民才有整体的意识，并且为了整体而活动，现代人作为个人本身就是不自由的——“市民的自由就是不需要普遍的事物，就是孤立的原则。但是，市民的自由（我们没有不同的两个词语来表达bourgeois[市民]和citoyen[公民]）是一个必然的环节，那是古代国家所不知道的。”\footnote{《哲学史讲演录》第2卷，K. L. 米希勒编辑，引自《全集》第14卷（1833年），第400页。（中译文参见《哲学史讲演录》第二卷，第364-365页。——译者）}

\section*{3}

然而，黑格尔直到晚年才完全洞见这个环节的历史必然性及其积极价值；因为毫无疑问，他只有在1820年左右才通过《法哲学原理》的市民社会部分“把握到”欧洲社会和政治制度内部的这种根本变化，并且这种“把握”恰恰来自于他对这种变化所做的系统的、富有成果的论述，而不像他在伯尔尼、法兰克福和耶拿开始经济学和政治学研究之后，仅仅对其进行断言而已。只要浏览一下他耶拿时期的文章、摘录和讲座，就会立即发现，根本而言，他从英国人和法国人的著作研究中所得到的事实，以及他对现代社会制度的理论洞见，同他对古代文化遗产的创造性接受漫无关联地混杂在一起。由此，如前述，市民(bourgeois)与罗马帝国时期并置，同时与希腊“政治生活”(politeuein)的“绝对伦理”相对立，从斯图亚特、斯密和李嘉图那里摘录的社会、需要、劳动和占有的经济领域获得了“相对伦理”的名称，并且作为古代悲剧和喜剧的背景出现，而对一般劳动制度的分析要么直接在“人民”和“政府”的标题之下进行，要么彼此独立，出现于人与自然的“实践”过程的框架之下。\footnote{对此，参见《批判的哲学杂志》第2卷，第2期（1802年），第79页。《政治学与法哲学文集》，G. 拉松编，第418页以下、464页以下，首先参见1803-1804年《耶拿讲座》（《耶拿实在哲学》第1卷，J. 荷夫迈斯特编，1932年，第220页及以下、236页及以下）和1805-1806年《耶拿讲座》（《耶拿实在哲学》第2卷，1931年，第197页及以下、231页及以下）。}除了他倾慕不已的古人鲜活的文化传统以外，黑格尔也直接指出了劳动和分工、机器、货币和商品、贫与富这些现代人的铁的事实，但未能深入它们的概念之中，或者至少未能将它们系统地结合起来。在耶拿时期，黑格尔似乎拒绝将这个凌乱世界的丰富内容结合起来，建立起他后来称之为人类世界的历史领域的“客观精神”的类别。另一方面，当他首次在1817年《哲学全书》初稿中试图结合古代和现代而建立这种内容方面的类别时，概念构造的古典化风格令他总是无法对现代状况做出到位的阐述，特别是在与该著1827年第2版进行对比时，这一点就非常明显。唯一论及这一主题的海德堡《哲学全书》第433节只谈到“普遍作品”及其特殊化为“诸等级”。尽管家庭表现为“个别性等级”，但是其作品为“特殊定在的需要，而其切近的目的是特殊的主体性，但是，要达到这一目的则以其他一切人的劳动为前提，同时也会影响到这一劳动”的领域，还不是后来的“市民社会”，而是“等级”，更确切地说，是（如其模糊隐约所言）“特殊性的等级”。\footnote{《哲学全书》（1817年），第433节。（中译文引自《哲学科学全书纲要》，第309页。——译者）对此，可以参考黑格尔于《法哲学原理》出版之后在《哲学全书》第2版（1827年）对这一段所做的重要改动（第513节以下）。这里，同在其他地方一样，由于“伦理”这一术语与近代家庭、市民社会和国家的发展联系起来而历史化了（第552节）。而在耶拿时期末，在《精神现象学》（1807年）中，对古代和现代的最早的历史图式化就已经出现，因此，罗马法状态便位于阿提卡共同体之后（在“真实的精神，伦理”部分），以及国家权力和财富的二元论表现为17、18世纪的特殊现象（在“自我异化了的精神，教化”部分）。}因此，对于社会领域，我们可以说，黑格尔正是通过近代经济与社会的概念结构与希腊城邦生活的概念结构之间的巨大差距，才达到了他对他自己历史地位的认识。正是由于旧政治学领域的古代传统概念无法应用于革命世纪的社会状况，导致黑格尔1820年左右构造出“市民社会”概念，作为国家和家庭之间的差异领域。因为在1820年《法哲学原理》面世以前，他既没有用过这个词语，也没有用过它的概念上的根本含义。因此，罗森茨威格提示的他在学生时代对舒尔策的《一切科学简论》摘录和弗格森的《文明社会史》是不正确的，因为在这两种情况下，讨论的都是旧的传统概念（尽管已非其本来面目）。在罗森茨威格明确提到的地方，舒尔策考察的是旧的意义上的市民社会，用他的话来说，就是“市民国家”的制度。因此，对舒尔策来说，从“市民社会的普遍概念”只是产生了“在上者的权力和司法权、臣服、奖惩之类的特殊概念”。\footnote{舒尔策：《一切科学简论》（1759年），第189-190页。}弗格森同样在与“政治社会”相同的意义上使用这个概念（第三部分，第6节），但在这个标题下，在广泛阐述（首先是斯巴达、雅典和罗马的）政治状况之余，也对“艺术和科学”进行了阐述。

如果这就解释了黑格尔的[市民社会]概念的外在起源的话（此诚为可疑），也依然有问题未获解决：黑格尔为何只在学生时代的摘录和早期神学政治片段中使用这个概念，\footnote{参见《早期神学著作集》（1907年），第41页以下、44、191页。在这些地方“市民社会”(bürgerliche Gesellschaft)依18世纪末的用法，与“国家”同义。就此而言，霍切瓦尔(Rolf K. Hočevar)在《青年黑格尔论等级制和代议制》(Stände und Repräsentation bei jungen Hegel, München 1968；即《慕尼黑政治学研究》第8卷)第9页第23条注释以及第201页第79条注释中对我的主张所提的反驳是站不住脚的。《早期神学著作集》第191页的这个句子特别明显地表明了这个词语的传统用法：“在市民社会的本性中……个人的权利变成了国家的权利。”这里，黑格尔的原型是门德尔松(Moses Mendelssohn)：《耶路撒冷》(Jerusalem, 1783)，引自莱比锡1869年版，第132页以下、146页。}其后却从未在其著作的其他地方用过它？在他的著作中，情况是这样的：即黑格尔与大约自1780年以来在克劳斯(Christian Jakob Kraus)、根茨(Friedrich Gentz)、洪堡(Wilhelm von Humboldt)、胡弗兰德(Hufeland)等人处出现的“societas civilis sive civitas”（市民社会或国家）的古老传统表达的消退甚至消解相一致，用“国家社会”(Staatsgesellschaft)来表述之。\footnote{对此，参见我的论文“市民社会”载《近代政治-社会概念词典》(Lexikon der politisch-sozialen Begriffe der Neuzeit)第2卷（1970年），O. 布伦纳、W. 康策、R. 科泽勒克(Koselleck)出版。}这个措辞非常清楚地表明，他不再使用旧的市民社会的概念，但是也还没有使用他自己的新概念。按照《纽伦堡哲学入门》，这个所谓的国家社会分为家庭和国家，前者是“自然社会”，后者是“人们通过法律关系形成的社会”。两者之间还没有后来所言的“市民社会”，因为“家庭的自然社会扩展为普遍的国家社会，它既是一个通过自然建立的联合，也是一个通过自由意志进入的联合”。\footnote{《纽伦堡哲学入门》，引自《全集》第18卷（1840年），卡尔·罗森克兰茨编，第47页以下、199页。[本段引文原文现收入《黑格尔全集》（历史考订版）第10卷第1分册，第359页。中译文参见《黑格尔全集》第10卷，第292—293页。——译者]}

然而，对于黑格尔的思想世界来说，这些处于《耶拿讲座》和海德堡《哲学全书》之间的纽伦堡中学预备课程，意义非常有限。由于在这些课程中“市民社会”了无踪影，因此它们就只是预兆性的。只有到了1817年以后，也就是抛弃了古代化的政治概念构造以后，黑格尔在耶拿时期获得的实际认识——在古代的政治遗产以外，随着近代“社会”和“市民”(bourgeois)，出现了一种与古典世界对立的不同现象——才成为对现代国家制度内部结构原则性的历史洞见。1820年，黑格尔将他1805-1806年在“个别劳动者”领域所见到的以及（这里，如果愿意的话，也总归是“古典的”）仅仅局限于这一领域的东西——“他（劳动的个人——笔者）在普遍物中有其无意识的实存；社会是他的自然，他依赖于社会的基本的、盲目的运动。这种运动从精神上和肉体上维系他，或扬弃他”\footnote{《耶拿实在哲学》第2卷（1931年），J. 荷夫迈斯特编，第231页。}——在概念上表达为“市民社会”以及对于市民社会来说同样典型的基本必然性。市民社会出现了市民的“市民自由”、他的部分依赖于“体系”的定在，在历史哲学讲座中，黑格尔谈到了这种定在。这样看来，黑格尔将“市民的”和“社会”组合成政治哲学的一个基本概念绝非偶然。这个概念，尽管外在地看，与亚里士多德的“政治共同体”(koinonia politike)、博丹、梅兰希顿或者沃尔夫的“市民社会”(societas civilis)和康德的“市民社会”的传统相一致，但其产生却恰恰是以与这个传统的完全断裂为前提的。就此而言，我们可以完全正当地说，一般而言，在黑格尔以前，甚至在他1820年以前的著作中，都不存在现代意义上的市民社会概念，只有通过这个概念，他才可能以一种影响重大却又必然的、富有成果的方式解决前面提到的“应用”问题。由此，在政治哲学的结构内，自博丹为主权概念奠基和卢梭提出公意以来，黑格尔做出了最为重要的改变。在他之后不久，这一改变即成为现代人的真正问题。

\section*{4}

黑格尔用“市民社会”对时代意识所造成的改变，绝不亚于现代革命之后果：通过政治集中于君主国家（更确切地说，革命国家），产生了一个去政治化了的社会，并且其重心转移到了经济上面。在同一时期，随着工业革命，这个社会在“国家经济学”（更确切地说，“国民经济学”）中经历了这种转变。唯有在这一过程中，欧洲社会的“政治”制度和“市民”制度始告分离。在此前古代政治的古典世界，“政治”制度和“市民”制度的意义相同：如托马斯·阿奎那所言的“communitas civili sive politica”（市民共同体或政治共同体），洛克所言的市民社会或政治社会。黑格尔自己也完全知道市民事务与政治事务的这种古老同一，在《法哲学原理》的一个地方还谈到了这两个词语在原始意义上的统一。他在第303节言及市民等级的、非政治的社会现代观念及其与旧的政治等级差别时，说道：“尽管在这些所谓的理论看来，一般的市民社会等级与政治意义上的等级相去甚远，然而语言却仍然保持了先前本来就有的这种统一。”\footnote{参见《全集》第8卷(1833年)，第303节，第398页。马克思自然会在其对黑格尔《法哲学原理》的批判性评论中注意到这句话，因此他立即写道：“因此，我们可以推论，现在这种统一已不复存在了。”（《黑格尔国家法批判》，1843年，引自《马克思恩格斯全集》，第一部分，第1卷第1册，第487页。）}在这里，黑格尔反对国家与社会对立的自由主义革命观念，但这并不能改变这一事实，即他在概念上比反对国家的自由主义运动还要尖锐，认识到了现代革命与“市民社会”的联系，并将其表达出来。因为黑格尔开始阐述市民社会的第一段，即《法哲学原理》第182节，便从孤立的个人、从19世纪早期在政治制度上得到了解放的社会（因而成为了前述意义上的“市民社会”）的私人这种事实出发，明确表达了当时出现的旧的市民社会传统的瓦解：“具体个人作为特殊的人自身就是目的。他作为各种需要的整体以及自然必然性和任意的混合体，构成市民社会的一个原则，——但这个特殊的个人本质上是与其他的这样的特殊性相关的，从而每个人都通过他人的中介，并且同时也完全只是作为通过另一个原则（即普遍性的形式）的中介，使自己有效并得到满足。”\footnote{同上书，第246页。下面的解释，参见我的论文“黑格尔《法哲学原理》中的传统与革命”，载《哲学研究杂志》(Zeitschrift f. philosophische Forschung)第16卷第2期(1962年)，特别是第218页以下，见本书第170页以下。}通过这一段话，之前已经解放了的西欧社会的功利原则及其出现在曼德维尔和斯密那里的经济模式、狄德罗和爱尔维修的个人利益、边沁和富兰克林的自利便以德国哲学的形式出现了。这里，结果也完全一样：在个人与个人之间，产生了一个独立的、基于利益的关系网络。正如黑格尔在下节所言，市民社会的两个原则“在其实现中的利己的目的”与它的活动条件，即利己的“普遍性”，建立起“一个在所有方面相互依赖的体系”，并且是这样的：“单个人的生活和福利及其权利的定在，都同所有人的生活、福利和权利交织在一起，建立在所有人的生活、福利和权利的基础上，并且唯有在这种联系中，单个人的生活和福利及其权利的定在才是现实的和有保障的。”\footnote{《法哲学原理》，《全集》第8卷(1833年)，第247页。在《哲学全书》第2版和第3版中，黑格尔简单地将市民社会称为“原子论体系”（参见《全集》第7卷第2册，L. 布曼[Boumann]编，1845年，第395页）。}要能够由此出发理解社会之为“市民社会”的描述，我们阅读第183节时，就必须对照第187节。第187节的第一句话是：“个体，作为这种国家的市民，就是以其自身的利益为目的的私人。”然而，这些私人性市民个人的单个利益所追求的目的建立在所有利益目的的内在联系之上，因此要达到单个利益的目的，这些个人就必须按照那种“普遍的方式”规定他们的意向和行为，从而使利益的联系、“社会”直接出现在他们的对面。同时，他们亦如黑格尔所言：“使自己成为这种社会联系之链上的一个环节。”\footnote{《全集》第8卷(1833年)，第251页。参见对这种联系作为一种“社会联系”的描述，它“是所有人都从其中得到他们的满足的普遍财富”。见柏林时期《哲学全书》，第524节（《全集》第7卷第2册，第395页）。}

正是这个普遍依赖的体系及其联系之链，发展出了有别于国家的、从需要和劳动的土壤中生长起来的并通过其活动不断产生的“私人”（第187节）之间的关系网络——“社会”（在该词的现代意义上）。黑格尔在其时代（1815年到1830年之间）的社会架构下（此时，社会尽管依然与等级相连，但已脱离了政治国家），详细阐述了社会的内部划分：经济上的需要的体系（第189-208节）、私法上的司法（第209-229节）以及通过警察和同业公会（第230-256节）实现的国家中的政治伦理整合。这个社会表现为“市民的”，因为它根本上已经逐渐演变为由私法调整的“作为私人的市民”的利益体系；它依然是政治的，实际上只是就它通过等级的和同业公会的团结，在与“警察”的联系中，自我运行在一个相对稳定的秩序体系当中而言的。

因此，我们必须认为黑格尔对市民社会的这种论述（这一论述并非如通常理解所认为的那样特别着眼于普鲁士，而是一般地着眼于1815年后从革命“倒退”的时代\footnote{甘斯在其关于《自然法和普遍法史》(Naturrecht und Universalrechtsgeschichte)的柏林讲座(1832-1833年冬季学期)导言中使用了这一表述。}）是整体性的，也是历史性的。可以说，它是一次伟大的尝试，既不将“社会”建构为基于劳动和享受的需要体系，也不将其建构为一种旧的欧洲社会安于旧的政治状况、安于统治和依附的自然因素，而是对革命以来欧洲人类世界得到解放的“社会”存在与其“政治的”、法的和伦理的秩序进行调解，并在调解中将它们提高到它的新概念上。因此，黑格尔的“市民社会”包含了后来的19世纪“社会”、社会学和国民经济学的对象、如他自己所言的“原子论的体系”，但这只是它的最底层；需要的体系尽管开始了社会的运动，但是社会运动要由“司法”进行规范和组织。因为，诚如黑格尔在第229节所言，在司法中，“市民社会（在市民社会中，理念迷失在特殊性之中，并且陷入内外分裂）回复到它的概念，即自在存在的普遍物与主观特殊性的统一”。\footnote{参见《全集》第8卷，第293页。}如果社会没有在法、伦理和政治上得到规治与团结，就不是“市民”社会。因此，黑格尔此处就司法而言的“概念”，同时也就是市民社会本身的概念。在黑格尔这里，这个概念在新的经济和私法成分（对社会来说，这些成分是构成性的）之上，依然保存了旧的伦理和政治的结构，尽管缩减为“警察”并局限于“同业公会”。因为，另一方面，社会的实体必然性，即警察和同业公会两者对它的“调整工作”（第236节）和“伦理化”（第253节，补充），便在同业公会中将古典世界的实质、“伦理的力量”（第145节）制度化了，并在警察中将“政治制度”（第269节）理性化了。但这也同时表明了黑格尔思想的历史穿透性，即他正好将警察和同业公会这两种秩序要素纳入他对市民社会的阐述之中。

对我们来说，将这两个概念相提并论不仅不恰当，而且更是不可理解：“同业公会”具有古代意味，“警察”则极具实际意味。而如《法哲学原理》所示，黑格尔通过它们描述的事情，并不能立马表明二者之间的内在亲和性及其与市民社会的变化形态之间的历史联系。要指出这种——对黑格尔的特殊概念来说恰好是构建性的——联系，我们就必须透过他的描述，回到直到18世纪为止的强大的古代政治学传统，追问其中的“警察”和“同业公会”究为何物，以及黑格尔本人还在多大程度上意识到它们的结构。首先要说明的是，《法哲学原理》意义上的警察概念跟我们的警察概念关系不大。这个概念最初随着17-18世纪近代国家的制度化和官僚化而得以接受，包含了国家对内部分化开来的社会进行的一般管理。就此而言，它是旧的“市民社会或国家”分化为“国家”和“社会”领域的概念上的对应物。更准确地说，它在管理的手段中标明了去政治化了的社会与政治国家之间的中介。“警察”本身就是变成了国家行政的旧政治，它将市民社会的政治制度及其统治技术表达了出来。\footnote{18世纪下半叶是“警察学”的伟大时期。1830年，青年莫尔(Robert von Mohl)仍以大部头的警察学论述开始其国家法研究。当时的著作有冯·尤斯蒂(J. H. G. von Justi)：《警察学原理》(Grundsätze der Polizeywissenschaft, 1756-1759)、冯·索南费尔斯(J. von Sonnenfels)：《警察原理》(Grundsätze der Polizey, 1765)、冯·普法伊费尔(J. F. von Pfeiffer)：《自然的、产生于社会的终极目的的普遍警察》(Natürliche, aus dem Endzweck der Gesellschaft entstehende allgemeine Polizey, 1779)、勒西希(K. G. Rössig)：《警察学教程》(Lehrbuch der Polizeywissenschaft, 1786)、荣格(J. H. Jung)：《国家警察学教程》(Lehrbuch der Staatspolizeywissenschaft, 1788)、洛茨(J. F. E. Lotz)：《论警察的概念》(Über den Begriff der Polizey, 1807)，等等。参见迈尔(H. Maier)近著《旧德国国家学和行政学（警察学）》(Die ältere deutsche Staats- und Verwaltungslehre [Polizeiwissenschaft], Neuwied/Berlin, 1966)。}当青年黑格尔在1805-1806年耶拿讲座中将现代“警察”同政治的起源、古典希腊哲学的“Politeia”（政制）联系起来时，还清楚地意识到这种联系：“这里，警察来自‘Politeia’（政制）、公共生活和统治、整体自身的行动，然而现在已经降格为整体提供各类公共安全、保护企业反对欺诈的行动。”\footnote{参见1805-1806年《耶拿讲座》，引自《耶拿实在哲学》第2卷，J. 荷夫迈斯特编(1931年)，第259页。}警察就是国家和社会之间的差异得以出现且不断得到调解的形式：它似乎是在市民社会的独立形态中唯一可能的“政治”结构，即行政管理，如《法哲学原理》第235节所言的“公共权力的监督和照管”。\footnote{《全集》第8卷(1833年)，第296页。参见与我们的解释相应的第249节对警察的规定：“作为一种保护和保障大量特殊目的和利益的外部秩序与机构，警察的职能首先是实现和维持包含于市民社会的特殊性之中的普遍物，使这些特殊目的和利益存在于这种普遍物之中。其次，它作为上级领导，也要照顾超出这个社会之外的利益（第246节）。”基于黑格尔这里的前提，洛伦茨·冯·施泰因在1850年后写下了他的重要理论，并且为国家管理的社会进行辩护（《行政学》，1869年及以后）。}作为“普遍物的保证力量”（如果愿意的话，也可以叫作市民社会的“概念”的保证力量），警察要与市民社会的“不可阻挡的活动”、“快速发展的人口和工业”、“财富的积累”（第243节）以及“大量群众下降到一定生活水平之下”（第244节）作斗争。这里，浮现在黑格尔眼前的——首先是英国（第245节注释）——是“产生了贱民”的“群众”。在他看来，这些人将会炸裂作为“市民的”（亦即伦理上有序、政治上受到约束的）社会的社会形式。以前，野蛮的原始民族、家奴、日工和手工艺人处于旧的市民社会以外，现在则是贱民。但在黑格尔这里，与政治学的哲学传统不同，贱民不再表现出这种传统由以获得其历史合法性的积极限制；相反，贱民的存在要求国家行政（即“警察”）去限制社会的“不可阻挡的活动”，并用市民社会去整合他（因此，黑格尔的市民社会还不是1840年左右的早期社会主义者所言的固化了的阶级社会、“资产阶级社会”或者“布尔乔亚社会”）。

市民社会的第二种制度，即同业公会，既与市民社会的解放有关，也与政治蜕化为警察有关。如果说警察是古典政治的“降格”形态的话，那么同业公会就会让人想到古老的家庭团体(Oikos-Verband)，想到“家庭社会”。随着市民社会的内部转型，家庭的“经济上的”和伦理-实体性的秩序失去了基础，因此黑格尔也就不再像康德那样讨论“家庭社会的法”，而只谈到在“作为婚姻的它的直接概念的形态”，在其“外部定在、所有权和财产”，在“子女教育和家庭解体”之中的“家庭”（第160节以下）。黑格尔《法哲学原理》表达的是崭新的、不再与家庭经济细胞相联的家庭概念，并且意识到了为“经济学”的市民社会做好了准备的历史变化。因为有别于政治学的古典传统（对这种传统来说，市民社会居于家庭社会之上），在黑格尔看来，家庭在“市民社会中是一个从属的东西，只构成基础：它不再有如此广泛的活动范围。反之，市民社会则是将人扯到自己这里来的巨大力量，它要求人们为它劳动，人们的一切都通过它，并依靠它而活动”。\footnote{参见《全集》第8卷，第238节，补充，第299页。}在现代，家庭社会的成员从家庭中被“扯了出来”，成为“市民社会之子”，市民社会用它自己的土地“代替了个人赖以为生的父辈土地”，使整个家庭的存在都依赖于它，听任“偶然性”支配。\footnote{同上书，第298页：“首先，家庭是实体性整体，它的职责在于照顾个体的特殊方面。一方面为他提供手段和技能，使他能够从普遍财富中获利；另一方面，在他失去工作能力时，维持他的生活和赡养。但是，市民社会把个体从这种联系中揪出来，使家庭成员之间彼此生疏，并承认他们都是自立的个人。”}

现在，黑格尔认为市民社会的运动原则，它在经济学领域中的“广泛活动”，首先在“营利等级”、市民等级中得到实现。农村大量地保存了旧的经济形式、家庭生活的“实体性”，而在城市中，“个人没有等级尊严，通过它的孤立而沦为利己的工商业者，他的生活和享受也极不稳定”。\footnote{参见《全集》第8卷，第250节和第253节，以及与第256节国家的推演有关的社会的城乡划分：“前者是市民工商业、完全沉入自身并且个别化的反思的所在地，后者则是以自然为基础的伦理的所在地”。}在黑格尔看来，同业公会就是这种持存者，是在“市民社会的劳动制度”及其运行之中的唯一稳定者，因此就是那些劳动中孤立起来的个体的“共同体”。因此，这种稳定物就不是一个固定的等级。因为在这里，在市民社会的劳动制度中，“由于等级不再存在”，所以没有人还在按其等级而生活。\footnote{同上书，第253节，第309页。}毋宁说，同业公会是旧的家庭共同体的延伸。如第252节所言，它作为“第二家庭”出现，以反对社会运动的偶然性和特殊性，“对于一般的、比较远离个人及其特殊急需的市民社会来说，这种地位有待于进一步规定”。在它的这种地位中存在着其实存的“伦理根据”。尽管控诉“贱民的产生与之相联的工商业阶级的奢侈浪费”以及“劳动的日益机械化”（第253节），这种伦理根据似乎向黑格尔证明了它的正当性。

显然，同业公会和警察相对于市民社会的这种“地位”（采用黑格尔的说法）是与18世纪以来市民社会在现代世界中变化了的地位相吻合的。因为如开头想要表明的那样，旧市民社会从自身出发，包含了一种与经济和政治领域的关系，它的形态即产生于这些领域。旧的市民社会坐落在家庭社会及其归属的经济之上，并与之相对照，生长在政治即公共领域的土壤，只有“经济上”独立的社会成员才能进入这一领域；相反，黑格尔第一次在原则上将现代市民社会作为主题，并将其提高到概念意识中。此种市民社会将革命以前的旧欧洲世界的实体性的经济和政治秩序结构降格为偶性，其所以如此，是因为它自己成为了私人生活和公共生活的实体。社会将其政治部分交给国家，仅仅从传统政治那里保留了警察，作为它自己基本运动的管理和秩序功能；随着工业企业的发展，旧家政学的“稳定的东西”、家庭团体解体了，因此市民社会试图将那些在城市的劳动分工中孤立存在的个人重新联合为一个共同体。这样，随着革命而得到解放的社会便在警察和同业公会中恢复了传统的伦理—政治因素，基于此，它也就能够将自己建立为一个“市民社会”。我们可以跟黑格尔一起说，它们是“市民社会的无组织分子绕之旋转的”枢轴。\footnote{《全集》第8卷，第255节，第310页。在黑格尔这里，这句话仅涉及家庭伦理（婚姻的神圣性）和同业公会（个人在同业公会中的尊严），但同时一般而言，也为市民社会第三部分奠定了基础。}通过政治学传统加给其制度的限制，黑格尔阻断了市民社会成为“需要的体系”所勾勒的现代经济社会的趋势。因为社会作为需要体系就是无组织的市民社会，由此也就不再是黑格尔所理解的并在《法哲学原理》中加以阐述的市民社会的概念了。黑格尔通过在概念上将现代的决定性的历史政治过程——社会与国家的分离以及社会中的传统和现代的交互关系——表达出来，便能够思考市民社会的现代形态，同时通过更加古老的结构去限定其实质力量——尽管只是一个早已过去了的时代。可以说，唯有在这里，在一个既进行中介同时又运动着的、突破了中介的中心，法哲学的“市民社会”才历史性地产生出来，并且具有了其真正的概念地位。

\part{法哲学与历史哲学}
\chapter{黑格尔《法哲学原理》中的传统与革命}

1833年5月，黑格尔逝世一年半以后，爱德华·甘斯出版了收入逝者友人版《黑格尔全集》中的《法哲学原理》第2版，并以对此书“未来命运”预兆般的暗示结束了他热情洋溢的前言：他发现“整部著作系由一块自由的金石所铸就”。作为黑格尔体系的一部分，甘斯说道，它将与整个体系共沉浮。也许，若干年后，它将与体系一道进入人们的观念、进入普遍的时代意识，届时它的“术语”将会消失，它的“深邃”将成为共同财富。但是，甘斯进一步总结道，其后，这种在此著中以及一般而言通过概念把握其时代的哲学将在哲学上过时，而仅具有历史意义；新的哲学发展已经在“已然改变的现实”中做好了准备：“这样，它的时代在哲学上已告终结，它成为历史的一部分。一种出自同一基本原理的新的哲学发展脱颖而出，以达到对现实变化的一种不同的理解。”\footnote{《全集》第8卷，前言，柏林，1833年，第VII页。}甘斯此言在不少方面都是预言性的。他道出此言时是黑格尔学派唯一的且依然非常虔诚的造反派。这段话预言了1830年代晚期和40年代黑格尔学派为了争夺这个体系的合法遗产而进行的争斗，也预言了——当然会令人想起1830年七月革命以及黑格尔对它的保留态度——1840年后现实的改变，最后还道出了黑格尔法哲学的特殊命运：在黑格尔法哲学出现后不久，它的时代就会在哲学上成为过去，然而它的历史影响才刚刚开始。甘斯趁“此书”出版之机所言的“未来命运”很快就在整个黑格尔体系上面得到应验。老年黑格尔主义、青年黑格尔主义和新黑格尔主义对法哲学的独特分析引人注目；事实上，这本书在1833年即已成为历史，而在随后的年代（首先是在洛伦茨·冯·施泰因和马克思这里）还“创造了”历史。

与《精神现象学》《逻辑学》《美学》《哲学史》和《历史哲学》在19世纪和20世纪所产生的巨大历史影响不同，《法哲学原理》和《自然哲学》的命运相同：它们受到同时代人的激烈排斥，随着实定的自然科学和历史法学的发展而陷入矛盾之中。对黑格尔法哲学的流俗成见——它迎合了普鲁士国家和1820年代的复辟时期——在一定程度上导致了从斯塔尔(Stahl)到埃德曼、海姆、卢格、罗森克兰茨、费舍尔(Fischer)这些批判者和评论家的空疏，1930年左右新黑格尔主义的哲学努力[首先可以在拉伦茨(Larenz)、杜尔凯特(Dulckeit)、布塞(Busse)和晚期宾德尔(Binder)那里见到]也是如此。\footnote{新黑格尔主义对《法哲学》的评论，见科勒(J. Kohler)发表在《法哲学和经济哲学档案》(Archiv für Recht- und Wirtschaftsphilosophie)第V卷(1912年)上面的论文；合集：宾德尔、布塞、拉伦茨《黑格尔法哲学导引》(Einführung in Hegels Rechtsphilosophie, 1931)；宾德尔(J. Binder)：《法哲学奠基》(Grundlegung zur Rechtsphilosophie, 1935)以及K. 拉伦茨和G. 杜尔凯特的作品，法哲学的新黑格尔主义开端为人们认为在黑格尔那里作为“国家中的人民共同体”而重新发现的“具体的共同体自身”[参见G. 杜尔凯特：《法概念和法形态》(Rechtsbegriff und Rechtsgestalt, 1936)，第17页]。黑格尔法哲学在19世纪和20世纪的影响史尚有待书写。}这里再度证实了甘斯所写前言的惊人准确性：前言当时就指出“此书的实际价值与对它的接受与传播极不相称。”\footnote{《全集》第8卷，前言，柏林，1833年，第V页。}众所周知，1840年后的政治运动进一步损害了对黑格尔法哲学的承认。它对时代的影响，不是来自对其意义的恰当评论，而是来自对它的批判。因此，法哲学的时代在哲学上已经终结，因为1840年后的时代不再是哲学的，而是政治的、经济的和历史的了。

毫无疑问，在法和国家领域，19世纪黑格尔的伟大辩护人非洛伦茨·冯·施泰因莫属。但是在他的《当代法国的社会主义和共产主义》（1842年）及其国家和社会理论的相同概念背后，它们所描述的社会和历史现实已经发生了多么巨大的变化啊！拉萨尔自称是黑格尔和甘斯的真正传人；就其哲学才能和法学教养而言，在黑格尔的学生中，他是1860年后唯一能够在法哲学基础上写出一部“政治学”的人。然而非常典型的是，他在他的第一部同时也是最后一部法哲学著述《既得权利的体系》（1861年）中宁愿投身于一个特殊的法学问题，向当时的经验主义和历史主义致敬。马克思觉得自己既不是黑格尔的辩护士，也不是他的传人（在这个词语的真正意义上），而是他的批判性的战胜者；因此，他就是在他的时代唯一如其所是地对待1821年《法哲学原理》的人：将它作为“唯一与正式的当代现实保持在同等水平上的德国历史”。\footnote{《〈黑格尔法哲学批判〉导言》，载《马克思恩格斯全集》（第1版），第一部分，第1卷第1册（1927年），第612页。（中译文引自《马克思恩格斯全集》第3卷，人民出版社2002年版，第205页。——译者）}青年马克思于1843年对第261-313节写下的《黑格尔国家法批判》（马克思通过这部著作彻底摆脱了经由革命改变了的政治形而上学传统），是19世纪黑格尔传承中对法哲学所做的唯一的、有时在历史方面还相当深刻的、与黑格尔的论述水平相当的评论。

显然，这部令人遗憾的未完成的评论性批判作品并未提出黑格尔法哲学在政治哲学传统中的地位问题，按其历史意图，它根本就不可能提出这个问题。与甘斯的预言一致，马克思的评论认为黑格尔法哲学仅仅属于历史，因此他的批判一开始就已经预设，法哲学的时代在哲学上已然过去。众所周知，批判的结果就是使哲学脱离法，而转向经济学。同时，这种公然宣称的非哲学依然深受这部最后的伟大的法的形而上学的影响，因为马克思不是在世界这部大书中，而是首先在黑格尔法哲学中学会看见并理解到了作为现代之特征的“对立”：国家与社会、国民与私人、政治生活与社会生活、中世纪与现代之间的对立。接下来的解释努力也恰好在其中见到了黑格尔的“深刻”：如马克思的评论所言，黑格尔“处处从各种规定（如我们各个邦所存在的那类规定）的对立开始，并且强调这种对立”。\footnote{《马克思恩格斯全集》（第1版）第一部分，第1卷第1册（1927年），第465页。[中译文引自《马克思恩格斯全集》第3卷，第69-70页。——译者]}对《法哲学原理》（首先是对其第三部分）的这种哲学与历史上的问题和形式的分析表明，这种对立一直深入形式的结构之中，并且整部著作阐述了一部在革命的基础上生长起来的非亚里士多德“政治学”的方案，其原理试图描绘和规定现代人的基础——他们的国家、社会、家庭和个人，同时又不割断与古典传统的联系。

\section*{1}

黑格尔给他的这部著作起了两个书名：《法哲学原理或自然法和国家学纲要》，这种做法乍看起来，不免令人侧目。第二个书名并非如人们所料，是一个副名；相反，通过添加“纲要”(im Grundrisse)二字，而成为法的“原理”(Grundlinien)的异名。由此，对“法哲学”形式上的强调便具有一种深意。

这一异名中的两个名称指出了黑格尔以前形而上学思想中的两门学科，即自然法和国家学。自然法是在近代欧洲（尤其是在17、18世纪）发展起来的，国家学则以“政治学”之名，属于欧洲的古老传统，并在直到沃尔夫为止的学院哲学中地位稳固。然而，旧政治学的本质恰好在于，其中并无自然法和国家学的分离。因为对于直到17世纪为止的欧洲思想（在德国则要推迟到18世纪中叶）来说，政治学是对人、公民或者政治的社会（或者共同体）在法律上进行规治的共同体进行全面阐述的理论，亚里士多德对于城邦作为政治共同体的论述乃是政治学的典范。\footnote{对此，参见布伦纳：《社会史的新道路》第30、207页，还有我在不久之后对市民社会所做的概念史研究。}由于城邦或国家作为政治共同体或者市民社会包含了自由市民（即基于其身份而享有政治权利的市民）在自由和法下的全部生活，因此非政治的和市民以外的人就表现为野蛮的，而“市民的”(civilis)一词便通过其与“自然的”(naturalis)的对立而获得了积极的含义。只有在16、17世纪，从马基雅维利、博丹和霍布斯开始，当现代国家从政治上组织起来的市民社会中解放出来时，国家学才从旧政治学中解放出来，从而18世纪出现了自然法和国家学之间的对立。此后，市民的(civilis)和自然的(naturalis)的意义便完全颠倒了过来。因为从历史—政治角度，而非按照其哲学内容来看，自然法表达了这样一种努力，即要么坚决限制现代国家对其本身原本就拥有的良好法秩序的社会的（旧的意义上的）“不合法的”干预，要么用一种新的自由和法的概念革命性地克服已经僵化为单纯的权力功能的国家政治。自然法与国家学之间的对立（或者如德国人所言，道德和政治之间的对立）在精神上引导了近代欧洲革命，并在此后与革命过程共进退。黑格尔法哲学以这种对立为前提，并在革命的历史背景下，如革命本身一样，想要通过哲学政治的努力达到对革命的克服。黑格尔在此著前面写上两个书名，作为其书名，正好表明了这一点。

“法哲学原理”是这部著作的第一书名，也是它的真正名称。我们看到，这个书名在第二个书名（“自然法和国家学”）前面。作为对现代国家的阐述，黑格尔的国家学想要成为一种突出意义上的“法哲学”。因此，在他这里，政治学，作为国家中人的社会生活的理论，便成了法哲学，因为他试图通过将前国家的自然法与（政治的权力功能为个人颁布的）法完全建立在自然—理性法——人之为人的自由——之上，扬弃二者之间的革命性对立。然而，对黑格尔来说，自由的这种自然—理性法（这里不详述其意义）具有一种历史内容。自从基督教进入世界历史起，这种内容就是灵魂在上帝面前的平等，从18世纪末开始，就是个人在革命国家面前的平等和自由。由于黑格尔理解到了革命的历史意义，并将革命的历史内容——所有人的自由和权利能力——接纳到他的政治思想中，因此对他来说，权利便在其与自由的联系中成为政治哲学的中心。随着革命的这一结果，革命以前的自然法和国家学之间的对立便失去了意义。此后，政治学的基础就只能是权利，但不是旧的市民社会的具有权利能力的市民的权利，而是权利能力的权利、追求其自由人的“概念”。

然而，黑格尔一方面想要在一种建立在权利之上的政治哲学中积极调和自然法和国家学之间的对立，另一方面他又回到了革命前的政治学传统。在这一传统中，无此种“对立”。如沃尔夫翻译这个名称时所言，政治学是关于“人的社会生活（特别是共同体）”的理论，此外别无他物。\footnote{克里斯蒂安·沃尔夫的德语版“政治学”全名为：《克里斯蒂安·沃尔夫致热爱真理者，为促进人类幸福对人的社会生活特别是共同体的理性思考》(Vernünftige Gedanken von dem gesellschaftlichen Leben der Menschen und insonderheit der gemeinen Wesen zur Beförderung der Glückseligkeit des menschlichen Geschlechts, den Liebhabern der Wahrheit mitgetheilet von Christian Wolff)，1721年第1版。}因此，说黑格尔法哲学的前提建立在普遍的法概念和自由要求之上（这两个成分表明它隶属于近代自然法），这是正确的：它作为国家理论（并且在国家理论中）发现它的基础。如同在亚里士多德那里，城邦是“真正的基础”，按照现实性来说是“首要的东西”，按照理性来说是人联合成社会的一切结合的“目的”。\footnote{黑格尔：《法哲学原理》，E. 甘斯出版(1833年)，第312页及以下。下面的引用均出自这一版本。}事实上，法哲学在很大程度上就是国家学，众所周知，从埃德曼和海姆、拉萨尔和罗森茨威格直到拉伦茨和杜尔凯特的新黑格尔主义对它的解释都集中于这一点。国家学这一名称必然令人想到旧政治学传统：在国家整体中阐述一族人民或者一个社会的政治秩序，黑格尔有意识地重拾这一传统。与古老的欧洲政治学传统相契合，黑格尔拒绝时人喜爱的“关于国家的哲学”及其目标：提出自己的“理论”，“世界上从未有过国家和国家制度，现在也还没有，但是现在——这个现在是永远继续下去的——似乎应该从头开始”。\footnote{同上书，第7页。（中译文引自《法哲学原理》，第4页。——译者）}他像传统一样，与当前的伦理、国家及其秩序紧密相连。同时，革命性的时代政治理论已经打破了传统意识，学院哲学以及康德的政治著作中非同寻常的辩护性论争清楚表明了这一点。这方面尤其典型的是革命后对“哲学与现实间的关系”的特别反思。国家的现实性及其实际的法和统治的形式既不是（如《法哲学原理》序言中所言）通过“建立彼岸物”而以乌托邦的方式摧毁掉，也不是在传统的政治或市民的科学或历史(scientia bzw. historia politica sive civilis)的意义上达到对其关系的简单陈述。相反，黑格尔的政治思想与现实的“肯定”关系意味着通过“理解当前和现实的东西”而“探究理性的东西”，因此本身就是一种现实与概念的“关系”，并且是一种“关于”国家与历史的哲学（如果人们想要违背黑格尔的话）。看似恢复传统的东西，实际上是一种通过革命破坏了的、其意义已然改变的传统。黑格尔的国家学正是在现代自然法的基础上生长起来，并且将要表明，政治在法中的中介，是一个既进行中介也在运动、逃离中介的中心。——由此，我们远远超出了与黑格尔法哲学的两个书名有关的法、自然法和国家学的概念分析，进至法哲学本身的体系问题及其与政治哲学传统的关系。

\section*{2}

可以假定，1820年黑格尔不仅在考虑到给他的柏林听众发一本讲座（“我循职进行法哲学授课”）提纲的需要时，认为有必要对《哲学全书》的相关部分进行单独阐述。\footnote{《法哲学原理》，第3页。（中译文参见《法哲学原理》，第1页。——译者）}相反，如果比较1817年海德堡版《哲学全书》“客观精神”相关章节与1821年《法哲学原理》，就会清楚地看到，它并非如柏林版序言所言，只是“对同一些基本概念所作的更为详尽尤其是更加系统的阐述”，而且还讨论了更多的东西。而通过附释澄清“正文的抽象内容”以及“更广泛地考虑当前流行的一些见解”的提示，也不能作为应当付印《法哲学原理》的“动因”。决定性的可能更是这一事实，即1817年后哲学体系的发展导致黑格尔超出《逻辑学》，进到“法的科学”。因为他在序言中说：“我已在我的《逻辑学》中详尽阐述了思辨知识的本性。”\footnote{同上书，第4页。（中译文参见《法哲学原理》，第2页。——译者）}就其阐述了“实践”知识的本质而言，《法哲学原理》正好与前者匹配。如黑格尔所强调的那样，法哲学正好关乎“科学”，尽管它并未总是在逻辑形式上非常清楚地显明“对象的具体、自身如此多样的属性”。实际上，在黑格尔本人生前出版的著作中，《法哲学原理》（1821年）与《逻辑学》（1812、1816年）彼此并列。如果将《精神现象学》视为体系（在1817年《哲学科学百科全书》中，这个体系总称为知识的总体）的导论，那么逻辑学和法哲学的体系撰构上的显著平行便会令人想到哲学的传统划分：理论哲学和实践哲学。这种划分存在于欧洲形而上学中，在18世纪，沃尔夫及其学派也在对亚里士多德学术传统的依赖中系统地坚持了这一划分。这里，实践哲学与逻辑学并列，人类灵魂的两方面的能力(facultas aminae)——认识能力和欲求能力(facultas cognoscitiva atque appetitiva)——被称为两者的基础(fundamentum)。自然法、政治学、伦理学的划分基于实践哲学的定义：“真正哲学的教导运用欲求能力求善避恶的这一部分，叫实践哲学。因此实践哲学就是指导欲求能力求善避恶的科学。”(Ea vero philosophiae pars, quae usum facultatis appetitivae in eligendo bono et fugiendo malo inculcat, philosophia practica dicitur. Est adeo philosophia practica scientia dirigendi facultatem appetitivam in eligendo bono et fugiendo malo.)\footnote{克里斯蒂安·沃尔夫：《理性哲学或逻辑学》（1740年），哲学概论之预备性讨论，第3章，第62节。}

特别的是，黑格尔没有在他的著作中援引过这个传统，因此需要考虑他在柏林时期对大学“完善的哲学课程”的意见。黑格尔认为，“完善的哲学课程”与明确、中肯、系统的阐述一样严肃和重要。这里的“大致”意思是说，“提供一种按照理论哲学讲座与实践哲学讲座的习惯划分而来的完善性。因此，我在一个学期讲授逻辑和形而上学，在同年度的另一学期讲授自然法和国家学，或是包含了伦理学或义务论意义上的法哲学”。\footnote{《柏林大学哲学系档案》，文献目录第2号第1卷，载《纽伦堡著作集》，J. 荷夫迈斯特出版，1938年，第XXIV页。}如《逻辑学》第1版序言中最尖锐地表达出来的那样，他确信旧形而上学已经终结。这一确信说明了黑格尔在哲学的古老划分上面已然断裂了的传统意识，因为前述引文不过是“大致”的回忆。公开表达出来的对于“大约25年来哲学思维方式在我们中间所遭受的完全改变”的意识、“精神的自我意识在这个时期对它自身所达到的立场”并不是偶然地导致他再次改变了他于1807年计划的《科学体系》的体系结构，并且放弃了精神现象学为一方面，逻辑学、自然哲学和精神哲学为另一方面的两分法。此后，黑格尔的哲学体系划分为他所言的“两门实在科学”，即自然哲学和精神哲学，逻辑学作为“纯粹科学”构成了它们的第一部分。\footnote{见《逻辑学》，第一部分，拉松编，1934年，序言，第7-8页。[中译文参见〔德〕黑格尔：《逻辑学》（上卷），杨一之译，商务印书馆1966年版，第5-6页。——译者]}

这里不能进一步讨论这个划分的问题。对我们的主题来说，根本在于，众所周知，法哲学与历史哲学一道从属于“精神哲学”第二部分即“客观精神”之下。客观精神的内容是人类世界的活动(Pragmata)：法和国家、经济和社会、政治和历史。因此，客观精神学说即以人与人之间，在家庭、市民社会和国家中有其基础和目的的那些人类行为为主题。这些行为就其本身的意义而言并不超越于人世之外，而是人世的“客观精神”，在欧洲形而上学传统中，它们是实践哲学（实践哲学首先是政治哲学和伦理学，因而是广义上的道德形而上学）的对象。与其对象的制度、在社会和政治上有序的人类世界相一致，直到18世纪为止，基于古代文化传统的学院哲学在实践哲学领域按照实践对象本身的结构和人的特定身份，区分了三门学科：伦理学、家政学和政治学。因为如沃尔夫所言，人的身份可以用两种方式来考察——“一个人要么在一定程度上是人，要么在一定程度上是市民。或者换而言之，一个人要么在一定程度上生活在人类社会或自然状态，要么在一定程度上生活在市民社会。因此实践哲学便以两种方式分为两个部分。”(vel quatenus est homo, vel quatenus est civis, aut, quod perinde est, vel quatenus vivit in societate generis humani seu statu naturali, vel quatenus vivit in societate civili. Ob duplicem hunc respectum philosophia practica in duas dispicitur partes.)\footnote{克里斯蒂安·沃尔夫：《理性哲学或逻辑学》(1740年)，哲学概论之预备性讨论，第3章，第63节。}这两个部分就是伦理学和政治学。伦理学的对象为人类的伦理状态，因此，它研究一般的道德存在者(entia moralia)和人之为人的地位。“对生活在自然状态或人类社会中的人进行思考的哲学部门，叫伦理学。因此，我们将伦理学定义为指导自然状态下（或者人就是他自己的主人、不服从别人的权力时）自由行动的科学。”(ea philosophiae pars, in qua homo consideratur tamquam vivens in statu naturali, seu in societate generis humani, Ethica appellatur. Quamobrem Ethicam difinimus per scientiam dirigendi actiones liberas in statu naturali, seu quatenus sui juris est homo, nulli alterius potestati subjectus.) 政治学的对象为市民社会(societas civilis)的法制度，因此不是类的伦理自然，而是通过法律规定的国家(res publica、Staat)的政治自然。它的主题不是作为人的人，而是作为市民的人，因此将他从他的伦理行为的普遍性和无规定性中拉出来，置于国家的市民法律的特殊性之下。“对生活在国家或者市民状态下的人进行思考的哲学部门，叫政治学。因此，政治学就是指导市民社会或者国家中自由行动的科学。”(Ea philosophiae pars, in qua homo consideratur tamquam vivens in Rep. seu statu civili, Politica vocatur. Est itaque Politica scientia dirigendi actiones liberas in societate civili seu Republica.)\footnote{克里斯蒂安·沃尔夫：《理性哲学或逻辑学》(1740年)，哲学概论之预备性讨论，第3章，第64、65节。}

在实践哲学这两门基本学科（政治学作为市民及其在市民社会中的权利和义务的理论，伦理学作为整个人类社会中的人的理论）以外，还有将人作为家庭成员来考察的家政学。如同亚里士多德《政治学》卷一，沃尔夫谈到了构成家庭之基础的“小型社会”(societates minores)：“先于国家或市民社会，有处于自然状态或者国家以外的小型社会，如夫妻社会、父权社会、主奴社会，它们通常来自同一家族。”(Dantur praeter Remp. seu societatem civilem societates aliae minores, quibus etiam in statu naturali seu extra Remp. fuisset locus, veluti societas conjugialis, paterna, herilis, quae simplicium nomine venire solent.)\footnote{克里斯蒂安·沃尔夫：《理性哲学或逻辑学》(1740年)，哲学概论之预备性讨论，第3章，第66节。}

在这个持续了好几百年的传统中，实践哲学的这种三分法不只是僵死的学院知识。它表达了革命前欧洲社会状况，反映了其政治、经济和伦理的风貌。随着17、18世纪革命，现代社会从国家中解放出来，这个传统瓦解了：先是在英格兰（霍布斯在政治学上是反亚里士多德派，洛克、弗格森和休谟已经完全站在新“社会”的基础上），然后在法国。伦理学与政治学分离，作为“道德哲学”，变成了它的敌人，家政学突破了其特有的限制，变成了18世纪的国家经济学和国民经济学——如黑格尔所言，变成了“在现代世界的基础上产生的若干科学之一”。\footnote{《法哲学原理》，第189节，第255页。（中译文参见《法哲学原理》，第204页。——译者）}在现代社会从国家中解放出来的过程中，以及与此相关的实践哲学（特别是旧政治学）的内容变化中，家政学中的革命发挥了决定性作用。对此，我们将在黑格尔这里具体指出来。有别于英格兰和法兰西，唯有在德国，18世纪仍然突出地保持了实践哲学的形而上学传统。这里，在沃尔夫及其学派，这一传统还在描绘着革命前的社会：家政学和政治学对应于家庭社会和市民社会的不同领域、国家和社会的统一，伦理与有效的法律相连，因此伦理学与政治学相连。\footnote{对此，参见青年洪堡对1785-1786年通俗哲学家恩格尔按照沃尔夫学派手册讲授的实践哲学的部分的笔记[《洪堡全集》(Gesammelte Schriften)，第一部分，第7卷第2分册，1908年，第363页以下]。}

对于当时西欧思想来说具有建构性的自然法与政治学的分离，在18世纪德国学院哲学中（除了大法学家普芬道夫和托马修斯以外）并未发挥决定性作用，这一点绝非偶然，因为自然法脱离了传统实践哲学的框架。西欧近代自然法的功能是与政府对立，保护“社会”的权利，使社会从政府解放出来，与此相对，沃尔夫将自然法的内容限定为一种普遍的实践哲学“理论”。自然法的定义在于它的规定之中：善行和恶行的科学，出现在欲望之中的人对善恶的内部知识。沃尔夫总结道：“自然法显然是实践哲学（即伦理学、政治学和家政学）的理论。因此，区分理论和实践并不难，自然法本身就能传达伦理学、家政学和政治学。”(Facile patet jus naturae esse theoriam philosophiae practicae, Ethicae scilicet, Politicae atque Oeconomicae. Quamobrem cum non opus sit, theoram a praxi distingui; jus naturae in ipsa Ethica, Oeconomica atque Politica tradi potest.)\footnote{克里斯蒂安·沃尔夫：《理性哲学或逻辑学》(1740年)，哲学概论之预备性讨论，第3章，第68节，补充。}沃尔夫只是形式上将自然法与实践哲学分开，它的内容则能够“传达”到社会的伦理、经济和政治的生活秩序之中去。

\section*{3}

现在，如果比较一下沃尔夫依赖亚里士多德和经院哲学最后一次在其庞大著作中建立起来的这个欧洲传统实践哲学的体系纲要与黑格尔政治形而上学的外部轮廓，仔细考察后就会惊奇地发现，尽管政治革命和18世纪下半叶开始的工业革命粉碎了传统实践哲学体系的整个结构，但是它的基础在黑格尔的《法哲学原理》中依旧宛然可见。概念的自我展开作为法哲学的辩证奠基（在黑格尔看来，在实践中，概念就是在其自由中的人的意志，自由就是意志的“概念”），导致他将这个领域划分为“抽象法”“道德”和“伦理”三个部分。抽象法和道德包含了人的前政治存在的理论。在抽象法中，可以认识到演绎为——以财产自由和职业自由为顶点的一——现代市民的私人权利的自然法；在道德中，可以认识到那种沉浸于主体的自我感之中的伦理学，通过18世纪晚期革命，形成了这种主体的自我感。在抽象法和道德中，人作为人，从政治状态和政治制度下解放出来，实现自己，表达自己：或是在与外部的关系中，与外部世界和周围世界的关系中，在所有权（第41-71节）、契约（第72-80节）和“不法”（第82-104节）等形式中，另一方面在与内部的关系中，在故意和责任（第105-118节）、意图和福利（第119-128节）、善和良心（第129-141节）等形式中。这里不可能对根本而言为探讨私法内容的18世纪自然法与黑格尔的抽象法（特别是私法思想的形而上学基础方面）进行具体比较。\footnote{首先参见《法哲学原理》的详细导言，黑格尔在意志的形而上学理论的基础上阐述了现代法权人格及其普遍自由要求的“概念”（第1-32节）。在黑格尔看来，不是“以一定身份而得到考察的人”作为不同于奴隶、农奴或“侨民”的“市民”，而是“人格性本身”构成法和所有权概念的基础（第40节）。}我们仅限于指出，对黑格尔来说，自然法与伦理学两者，都是人类世界伦理—政治状态的抽象物：作为私法的自然法，标明为“抽象法”，是抽象人格性的法；18世纪主观化了的伦理学，作为道德，是同样抽象的主观性的道德。作为人的人只有在作为市民的人与作为家庭、社会和国家之内的伦理—政治的存在中，才“在实践上”实现其自身。

相比较而言，抽象法和道德探讨了传统政治哲学所不知道的现代内容——想想人在外部世界的实现（第41节以下）、在劳动和使用中对人格的所有权要求的奠基（第56、59节以下）、所有权和法权人格的等同（第51节）以及人格的主体存在、作为“道德的概念”的“意志对它本身的内部关系”——，伦理则明确回到旧的欧洲伦常传统的立场。在黑格尔这里，伦理指的是个人与特定人民和国家的“伦理力量”（第145页）及“必然关系”（第148节）的同一、古代市民的“习俗”(Ethos)。伦理学和政治学同样地包含于这种习俗之中，以保存法、伦理和国家的古老统一。黑格尔在其讨论实践哲学主题的第一篇文章《论自然法的科学探讨方式》（1802年）中写道：“个人的伦理，以及反过来，个人伦理的本质绝对是实在的（因而就是普遍的）绝对伦理。个人伦理是整个体系的脉动，并且自身就是整个体系。这里我们注意到语言的暗示。这种暗示尽管在以前遭到人们的拒绝，但从前面所言来看，它是完全正当的。这就是：绝对伦理的自然中即有一种普遍物或者伦常，因此，表示‘伦理’的希腊词，还有德语词，都卓越地表明了伦理的这种自然，而最近的伦理体系由于以一种自为存在和个别性为原则，就无法做到使用这些词语而不言明它们的联系。”语言的这种内在暗示如此重要，以至于18世纪的体系为了标明“其本质”，即个别化了的个人，“就不能滥用这些语词，而是采用了‘道德’一词。尽管从其词源上看，这个词语也指向了同样的方向，但因为它是一个新造词，就不能如此直接地反对其蹩脚的含义”。\footnote{见《批判的哲学杂志》第2卷第3册(1803年)，第1-2页。}

文献中反复提到黑格尔国家概念的“古代化”手法，同时毫无疑问，古代城邦理想对于他拒绝将个人从伦理力量的束缚中革命性地解放出来，对于他将国家制度和自由意识、伦理和法等同起来，是多么重要。因此，就形式而言，法哲学第三篇与古代实践哲学第三部分相同：“伦理”代表了伦理学和政治学的统一，对于伦理法的国家制度的传统理论——国家或者市民社会(civitas sive societas civilis)的理论——来说，这种统一是根本性的。黑格尔在他自己的“政治哲学”中并不是徒然谈到“客观的”而非“道德的”“伦理义务论”与“将在接下来的第三部分系统阐述的伦理必然性圆圈”（即国家中的“市民”关系）之间的同一性。\footnote{《法哲学原理》，第148节。一种“内在、一贯的义务论”（即旧义上的伦理学）“不外乎阐述由于自由的理念而是必然的，因而在其在国家的整个范围内是现实的那些关系”（第213-214节）。（中译文参见《法哲学原理》，第167页。——译者）}

由于黑格尔重拾古代传统，因此将现代政治和工业革命的内容包含于这种对伦理概念的古代理解之中，也就更加令人惊异。因为黑格尔《法哲学原理》的吊诡正好在于，它在此名义下记录了革命时代的基本趋势，并且在政治哲学的体系中做出了自博丹的主权概念和卢梭的公意以来最深刻的改变。唯有在这里，青年黑格尔广泛的国民经济学研究和政治研究才引人注目地获得了体系性成果。在耶拿时期的文章和讲座中，有——乍看之下令人咋舌的一一从研究英国和法国的著作获得的对于革命过程中把握到的欧洲工业和政治世界的制度的洞见，这种洞见根本上与古代文化传统无关，但与之并列。资产者被搬到了罗马帝国时期，社会的经济领域、需要、劳动和占有，被称作与希腊政治活动(politeuein)的绝对伦理相对立的“相对伦理”，古代悲剧在现代相对伦理中上演，法哲学所认识到的对“市民社会的劳动方式”的分析要么在个人与自然相关的实践过程框架下进行，要么尴尬地出现在“人民”“政府”“国家”等标题之下。\footnote{参见《批判的哲学杂志》第2卷第2册(1802年)，第79页。收入《政治学与法哲学文集》，G. 拉松编，第418页以下、464页以下，特别是参见1803-1804年《耶拿讲座》（《实在哲学》第1卷，J. 荷夫迈斯特编，1932年，第220页以下和236页以下）和1805-1806年《耶拿讲座》（《实在哲学》第2卷，1931年，第197页以下、213页以下和231页以下）。}实践哲学的旧结构再也不能容纳18世纪末政治、社会和经济的变化，这首先可以在1803-1804年和1805-1806年耶拿讲座的两份草案中看出来。这里，裂变后的世界的丰富内容并没有得到合理的体系化处理。另一方面，当黑格尔试图在1817年《哲学全书》中第一次对精神的“客观”内容进行严格体系化时，古式风格仍然使他无法对现代关系做出恰当的阐述。\footnote{唯一阐述这一主题的《海德堡哲学全书》第433节只谈到“普遍作品”及其在“诸等级”中的特殊化，其中家庭作为“个别性等级”出现。“特殊的”等级——“其作品是特殊定在的需要，而其切近的目的是特殊的主体性，但是要达到这一目的则以其他一切人的劳动为前提，同时也会影响到这一劳动”——还不是“市民社会”而只是“等级”。比较一下黑格尔在《法哲学原理》出版后在《哲学全书》第2版(1827年)对这一部分所做的全面改动（第513节以下）。在这里和其他地方，黑格尔将“伦理”术语历史化了，并将其与近代家庭、市民社会和国家的发展联系起来（第552节）。}只有1821年法哲学才能将革命的结果纳入实践政治哲学的构造，但因此也就同时摧毁了其传统的体系构架。

首先，当黑格尔将伦理，即人类世界的伦理—政治制度，它的“现实精神”，理解为家庭、市民社会和国家的著名三分时，他似乎回到了古代的欧洲传统。因为家庭和国家乃是大革命以前的政治哲学从亚里士多德（《政治学》卷第1卷，第1-2章；《尼各马可伦理学》第8卷，第14章，1162a17以下；《欧德谟伦理学》第7卷，第10章，1242a23以下）那里继受而来的主题，二者构成了人类共同生活的基础，其中国家(polis或civitas sive res publica)与社会等同，被称作“市民社会”(koinonia politike或societas civilis)，以区别于家庭的“家庭社会”(societas domestica)。\footnote{与“社会的”与“政治的”（或者“市民的”与“国家的”）之间的现代区分相对立，古代政治哲学的基本区分为“oeconomicus”（经济的）或“domesticus”（家庭的）与“civilis”（市民的）之间的区分，因此市民社会表现为公共—政治社会。}事实上，外部结构上的相似令人惊异，并且《法哲学原理》中黑格尔思想的历史来源最令人信服地表现在该书最后面的这个最重要部分对这三个概念的形式接受上。另一方面，真正说来，只有理解到这三个概念获得了完全不同的意义、这三个词语获得了全新的内容时，才能敏锐、清楚地把握黑格尔与传统的断裂。因为黑格尔法哲学所言的家庭在结构上与古老的家庭团体大相径庭，市民社会也完全不同于旧欧洲的市民社会，因此黑格尔的国家也就不同于传统的国家模式(Civitas-Modell)，因为它已经从自身中区分出“市民社会的坚实体系”。\footnote{《历史哲学》讲座导言中的这一说法，参见《历史中的理性》，J. 荷夫迈斯特编，1955年，第209页。}《法哲学原理》第三部分结构的伟大正在于，通过其思想的历史实质（即黑格尔将现代革命的各种趋势接受到法哲学之中），黑格尔对这些概念的历史来源进行了改造，并与之断裂。此后，在黑格尔伦理纲要中的传统与革命之携手同行，就应当在三个概念的结构中得到具体表达；在这三个概念中，国家和市民社会便在它们的对立中成为现代政治哲学的真正中心。

\section*{4}

黑格尔在第157节以一个简要的预先规定开始其详细的阐述。“现实精神”、人类世界伦理的实际本质是：

\begin{quote}
第一，直接的或自然的伦理精神——家庭；这种实体性向前推移，丧失其统一性，进而分裂，过渡到相对物的立场，因此就成为——

第二，市民社会，这是各个成员作为独立的个人在一个形式普遍性之中的联合，这种联合是通过成员的需要，通过作为保障人格和财产手段的法制度，以及通过维护他们的特殊利益与共同利益的外部秩序而建立的。这个外部国家——

第三，返回并凝聚为国家制度，它是实体性的普遍物的目的和现实性，也是致力于该普遍物的公共生活的目的和现实性。
\end{quote}

这一描述的引人注目之处在于，“市民社会”在这里即已占据相对较大的篇幅。家庭和国家代表并保存了旧的意义上的伦理、它们的“实体性”，市民社会则是“家庭统一性的丧失”。黑格尔引入这个概念的方式表明他显然摆脱了传统伦常，而市民社会的“矛盾”则使他在进行新的描述时对内容做出了详尽的规定。

老年黑格尔派甘斯、魏塞、H. F. W. 欣里希斯(Hinrichs)、亨宁(Leopold von Henning)已经认识到[黑格尔]将（现代意义上的）“市民社会”引入政治哲学乃是一个重大成就。毫无疑问，只有到了1820年，黑格尔才随着这个概念“理解”到了欧洲社会政治制度内部的这些基本改变，而不只是像他开始政治研究时那样进行断言。从亚里士多德到康德的法哲学传统将国家称作市民社会，因为人类社会本身就是在政治上（在自由市民的权利能力和在等级特权中）和在社会上（在家庭的经济上的重要地位）受到规治的。黑格尔则将国家的政治领域从已经变成“市民的”社会领域中区分出来。由此，“市民的”(bürgerlich)便获得了排他性的“社会的”含义，不再如18世纪那样作为“政治的”同义词使用：它指的是在绝对国家私人化了的现代市民——作为布尔乔亚的市民——的单纯“社会”地位。第190节补充用法语术语说明了这一点。这个从其政治—法上的含义解放出来的市民概念与同样解放出来的社会联系起来。它们的政治实体消解于它们两者通过现代革命而获得的社会功能之中。如此看来，黑格尔将“市民的”和“社会”结合为一个政治范畴绝非偶然。这个概念前所未有，在他自己1820年以前的著作中也未曾见到。因此，罗森茨威格所说的黑格尔对他所知道的舒尔策的《一切科学简论》（1759年）和弗格森的《文明社会史论》（1767年）的摘抄的提示就是错误的，因为此二人涉及的都是旧的传统概念（尽管形式上有所减弱）。\footnote{罗森茨威格：《黑格尔与国家》第2卷，1920年，第118页。在舒尔策这里，罗森克兰茨提到的地方在旧的意义上考察市民社会，如他所言，一个“市民国家”的社会制度（《一切科学简论》[1759年]，第189-190页）。弗格森的《文明社会史》（1767年）同样在与“政治社会”相同的意义上使用这个概念（第三部分，第6章），但在此名下除了广泛讨论政治关系（首先是斯巴达、雅典和罗马）外，还讨论了“艺术和科学”。}黑格尔从未在他的著作中接受过出现于从托马斯·阿奎那和大阿尔伯特到博丹、霍布斯、洛克和康德之中，并将革命以前的市民社会的制度表达出来的政治哲学的传统公式\footnote{参见“‘市民社会’概念及其历史起源问题”，本书第141页以下。}：“civitas sive societas civilis sive res publica latius sic dicta”（国家或者市民社会或者广义上的国家），一般而言，他也不再理会它的这种意思，就如法哲学似乎提到——从新的市民社会的立场出发——这种等同的一段话所言：“如果把国家想象为各个不同的人的统一，亦即仅仅是共同性的统一，其所想象的只是指市民社会的规定而言。许多现代的国家法学者都不能对国家提出除此以外任何其他看法。”\footnote{《法哲学原理》，第182节，补充。（中译文引自《法哲学原理》，第197页。——译者）}这里，与传统的断裂经典地表达了出来，然而另一方面，黑格尔确实完全知道旧市民社会与政治国家的同一，而且有意思的是，他是从语言出发，并且基于当时大体上还是受到等级束缚的社会[而获得这种知识的]。“尽管在这些所谓的理论看来，市民社会的一般等级和政治意义上的等级是绝不相同的东西，但语言仍然保持了以前就存在的二者之间的结合。”\footnote{同上书，第303节，第398页。（中译文引自《法哲学原理》，第323页。——译者）}同样，《法哲学原理》“市民社会”部分第一段就从实际出发、从由政治制度中解放出来的社会（即前述意义上的“市民社会”）的私人—市民出发，清楚表达了对旧概念的传统的完全抛弃：“具体个人作为特殊的人本身就是目的。他作为各种需要的整体以及自然必然性和任意的混合体，构成市民社会的一个原则，——但这个特殊的个人本质上是与其他的这样的特殊性相关的，从而每个人都通过他人的中介，并且同时完全只是作为经过了另一个原则（即普遍性的形式）的中介，使自己有效并得到满足。”\footnote{《法哲学原理》，第182节，第246页。对“市民社会”的最好阐述仍然是由罗森茨威格（《黑格尔与国家》第2卷，1920年，第118页以下）和卡尔·洛维特：《从黑格尔到尼采》（1941年）第二部分第1章第2节做出来的。此外，也可以比较克罗纳(R. Kroner)发表在《应用社会学档案》(Archiv für angewandte Soziologie)第4卷(1931年)上面的文章和里特尔的著作《黑格尔与法国大革命》对市民社会作为现代“劳动社会”的思辨阐述。}这里，获得解放的西欧社会的功利原则及其经济模式、狄德罗和爱尔维修的个人利益、边沁和富兰克林的“自利”便以德国哲学的形式出现了。于是，结果也完全一样：在个人与个人之间，产生了一个独立的、基于利益的关系网络。黑格尔在第183节写道：“利己的目的在其实现中……建立起一个在所有方面相互依赖的体系，单个人的生活和福利及其权利的定在，都同所有人的生活、福利和权利交织在一起，建立在所有人的生活、福利和权利的基础上，并且唯有在这种联系中，单个人的生活和福利及其权利的定在才是现实的和有保障的。”个体，“作为这种国家的市民，就是以其自身的利益为目的的私人”（第187节），基于这个“市民社会”特有的“必然性”（第186节），他们使自己成为其强有力的联系之“链上的一个环节”。

这个对新社会来说是典型的必然性(1)出现在其最低级阶段，即作为自然（“a. 需要及其满足的方式”）、劳动（“b. 劳动的方式”）和占有（“c. 财富”）的需要的体系（第189-208节）。这些成分以它们的方式在“社会化”的形式中产生了一种自身的必然性：无限多样的需要及其满足的手段（第190-195节）；劳动分工，“人的依赖和相互关系”变成了“完全的必然性”（第198节），同时劳动“机械化”，人可以“走开”，机器登场（第198节）；社会生产的“普遍财富”分配为特殊财富及其结果，“人的设定起来的不平等”（在黑格尔看来，“在市民社会中不但不扬弃不平等，反而从精神中产生它”[第200节]），以及社会的“全部联系”区分为各等级（第201节以下）。这种必然性也(2)支配了市民社会的上面的层次，在这一层次，在“司法”中，市民社会似乎回到了以前的制度要素，即法（第209-229节）。然而，这种司法之法本质上是私法，即适用于“市民社会中无限零星和复杂的所有权和契约的关系和种类的素材”（第213节）的“市民”法。这个领域的必然性在于：(1)“法的设定的存在”（第212节以下）；(2)其内容上的即无限多的案例以及特定的个别案例的量的特性，这些个案加剧了实定性即其作为设定的存在的特性（第214节）；(3)私人间的相互关系的无限变化的性质，这种性质限定了法，并将其限制于法律。只有第三阶段，“警察和同业公会”（第230-256节），才能将必然性“伦理化”，这种必然性像冰冷的古代命运一样支配着市民社会的个人。尽管“警察”——这是古代社会政治制度的一个残余概念——局限于外部“秩序的范围”以及平衡“无意识的必然性”（第231节以下、236节以下），但在市民工商业的同业公会，“伦理物作为内在物回到了市民社会中”（第253节以下）。

与设定起来的法相比，警察作为“普遍物的保证力量”，更加必然地同社会的特殊性不停地斗争。在即将进入国家的政治领域之前，黑格尔在“需要的体系”中转向现代社会经济的自然基础之后，便在这一部分第243-248节将工业革命的趋势纳入市民社会的概念中。此后，工业即作为“市民社会的劳动组织”而被谈及（第251节）。市民社会再也不像在沃尔夫和政治传统中那样，由“许多家庭”组成，\footnote{克里斯蒂安·沃尔夫：《克里斯蒂安·沃尔夫致热爱真理者，为促进人类幸福对人的社会生活特别是共同体的理性思考》，第214节。}而是（如果在其“不可阻挡的活动”中去理解的话）由全新的成分构成：富裕和贫困、工业和无产阶级大众、人口增长和殖民扩张。市民社会(1)“在它自身”带来了“快速发展的人口和工业”，(2)导致“大量群众下降到一定生活水平之下”，产生了所谓的“贱民”，(3)陷入一种“辩证法”，这种辩证法表现在这一事实中：“尽管财富过剩，市民社会总是不够富足。也就是说，它所占有而属于它所有的财产如果用来防止过分贫困和贱民的产生，总是不够的。”\footnote{《法哲学原理》，第245节。“国家容纳市民社会”，原文为“In-Korporation der Bürgerlichen Gesellschaft vom Staate her”，即“从国家出发吸纳市民社会”，意为在承认市民社会相对独立的前提下，将其纳入国家之下，以消除前述市民社会的“辩证法”。——译者}在黑格尔看来，由这种辩证法，对市民社会来说，只有两条出路：要么是其内部的一条内在的出路，市民社会因其矛盾趋势，而走出自身，进行殖民（第248节）。这条道路只是拖延，并不解决问题。要么是另一条——人们可能称之为——“超越的”道路，即《法哲学原理》第三部分（第257节以后）以之作结的国家容纳市民社会。然而，如1840年后的事件所表明的，这条道路并未扬弃前述辩证法，而是跳过了它。如果扩张达到了可以利用的空间的极限，殖民的道路就会对社会内部的矛盾力量无能为力，而当国家面对时代的可能性而无所作为时，就会作为乌托邦而在时代的现实性面前撞得粉身碎骨。

无论如何，黑格尔在《法哲学原理》的这个部分（第243-248节）和《历史哲学讲演录》的一个地方具体指明了市民社会是一个时代概念。黑格尔在他关于亚里士多德政治学的演讲结束时说道，“我们近代国家的抽象权利把个人孤立起来，允许他如此”，继而建立起一种必然的联系，以至于“真正说来，没有人具有整体的意识，或者为了整体而活动；他为整体而工作，但是不知道他是怎样为它工作的，他只关心他个人的保护”。\footnote{《哲学史讲演录》第2卷，K. L. 米希勒(Michelet)编，1833年，《全集》第14卷，第400页。（中译文参见《哲学史讲演录》第二卷，第364页。——译者）}黑格尔在这里用近在眼前的现代工厂形象说明现代个人、国家和社会的结构划分：“这是一种分工的活动，在这种活动中，每个人只占一份；正如在一个工厂里面，没有人自己单独制造一件完整的产品，每个人都只制造产品的一部分，不懂得制造其他各部分的技能，只有少数几个人才把各部分装配成一件产品。”只有自由的人民才有整体的意识并为它而活动，现代人作为个人本身就是不自由的——“市民的自由就是不需要普遍的事物，就是孤立的原则”（强调字体为黑格尔所加）。

《法哲学原理》第305-308节（在“伦理”第三部分，关于国家的章节）更加明确地规定了这里明言的这种时代上的联系。这几节一方面在里尔之前30年就论述了市民社会中流动的和不流动的社会力量，即等级要素的两个方面：建立在不动的自然基础之上的贵族等级和在工商业中努力拼搏的市民等级。黑格尔认为，市民等级构成了“市民社会的流动方面”（第308节）。另一方面，在将社会等级提升至国家时，这几节也将当时的两院制（世袭的贵族院和民选的代议院）固定了下来：“私人等级在立法权的等级要素中获得政治意义和政治效能。所以，这种私人等级既不是简单的不可分解的集合体，也不是分裂为许多原子的群体，而只能是它现在这个样子，就是说，它分裂为两个等级：一个等级建立在实体性的关系上，另一个等级则建立在特殊需要和以这些需要为中介的劳动上。”\footnote{《法哲学原理》，第303节，第397页。（中译文引自《法哲学原理》，第322页。——译者）}再有意思的是，黑格尔这里在用一个可以说是“政治上”修正了的市民的等级社会概念，去反对社会可能消解为“简单的不可分解的原子”，以及“许多原子的群体”，同时也反对国家的政治形式化。就等级架构上而言，市民社会依然是政治社会，并且正好处于解放了的孤立个人与国家组织之间。由此，在黑格尔这里，市民社会就不仅仅是在英国国民经济学推动下对国家学所做的一个补充，而是正好处于这两个过程的交合处：一是旧市民社会的去政治化过程，二是国家从市民社会中解放出来。因此《法哲学原理》第308节便谈到了“市民社会流动部分”的议员：“由于这些议员是市民社会派出来的，所以就直接得出这样一个结论：市民社会在派议员时是以它的本来面目出现的，就是说，它不是一群原子式地分散的单个人，不是仅仅为了完成单一的和临时的活动才一时凑合起来而事后则没有任何进一步联系的人们；相反地，它是一个分为许多有组织的协会、自治团体、同业公会的整体，这些团体因此才获得政治联系。”\footnote{《法哲学原理》，第308节，第400页以下。（中译文引自《法哲学原理》，第325-326页。——译者）}黑格尔这一概念的历史实质及其某种意义上的吊诡、双重含义首先表明，随着这一概念，从法国革命开始的市民社会与政治生活的分离尽管应当被扬弃，然而同时由英国开始的经济革命却要接受到这一概念之中。黑格尔市民社会的这种复杂的双重时代关联，其不仅指向过去和未来，而且指向现代革命的两股潮流，指向法国和英国，这一点构成了它的历史现实性，并且成为它在下一代发生政治影响的基础。

\section*{5}

即使我们并不接受里特尔新近出版的黑格尔解读\footnote{《黑格尔与法国大革命》，美茵河畔法兰克福，1965年新版。}的所有要点，那么市民社会在法哲学的概念结构中一再出现——作为抽象法的具体基础（第209节以下），作为主观道德的广泛的活动领域（在贫困领域，第242节），作为消失在其两极分化之中的“伦理事物的现象世界”以及作为伦理的“沦丧”（第182、184和229节），作为个体之“子”和家庭的实体性基础（第238节），最后作为国家的“前提”（第182节）——这一事实也就证明了他的这种论点的正确性：自1820年起，市民社会即“进入哲学及其政治理论的中心”。\footnote{《黑格尔与法国大革命》，第57页以下。}更加具体地说，市民社会尤其处于“伦理”的三个部分的中间，并且表现为将家庭和国家连接起来的部分，这在欧洲政治哲学传统中是前所未有的。由于这里不能对这个历史问题进行深究，因此就仅限于提到托马修斯在18世纪初如何界定市民社会。他在《政治明智短论》中写道：“然而人的社会本身要么是市民社会，要么是家庭社会。家庭社会构成市民社会的基础，因为这里市民社会的意思不外乎许多家庭社会及其人口的联合，而处于一个普遍的统治之下。”\footnote{托马修斯：《政治明智短论》，1705年，第204页以下。}这个定义完全处于旧传统之下，并且像《普鲁士国家的普通邦法》(Allgemeines Landrecht für die Preußischen Staaten, 1794)这样一部18世纪末制定的贴近当时时代社会现实的法典一样，当它在第一部分第2节宣称“市民社会由更多更小的、通过自然或法律（或是通过两者）结合起来的社会和等级组成”\footnote{《普鲁士国家的普通邦法》，第一部分，第1章，第2节。}时，也同样与传统相去不远。

与传统唯一的但却重大的区别也许是用诸等级来补充“社会”——沃尔夫的小型社会。在《普通邦法》中，市民社会仍然表现为政治社会、国家，其根据在于至少理论上保持不变的家庭、家庭社会的传统地位；相反，市民社会则以其政治公共领域的结构与之形成鲜明对照。因此，《普通邦法》下一节便立即规定：“夫妻之间、父母与子女之间的联合构成家庭社会。仆人也要算入家庭社会。”\footnote{《普鲁士国家的普通邦法》第一部分，第1章，第3节。}在这一点上，《普通邦法》的作者与同一时期康德的《道德形而上学》（1797年）第一部分即《法权论》完全一致。在康德的《法权论》中，市民社会同样保留下来，而家庭社会法也有其自己的部分。\footnote{《法权论》，第22节，见《康德全集》（科学院版）第6卷，第276页。在康德看来“家庭制度”的特征是“物的方式的人格权”，即“占有一个作为物的外部对象，并将这个对象作为一个人格来使用”（第276页）。通过这种法权关系，康德将家庭团体的古老“家政”结构、主奴关系（第283页）固定下来。对此，参见黑格尔在《法哲学原理》第40节对这种“物的方式的人格权”的批判。}然而，一旦“国民经济学”的发展突破了这种“家庭社会”的限制（如在研读过斯图亚特和斯密的黑格尔这里），社会便获得了指派给“oeconomia”（家政）的功能。相反，“societas civilis”作为“市民社会”便获得了新的地位，随着工业革命开始，市民社会的实体必然性支配了个人及其“小型社会”，而家庭社会则相应地解体。

我们可以在黑格尔的家庭概念上准确地看出这种交互影响。作为“伦理”的第一阶段，家庭不是一个拥有不同地位的成员的“社会”，它不是由丈夫和妻子、父母和子女的“小型社会”构成，而且绝对不会由主人和奴隶构成；相反，它是一个法律上的“人格”，“在所有物中”，而非在家庭中，有其外部实在性。\footnote{《法哲学原理》，第169节：“家庭作为人将在所有物中具有它的外部现实性。它只有在作为财富的所有物中才具有它的实体性的人格性的定在。”}不同于康德，在黑格尔这里，传统欧洲家庭的“经济”性消失了，并被18世纪晚期的“情感性”家庭概念所取代。\footnote{布伦纳在其论文《“完整家庭”和传统欧洲的“经济学”》(Das "Ganze Haus" und die alteuropäische "Ökonomik")中使用了这一表述，该文后被收入《社会史的新道路》（1956年），第44页。}《法哲学原理》第158节规定：“作为精神的直接实体性，家庭以精神的自我感觉到的统一性——爱——为规定，因此这里情感就是，精神的个体性在这种作为自在自为地存在着的本质性的统一性中自我意识到，自己在其中不是作为自为的个人，而是作为成员而存在。”

“必然性”作为市民社会的基本规定，消解了作为基本经济单位的家庭。我们看到，市民社会使其结构交织起来，结合为一个实体性的、经济的和法的整体。个人尽管还是市民社会的“成员”，然而首先不再是参与到其实体之中，而是作为“独立的”个人与市民社会的运动相对立。因为第238节说，唯有家庭才“首先”是“实体性整体”，它的职责在于“照料个人的特殊方面，既要考虑到他的手段和技能，使他能够从普遍财富中谋生，又要考虑到他在丧失工作能力时的生活和照顾”。然而实际上，如黑格尔所言，家庭“在市民社会中是从属的东西，它只构成基础：它的活动范围并不如此广泛。相反，市民社会才是惊人的力量，它把人扯到自己这里来，要求他为它工作，要求他的一切都要通过它，并且通过它而行动”。\footnote{《法哲学原理》，第238节，补充，第299页。（中译文参见《法哲学原理》，第241页。——译者）}这里，市民社会的必然性表现为它之所是：居于家庭的伦理实体和个人生活之上的“惊人的力量”，它使个人彼此“生疏”，只承认他们是“独立的个人”：“但是，市民社会把个人从这种联系中揪出，使家庭成员彼此生疏，并承认他们是独立的个人；此外，市民社会又用它自己的土地取代个人赖以为生的外部的无机自然和父辈的土地，使整个家庭的存在都依赖于它，听从偶然性的支配。由此，个人就成为市民社会的子女，市民社会要求于他，而他也对市民社会拥有权利。”\footnote{《法哲学原理》，第238节。（中译文参见《法哲学原理》，第241页。——译者）}这些地方最清楚地表达了黑格尔关于家庭团体中“小型社会”的解体及其与市民社会概念之间的联系的看法。

于是，黑格尔的市民社会概念涉及传统意义上的“家庭”社会的结构变化，同时另一方面也关系到“政治”国家的地位变化。这里，政治学不再是市民科学(scientia civilis)，即在其政治制度中对市民社会进行阐述，而是与之相对的“国家科学”。在欧洲旧传统中，国家和社会在市民社会的关系概念中交织在一起，在黑格尔这里，它们才开始被“设定”到“关系”之中，并且彼此处于一种特殊的“差别”之中：“市民社会是处于家庭和国家之间的差别[阶段]，尽管其形成要晚于国家。因为它作为差别，以国家为前提，它要存在，前面就必须要有一个独立的国家。”\footnote{同上书，第182节，补充，第246页。（中译文参见《法哲学原理》，第196页。——译者）}黑格尔将现代的国家制度的现实性纳入神一般静止的、实现了其目的的、自足的古代城邦理念的传统形式之中，在其内部解决了国家与社会之间的二元论。因为黑格尔明言，这个国家乃是“现代国家”，其本质在于，“普遍物是同特殊性的完全自由和私人福利相结合的，所以家庭和市民社会的利益必须集中于国家；但是，特殊性也必须保持其权利，如果没有特殊性自己的知识和意志，那么目的的普遍性也就不能前进”。\footnote{同上书，第260节，补充，第322页。（中译文参见《法哲学原理》，第261页。——译者）}

在黑格尔这里，国家不仅“设定”了市民社会，而且相反，还为其自身的制度现实性预设了市民社会的存在。柏林历史哲学讲座对北美政治状况所作的令人印象深刻的经验分析表明了这一点。这里，黑格尔得出了这样的结论：现在北美仍然不是一个完全发展了的国家，因为它还没有市民社会。在北美，这种发展受到西部土地占领的阻碍，人口的“不断迁移”阻碍了城镇化进程：“只有像在欧洲一样，农业的简单增长受阻，居民才会不再走向田地，而是涌向城市工商业，形成紧密的市民社会体系，满足有机国家的需要。”\footnote{《历史哲学》，导论，引自《历史中的理性》，1955年，第209页。}这里，如布伦奇利在19世纪晚期所言，市民社会已经带有了“第三等级概念”的特征。

黑格尔与传统的断裂并不只是表现在这一点上面，即政治学的对象分化为国家和市民社会的差别。它更是通过这种现代“国家”与历史的关系而表现出来。因为它的前提并不仅仅基于市民社会的动力学，而是同时向上延伸到历史之流中。国家理念在其臻于完成的顶点表现为“类以及与个别国家相对的绝对权力、在世界历史过程中赋予自己现实性的精神”。\footnote{《法哲学原理》，第259节。（中译文参见《法哲学原理》，第259页。——译者）}国家不再是传统的政治状态概念，不再是贯穿传统政治思想之始终的国家或市民社会的静止不变的自然模型；相反，历史作为新的上司高踞其上。在历史中，“伦理的整体自身、国家的独立性”——这里我们必须替黑格尔补充完整：再次——“委诸偶然性之下”。\footnote{同上书，第340节。（中译文参见《法哲学原理》，第351页。——译者）}对于黑格尔《法哲学原理》的这种结尾，甘斯言道，“真正说来”，其中包含了“此书最重要的价值”。一般而言，这种说法道出了他那一代人的心声和19世纪观念。在甘斯看来，这个结尾乃是一幅“壮丽的景象”，从国家的高度出发，我们看到众多国家“卷入历史的世界海洋”；同时，他在当时率先指出，黑格尔在《法哲学原理》末尾给出的“历史发展概要只是对这一领域中那些更加重要的关切的一个预告”。\footnote{《法哲学原理》，前言，第VIII-IX页。}

\section*{6}

如果我们全面考察黑格尔法哲学的概念结构及其建构原则，那就会说，作为革命的结果和遗产，市民社会和历史这两个要素耸立其中。两者同时改变了自古代以来一直有效的政治哲学传统所阐述的家庭和国家的内部制度。正如家庭居于作为家庭成员的个人与市民社会“之间”一样，国家运行于市民社会与历史之间。在黑格尔这里，正是在古代实践哲学体系的建筑术上的起源，产生了法的富于艺术性的辩证结构。最后，法哲学作为“中介体系”由此出发获得了富有历史内容的独特财富，这种财富使它在原则上与卢梭的《社会契约论》、康德的《法权论》和费希特的《自然法基础》等政治著作区分开来。

黑格尔在古代哲学史讲座中言及柏拉图时，反对对《理想国》的流行成见，说柏拉图“不是一个玩弄抽象理论和抽象原则的人，他的真实精神曾经认识了并表述了真实的事物。这不能是别的，而只能是他生活于其中的世界的真实事物，也只能是那唯一很好地活在他本人和希腊里面的[时代]精神的真实事物”。在更深的意义上，这种说法也适用于他本人。事实上，他用《法哲学原理》以现代世界的“实体性方式”描述了这个世界的真正“伦理”。他用来结束柏拉图讲座的这一部分的那些话语，也能很好地放在作为最后的形而上学的哲学“政治学”的法哲学的开头：“没有人能够跳出他的时代，他的时代的精神也就是他的精神；但问题在于认识到时代精神的具体内容。”\footnote{《哲学史讲演录》，《全集》第14卷，第275页。（中译文引自《哲学史讲演录》第二卷，第248页。——译者）}
\chapter{黑格尔历史哲学中的进步与辩证法}

\section*{1}

在《历史哲学讲演录》导言中，黑格尔精练地、简明扼要地将世界历史界定为“自由意识的进步”，并补充道：“【……】一种我们必须在其必然性中予以认识的进步。”因此，讲座的主题就是认识世界历史中的必然进步——进步的理论构成历史哲学的核心，自由理念则是这个理论的基准点和线索。\footnote{此文系在1969年巴黎举办的第七届黑格尔国际大会上的谈话。}除了《法哲学原理》序言的那个著名的模糊表达（“凡是现实的，都是合乎理性的；凡是合乎理性的，都是现实的”）以外，黑格尔哲学中没有第二个命题像历史哲学的这个基本命题那样，被人频繁引用，同时也被彻底误解。这两个命题的影响（同时下面也会讲到它们的联系）清楚地表明，无论是显而易见的晦涩，还是清楚明白的表述，都会使人对哲学命题产生误解和曲解。

事实上，与先前所有世纪相比，19世纪在科学技术以及与之相关的对自然的支配上都取得了无可比拟的进步，对这样一个时代来说，还有什么能比黑格尔关于自由意识的进步的历史公式更加令人信服呢？如果说，在黑格尔看来，历史哲学的任务在于认识这种进步的必然性\footnote{指前面所言及的“历史哲学的任务在于发现历史进步的必然性或规律”。——译者}，则它的追求便与19世纪实证主义学派的诉求并无二致：发现这种所有人有目共睹的历史进步背后世界历史的各个阶段或者各个时代有规则的更替规律，这种规律决定了历史事件的进程及其发展趋势。就像18世纪的杜尔哥(Turgot)和19世纪的孔德一样，黑格尔也接过该定义，提出了一种世界历史的三阶段规律，作为他的世界历史规划的基础：按照自由的规则，最初一个人是自由的，然后一些人是自由的，最后所有人都是自由的，由此，世界历史就是一个由东方世界、希腊—罗马世界和日耳曼世界或现代世界构成的时间序列。在这个自由进展的理念中，黑格尔似乎完全接近了他的时代的那种进步的乐观主义。这种乐观主义认为，所有人的自由都在其中得到实现的市民社会就是一切可能世界中最好的世界。由此，这种神正论便经历了一种莱布尼茨从未梦想过的扩张。通过认识肯定物、认识社会整体的向前发展，世界的负面物（在莱布尼茨这里是个体的恶）得到了辩护：历史的进步不是假象、不是欺骗，世界精神绝不撒谎。

然而，对黑格尔公式的误解并不在于，它似乎用一种自由的量的进展去表现世界历史的进程，而在于它以现时代的进步为尺度去打量这一进程。换言之，从最新的进展出发，将以往的历程构想为一个必然通向这一进展的发展过程。这种时间序列上的定位“预设了，在世界历史中，唯有后果丰富者方才有效，世界事件的前后序列须按成果的理性去评定”。\footnote{参见洛维特：《从黑格尔到尼采》，第3版，苏黎世，1953年，第238页。}19世纪下半叶，达尔文的进化论在自然领域中为这种成果评价提供了坚定的支持，由此，这种成果评价就构成了黑格尔的进步观念为何如此流行的原因。黑格尔的进步观使按照“强者的权利”在与封建统治势力斗争中获胜的市民阶级心安理得，就像随着工业生产力的解放，无产阶级出现为市民等级的世界历史的对手一样。然而，自由意识的进步理念却变成了市民意识的进步的意识形态。在达尔文主义进化论和实证主义社会学的影响下，历史变成了一个掩盖了历史主体自身的作用及其活动意义的线性的自然过程。\footnote{卢卡奇：《历史与阶级意识》，柏林，1923年，第215页。也参见梅洛-庞蒂：《辩证法的历险》，法兰克福，1968年，第44页以下。}进步在一定程度上抽身隐退，变成了需要依赖的事实，或是一种人们必须服从的规律。

这个偶像化了的进步概念为社会主义理论所克服，乃是19世纪历史的悖论之一。恩格斯在其著作《路德维希·费尔巴哈与德国古典哲学的终结》（1888年）中对黑格尔历史哲学的解释实现了这种突破。恩格斯与黑格尔一致，也承认自然和历史的发展有着根本的区别。自然中，只有无意识的、盲目的力量彼此发生作用，不会出现有意愿的、有意识的目的，而在社会历史中，则是有意识的、经过思虑和凭激情行动的、追求一定目的的人。但在恩格斯看来，这种区别并不能改变下述事实，即“历史过程”是受内在的普遍“规律”支配的，问题只是在于发现此种规律。\footnote{《马克思恩格斯全集》第21卷，柏林，1962年，第296页以下。（中译文参见《马克思恩格斯选集》第四卷，人民出版社2012年版，第253-254页。——译者）}这里，取代自由意识中进步理念的是一种自然历史发展的理论，这种理论根本而言是实证主义的。在这种理论中，进步变成了规律，就像支配物体在空间中的运动规律一样，它也支配了人在时间中的行为。诚如马克思所言，时间就是历史的空间。

进步思想的这种改变一直影响至今，其结果是，黑格尔的进步观念与辩证法之间的原始联系消失不见了。我的阐述旨在通过对黑格尔历史哲学中的历史运动模式及其时间结构进行解释，恢复这种联系。我可以引述一名现代经典诗人的证言，这位诗人——在同代诗人中，这是一个独一无二的现象——同时又是一个杰出的辩证法家和黑格尔哲学的爱好者。我指的是布莱希特(Bertold Brecht)。在他不久前出版的遗存下来的哲学笔记中，对社会主义者们的“莽撞的进步概念”做出了这样的简洁而又难以捉摸的表述：“‘进步’这个概念具有非常愉快的政治性质，但它对‘辩证法’的概念却有不利的后果。”\footnote{布莱希特：《政治与社会文集》(Schriften zur Politik und Gesellschaft)第1卷，W. 黑希特(Hecht)出版，柏林，1968年，第250页。}针对线性进步导致了相反方向的反驳，布莱希特以其独特的方式和辩证法家的态度，证明基于“普遍”规律的辩证法理论于历史实践是不利的，甚至是荒唐的。他写道：“如果人们向这种简单然而动人的观点的信徒指出，他们的看法就像那些看手相的人的看法一样，认为在看完手相确定即将到来的事情之后，能够阻止这些事件的发生，那么这些信徒就会郁郁不乐，咕哝发火。”柏拉图对话中那些无所不知、自以为是的智术师被他们自己的武器打败以后，也是这般束手无策。实际上，辩证法并不是普遍的、包含了自然和历史的规律科学，而是如布莱希特所言：“一种思维方法，或者不如说是一系列彼此相连的可理解的方法，它们允许将一些僵化的观念消解掉，用实践去反对处于统治地位的意识形态。”\footnote{布莱希特，《政治与社会文集》第1卷，第251页。}

我想要证明，黑格尔的进步理论有别于后来的进步的意识形态，它是要使辩证法的这种双重观点生效。

\section*{2}

在我们着手证明之前，似乎有必要首先简述一下黑格尔历史哲学赖以产生的历史前提。现代进步概念的历史始于16世纪，一开始是不停地发展，然后到启蒙运动时期导致了一系列二律背反。这些二律背反构成了黑格尔为此概念进行辩证奠基的起点和出发点。现代进步观念起源于城市商业和制造业等级的解放理想。随着培根、笛卡尔、霍布斯和莱布尼茨对科学理论和行动理论提出新的筹划，现代进步概念也就产生出来。这种新的理论不再像经院主义—亚里士多德哲学那样将它们的目标预设在一个最高的善或者最高的知识中，而是认为它们自身原则上是开放的，而不是封闭的。自然科学构成了这种理论的典范，其整齐划一的直线运动的概念预设了自然界并无固定的目的和目标，也预设了运动相对于静止的优先性。正如自然的科学一样，其实践上的推论，即通过劳动和行动控制自然，也不是完美的、自身封闭的，而是一个迈向未来的过程。\footnote{参见培根：《新工具》第一卷，1620年，80-84。}现在，科学和行动的过程性成为伦理学和人类学的基本概念。霍布斯用一种善(bonum)代替了旧道德哲学著作中的至善(summum bonum)，它不像后者那样在于完满和静止，而在于无休无止地从满足愿望到满足愿望的行进。\footnote{霍布斯：《利维坦》（1651年），第一部分，第11章。}近代市民社会于其产生时即被理解，它的至善包含了它对幸福的解放性要求。莱布尼茨表达了同样的观念，他理解的幸福不是尽善尽美的快乐，而是向新的完善和新的快乐不断迈进。\footnote{莱布尼茨：《自然和恩典的理性原则》(Vernunftprinzipien der Natur und der Gnade, 1714)，第18节。}这个无限进步的概念成为18世纪启蒙运动理解人的自然和人的历史的关键。莱布尼茨引入的关键词叫可完善性，它用自己的手段摧毁了传统的自然目的论，并在本体论上使市民阶级的解放要求合法化。按照这种将进步提高到宇宙论层次的理论，事物并不因其符合亚里士多德的实体和本质的本体论要求永恒不变就是完美的，而是相反：事物的可变性构成它们在一种更高意义上可以称之为完美的根据。完善只有作为可完善性（也就是说，作为完善的过程），才是可能化。

这里，进步思想达到了其最宽广的范围，由此产生了那些贯穿了启蒙运动历史哲学思想的二律背反。第一个二律背反是在【进步】概念自身中的二律背反，这个概念一方面似乎包含了一个目的——万物的完善，等等，然而另一方面却将此目的表现为原则上是不可达到的。由此产生了第二个二律背反，它是在运用这个概念时产生的。换言之，我们要么假定进步只对个人有意义，要么想要反过来证明，只有类、人类才能达到完善的目标。这个二律背反在18世纪文学上、在门德尔松和莱辛争论中得到了解决。\footnote{参见门德尔松《耶路撒冷》，《全集》第5卷，莱比锡，1844年，第120页；莱辛，《人类的教育》，《全集》第8卷，柏林，1956年，第590页以下。}第三个二律背反形式上看似乎重复了第二个二律背反的立场，但内容上却存在着原则性差异。它接受了曼德维尔和卢梭对市民社会经济、文化和政治矛盾的发现，由此，在18世纪下半叶，解放的理想便进到了批判阶段：社会可完善性的原因不是德性，而是恶德，其结果不是平等，而是不平等，其前提不是贫富均衡，而是贫富对立。早期启蒙的进步乐观主义证明其自身是一个为了保证概念自身得到应用而必须抛弃的幻想。这出现在卢梭这里。在他看来，所有超出人类自然状态的进步，既是人类文化和理性完善的脚步，也同样是人类的衰落。\footnote{《论人类不平等的起源》，《卢梭选集》，巴黎，无出版年份，第65页。}康德得出了正好相反的结论：对人类来说，人走出自然状态的过程，从单纯动物的野蛮状态过渡到人道状态，乃是从坏到好的进步，对个人来说，则并非如此。\footnote{康德，《人类历史起源臆测》（1786年），对第一部分的解说。（“从坏到好的进步”里德尔原文为“Fortschritt vom Bessern zum Schlechteren”，兹据康德原文更正。康德原文见Kants Werke. Band VIII. Akademie-Textausgabe, Berlin 1968, p. 115。——译者）}

第三个二律背反的决定性观点是由卢梭和康德阐述的一种进步辩证法的观念，这种观念使他们二人发现，历史的运动结构自身就是辩证的，是一种进步和回归的运动。这种观点构成了黑格尔的直接前提。然而卢梭对于在市民社会及其文化的基础上解决进步的辩证法仍然心存疑虑，并且着眼于一种依然与自然相契合的社会形式，筹划其《社会契约论》的模式，康德则将进步辩证法作为其历史哲学的出发点。这种历史哲学不再是将此领域以前的努力变成空洞无益的游戏的种种虚构。在特别严肃的卢梭那里，历史根本而言是一部哲理小说，因为回归的运动与退步的运动并无本质区别。康德用一种进步理论回答卢梭，认为进步就是人类所有原始禀赋的发展，并将通过经验后天认识到的现代市民社会的对抗性——人的非社会的社会性——提升为历史运动的动力，同时又按照一条先天的线索，将其提升为“人类合法秩序的原因”。通过证明人类行为中的有规律的进程及其可能同意组建一个社会（在这个社会中，每个人的自由与所有人的自由都能按照一条普遍法则而共存），解放的模式便得到了辩护。康德在其《世界公民观点之下的普遍历史观念》（1784年）第五个命题中说：“自然迫使人要去解决的人类的最大问题，就是建立一个普遍的照管权利的市民社会。”\footnote{康德原文见Kants Werke, Band VIII, p. 22。（中译文参见〔德〕康德：《历史理性批判文集》，何兆武译，商务印书馆1990年版，第8页。——译者）}康德历史建构中的先天线索，就是在权利思想中得到保证的自由，而不是至善（或者德性和幸福）的实现。历史哲学中的这个哥白尼式转向在于，偏离旧的神学和伦理-教育的历史思辨，转向现代市民社会及其历史，现在，这个历史需要一个新的、非形而上学的神正论。康德高呼：“因为在没有理性的自然界之中我们赞美造化的光荣与智慧并且把它引入我们的思考，究竟又有什么用处呢；假如在至高无上的智慧的伟大舞台上，包含有其全部的目的在内的那一部分——亦即全人类的历史——竟然始终不外是一场不断地和它相反的抗议的话？这样一种看法就会迫使我们不得不满怀委屈地把我们的视线从它的身上转移开来，并且当我们在其中永远也找不到一个完全合理的目标而告绝望的时候，就会引导我们去希望它只能是在另外一个世界里了。”\footnote{《康德全集》（科学院版）第8卷，第30页。（中译文引自《历史理性批判文集》，第20页。——译者）}

\section*{3}

黑格尔的历史哲学由此出发，我的主张是，只有它才完成了对哥白尼式转向后果的思考，并将18世纪进步理论的那些二律背反抛到了身后。黑格尔同意康德，认为世界历史的基础不是幸福或者德性与人道的教育，而是在自由与权利关系之中的自由。黑格尔称莱辛和赫尔德的“人类教育”观念是“富于精神的”，但是这种观念“只是远远地”触动他，并非“此处所谈”。\footnote{黑格尔：《历史中的理性》，第150页。参见《法哲学原理》，第343节，附释，《全集》（百年纪念版）第7卷，第447页以下。}他与卢梭、康德一致，认为可完善性概念本身无法满足，在应用到历史中时是没有标准的：“在本讲座中，一般说来，进步具有量的形式。更多的知识，更好的教养，只能是这样的比较；可以一直这样说下去，而不言及任何规定性，说到某种质的东西。”\footnote{黑格尔：《历史中的理性》，第150页。}然而，比一致更加重要的是在康德和黑格尔的历史进步观之间存在的区别。我将其总结为五点，其中第一点是决定性差异，它决定了其余几点。

1. 不同于康德及其18世纪前人，在黑格尔看来，一个普遍的照管权利的市民社会的理念不再是悬设，而是现代世界的现实性，也就是那个在法国革命中作为世界历史的原则而使所有人的自由都变得有效并得到承认的“宪政”国家的原则。在它面前，作为历史目标（同时历史也从来就不能达到这一目的）的进步的二律背反已然失效。

2. 在康德这里，进步的过程有两个要件，一是在历史之前即已存在于其最终形式之中的自然的计划，二是带有作为质料的激情的人的自然，在这种质料中，前述形式得以实现其自身。\footnote{我在这里使用了科林伍德的一个中肯的表达，见R. G. 科林伍德：《历史哲学》，斯图加特，1955年，第119页。}黑格尔则将进步置于历史自身当中，由此，历史即作为自由意识的进步而得到理解。在黑格尔看来，一般而言，历史哲学的问题并不在于能否看出自然或者神圣天意的计划，或是将这种计划作为外在地考察经验历史的先天线索，而是在于一次性的历史过程是否包含了一种为其辩护的意义——这个问题的解答存在于时间和概念的辩证法当中，对此下面还会言及。

3. 与康德及其18世纪前人相对，在黑格尔这里，在理性支配的未来与无理性的过去之间并不存在一道鸿沟，因为对他来说，进步的目标不在未来，而在当下。\footnote{参见上书，第111页。}从当下出发，可以将历史回头构建为一个连续向前发展的过程。可以说，这里有着黑格尔历史哲学的限度，但也是其伟大之所在：通过历史连续性的思想，将历史中的进步合法化了。

4. 在康德及其前人那里，作为进步的历史活动于个人与类的关系的二律背反之中，但在黑格尔看来，这种关系就其本身及其与历史的关联而言都只是一种抽象物。就其本身而言是抽象的，因为在个体的自知、它的自我意识(Ich-Bewusstsein)中，类是同时得到思考的，并且其中也就有“我即我们”，或是黑格尔以“精神”之名为解释历史的进步奠定基础的那种原始的综合性的总体意识，在此总体意识中，个体和类两者达到了自由法意识；在与历史的关联中是抽象的，因为个体在其要素中本质上就是“普遍的自然”，即一个特定的民族，而类是一种严格意义上的历史的类，也就是那种世界精神的形态，它从各个民族和国家的过程中作为它们的历史而产生出来。\footnote{参见黑格尔，《法哲学原理》，第340节，第446页。《哲学全书》，第548节，《全集》（百年纪念版）第10卷，第426页以下。}

5. 尽管康德在回答“人类是否应该理解为不断地朝向改善而进步”时，\footnote{参见康德，《系科之争》（1798年），第二篇，1—10。}为了未来而先天地考虑了历史的问题——据此，讲史者制造并安排了他事先预告的那些历史事件——但他并未为了当下和过去而考虑存在于这种观念中的实践观点、人的创造性活动在历史中的作用：在一定程度上，实践被自然的计划消解掉了。相反，尽管黑格尔这里也有与康德的自然计划类似的“理性的狡计”学说，但他将历史、精神的王国视为人类劳动的产物和结果：“精神的领域由人产生。人能对神的国度提出各种各样的见解，因此总有一个精神的王国在人世中得到实现，并被人设定于实存之中。”\footnote{黑格尔：《历史中的理性》，第50页。}在字面意义上，历史中的进步由人所拟就：“然而，事实上，我们之所是，同时就是历史的【……】属于我们的、属于当今世界的理性的财富，不是直接产生的，也不是成长于当下的土壤；相反，其本质在于，它是一种遗产，进一步说，它是劳动的结果，并且是人类历代祖祖辈辈劳动的结果。”\footnote{黑格尔：《哲学史讲演录》第1卷，《全集》(百年纪念版)第17卷，第28页。}

\section*{4}

这里，我不详细讨论第二至五点，而只论述第一点。它最好地说明了我所言的“黑格尔是康德引入历史哲学中的哥白尼式转向的完成者”的主张。为此目的，我对康德和黑格尔与法国大革命之间的关系进行了比较。这个比较不是外在的表态，而是要突出这场革命对于他们的历史哲学思想的价值和意义。康德在18世纪90年代开始将人类统一进步的理念看作自然的单纯意图，以及在一定程度上与历史过程并列的想象的历史目的，并且已经致力于将人类统一进步的理念与历史过程联系起来。与《纯粹理性批判》中的图式论问题——应当思考单纯的知性概念如何运用于可能的经验对象——相类似，现在历史哲学也产生了在理念中设想的、通过自然意图予以保证的历史目的与实际的历史过程之间的关系问题。

康德用历史符号理论解决这个理念在经验中的图式化问题。这个问题在于，找出一个历史事件，它在时间上是不确定的，却暗示了人类作为其向改善推进原因的特性和能力，由此即可得出到达那种历史目标之可能性的结论——“然后这一结论还要能够这样地扩大到以往时代的历史（即它永远是在前进的），以至于那个事件的本身并不必须被看成是这种历史的原因，而是必须被看成只不过是一种示意、是一种历史符号(Signum rememorativum, demonstrativum, prognostikon)【……】”。\footnote{康德，《系科之争》，第二篇，5。（中译文引自〔德〕康德，《论教育学》，赵鹏、何兆武译，上海人民出版社2005年版，第116页。——译者）}对康德来说，这种历史符号——这里他的进步理论的二律背反又重新开始了——不是法国大革命，而是虽未参加然而同情法国革命的公众的道德反应。黑格尔在历史哲学讲座中报道了法国革命那些观众的热情和崇高情感；但在他看来，真正的历史符号乃是法国革命事件本身，其中包含了一种对于以后和以往的意义：对于以后具有意义，是因为现在，权利和自由的思想在历史上第一次被提高为政治制度的基础，“从此以后，一切都应该建立在这个基础之上”；对于以往具有意义，是因为自从阿那克萨哥拉将理性提高为宇宙的原则以来，哲学（再次）出现在这种思想的符号中。同时，这是从理性作为宇宙原则到理性作为历史原则的何种飞跃和迈进啊！黑格尔这样说明这种进展：“自从日丽中天、众星绕它运行以来，还从未见过此种事情：人立于头脑（即思想）上面，按照思想构筑现实。阿那克萨哥拉最早说过，努斯($\nu o\tilde{\upsilon}\varsigma$)统治世界；然而直到现在，人们才认识到，思想应该统治精神的现实。”\footnote{黑格尔：《世界历史哲学讲演录》第4卷，G. 拉松出版，莱比锡，1920年，第926页。}黑格尔以这种方式将法国革命本身理解为历史符号和权利与自由思想的现实化，由此，康德进步理念的图式化问题便获得了不同的解答。黑格尔用时间与概念、自我意识与时间的辩证法取代了那种图式化，由此，他的思想便实现了向现代世界以及发生在现代世界基础上的历史过程的突破。

哲学中揭橥自由之基础的事件，作为时间中的事件，发生在法国革命之前。卢梭在自由意志的思想中发现了一切权利的实体基础，康德在“我思”、意识的先验统一中发现了一切经验的最高条件和所有思想规定的起源。这两个事件共同构成了黑格尔哲学的基本经验。黑格尔哲学从与自身同一的自我的思想出发，同时结合费希特和早期谢林的思想对其进行了解释。自我不是空洞的统一性；相反，按照概念的系统构造，它是三个自身区分开来的环节的统一。首先，在康德的“我思”的意义上，是纯粹自身与自身发生关联的统一性，通过抽除所有内容，这种统一性回到与自身的无限制等同的自由之中，并将自己规定为普遍性；其次，是绝对规定了的存在，此种存在与他者对立，并排斥他者，此即个别性或者个人的人格性；第三，是绝对的普遍性，它同时又是直接的绝对个别化——自我的具体结构，黑格尔将它与概念的结构等同起来：普遍物与个体合而为一。\footnote{黑格尔：《逻辑学》，第二部分，莱比锡，1951年，第219页。}

“精神”就是这个结构向——黑格尔用来说明社会历史世界及其变化的——那个总体意识的辩证展开。然而，并不是“精神”的概念才包含了与时间的联系；相反，概念的概念、自我亦然——这是超出康德以及之前所有哲学的决定性步伐。在康德那里，尽管“我思”必然伴随时间性的表象过程，但它自身却在时间之外；对黑格尔来说，自我根本上就在时间之中：“自我在时间之中，时间就是主体自身的存在。”\footnote{黑格尔，《美学》，第三部分，《全集》（百年纪念版）第14卷，第151页。}自我不是从它与时间的关系出发规定自身；相反，它从时间与自我或概念的关系出发规定时间：时间就是“定在着的概念自身”\footnote{黑格尔，《哲学全书》，第258节，注释，《全集》（百年纪念版）第9卷，第258页；《精神现象学》，《全集》（百年纪念版）第2卷，第44、612页。}，或如黑格尔所言，时间“与自我（即纯粹自我意识的自我）是同一个原则”。在自我与时间的辩证法中，黑格尔实现了体系向历史的突破，一笔勾销了康德的图式化问题。康德则还想要用图式化去克服独断论形而上学的隐秘神话学以及与之相应的道德恐怖主义（人类发现自己不断朝着更恶倒退）、幸福主义（人类不断朝向改善前进）和阿布德拉主义（人类发现自己永远处于静止状态，愚蠢地忙忙碌碌乃是人类的特性）的荒谬历史主张。\footnote{参见康德，《系科之争》，第二篇，3，a-c。}

自我与时间的辩证关系使得在历史中指出康德将其适用于自然经验的那种联系和统一性成为可能。自我并不是时间在上面随意涂画的白板，也不再如那些独断的历史观念所见（康德的历史构想部分地也是如此），是那些相互交错的事件和意见的漠不相关的、没有自我的中途点\footnote{参见狄尔泰，《施莱尔马赫传》(Das Leben Schleiermachers)第2卷，2，哥廷根/柏林，1966年，第538页。}，而是时间和历史的统一，并在这种统一中自己出场。黑格尔把自我的这种出场叫作自我的自由。时间与自我（即纯粹自我意识的自我）是同一个原则，而纯粹自我意识同时就是自由的原则。\footnote{黑格尔：《哲学全书》，第258节，注释，《全集》（百年纪念版）第9卷，第79页。}

黑格尔按照自由意识进步理念的线索，基于这种自我、时间和自由的历史辩证法，对历史进行解释。这种辩证法将黑格尔哲学的基本经验融入其哲学体系之中。现在，在法国革命的历史事件中得到实现的市民社会的理性理想，就能通过实际发生的历史而得到辩护。我在其中看到了由康德引入的历史哲学中的哥白尼式转向的完成。这种通过时间性事件所做的辩护并不意味着，当黑格尔在前后相继的世界事件的尺度上按照成果理性对进步进行量度时，将历史化为一个线性的过程。情况恰好相反。当黑格尔预设的市民社会的理性理想在国家中成为现实，而现代世界的社会历史现实变得理性时，则自由的意识也就不再归结为单义的进步运动，而是归结为双义的进步和回归的运动。悖论的表达就是：真正的进步、它的完成就是回归——黑格尔的理论用它实现了精神的体系概念，即精神在它自己的别处也是自在的或是自由的：“理性认识真实的东西，自在自为的存在者，这种东西是没有限制的。精神的概念就是回归自身，以自己为对象；因此，前进就不是所向无定地向无限物进发，而是有一个目的，即回归自身。”\footnote{黑格尔：《历史中的理性》，第181页。参见同书，第58、65、67、70页以下。}回归的结构既不是说，进步在历史中到达了终点——黑格尔排除了这种看法，因为从地理上看，自由的实现，作为世界历史的目标，仅仅出现世界的一个角落，并且“从昨天”才开始\footnote{参见黑格尔：《法哲学原理》，第62节，附释。}，同时它也没有预设，历史运动就是重复——同时代人的复辟企图意义上的或是更加古老的、直到维科为止一直有效的政治制度的循环论意义上的——过去的社会形式。黑格尔拒斥浪漫派国家理论中的中世纪教义，称之为“我们时代的无聊，想把其中的精髓变成口号”。\footnote{黑格尔：《世界历史哲学》第4卷，第840页。}尽管他想要在他自己的时代和晚期罗马帝制时代之间确定一些共同的结构，但他绝不认为在历史衰落时期之间具有“重大相似性”，而在从福尔格拉夫(Karl Vollgraff)到拉绍尔克斯(Ernst von Lasaulx)和布克哈特再到斯宾格勒的19世纪保守派历史理论中，这种相似性发挥了显著作用。

这样，内在于历史运动之中的进步，就像历史运动一样，是不可逆的。这并不是说，进步就是直线式的，以一种直接的、无对立的自然过程的方式行进的。如同历史中的所有变化一样，每一进步都是矛盾的环节——有和无、产生和消逝、进展和衰落——的统一。至关重要的是，黑格尔将变化思考为一个缓慢的、没有断裂的过程：“用变化的渐进性来理解产生和消逝，就是同语反复所特有的无聊；那意味着：正在产生或者消逝的东西，预先就已经是现成的了，而变化则成了外在区别的简单改变，这样实际上就是同语反复。”\footnote{黑格尔：《逻辑学》，《全集》（百年纪念版）第4卷，第461页。（中译文引自《逻辑学》上卷，第405页。译文略有改动。——译者）参见伽兰第(R. Garaudy)：《上帝死了。对黑格尔的一个研究》(Gott ist tot. Eine Studien über Hegel)，柏林，1965年，第386页以下。}因此，进步首先是在历史运动的关节点和转折点上、在时间与自我发生分离的时期得到实现的。在这种时期，一个个人或一个民族的精神不再与其制度相一致，制度也不再与精神一致。于此，可以区分出进步的三种模式，据黑格尔看来，它们都在历史中出现，同样必然，同样正当：革命的模式、改革的模式以及最后作为世界精神的终极理由(ultima ratio)——一个民族被排除在历史进步民族的序列之外：“如果一个国家的法制所表示的真理已经不符合于它的自在本性，那么它的意识或概念与它的现实性就存在着差别，它的民族精神也就是一个分裂了的存在。有两种情况可以发生：这民族或者由于一个内部的强力的爆破，粉碎了那现行有效的法律制度，或者较平静地、较缓慢地改变那现行有效的但却已不复是真的伦理、已不能表现民族精神的法律制度。或者（第三种模式即在其中）一个民族缺乏理智和力量来作这种改变，因而停留在较低劣的法律制度上；或者另一个民族完成了它的较高级的法制，因而成为一个较卓越的民族，与之相对，而前一个民族必定会不再为一个民族，并受制于这个较卓越的民族。”\footnote{黑格尔：《哲学史讲演录》第2卷，《全集》（百年纪念版）第18卷，第276页以下。（中译文引自《哲学史讲演录》第二卷，第249-250页。——译者）}第三种模式包含了作为世界历史的现实历程出现在《法哲学原理》结尾的那些有限的民族精神的那种“辩证发展现象”。这是历史哲学在哲学体系中的位置。如果我们想要在黑格尔的体系和历史统一的视角下理解康德引入的方法论转向的意义，就必须说明黑格尔历史哲学的这种位置。

\section*{5}

在黑格尔的阐述中，国家的普遍理念在其完成的顶点下降为在历史的活动中彼此对立的“四种”特殊的“国家”。现代世界的国家与市民社会处于差异的关系之中，自己平衡了市民社会的内部差异。它不再是传统的政治状态概念，也不是古典政治学和近代自然法理论的不变的和无历史的自然模式或契约模式；相反，历史——作为实践哲学的新的上司——高踞其上。在历史中，“伦理整体本身、国家的独立性都被委诸偶然性之下”（第340节）。黑格尔在《法哲学原理》体系最后引入的历史，就是从霍布斯到卢梭的自然法学说于其体系开端处设定的那种自然状态理念的现实性。在黑格尔这里，自然法理论家们设想的从自然状态到市民社会并在与国家的同一中静止下来的运动，不是始于国家同市民社会的关系，而是始于与其他国家的关系。这种自然状态不是虚构的状态，而是现实的状态，是历史的运动。《法哲学原理》将这个运动接受到其体系之中，从而摆脱了法和社会的抽象自然理论。在世界历史中，“普遍精神”的定在的要素，作为法概念自我展开的起点和终点，就是“精神性的现实性”，在这种现实性中，人和物的“自然”就不是固定不变的规律，而是黑格尔思考为自由的概念，它“在其内在性和外在性的全部范围内”实现其自身。在这个与法哲学体系结束和历史哲学开端同时发生的运动的辩证法中，时间、概念与自由统一起来了。历史作为自由意识中的进步，就是对于在现代宪法国家中得到实现的那种自由的辩护。

由此出发，即可明白黑格尔历史哲学提出的非同寻常的要求。历史哲学的主题和对象不再是对历史过程本身无甚兴趣、根据一条先天的线索写就的一种可能的历史的理念，而是可以理解为实在性的历史本身。这种要求并非现在才被人们斥为幻想，而我也绝非想要以一个正统黑格尔派的态度着力为其辩护。最后，在我为其辩护时，也只想提出一点意见，这首先是为了质疑那些罕见一致地臆想自己正在接近一个后历史时代的现代学派——结构主义、实证主义和精神科学的解释学——对解放的批判的历史概念的抹黑。事实上，黑格尔的辩证历史解释提出的要求可以简化为列维—斯特劳斯\footnote{参见列维—斯特劳斯：《野性的思维》，巴黎，1962年，德文版，美茵河畔法兰克福，1968年，第292页。}拿来解释萨特的辩证理性批判问题的那个表达式：在何种条件下法国革命的神话是可能的？但是人们可以从黑格尔哲学得到一个标准（这在萨特这里也许没那么容易），从而不可能将历史和神话混同起来。并且我认为，这个标准存在于黑格尔确定的概念（或自我）与时间的关系中。正如时间的结构必然以自我的结构为前提，自我也同样以时间为前提：在严格意义上，革命是时间中的事件，是历史符号，它作为这样的符号意味着进步，而不是重复同一结构。古典政治（我用它表示其哲学维度。在我看来，在此维度中，必然要讨论神话和历史问题）要么回溯到神话（如柏拉图从《理想国》到《蒂迈欧》的宇宙论及其关于时间的有规律的轮回学说的过渡），要么回溯到循环的时间结构的神话考古学（如亚里士多德《政治学》第1卷），以巩固城邦的存在；现代自然法则在源自历史的自然状态的学说中进行虚构，创造出一部哲理小说，以使国家作为市民社会的生成过程能够止步于国家之中——已经有了一部历史，然而再也没有历史了。黑格尔让他的《法哲学原理》体系纲要结束于对市民社会和国家之实在性的历史展望中，在这种实在性面前，市民社会和国家都“只是作为观念物而存在”。\footnote{黑格尔，《法哲学原理》，《全集》（百年纪念版）第7卷，第446页。（原引文引自《法哲学原理》第341节，而非第340节。中译文参见《法哲学原理》，第351页。——译者）}出现于历史中的辩证法确实是一种辩护的辩证法，但也是一种克服的辩证法。黑格尔说过，唯有当下存在，以前和以后都不存在；\footnote{黑格尔，《哲学全书》，第二部分，第259节，补充，《全集》（百年纪念版）第9卷，第86页。}然而他还补充道，具体的当下是过去的结果，并且孕育了未来。真正的当下是概念把握到的时间整体。

\chapter*{各章来源}
\addcontentsline{toc}{chapter}{各章来源} % 手动加入目录
\markboth{各章来源}{各章来源} % 设置页眉

“客观精神与实践哲学”，在1968年里尔举办的第一届国际黑格尔大会所作演讲的扩充稿。

“制度中的辩证法”，在帕多瓦大学的演讲（1976年），未刊稿。

“自由法则与自然的统治”，摘自《体系与历史》，美茵河畔法兰克福，1973年。

“黑格尔对自然法的批判”，载《黑格尔研究》1967年第4卷。

“对国民经济学的接受”，尚未出版的黑格尔法哲学和国家哲学中的一章的扩充稿。

“黑格尔《法哲学原理》中的传统与革命”，载《哲学研究杂志》1962年第16卷。

“‘市民社会’概念及其历史起源问题”，载《法哲学与社会哲学档案》1962年第48卷。

“黑格尔历史哲学中的进步与辩证法”，摘自《体系与历史》，美茵河畔法兰克福，1973年。
\chapter*{译后记}
\addcontentsline{toc}{chapter}{译后记} % 手动加入目录
\markboth{译后记}{译后记} % 设置页眉

本书作者为当代德国著名哲学家曼弗雷德·里德尔(Manfred Riedel，1936—2009)。里德尔于1936年5月出生在莱比锡的一个乡村，1954年进入莱比锡大学，1957年转入海德堡大学，曾先后师从恩斯特·布洛赫、卡尔·洛维特、伽达默尔等当代德国哲学大师。1960年，他在洛维特的指导下以《黑格尔思想中的理论与实践》一文获得博士学位。1968年，他完成了教授资格论文《市民社会：一个古典政治学和现代自然法范畴》。他曾先后在海德堡大学、马堡大学、萨尔布吕肯大学任教，1970年成为埃朗根-纽伦堡大学哲学正教授，1980—1981年担任纽约社会研究新学院休斯(Theodor-Heuss)客座教授，并先后在都灵、罗马、威尼斯等地担任客座教授。1992—1993年，他担任耶拿大学席勒讲座教授，1992年获得哈勒-维滕堡大学的实践哲学教席，并于2004年荣休。此外，里德尔曾于1991年至2003年间担任德国海德格尔协会主席，并从2005年起担任佛罗伦萨意大利人文科学院院士。里德尔的研究涉及康德、黑格尔、尼采、狄尔泰和海德格尔等人，从古代、中世纪到欧洲近代的市民社会史，构成其研究的一个重点，他的《在传统与革命之间：黑格尔法哲学研究》即为这方面研究的一部力作，也是当代德国学界黑格尔法哲学研究方面的一部经典作品。

该书第1版于1969年以《黑格尔法哲学研究》(Studien zu Hegels Rechtsphilosophie)之名出版(Frankfurt am Main: Suhrkamp, 1969)，次年刊行第2版。1982年，里德尔将原书内容稍加扩充，并加上主标题“在传统与革命之间”（原标题则保留为副标题）后，由位于斯图加特的克莱特-科塔(Klett-Cotta)出版社出版，是为第3版。自第1版出版以来，先后有意大利文、日文、西班牙文、韩文、英文等诸种译本问世，其影响也随之扩展到国际学界。当原著出版50周年之际，其中译本能够面世，固然有些恨晚，但也不能不说是一件幸事。

本书的翻译出版首先要感谢商务印书馆“政治哲学名著译丛”主编吴彦博士的盛情邀请。此次翻译根据的是里德尔原书1982年扩充版，并且参考了赖特(W. Wright)的英译本(Manfred Riedel, Between Tradition and Revolution: The Hegelian Transformation of Political Philosophy, Cambridge University Press 1984)。翻译的具体分工是：前言部分由两位译者共同完成，第一、第二部分的翻译由黄钰洲承担，并由朱学平进行了校对；第三、第四部分的翻译由朱学平承担。“术语索引”和“人名索引”由黄钰洲录入。需要指出的是，我国罗马法学者普遍将“jus civile”译为“市民法”，本书也同样将“societas civilis”译为“市民社会”，其含义是“公民社会”，而非黑格尔意义上的“市民社会”。由于译者学养和语言水平有限，译文中必然存在错误不当之处，对此，译者真诚期待读者和学界同人不吝批评指正。最后，我们也衷心感谢商务印书馆的编校者为本书出版所付出的辛勤劳动！

\begin{flushright}
译者\\
2019年1月
\end{flushright}

\end{document}